% ============================================================================
%  main.tex
% ----------------------------------------------------------------------------
%  "C, Memory, and Systems" - A Definitive Course in C Programming,
%  Data Structures, and Systems Architecture.
%
%  Pedagogically inspired by the C course of Salvatore Sanfilippo (antirez):
%  start from the raw realities of memory, bytes and addresses; build real,
%  production-shaped software; make every memory operation explicit.
%
%  Compile with:  ./build.sh        (runs pdflatex twice, opens zathura)
% ============================================================================

\documentclass[11pt,oneside]{book}

% ---------------------------------------------------------------------------
%  Packages
% ---------------------------------------------------------------------------
\usepackage[T1]{fontenc}
\usepackage[utf8]{inputenc}
\usepackage{lmodern}

\usepackage[margin=2.4cm, headheight=15pt, footskip=1.2cm]{geometry}

\usepackage[table]{xcolor}
\usepackage{amsmath,amssymb}
\usepackage{graphicx}
\usepackage{microtype}
\usepackage{setspace}

\usepackage{fancyhdr}
\usepackage{enumitem}
\usepackage{titlesec}

\usepackage{tikz}
\usepackage{ifthen}
\usepackage{listings}
\usepackage{fancyvrb}

\usepackage[many]{tcolorbox}
\usepackage[hidelinks]{hyperref}

% ---------------------------------------------------------------------------
%  Colors (palette tuned for printed and screen reading)
% ---------------------------------------------------------------------------
\definecolor{cbk}{HTML}{3A3A3A}         % base text
\definecolor{ccmt}{HTML}{6A9955}        % comments
\definecolor{ckw}{HTML}{0000B0}         % keywords
\definecolor{cstr}{HTML}{A31515}        % strings
\definecolor{cnum}{HTML}{098658}        % numbers
\definecolor{cbg}{HTML}{F7F7F5}         % code background
\definecolor{cframe}{HTML}{C8C8C4}      % code frame
\definecolor{cins}{HTML}{E8F0FE}        % insight box background
\definecolor{cwar}{HTML}{FDF1E0}        % warning box background

% ---------------------------------------------------------------------------
%  Code listings: an explicit, readable C style.
%
%  We deliberately do NOT shade clever one-liners. Every snippet here uses
%  explicit statements, named temporaries and separate lines for pointer
%  moves, so the reader can trace memory exactly.
% ---------------------------------------------------------------------------
\lstdefinestyle{csrc}{%
  language=C,
  basicstyle=\ttfamily\small\color{cbk},
  keywordstyle=\color{ckw}\bfseries,
  commentstyle=\color{ccmt}\itshape,
  stringstyle=\color{cstr},
  numberstyle=\tiny\color{gray},
  showstringspaces=false,
  tabsize=4,
  breaklines=true,
  breakatwhitespace=true,
  frame=single,
  rulecolor=\color{cframe},
  backgroundcolor=\color{cbg},
  captionpos=b,
  numbers=left,
  numbersep=8pt,
  xleftmargin=1.6em,
  framexleftmargin=1.6em,
  aboveskip=1.2em,
  belowskip=1.2em,
  morekeywords={size_t,ssize_t,uint8_t,uint16_t,uint32_t,uint64_t,int8_t,
                int16_t,int32_t,int64_t,offsetof,NULL,malloc,calloc,realloc,
                free,sizeof,struct,typedef,enum,union,const,volatile,restrict,
                static,inline,printf,fprintf,sprintf,snprintf,scanf,fgets,
                puts,exit,abort,assert,memcpy,memset,memmove,strlen,strcpy,
                strncpy,strdup,strcmp,strncmp,strtol,errno,perror}
}

\lstset{style=csrc}

% ---------------------------------------------------------------------------
%  Verbatim: ASCII diagrams must never be split across a page boundary,
%  or they render as "half broken" art that jumps to the next page.
% ---------------------------------------------------------------------------
\fvset{samepage=true}
\RecustomVerbatimEnvironment{verbatim}{Verbatim}{}

% ---------------------------------------------------------------------------
%  Boxes for pedagogy: insights and rules of thumb.
% ---------------------------------------------------------------------------
\newtcolorbox{insight}[1][]{%
  enhanced, breakable, colback=cins, colframe=blue!55!black, boxrule=0.6pt,
  arc=2mm, left=3mm, right=3mm, top=2mm, bottom=2mm,
  title={\textbf{\sffamily Insight}}\ifthenelse{\equal{#1}{}}{}{ --- \textbf{#1}},
  fonttitle=\normalfont, attach boxed title to top left={yshift=-2mm, xshift=3mm},
  colbacktitle=cins, coltitle=blue!55!black, boxed title style={empty}
}

\newtcolorbox{ruleofthumb}[1][]{%
  enhanced, breakable, colback=cwar, colframe=orange!80!black, boxrule=0.6pt,
  arc=2mm, left=3mm, right=3mm, top=2mm, bottom=2mm,
  title={\textbf{\sffamily Rule of Thumb}}\ifthenelse{\equal{#1}{}}{}{ --- \textbf{#1}},
  fonttitle=\normalfont, attach boxed title to top left={yshift=-2mm, xshift=3mm},
  colbacktitle=cwar, coltitle=orange!80!black, boxed title style={empty}
}

\newtcolorbox{osconcept}{%
  enhanced, breakable, colback=blue!5, colframe=blue!60!black, boxrule=0.6pt,
  arc=2mm, left=3mm, right=3mm, top=2mm, bottom=2mm,
  title={\textbf{\sffamily Why this is an OS problem}},
  fonttitle=\normalfont, attach boxed title to top left={yshift=-2mm, xshift=3mm},
  colbacktitle=blue!5, coltitle=blue!60!black, boxed title style={empty}
}

\definecolor{cpurp}{HTML}{5B2D8E}
\newtcolorbox{researchnote}[1][]{%
  enhanced, breakable, colback=purple!5, colframe=cpurp, boxrule=0.6pt,
  arc=2mm, left=3mm, right=3mm, top=2mm, bottom=2mm,
  title={\textbf{\sffamily Research Corner}}\ifthenelse{\equal{#1}{}}{}{ --- \textbf{#1}},
  fonttitle=\normalfont, attach boxed title to top left={yshift=-2mm, xshift=3mm},
  colbacktitle=purple!5, coltitle=cpurp, boxed title style={empty}
}

\newcommand{\paperref}[2]{%
  \par\medskip\noindent{\small\textbf{Reference:}\ #1\ ---\ \url{#2}}\par
}

\newcommand{\code}[1]{\texttt{\small #1}}

% ---------------------------------------------------------------------------
%  Page style
% ---------------------------------------------------------------------------
\pagestyle{fancy}
\fancyhf{}
\fancyhead[L]{\small\textit{C, Memory, and Systems}}
\fancyhead[R]{\small\textit{\nouppercase{\leftmark}}}
\renewcommand{\headrulewidth}{0.4pt}
\fancyfoot[C]{\small\thepage}

\titleformat{\chapter}[display]
  {\normalfont\huge\bfseries}{\chaptertitlename\ \thechapter}{18pt}{\Huge}
\titlespacing*{\chapter}{0pt}{-20pt}{18pt}

% ---------------------------------------------------------------------------
%  Title page
% ---------------------------------------------------------------------------
\title{\vspace{-2.5cm}\Huge\textbf{C, Memory, and Systems}\\[0.8em]
       \Large A Definitive Course in C Programming, Data Structures,\\
       and Systems Architecture\\[1.2em]
       \normalsize\emph{Beginning from bytes and addresses, building toward databases,}
       \newline\normalsize\emph{network servers, and virtual machines.}}
\author{The Open Book of the Machine}
\date{}

% ---------------------------------------------------------------------------
%  Document
% ---------------------------------------------------------------------------
\begin{document}

\frontmatter
\maketitle
\tableofcontents

\mainmatter

% Every chapter is guarded so the project compiles even while later
% chapters are still being written. The full directive list for all
% seventeen chapters lives here.
\IfFileExists{chapter1.tex}{% ============================================================================
%  chapter1.tex
%  The Foundations of Memory and Pointers
% ============================================================================

\chapter{The Foundations of Memory and Pointers}

\section{The Machine Underneath}

Before we talk about C, we must talk about the machine. Everything a
computer does reduces to a staggeringly simple model: there is a large
array of bytes, each byte has a number (its \emph{address}), and the
processor runs a loop that reads instructions from memory and acts on
them. The CPU has no understanding of ``variables'', ``strings'', or
``objects''. It understands \emph{addresses} and \emph{bytes}. When we
write C, we are simply describing, with a thin layer of sugar, what the
machine will actually do.

This is the heart of the Sanfilippo approach to C: we never hide the
machine from the student. If you understand bytes, addresses, and
lifetimes, then every higher-level concept---data structures, network
servers, databases, virtual machines---is just a disciplined arrangement
of these three ideas. If you skip this foundation, every later chapter
will feel like magic, and magic breaks in production.

Let us begin with the smallest unit of storage the CPU can address
directly: the \emph{byte}.

\begin{verbatim}
   A byte is 8 bits. Each bit is a switch: 0 (off) or 1 (on).

   Address 0x7ffe29a00000:

        +---+---+---+---+---+---+---+---+
        | 0 | 1 | 0 | 0 | 0 | 0 | 1 | 1 |   = 0x43 = 67 = ASCII 'C'
        +---+---+---+---+---+---+---+---+
         bit7 bit6 bit5 bit4 bit3 bit2 bit1 bit0

        bit0 has value 1 (2^0), bit1 has value 2 (2^1),
        bit2 has value 4 (2^2), and so on, up to bit7 = 128.
        This byte represents the number 64 + 2 + 1 = 67.
\end{verbatim}

Every byte in the address space has a unique address. On x86-64 the
implemented virtual address space is actually numbered from $0$ to
$2^{48}-1$ (or to $2^{57}-1$ with five-level paging); the 64-bit register width
would allow $2^{64}$, but the hardware does not implement that full range. Either
way, it is so large that we will never fill it. We write addresses in hexadecimal
because it maps cleanly onto groups of four bits, and because writing
$2^{64}$ as decimal numbers is a waste of human attention.

\begin{insight}[Why hexadecimal]
A byte is two hexadecimal digits. One hex digit is four bits---a
\emph{nibble}, half a byte. That is
why \code{0x43} reads like a byte: \code{0x} is the prefix, \code{4} is
the high nibble, \code{3} is the low nibble. When you see a memory
dump, you can mentally group it into bytes by pairs of digits.
\end{insight}

\subsection{Bits, Bytes, and Words}

Three sizes will keep recurring, so let us settle them now. A
\emph{bit} is one switch, 0 or 1. A \emph{byte} is eight bits and is
the smallest addressable unit. A \emph{word} is the machine's
\emph{natural unit of data}: the size of a general-purpose register,
and the size the CPU prefers to load and store in one operation.

\begin{verbatim}
   The three units, drawn to scale:

     a bit:        a byte (8 bits):        a word (a register):
                   1 0 1 1 0 0 1 0
     +---+          +---+---+---+---+---+---+---+---+       +---------------+
     | 1 |          | 1 | 0 | 1 | 1 | 0 | 0 | 1 | 0 |       |   0x00000043  |
     +---+          +---+---+---+---+---+---+---+---+       +---------------+
                                            ^
                                            ^   four bytes grouped:
                                                the natural 64-bit word
                                                = 0x0000000000000043

   On a 64-bit machine, one word is 8 bytes and holds up to 2^64
   different values -- enough to address every byte in the address
   space. That is no coincidence: a POINTER is exactly one word.
\end{verbatim}

The word size is the machine's identity, and it has consequences you
will feel all book long:

\begin{itemize}
  \item A \emph{pointer} is one word. On a 32-bit machine pointers are
    4 bytes; on a 64-bit machine they are 8. That is why
    \code{sizeof(int *)} printed 8 in Chapter~1's table, and why the
    address space is $2^{64}$ bytes on modern systems.
  \item The CPU likes \emph{aligned} access: reading a word from an
    address divisible by the word size is fastest; straddling a
    boundary costs extra. Alignment, and the padding it forces into
    structs, is the subject of Chapter~3.
  \item The word survives in our vocabulary: a \emph{double word} is
    two words, hex dumps are laid out in words, and ``word size'' is
    how we say ``how wide is this machine?''.
\end{itemize}

\begin{verbatim}
   Quick sanity check -- every claim above in one line:

     64-bit machine   word = 8 bytes   pointer = 8 bytes
     32-bit machine   word = 4 bytes   pointer = 4 bytes
     sizeof(int *) always equals sizeof(void *) on any one machine:
     both are one word, whatever that word is.
\end{verbatim}

\begin{insight}[Why we say "byte" and not "word" for addresses]
The CPU addresses memory by \emph{byte}, not by word: consecutive
addresses differ by one byte. Words are built from bytes at aligned
addresses. So while a register moves one word at a time, an address
still names one byte---and pointer arithmetic from Chapter~1 works in
units of the pointed-to type, not in words. Both facts are true at
once: bytes are addressable, words are natural.
\end{insight}

\subsection{Registers, Cache, and RAM: The Memory Hierarchy}

Between the CPU and the huge, slow RAM we have been assuming sits
something your code never sees but that decides how fast it runs:
the \emph{memory hierarchy}. Access to a byte is not uniform. The CPU
keeps a handful of registers (fastest, tiny), then one or more
\emph{caches} (fast, medium-sized), then RAM (slower, huge), then
disk (far slower, enormous). Each level holds a copy of data that the
next level up is using, in blocks.

\begin{verbatim}
   The hierarchy, with the "distance from the CPU":

     level      size      speed        managed by
     -------    -------   ----------   ---------------
     registers   ~100s B   ~1 ns        the program (you!)
     L1 cache     ~32 KB    ~1 ns       the CPU, automatically
     L2 cache    ~256 KB    ~3 ns       the CPU, automatically
     L3 cache    ~8 MB      ~10 ns      the CPU, automatically
     RAM         ~GB        ~100 ns     the OS (the page model)
     disk        ~TB        ~ms         the OS (filesystems)

   (Numbers are indicative; what matters is the SHAPE: each level is
    bigger and slower, and faster levels cache the slower ones.)
\end{verbatim}

A \emph{cache} is a small, fast store that keeps recently used data so
the CPU rarely has to go to the slower level. When the CPU wants a
word, it looks in the cache first; only a \emph{cache miss} sends it
to RAM. Data moves between cache and RAM in fixed-size units called
\emph{cache lines}---typically 64 bytes. This one fact explains more
systems performance than any other:

\begin{verbatim}
   Why cache lines matter to your C code:

     two adjacent ints:  cache line = [ 42 ][ 43 ]...  one fetch loads
     both. A struct's fields arrive together (ch.3's AoS).

     two ints 1000 bytes apart: two cache lines, two fetches. If an
     algorithm walks field-by-field across many structs, it drags
     whole cache lines it never uses.

   Consequence: the COST of an operation depends on which cache lines
   it touches, not just how many instructions it runs. This is why
   Chapter 3's "array of structs vs struct of arrays" decision, and
   Chapter 4's unrolled lists, are performance choices, not style.
\end{verbatim}

The hierarchy also explains the numbers this book quotes. ``A cache
miss costs about 100 times a cache hit'' is the real reason aligned,
sequential, localized data wins. And the page model of the OS sits one
level down: RAM caches disk, exactly as L1 caches RAM. Same shape,
different scale.

\begin{insight}[Cache is invisible and irresistible]
The CPU manages the caches with no help from your program. You cannot
store into L1 ``by hand.'' But you \emph{can} arrange your data so
that the automatic caching works for you---contiguous, aligned, walked
in order. That arrangement is most of what ``performance engineering''
means, and it runs through every later chapter.
\end{insight}

\section{Sizes of the Primitive Types}

C's primitive types are contracts about the interpretation of a range of
bytes. The type tells the compiler (a) how many bytes the object
occupies, and (b) how to interpret those bytes. The standard defines
\emph{minimum} sizes, but on our target machines---a modern 64--bit
x86-64 or ARM64 system---the sizes are stable:

\begin{verbatim}
   type            size (bytes)   typical range
   --------------  ------------   ------------------------------------
   char            1              -128 .. 127           (signed)
   unsigned char   1                0 .. 255
   short           2             -32768 .. 32767
   unsigned short  2                0 .. 65535
   int             4              -2147483648 .. 2147483647
   unsigned int    4                0 .. 4294967295
   long            8              (on 64-bit Linux: 8 bytes)
   unsigned long   8
   float           4              ~7 decimal digits of precision
   double          8              ~15 decimal digits of precision
   pointer         8              (8 bytes, an address in the space)
   --------------  ------------   ------------------------------------
\end{verbatim}

Do not trust your memory of this table on every platform: in C, you
always ask the compiler for the true size at build time. That is the
job of \code{sizeof}.

\begin{lstlisting}[caption={Sizes as reported by the compiler itself.}]
#include <stdio.h>

int main(void) {
    printf("char                = %zu bytes\n", sizeof(char));
    printf("short               = %zu bytes\n", sizeof(short));
    printf("int                 = %zu bytes\n", sizeof(int));
    printf("long                = %zu bytes\n", sizeof(long));
    printf("double              = %zu bytes\n", sizeof(double));
    printf("int pointer         = %zu bytes\n", sizeof(int *));
    printf("void pointer        = %zu bytes\n", sizeof(void *));
    return 0;
}
\end{lstlisting}

Notice the format specifier \code{\%zu}: the letter \code{z} says
``the argument is a \code{size\_t}'', which is exactly the unsigned
type \code{sizeof} yields. You will meet \code{size\_t} constantly;
it is the type of every count in the standard library.

\begin{ruleofthumb}[The Loop Counter Rule]
In this course we use \code{size\_t} \emph{only} for loop counters that
iterate over arrays and buffers. Everywhere else---indices computed from
user input, signed arithmetic, difference between addresses---we use
plain \code{int} or \code{long}. This keeps signs explicit and avoids
the silent traps of unsigned wraparound when you subtract one index from
another.
\end{ruleofthumb}

\section{The Address Space of a Process}

When the operating system launches your program, it does not hand your
process raw physical RAM. It constructs a \emph{virtual address space}:
a private, contiguous-looking map of addresses in which your program
lives. Each process gets the illusion of owning the whole machine. The
map is laid out in a fixed order, from low addresses to high addresses:

\begin{verbatim}
   +--------------------------------------------------+  high addresses
   |                                                  |  0x7fffffffffff
   |                     STACK                        |
   |  local variables, return addresses, arguments    |
   |  grows DOWN toward lower addresses               |
   |  |                                               |
   |  v                                               |
   +--------------------------------------------------+
   |              (free region; this is where          |
   |               memory mappings and shared libs    |
   |               get placed)                        |
   +--------------------------------------------------+
   |  ^                                               |
   |  |                                               |
   |                     HEAP                         |
   |  memory handed out by malloc()/calloc()/free()  |
   |  grows UP toward higher addresses                |
   +--------------------------------------------------+
   |                     BSS                          |
   |  uninitialized globals and statics (zeroed)     |
   +--------------------------------------------------+
   |                     DATA                         |
   |  initialized globals, statics, string literals  |
   +--------------------------------------------------+
   |                     TEXT                         |
   |  the machine code of your program                |
   +--------------------------------------------------+  low addresses
   |                     0x0  (unmapped, fault)       |
   +--------------------------------------------------+
\end{verbatim}

This diagram is worth memorizing; we will return to it in every later
chapter. Three regions dominate our mental model:

\begin{description}
  \item[TEXT] Read-only. The instructions themselves. If you try to write
    here, the CPU raises a fault.
  \item[DATA/BSS] Globals and statics. BSS is not stored in the binary;
    the OS simply promises to give you zeroed memory at startup.
  \item[STACK] Automatic (local) variables, function call frames, return
    addresses. It grows \emph{down}.
  \item[HEAP] Everything you allocate at runtime with \code{malloc}.
    It grows \emph{up}.
\end{description}

The stack and the heap grow toward each other. This is why the phrase
``stack overflow'' exists: a runaway recursive function walks the stack
down until it collides with something it may not touch, and the
operating system kills the process. Not a metaphor; a literal collision
of addresses. (A \emph{recursive} function is simply one that calls
itself; each call pushes another frame onto the stack, and one that
never stops pushing is what ``overflows'' it. We write real recursive
functions---tree traversal, directory walking---starting in
Chapter~4.)

\subsection{Virtual Memory and Pages}

Two more words from the diagram need exact definitions now, because
they will carry every chapter from here on.

\emph{Virtual memory} is the mechanism behind the ``virtual address
space'' just described. Your program's addresses---the \code{0x7ffc}
numbers you printed in Chapter~1---are \emph{virtual}: they are not
physical RAM addresses at all. The kernel keeps a translation table,
the \emph{page table}, and the CPU consults it on every access,
converting your virtual address into the physical address where the
byte actually lives. This gives every process a private, huge address
space, \emph{isolation} (your page table never contains another
process's physical pages), and \emph{flexibility} (a virtual region can
be backed by RAM, by a file, or by nothing yet).

A \emph{page} is the unit of that machinery. Virtual memory is not
managed byte by byte; it is chopped into fixed-size chunks---4~KB on
our machines---and each chunk is mapped, protected, and shared as a
whole.

\begin{verbatim}
   Virtual memory, drawn:

   YOUR address space (virtual)        PHYSICAL RAM        the page table
   +-----------------+                +-----------+        maps virtual
   | stack (pages)   |                |           |        page -> physical
   | heap  (pages)   |  ---page---->  |           |        frame
   | data  (pages)   |  table lookup  |           |
   | text  (pages)   |                |           |
   +-----------------+                +-----------+

   a virtual address: 0x7ffc1004
        |
        | split by the CPU into:  page number  0x7ffc   +  offset 0x004
        |                         (look up in the page table)
        v
   physical address: the RAM cell that holds the byte

   - A page that is "not present" in the page table is not in RAM;
     touching it makes the CPU raise a PAGE FAULT, and the kernel
     decides what that means (load the page, or kill the process).
   - This is why the address space can look far larger than RAM: most
     pages are simply not present, and RAM holds only what is in use.
   - Each process has its OWN page table, which is what makes two
     processes using the same virtual address able to hold different
     data -- and what makes process A unable to read process B.
\end{verbatim}

The page explains several things you have already seen. A 4096-byte
buffer in Chapter~5 is exactly one page. The stack that ``overflows''
is a set of pages the kernel maps on demand and revokes on fault. And
the mmap machinery of Chapter~12 and the allocators of Chapter~13 are
both, at bottom, negotiations over pages.

\begin{verbatim}
   One page, drawn at scale (4 KB = 4096 bytes = 512 eight-byte words):

     +--------------------------------+  address P (page-aligned)
     |                                |
     |         4096 bytes             |
     |         (this page)            |
     |                                |
     +--------------------------------+  address P + 4095

   Pages are PAGE-ALIGNED: a page's address is a multiple of 4096.
   That is why malloc() returns 16-byte-aligned pointers but the heap
   itself is built from page-sized regions.
\end{verbatim}

\begin{insight}[A page is the OS's atom]
The OS never manages memory in bytes; it manages in pages. Everything
that costs you performance---a page fault, a copy-on-write, a
shared-memory mapping, a demand-paged file---happens page by page, at
4~KB granularity, and nothing you do in user space is ever smaller
than that at the OS level. When a program feels "slow on first
access," the page is the culprit.
\end{insight}

\section{Addresses and Pointers}

A \emph{pointer} is a variable whose value is an address. The type of a
pointer records two things: the kind of data stored at the destination,
and the fact that the value is an address. The type is a contract; the
value is just an integer that happens to be a number of a memory cell.

\begin{lstlisting}[caption={A pointer is a variable holding an address.}]
#include <stdio.h>

int main(void) {
    int value = 42;
    int *ptr;                 /* "ptr will hold the address of an int" */

    ptr = &value;             /* &value: take the address of value */

    printf("value  = %d\n", value);
    printf("&value = %p\n", (void *) &value);
    printf("ptr    = %p\n", (void *) ptr);
    printf("*ptr   = %d\n", *ptr);

    return 0;
}
\end{lstlisting}

The \code{\&} operator, called \emph{address-of}, produces the address
of the variable that follows it. The \code{*} operator, called
\emph{dereference}, goes from the pointer to the object it points at.
They are inverses: \code{*(\&value)} is \code{value}.

Let us draw the state of memory at the moment just before the final
\code{printf}:

\begin{verbatim}
   Suppose the compiler placed:
     value at address 0x7ffc1004
     ptr   at address 0x7ffc1008

   We depict each memory cell as a box containing its value.
   A pointer is drawn as an arrow to the cell it targets.

              address       bytes               interpretation
   value      0x7ffc1004    [ 42 ]              an int holding 42
   ptr        0x7ffc1008    [ 0x7ffc1004 ]      an int* holding value's address

   Box-and-arrow view:

        +------------------+          +------------------+
        |       ptr        |          |      value       |
        |  0x7ffc1008      |--------->|      42          |
        |  (holds address) |          |  (an int)        |
        +------------------+          +------------------+

   And in full:

        +---------------------------------------------+
        |  ptr = &value                               |
        +---------------------------------------------+
        |  "ptr"                |  "value"            |
        |  addr 0x7ffc1008      |  addr 0x7ffc1004    |
        |  value 0x7ffc1004 --->|  value 42           |
        +---------------------------------------------+
\end{verbatim}

The arrow is the single most important diagram in this book. When you
read C code, your job is to draw the arrows. When you write C code,
your job is to know exactly which arrows exist and where they point.
A C program is only correct if its arrows and its ownership rules are
correct; everything else follows.

\section{Dereferencing Is Reading or Writing Through an Arrow}

Dereferencing is a two-way street. It can \emph{read} the cell the arrow
points to, or it can \emph{write} through it. A dereference on the left
side of an assignment writes to the pointed-to cell.

\begin{lstlisting}[caption={Reading and writing through a pointer.}]
#include <stdio.h>

int main(void) {
    int value = 42;
    int *ptr = &value;

    /* Read: fetch the int that ptr points at. */
    printf("before: %d\n", *ptr);

    /* Write: store a new int INTO the cell ptr points at. */
    *ptr = 100;

    printf("after:  value = %d, *ptr = %d\n", value, *ptr);

    /* value was modified, because ptr points at value. */
    return 0;
}
\end{lstlisting}

\begin{verbatim}
   Before:                      After *ptr = 100;

     ptr ---> [ 42 ]              ptr ---> [ 100 ]
                 value                       value
\end{verbatim}

There is nothing magic here. \code{*ptr = 100} means: ``go to the
address stored in \code{ptr}, and write the bytes of \code{100} there.''
The C compiler turns it into one store instruction. The pointer is just
the messenger; the object it points at is what actually changes.

\begin{insight}[Pointers are how C returns multiple values]
Functions return one value. If you want a function to deliver two
results, you pass pointers to the caller's variables and the function
writes through them. This is exactly how \code{scanf} works, and how
we will write our own data structure APIs in later chapters:
\code{int list\_pop(struct list *l, int *out\_value);}
\end{insight}

\section{Pointer Arithmetic}

A pointer plus an integer steps over whole \emph{elements}, not bytes.
If \code{int *p} holds an address, then \code{p + 1} holds the address
of the next \code{int} object, which is 4 bytes later on our machine.
The compiler knows the size of the pointed-to type, so it does the
multiplication for you.

This is the C idiom that builds every array algorithm, and it is the
place where ``clever'' code tends to hide. In this course we refuse
obfuscation: we move pointers on explicit lines, we keep the current
element in a named variable, and we only ever use \code{size\_t} as the
loop counter. Compare the same loop written two ways.

The style you will find in legacy code (and, alas, in many textbooks):

\begin{lstlisting}[caption={Legacy style -- compact, but opaque.}]
/* Copies the string src into the buffer dst. Obscured, hard to debug. */
void copy_string(char *dst, const char *src) {
    while (*dst++ = *src++);   /* we do NOT teach this */
}
\end{lstlisting}

That line does four operations inside a test. If a bug appears, you
cannot step through it one move at a time. The style we use:

\begin{lstlisting}[caption={Our explicit style: one move per line, named pointers.}]
#include <stddef.h>
#include <string.h>

/* Copy src into dst, including the terminating NUL byte. */
/* Precondition: dst has room for strlen(src) + 1 bytes.     */
void copy_string(char *dst, const char *src) {
    size_t i;
    size_t len;

    len = strlen(src);

    for (i = 0; i <= len; i++) {
        /* Copy one byte: the element at position i of src. */
        dst[i] = src[i];
    }
}
\end{lstlisting}

Notice what we gained: every step is a single, auditable statement.
We can put a breakpoint on any line, print \code{i} and \code{len},
and inspect the bytes one at a time. The clarity costs almost nothing
in speed; a modern compiler turns the explicit loop into the same
optimized code it would generate for the clever version. Clarity is
nearly free. We take it.

\section{Arrays Are Contiguous Bytes; Pointers Point at Elements}

An array of \code{N} elements of type \code{T} is simply \code{N}
objects of type \code{T} laid out consecutively in memory. There is no
header, no length, no metadata. The address of the array is the address
of its first element.

\begin{verbatim}
   int numbers[4] = {10, 20, 30, 40};

   A 4-byte int is drawn as one box. The array occupies 16 consecutive
   bytes:

      numbers +----+----+----+----+
              | 10 | 20 | 30 | 40 |
              +----+----+----+----+
      address: 1000 1004 1008 1012

      numbers            == &numbers[0]   == 1000
      numbers + 1        == &numbers[1]   == 1004
      numbers + 2        == &numbers[2]   == 1008

      *numbers           == 10
      *(numbers + 2)     == 30            == numbers[2]
\end{verbatim}

The expression \code{numbers[i]} is defined to mean
\code{*(numbers + i)}. The name of an array, used in an expression,
\emph{decays} into a pointer to its first element. This is why
passing an array to a function passes the address, not a copy: any
write inside the function modifies the caller's array.

\begin{lstlisting}[caption={Arrays decay to pointers when passed to functions.}]
#include <stdio.h>

/* Fill a buffer of n ints with the values 0..n-1. */
void fill_seq(int *buf, size_t n) {
    size_t i;

    for (i = 0; i < n; i++) {
        buf[i] = (int) i;
    }
}

int main(void) {
    int numbers[4] = {0, 0, 0, 0};
    size_t i;

    fill_seq(numbers, 4);           /* "numbers" is &numbers[0] here */

    for (i = 0; i < 4; i++) {
        printf("numbers[%zu] = %d\n", i, numbers[i]);
    }

    return 0;
}
\end{lstlisting}

Note the parameter type \code{int *buf}. In the declaration of a
function parameter, \code{int buf[]} and \code{int *buf} are identical.
We write the pointer form, because it tells the truth about what is
passed.

\section{Stack vs.\ Heap: Who Owns the Memory?}

This distinction is the pivot of all of C. Two places hold the data
your program manipulates, and they differ in \emph{lifetime} and
\emph{management}.

\subsection{The Stack}

Local variables live on the stack. When a function is called, the
machine pushes a \emph{stack frame} containing its locals and return
address; when the function returns, that frame is popped and its
storage is, effectively, gone. Stack memory is:

\begin{itemize}
  \item allocated and freed automatically (by the call/return protocol),
  \item extremely fast,
  \item limited in size (typically ~8 MB on Linux),
  \item invalid the moment the function returns.
\end{itemize}

\begin{lstlisting}[caption={A classic lifetime bug: returning a pointer to a stack local.}]
#include <stdio.h>

int *bad_function(void) {
    int local = 42;

    return &local;   /* ERROR: local dies when bad_function returns */
}

int main(void) {
    int *p = bad_function();

    /* p dangles: it names memory that the stack has already reclaimed. */
    printf("%d\n", *p);   /* undefined behavior */
    return 0;
}
\end{lstlisting}

\begin{verbatim}
   The stack during main() -> bad_function() -> return:

      before call:     during bad_function():    after return:
      +--------------+  +--------------+          +--------------+
      | main frame   |  | main frame   |          | main frame   |
      |  p = ?       |  |  p = ?       |   p ---> |  p = 0x...   |
      +--------------+  +--------------+          +--------------+
      (stack top)       | bad frame    |          (frame popped)
                        |  local = 42  |<-- &local is now garbage
                        +--------------+
\end{verbatim}

The address value survives; the storage behind it does not. We call
\code{p} a \emph{dangling pointer}. No compiler can reliably catch
every use of one; only discipline can. The rule:

\begin{ruleofthumb}[Heap for anything that outlives its creator]
If data must be returned from a function, or shared between functions,
or sized by values known only at runtime, it goes on the \emph{heap},
allocated with \code{malloc}, and freed with \code{free} by whichever
party becomes its owner.
\end{ruleofthumb}

\subsection{The Heap}

The heap is managed explicitly. \code{malloc(size)} asks the OS (via a
bigger allocator such as \code{glibc}) for \code{size} bytes and
returns a pointer to them, or \code{NULL} if it cannot. The bytes are
\emph{yours} until you call \code{free}. Nothing is reclaimed
automatically. This is the price of lifetime control, and it is a fair
price: with it we build linked lists, hash tables, trees, and whole
servers, all with deterministic behavior.

\begin{lstlisting}[caption={A correct heap round-trip.}]
#include <stdio.h>
#include <stdlib.h>
#include <string.h>

int main(void) {
    char *buffer;
    size_t len;

    len = 32;
    buffer = malloc(len);
    if (buffer == NULL) {
        fprintf(stderr, "out of memory\n");
        exit(1);
    }

    strcpy(buffer, "hello from the heap");

    printf("%s\n", buffer);

    free(buffer);          /* hand the bytes back */

    /* After free, buffer is a dangling pointer; do not dereference it. */
    return 0;
}
\end{lstlisting}

The \code{== NULL} check is not optional. \code{malloc} fails when the
system cannot honor the request, and a good program treats that
possibility as part of its normal control flow.

\section{Pointers to Pointers}

Sometimes the thing a pointer points at is, itself, a pointer. A
\code{char **argv} is a pointer to an array of strings; a function that
must \emph{reassign} a caller's pointer takes a pointer to that pointer.
This is how we implement functions that ``return'' newly allocated
objects through their arguments.

\begin{lstlisting}[caption={Writing a new pointer into a caller's variable.}]
#include <stdio.h>
#include <stdlib.h>

/* Allocates a small string and stores its address in *out. */
void make_message(char **out) {
    char *s = malloc(16);

    if (s == NULL) {
        *out = NULL;
        return;
    }
    s[0] = 'O'; s[1] = 'K'; s[2] = '\0';

    *out = s;          /* write the pointer into the caller's variable */
}

int main(void) {
    char *msg = NULL;

    make_message(&msg);      /* &msg is a char** */

    if (msg != NULL) {
        printf("%s\n", msg);
        free(msg);
    }

    return 0;
}
\end{lstlisting}

\begin{verbatim}
   Before make_message(&msg):        After *out = s:

     msg  [ NULL ]                     msg  [ 0x.. ] ---> [ O ][ K ][ 0 ]
     out  [ &msg ]                     out  [ &msg ]
\end{verbatim}

The pattern is always the same: the parameter is a pointer to the
variable we want the function to modify, and inside the function we
dereference it to perform the write.

\section{NULL: The Pointer That Points at Nothing}

\code{NULL} is the address 0, and by convention it means ``this pointer
is not currently valid''. It is the only address you may
\emph{compare} against but never \emph{dereference}. The processor does
not map address 0, so dereferencing a \code{NULL} pointer traps
immediately with a segmentation fault, which is a favor: the bug
surfaces exactly where it happens, instead of corrupting data silently
a million instructions later.

\begin{verbatim}
   ptr ---> [ 42 ]       valid pointer, points at a real int

   ptr = NULL;

   ptr ---> 0x0          NOT a real object.
   Dereferencing it:    *ptr  ==>  SIGSEGV (segmentation fault)
\end{verbatim}

Use \code{NULL} honestly: every allocation is checked against it,
every function that may fail returns it, and every pointer in a
structure that is ``not set'' holds it. \code{NULL} is the C equivalent
of ``no object here'', and it costs nothing to check but a comparison.

\section{Undefined Behavior: The Abyss}

The C standard reserves the right to do \emph{anything} when your
program crosses certain lines: dereferencing a dangling or \code{NULL}
pointer, reading an uninitialized variable, overflowing a signed
integer, indexing an array out of bounds. The compiler may assume you
never do these things, and optimize accordingly. When you violate the
contract, the program does not ``throw an exception''; it simply has no
guaranteed behavior. It might work today and break tomorrow, or under a
different compiler flag, or when the array happens to sit next to
something important.

This is why the discipline in this course is non-negotiable:

\begin{itemize}
  \item Every pointer is either valid, \code{NULL}, or untouched.
  \item Every array index is checked against a length you carry.
  \item Every \code{malloc} result is checked.
  \item Every byte a function writes is justified by a size argument.
\end{itemize}

\section{Command-Line Arguments: argc and argv}

Almost every real program is launched with arguments: \code{./tool
input.txt -v}. Those arguments arrive in \code{main} through two
parameters, and since every later chapter's practical programs use
them, we define them once here.

\begin{lstlisting}[caption={main receives its command line as argc and argv.}]
#include <stdio.h>

int main(int argc, char *argv[]) {
    int i;

    printf("program name: %s\n", argv[0]);
    printf("number of arguments: %d\n", argc);

    for (i = 1; i < argc; i++) {
        printf("  argv[%d] = %s\n", i, argv[i]);
    }

    return 0;
}
\end{lstlisting}

\begin{verbatim}
   Running:   ./prog alpha "two words"

   argc = 3           (the count of strings in argv)
   argv[0] = "./prog"       the program's own name
   argv[1] = "alpha"
   argv[2] = "two words"    (quotes group a string with a space)

   main's parameters:
     argc   an int: how many strings are in argv (always >= 1)
     argv   an array of char* strings; argv[argc] is always NULL

   - The SHELL (or your parent process) builds this array and passes
     it to the new program at exec time (ch.5). The kernel copies it
     onto the new program's stack.
   - argv[0] is not an argument the user typed; it is the name the
     program was invoked as. Arguments the user typed begin at argv[1].
   - That is why every practical in this book guards with
     "if (argc < 2) { fprintf(stderr, \"usage...\"); }": it is
     checking that the user actually supplied argv[1].
\end{verbatim}

The two spellings \code{char *argv[]} and \code{char **argv} are
identical here, by the array-decay rule of Section~1.8. We write
\code{char *argv[]} because it says what it is: an array of strings.
And the strings themselves---the \code{argv[i]} values---are ordinary
\code{char*} pointers into a region the kernel placed on the stack
before \code{main} ran.

\section{Practical Exercise: A Memory Inspector}
Let us consolidate the chapter into a small, real tool. It prints the
layout of a few variables and then hex-dumps a region of memory so you
can literally \emph{see} bytes, addresses, and pointer values. Run it;
look at the output until it stops being surprising.

\begin{lstlisting}[caption={memory\_inspect.c: see bytes, addresses, and pointers.}]
#include <stdio.h>
#include <stdlib.h>

/* Prints the raw bytes of the n bytes starting at p. */
void hex_dump(const unsigned char *p, size_t n) {
    size_t i;

    for (i = 0; i < n; i++) {
        /* Print byte number i, then a separator between groups of 8. */
        printf("%02x", p[i]);
        if ((i % 8) == 7) {
            printf("  ");
        } else {
            printf(" ");
        }
    }
    printf("\n");
}

int main(void) {
    int value = 0x11223344;   /* one int, four bytes we can recognize */
    int *ptr = &value;
    char text[] = "abc";      /* four bytes: 'a','b','c','\0'         */

    printf("=== addresses ===\n");
    printf("&value     = %p\n", (void *) &value);
    printf("&ptr       = %p\n", (void *) &ptr);
    printf("&text      = %p\n", (void *) &text);
    printf("ptr        = %p\n", (void *) ptr);

    printf("\n=== sizes ===\n");
    printf("sizeof(int)     = %zu\n", sizeof(int));
    printf("sizeof(int *)   = %zu\n", sizeof(int *));
    printf("sizeof(text)    = %zu\n", sizeof(text));

    printf("\n=== the int as raw bytes ===\n");
    hex_dump((unsigned char *) &value, sizeof(value));

    printf("\n=== the pointer as raw bytes ===\n");
    hex_dump((unsigned char *) &ptr, sizeof(ptr));

    printf("\n=== the string (with its NUL terminator) ===\n");
    hex_dump((unsigned char *) text, sizeof(text));

    printf("\n=== stepping over an int with pointer arithmetic ===\n");
    {
        /* p moves from value to each byte-sized window inside it. */
        unsigned char *q = (unsigned char *) &value;
        size_t i;

        for (i = 0; i < sizeof(value); i++) {
            printf("byte %zu at %p = 0x%02x\n", i, (void *) &q[i], q[i]);
        }
    }

    return 0;
}
\end{lstlisting}

If your machine is little-endian (x86-64 and ARM64 both are), the int
\code{0x11223344} is stored as the byte sequence \code{44 33 22 11}:
the least significant byte comes first in memory. The pointer's raw
bytes encode the address, and you will recognize them as the value
\code{ptr} printed in the addresses section.

\begin{insight}[Endianness]
Endianness is the order in which the bytes of a multi-byte number are
stored. Little-endian machines store the least significant byte first.
Most of the time you will not care, but the moment you write a file
format, a network protocol, or a virtual machine, you will care a great
deal. We will revisit endianness when we serialize data to disk and to
the network in later chapters.
\end{insight}

\begin{researchnote}[Memory allocators and custom allocators]
All through this chapter we told you to trust \code{malloc} and check
its result. Researchers have actually measured whether programs benefit
from \emph{custom allocators}---hand-written replacements for
\code{malloc} tuned for one program. ``Recent research'' here means a
careful empirical study by Emery Berger and colleagues, and its
conclusion is an excellent corrective to a very common habit: writing a
custom allocator is usually wasted effort.

\begin{verbatim}
   The study:  Berger, Zorn, and McKinley studied eight applications that use
   custom allocators and compared each one, in wall-clock time and in
   memory, against the Lea allocator (the general-purpose malloc of the
   day). Their finding, in one line:

       "Custom allocators are rarely effective: for most programs a
        general-purpose malloc performs as well as or better than the
        custom allocator; only region-based allocators (arenas)
        consistently show gains, of up to 44%."

   (That is a brief paraphrase of the paper's findings, condensed for the book.)
\end{verbatim}

What makes the study trustworthy, and worth quoting: it is an
\emph{empirical} result (measured, not argued), it was published in a
peer-reviewed venue, and it runs against the folklore that every
performance-sensitive program needs a hand-tuned allocator. The lesson
for this book is simple: use \code{malloc}, use \code{free}, check
failure, and profile \emph{before} you ever consider writing a custom
allocator. We will revisit this when we design our own allocators in
Chapter~13, where the goal will be \emph{policy} (object pools, arenas)
rather than raw speed.

\paperref{E.~D.~Berger, B.~G.~Zorn, K.~S.~McKinley, ``Reconsidering
Custom Memory Allocation,'' ACM OOPSLA 2002}{https://people.cs.umass.edu/~emery/pubs/berger-oopsla2002.pdf}
\end{researchnote}

\section{Structures in Bytes}

Every structure in this book---variable, struct, packet, inode---is,
at bottom, a fixed arrangement of bytes. This section starts the habit
of reading that arrangement; every later chapter ends with its own
byte map. On our 64-bit little-endian machine, the scalars look like
this:

\begin{verbatim}
   type         size   bytes (little-endian storage)
   ---------    ----   ------------------------------------------------
   char         1      01
   short        2      01 00
   int          4      01 00 00 00
   long         8      01 00 00 00 00 00 00 00
   float        4      00 00 80 3f            (the float value 1.0)
   double       8      00 00 00 00 00 00 f0 3f   (the double value 1.0)
   pointer      8      e6 3c a2 29 a0 00 00 00   (some address)
\end{verbatim}

Little-endian means byte 0 is the least significant byte, so the value
\code{1} is stored as \code{01 00 00 00} and \code{0x11223344} as
\code{44 33 22 11}. Reading a memory dump is reading these sequences in
groups of the type's size. A stack frame is bytes too: on x86-64 the
\code{call} instruction pushes an 8-byte return address, the callee
saves the old frame pointer, then makes room for its locals:

\begin{verbatim}
   one stack frame (the stack grows downward):

   +-------------------------------+
   | caller's frame                |
   +-------------------------------+
   | return address          (8B)  |  pushed by the call instruction
   +-------------------------------+
   | saved frame pointer     (8B)  |  pushed by "push rbp"
   +-------------------------------+  <- rbp
   | locals (e.g. int value, 4B)   |
   | and saved registers           |
   +-------------------------------+  <- rsp
\end{verbatim}

\code{\&value} is a byte address inside that range---which is why the
frame diagram of Chapter~1 and the byte diagram here are the same
picture.

\section{Summary of the Mental Model}

At the end of this chapter you should be able to answer, for any
snippet of C you read: \emph{which variable lives where, how many bytes
it occupies, what pointer points at what, and who owns the memory.}

\begin{verbatim}
   The four questions of every C program:

   1. Address : where do the bytes live?       (stack / heap / globals)
   2. Size    : how many bytes?                (sizeof gives the truth)
   3. Arrows  : which pointers point where?    (draw the diagram)
   4. Owner   : who calls free()?              (exactly one party)
\end{verbatim}

If you can draw the memory diagram for a program, you understand the
program. If you cannot, you do not---regardless of how many tests pass.
This habit, more than any feature of the language, is what separates
system programmers from everyone else.

In the next chapter we open the machine's fastest memory: the
\emph{registers}, built from flip-flops---bistable loops of
transistors that hold a bit---and descend from the register image to
the electronics and chemistry that make it retain its value. The
chapter after that combines scalars into \emph{structures}, and
learns why the machine, in its demand for aligned access, quietly
inserts invisible bytes between our fields.
}{}
\IfFileExists{chapter2.tex}{% ============================================================================
%  chapter2.tex
%  Registers: From the Register Image Down to the Atom
% ============================================================================

\chapter{Registers: The Bistable Heart of the Machine}

\section{A Snapshot of the Whole Machine}

A process is running. It has just finished one instruction and is
about to start the next when the operating system decides that
another process deserves the CPU. The machine must not lose the
process that was running---a few milliseconds from now it will come
back and finish exactly where it left off.

So the CPU does something extraordinary: it photographs itself. It
copies into memory everything that defines ``where this computation
stands''---which instruction runs next, where the stack points, the
value of every local the compiler is juggling---and later, when the
process returns, it restores the photograph and the computation
continues as if nothing had happened. Chapter~5 will call this
snapshot the \emph{register image}, and Chapter~11 will show a thread
switch saving and restoring it, byte by byte.

\begin{verbatim}
   The register image of an x86-64 process, mid-context-switch:

   RAX  0x000000000000003f     R9   0x00007ffc1004d0
   RCX  0x0000000000000001     R10  0x0000000000000000
   RDX  0x0000000000000000     R11  0x0000000000000246
   RBX  0x0000000000000000     R12  0x0000000000000000
   RSP  0x00007ffc1004f0       R13  0x0000000000000000
   RBP  0x00007ffc1004f0       R14  0x0000000000000000
   RSI  0x00007ffc1004d0       R15  0x0000000000000000
   RDI  0x0000000000000001
   RIP  0x0000000000401167     (the next instruction)
   RFLAGS 0x0000000000000202   (carry, zero, sign, ...)
\end{verbatim}

Sixteen general-purpose registers at eight bytes each, plus the
instruction pointer and the flags: 128 bytes of ordinary, unglamorous
data. And yet those 128 bytes are the whole machine, suspended in
mid-thought. Every computation that will ever run on this CPU spends
its working life inside them.

\begin{osconcept}
None of this saving and restoring happens on its own. The decision to
interrupt a running process, the choice of \emph{which} process runs
next, and the obligation to put the old process's 128 bytes away and
lift another's back out are all made by the \emph{scheduler}, a piece
of the kernel. A thread switch is, in the most literal sense, an
exchange of register images; Chapter~11 will write that exchange by
hand. Registers are hardware, but their contents are handed between
processes only because the OS supervises the transfer.
\end{osconcept}

That is the image this chapter is about. By its end you will know,
bit by bit, what it is made of---because underneath those 128 bytes
lies a ladder of smaller and smaller things. A register. A
flip-flop. Two logic gates. A pair of transistors. And finally, a few
billion silicon atoms, a handful of which have been swapped for
dopants. Let us climb down.

\section{The Register: One Word, Sixty-Four Flip-Flops}
\label{sec:register}

The image is sixteen words. A \emph{word} is the machine's natural
unit (Chapter~1), and each word in the image lives in a
\emph{register}---a small store on the CPU die itself, nanoseconds
away, managed entirely by your program.

A register is not mysterious. It is a row of storage cells, one per
bit: sixty-four cells for a 64-bit register. Each cell is a
\emph{flip-flop}, a circuit that can remember a single 0 or 1, and
all sixty-four flip-flops in a register share one \emph{clock} line
so that the whole word changes at the same instant. When the compiler
emits \code{mov rax, 42}, the CPU drives 42 onto the data lines and
pulses the clock; on that pulse every flip-flop captures its bit, and
\code{rax} now holds 42.

\begin{verbatim}
   A 64-bit register is 64 flip-flops and one clock:

       D0 ----->[D Q]----- Q0
       D1 ----->[D Q]----- Q1
       ...      [D Q]
       D63----->[D Q]----- Q63

                 +-------+
       clock ----> all 64: on the rising edge each captures its D.
                 +-------+
\end{verbatim}

The shared clock is what makes a register \emph{one object} rather
than sixty-four coincidences: the whole word is latched in a single
tick. That is why a register has the byte size of a word, why the
register image is laid out in 8-byte units, and why your program can
treat a register as an indivisible thing.

The sixteen registers are not interchangeable. The machine gives each
one a personality, and the compiler and the hardware agree on it:

\begin{verbatim}
   RAX   accumulator: arithmetic results, function return values
   RCX   counted loops; the first "loop counter" a machine ever had
   RDX   data/remainder; the high half of wide arithmetic
   RBX   a callee-saved scratch register
   RSP   the stack pointer: where the top of the stack is
   RBP   the base/frame pointer: where the current frame starts
   RSI   source index: string ops, second function argument
   RDI   destination index: string ops, first function argument
   R8..R15  a general-purpose pool added by the 64-bit extension
\end{verbatim}

Most of these roles are convention, not law---the CPU will happily do
arithmetic in \code{r15}---but the convention is so deeply baked into
compilers, debuggers, and system call interfaces that your code
should treat it as the layout of the machine. \code{rsp} and
\code{rbp} in particular are \emph{hard} convention: Chapter~6 will
show the stack machinery that depends on them.

\begin{insight}[Why sixteen?]
The number of registers is an accident of instruction encoding. Each
instruction that names two registers must fit those names into a few
bits; \code{R8}--\code{R15} came with a 64-bit extension byte added
to that encoding. A machine with 32 or 128 registers is perfectly
possible---and modern out-of-order CPUs secretly have far more---but
\emph{architecturally}, the number is whatever the instruction format
can address.
\end{insight}

Now the question this chapter descends toward: what is a flip-flop,
that it can hold a bit?

\section{Down One Level: The Flip-Flop}
\label{sec:flipflop}

A flip-flop's ancestor is the \emph{SR latch}, and the SR latch is
almost absurdly simple: two logic gates wired so that each gate's
output is one of the other gate's inputs. With no inputs asserted,
the circuit is just a loop. And that loop is a memory.

\begin{verbatim}
   The SR latch: two NOR gates, cross-coupled.

         S -------------------------+
                                     v
                        +-----------------------+
                        |          NOR          |
                        |                       |----------- Q
                        +-----------------------+
                              ^
                              |   (Qbar feeds back)
                        +-----------------------+
                        |          NOR          |
                        |                       |----------- Qbar
                        +-----------------------+
                              ^
         R --------------------+
\end{verbatim}

This is the drawing you will find in Wikipedia's flip-flop article
and in every digital-logic textbook; in vacuum-tube form it dates to
the 1918 Eccles--Jordan trigger relay. A NOR gate outputs 1 exactly
when both its inputs are 0, which gives the latch a tiny truth table:

\begin{verbatim}
   S   R   action
   --- --- -----------------------------------
   1   0   set     : Q = 1 (the latch remembers 1)
   0   1   reset   : Q = 0 (the latch remembers 0)
   0   0   hold    : Q stays wherever it was
   1   1   forbidden: both outputs collapse; never drive this
\end{verbatim}

Force S to 1 and Q rises; the lower NOR sees that high and its output
falls; and when you let go of S, the loop is on its own---and it
stays exactly where it is, because each gate holds the other in
place. That is the entire mechanism of computer memory in one
sentence: \emph{a loop that reinforces itself}.

A bare latch is transparent---a stray signal can flip it at any time.
A CPU wants its registers to change only at agreed instants, so the
flip-flops it uses are \emph{edge-triggered D flip-flops}: one data
input, one clock, and the rule that the output changes only on the
clock's rising edge, and only to whatever the data input held at that
exact moment. A D flip-flop is two latches in series---master and
slave, opened on opposite clock phases---so it samples, then holds,
with no window for accidental change. Sixty-four of them sharing one
clock is the register of Section~\ref{sec:register}.

\begin{verbatim}
   A D flip-flop as master and slave, one data path:

         D ---->[master latch]----[slave latch]----> Q
                       |                  |
                       +--------+---------+
                                |
                                v
                          clock (inverted for one of them)
   On the rising edge the master opens and samples D; the slave is
   locked. On the falling edge the master locks and the slave opens,
   copying the sampled value to Q. The output therefore changes only
   around the clock edge, never in between.
\end{verbatim}

\section{Why the Loop Holds: Bistability and the Math of Feedback}
\label{sec:bistable}

The latch holds because it is \emph{bistable}: it has exactly two
stable states, and every other state is shoved toward one of them.
The mathematics is clean, and worth holding in your hands.

Model each NOR gate as a sharp inverting element. With $Q$ and
$\bar{Q}$ the two stored nodes, the pair of gates imposes

\begin{equation}
Q = \overline{S \lor \bar{Q}}, \qquad \bar{Q} = \overline{R \lor Q},
\end{equation}

where $\lor$ is logical OR and the overbar is negation. In the hold
condition ($S = R = 0$) this collapses to

\begin{equation}
Q = \overline{\bar{Q}}, \qquad \bar{Q} = \overline{Q},
\end{equation}

which has exactly two solutions: $Q=1,\ \bar{Q}=0$ and $Q=0,\
\bar{Q}=1$. A third, formal solution---$Q$ and $\bar{Q}$ both
halfway---exists on paper but not in practice: it is an
\emph{unstable} equilibrium. The tiniest noise tips the loop, and
positive feedback then races it all the way to one of the stable
corners. That razor's edge is called \emph{metastability}, and it is
why real hardware occasionally hiccups when two signals race.

The analogue view says the same thing with a picture. An inverter's
output $v_{out}$ is a function $f$ of its input; two inverters in a
loop must satisfy $v_{out} = f\bigl(f(v_{out})\bigr)$. Because $f$
is inverting and steeper than $-1$ at its midpoint, the composition
crosses the diagonal three times: the two stable points and the
unstable middle.

\begin{verbatim}
   v_out
   1 |  x                     stable: Q = 1
     |    .
     |     .     the middle diagonal crossing is the
     |       .        .       metastable point: unstable
     |        .    .
     |           .  x         stable: Q = 0
   0 +----------------------------- v_in
\end{verbatim}

That is the mathematical reason a register does not ``half-remember'':
the loop has no resting place between 0 and 1, so it cannot rest
there. A flip-flop is \emph{static} memory---it holds its two stable
voltages as long as power flows, needing no help. There is a second,
dynamic family, the DRAM capacitor, whose retention is a race against
leakage and whose math is even more explicit. Both families are built
from the same two ingredients, which is the next level down.

\section{The Clock: Making a Million Flip-Flops Agree}
\label{sec:clock}

Every register in the CPU is edge-triggered, and edge-triggering only
means something if every edge arrives at the same instant across the
whole die. That agreement is the job of the \emph{clock}: a single
square wave, distributed to every flip-flop, that ticks once per
cycle. At roughly 3--5 billion ticks a second in a modern CPU, each
tick lasts about a third of a nanosecond.

\begin{verbatim}
   The clock, drawn against time (t is running to the right):

   clock   ____|----|____|----|____|----|____|----|____
                  ^         ^         ^         ^
              rising edges: on each one, every flip-flop in the
              CPU samples its D input and changes exactly then.
\end{verbatim}

A synchronous machine is a heartbeat: flip-flops remember, and the
combinational logic \emph{between} registers---gates wired in a
straight line, no loops---does the computing during the quiet part of
each cycle. One tick to sample, one interval to compute, one tick to
latch the answer, and the whole machine advances in lockstep. This is
why your program sees registers change ``all at once'' and never
``in mid-thought'': the clock simply does not allow half-updated
words.

\subsection{Setup and Hold: The Flip-Flop's Polite Window}

A flip-flop samples on the edge, but it needs a moment of peace around
that edge. The data must be settled and steady for a short time
\emph{before} the edge---the \emph{setup time}, $t_{su}$---and must
stay put for a short time \emph{after} it---the \emph{hold time},
$t_{h}$. These are physical facts about the transistors inside the
flip-flop, on the order of tens of picoseconds.

\begin{verbatim}
   The window a flip-flop demands around a rising edge:

   clock    ____|----|________|----|____
                        ^
   D input  ____|_______________|________
             <-- stable high --><-- stable low
                        |<- t_su ->|<- t_h ->|
                        D must not change inside this window.

   Violate it, and the flip-flop does not guarantee a clean 0 or 1:
   it may sit halfway, decide randomly, or take a while to decide.
\end{verbatim}

Setup and hold are the price of agreeing on \emph{when}. The clock
must arrive at every flip-flop at nearly the same time (a skew of a
few dozen picoseconds is acceptable), the combinational logic must
finish within the cycle, and no signal may sneak through a latch
while its data is still being written. Engineers design for these
rules with static timing analysis; your program never sees them,
because the CPU is built to meet its own deadlines by construction.

\subsection{Metastability in the Wild: Crossing Asynchronous Borders}

Metastability is not an academic footnote. It appears the moment
something \emph{asynchronous} touches the clocked machine: a keyboard
key, a network packet, a disk interrupt. Those signals do not march to
the CPU's clock, so when a flip-flop samples them, the sample can land
right inside the forbidden window---and the flip-flop goes metastable,
hesitating between 0 and 1 for a while before tumbling to one side.

The trick is to not care which side it falls on. A
\emph{synchronizer} is two flip-flops in a row fed by the same clock:
the first may go metastable, but it has a whole clock cycle to settle,
and the second flip-flop samples it \emph{after} that, by which time
it is a solid 0 or 1. We lose a cycle of latency; we gain a guarantee
that the value the logic sees is a clean bit.

\begin{verbatim}
   Two flip-flops turning an async input into a safe signal:

        async in --->[D Q]---->[D Q]----> into the clocked logic
                      flop 1    flop 2
                         |         |
                         +----+----+
                              |
                              v
                          the one clock (both flops)

   flop 1 has one full clock to resolve; flop 2 then samples a
   settled value. The value is "arbitrary," but it is never
   mid-flight, and that is what the rest of the machine needs.
\end{verbatim}

That is why hardware ``occasionally hiccups'' in the face of racing
signals: no synchronizer is perfect, and with enough samples the
first flop occasionally takes longer than a cycle to decide. The
failure rate is not zero---it is merely small---and how small is a
measured engineering quantity. It is also exactly why every external
interrupt in the kernel arrives through hardware that has already
been synchronized before it ever reaches a register.

\begin{researchnote}[Metastability is measured, not guessed]
The frightening part of metastability is that a synchronizer can
theoretically sit unresolved forever. In practice it cannot, but the
failure rate is governed by a stark formula. For a two-flip-flop
synchronizer given a resolution time $t_r$, the mean time between
failures is roughly

\begin{equation}
\text{MTBF} \approx \frac{e^{t_r / \tau}}{T_0\, f_\text{clk}\, f_\text{data}},
\end{equation}

where $f_\text{clk}$ and $f_\text{data}$ are the clock and data
frequencies and $T_0$ and $\tau$ are constants of the flip-flop
technology. The exponential in the numerator is the whole story:
each extra clock cycle of resolution time buys an astronomically
larger MTBF. Thomas Chaney and Charles Molnar proved the phenomenon
was fundamental, not a manufacturing flaw, in 1973, and the formula
appears in nearly every serious digital-design text since.

\paperref{T.~J.~Chaney, C.~E.~Molnar, ``Anomalous Behavior of
Synchronizer and Arbiter Circuits,'' IEEE Transactions on Computers,
1973}{https://ieeexplore.ieee.org/document/1672263}
\end{researchnote}

\begin{ruleofthumb}[Synchronize before you trust]
In software terms, treat every value that arrives from outside the
clocked core---a pressed key, a network byte, a changed file---as
asynchronous and therefore suspect. The hardware synchronizer turns a
``maybe metastable'' bit into a ``settled, but arbitrary'' bit; your
program's job is to re-validate that arbitrary value (parse it, check
it) before trusting it. The register image itself never needs this
treatment: everything in it lives on the clock.
\end{ruleofthumb}

\section{Down Another Level: Transistors and Capacitors}
\label{sec:transistors}

The gates and inverters above are not symbols; each is built from
\emph{field-effect transistors}, which are voltage-controlled
switches. A voltage on the gate opens or closes a channel for
current. A CMOS inverter is two transistors: an n-type one that
conducts when its input is high and pulls the output low, and a
p-type one that conducts when the input is low and pulls the output
high.

\begin{verbatim}
   A CMOS inverter: one p-type and one n-type transistor.

           Vdd
            |
           [P]    p-type: on when input is 0
            |----- output
           [N]    n-type: on when input is 1
            |
          GND

   input = 1 -> N on, P off -> output 0
   input = 0 -> P on, N off -> output 1
\end{verbatim}

A static RAM (SRAM) cell---the flip-flop of Section~\ref{sec:flipflop}
drawn in silicon---is six transistors: two cross-coupled inverters
(four transistors) that hold the bit, plus two \emph{access}
transistors that connect the cell to the bit lines when the row's
word line is high. Read or write, the cell is a tiny latch, exactly
the loop of the previous section.

\begin{verbatim}
   The 6-transistor SRAM cell: the two inverters, plus access.

            Vdd                     Vdd
             |                       |
            [P1]                    [P2]     <- pull-up transistors
             |                       |
             +--- Q ----+---- Qbar --+
             |          |            |
            [N1]       [A1]        [A2]     <- access transistors
             |          |            |
           GND     bit line     bit line bar
                         \        /
                          +------+--------- word line (their gates)

   word line low  : A1, A2 open   -> the cell floats, holding Q
   word line high : A1, A2 closed -> bit lines read or write Q
   The two inverters actively pump Q and Qbar to the rails, which is
   why SRAM never needs refreshing.
\end{verbatim}

A dynamic RAM (DRAM) cell is leaner: one access transistor and one
\emph{capacitor}. The bit is simply whether the capacitor holds
charge. That frugality is why DRAM is cheap and dense---and its cost
is that the charge leaks away and must be \emph{refreshed} on a
timer.

\begin{verbatim}
   A DRAM cell: one transistor, one capacitor.

        word line
            |
            v
           [T]     the access transistor
            |
   bit line ---+====   the storage capacitor (the bit)
               |
              gnd
\end{verbatim}

Here is the retention math, concrete and vivid. A capacitor stores
charge $Q = C\,V$. A leakage current $I$ drains it at

\begin{equation}
\frac{dV}{dt} = -\frac{I}{C},
\end{equation}

and if the leak behaves like a resistance $R$, the voltage decays as

\begin{equation}
V(t) = V_0\,e^{-t/RC}.
\end{equation}

Take a modern cell: roughly $C \approx 25$ fF and a standby leakage
of about a picoampere. Then $dV/dt \approx 40$ V/ms, so a stored
voltage of about a volt falls below the sensing threshold within ten
to twenty milliseconds. That is why the JEDEC standard insists every
DRAM row be refreshed at least every 64 ms---the memory must be
re-read and topped up before its bits fade into noise. The register's
flip-flops need no such clock: their inverters actively pump the
stored node back to the rail, which is why the register image holds
its values while the power is on.

\begin{verbatim}
   Two families, one difference:

                         holds the bit        refreshed?
   register flip-flop    active inverter      never (static)
   SRAM cell             active inverter      never (static)
   DRAM cell             a passive capacitor  every ~64 ms

   Static families actively pump their bit; dynamic families store
   charge and fight leakage. Everything else is the same loop.
\end{verbatim}

\begin{insight}[The cache hierarchy is this section in disguise]
Chapter~1's memory hierarchy is a ladder of exactly these cells.
L1/L2/L3 cache is SRAM---cross-coupled inverters, no refresh, fast
but expensive. Main memory is DRAM---one transistor and a capacitor,
cheap and dense, refreshed continuously by the memory controller.
And the register file is the fastest SRAM of all, right at the top
of the ladder. Every ``level'' you met in Chapter~1 is one of the two
circuits drawn here, scaled up.
\end{insight}

\section{The Bottom: Chemistry, Silicon, and a Handful of Electrons}
\label{sec:chemistry}

What are those transistors and capacitors made of? Silicon, arranged
in a crystal, with a few of its atoms deliberately replaced. This is
the chemistry underneath every bit your program has ever stored.

Pure silicon is a semiconductor. Each atom shares four valence
electrons in a perfect lattice, so at room temperature almost nothing
conducts. \emph{Doping} changes that by substituting a handful of
atoms:

\begin{itemize}
  \item \emph{Phosphorus} brings five valence electrons and donates
    a free one: the material becomes \emph{n-type}, full of mobile
    electrons.
  \item \emph{Boron} brings three valence electrons and leaves a
    missing bond---a \emph{hole} that moves like a positive
    charge---and the material becomes \emph{p-type}.
\end{itemize}

A transistor is a p-type body with n-type source and drain, a few
nanometres of silicon dioxide on top, and a metal gate. Raise the
gate voltage and its electric field pulls electrons to the surface,
forming an \emph{inversion channel} that connects source to drain.
No field, no channel, no current: that is the voltage-controlled
switch of the previous section, made of silicon and a sprinkling of
boron and phosphorus atoms. At doping levels near $10^{17}$ atoms per
cubic centimetre, roughly one silicon atom in a million is a dopant;
the ``free'' electrons that carry a signal are donations from those
rare, chosen atoms.

\begin{verbatim}
   The chemistry of a switch:

        gate
          |
         SiO2    (a nanometre of glass)
     n+  | p-type | n+
     S   | channel| D
         |  (electrons pulled in by the gate's field)

     electrons are the carriers; holes are their absence.
\end{verbatim}

Leakage, the enemy of the DRAM capacitor, is a chemistry story too.
An ``off'' transistor still lets a few carriers slip through the
barrier, and subthreshold leakage grows exponentially as devices
shrink. A modern DRAM cell stores its bit as only tens of thousands
of electrons; every time the cell shrinks it holds fewer, leaks
faster, and must be refreshed more aggressively. The famous scaling
rule---R.~H.~Dennard's 1974 prediction that smaller transistors would
stay fast and cool---broke precisely because of this chemistry of
tunnelling and leakage. That is why the industry moved to new
dielectrics, finFETs, and three-dimensional transistors: the deepest
detail on this ladder is where Moore's law meets its limit.

\begin{verbatim}
   The full descent, one ladder:

   register  -> 64 flip-flops sharing a clock
   flip-flop -> two cross-coupled inverting gates, edge-triggered
   gate      -> a pair of transistors (CMOS inverter, NOR, ...)
   transistor-> a sliver of doped silicon with a field-effect switch
   silicon   -> a crystal of atoms, a handful doped with P or B
   bit       -> the presence or absence of electrons at one node
\end{verbatim}

\section{Back Up: The Image, Understood}
\label{sec:backup}

Let us return to the opening scene. The process is interrupted; the
CPU photographs itself; 128 bytes of register image go into memory.
Now you know what those bytes are:

\begin{verbatim}
   one bit      = one flip-flop = two cross-coupled gates
                = four to six transistors = a loop of doped silicon
   one register = 64 flip-flops sharing a clock   (8 bytes)
   the image    = 16 registers + RIP + RFLAGS     (128 bytes)
\end{verbatim}

The register image is not a metaphor. It is 128 bytes of bistable
loops---each loop a pair of gates, each gate a pair of transistors,
each transistor a sliver of doped silicon---sitting on the CPU die
and holding the exact state of a computation because each loop
actively refuses to change. That refusal, powered by positive
feedback and paid for elsewhere with leakage and refresh, is the
deepest reason the data structures, processes, and databases of this
book work at all.

\begin{insight}[Every byte you store is this story]
The 128-byte register image is the fastest memory there is, and its
cousins---L1 cache, SRAM, DRAM, disk---are the same story at larger
scale. Cache is SRAM: more cross-coupled inverters, still static,
still loop-held. RAM is DRAM: capacitors racing leakage, refreshed
64 times a second. The hierarchy of Chapter~1 is, from the bottom up,
a ladder of the very same physics you just descended.
\end{insight}

\section{Registers at Work: The Datapath and the ALU}
\label{sec:datapath}

The descent is complete, so let us watch the image \emph{work}. A
register's job is not to sit still; it is to feed a computation. The
machine that does the feeding is the \emph{datapath}: the register
file (all sixteen registers in one bank), the arithmetic/logic unit
(the \emph{ALU}), and the wires that connect them. Every instruction
you will ever compile is a script for this small, closed system.

\begin{verbatim}
   The instruction loop of one add rax, rbx, step by step:

        +--------+      +--------+      +--------+
        |  RIP   |      |  RAX   |      |  RBX   |   register file
        +--------+      +--------+      +--------+
           |    ^           ^               ^
           |    |           |               |
           |    |       read port A     read port B
           |    |
           |  +-----------------------------+
           |  |          control logic      |  "add"
           |  +-----------------------------+
           |            |          |
           |            v          v
           |        +---------------+
           |        |     ALU       |  computes RAX + RBX
           |        +---------------+
           |            |
           |      +-----+------------+-----------------+
           |      |                  |                 |
           |      v                  v                 v
           |  write port -> RAX   RFLAGS           result discarded
           |                          (CF, ZF, SF, OF set)
           +--- RIP advances past the instruction; repeat forever
\end{verbatim}

1. \textbf{Fetch.} The instruction pointer \code{rip} holds the
address of the next instruction's bytes in memory. The CPU reads
those bytes and advances \code{rip} past them.
2. \textbf{Decode.} The control logic works out what the bytes mean:
``add the value in \code{rbx} to the value in \code{rax}.''
3. \textbf{Read.} The register file places \code{rax} on read port
A and \code{rbx} on read port B. Two reads in the same tick is what
the two ports are \emph{for}.
4. \textbf{Execute.} The ALU computes the sum and, as a side effect,
latches the result's meaning into \code{rflags}: whether it was zero
(\code{ZF}), negative (\code{SF}), carried out (\code{CF}), or
overflowed (\code{OF}).
5. \textbf{Write.} On the clock edge, the result is written back
into \code{rax} through the write port, and the whole loop repeats.

Three registers you already know appear here in their true roles.
\code{rip} is the machine's cursor into the program; \code{rflags}
is the machine's memory of what the last computation \emph{meant};
and the register file is the working surface the ALU reads from and
writes to. An add is a fetch, a decode, a read, an add, a write---and
a handful of clock cycles, in a machine that does billions of them a
second.

\begin{insight}[Architectural vs.\ physical registers]
The sixteen registers your program sees are the \emph{architectural}
register file: the ones Chapter~5 saves and Chapter~11 restores. A
modern out-of-order CPU keeps a much larger \emph{physical} register
file underneath---hundreds of cells---and renames the architectural
names onto physical cells so that independent instructions can use
the same \code{rax} in the same tick without colliding. Your program
only ever sees the sixteen names; the renaming is invisible magic
that makes the datapath faster. The image you save is still 128
bytes, because the contract with your program is the architectural
file.
\end{insight}

\section{RFLAGS, Bit by Bit: How a Comparison Decides a Branch}
\label{sec:rflags}

The flags are easy to ignore and impossible to avoid. Every
arithmetic instruction sets them; every conditional branch reads
them; and compilers generate code that lives and dies by them. In the
register image they are a single 64-bit word, but only a handful of
the bits carry meaning:

\begin{verbatim}
   The meaningful bits of RFLAGS (x86-64):

   bit   name  meaning
   ---   ----  ------------------------------------------------
   0     CF    carry: unsigned overflow from the last operation
   2     PF    parity: the low byte of the result had even parity
   4     AF    auxiliary carry: bit 3 carried (BCD arithmetic)
   6     ZF    zero: the result was exactly zero
   7     SF    sign: the result is negative (its top bit is set)
   8     TF    trap: single-step, used by debuggers
   9     IF    interrupt enable: maskable interrupts are allowed
   10    DF    direction: string instructions move up or down
   11    OF    overflow: signed overflow from the last operation
\end{verbatim}

The six arithmetic flags---\code{CF, PF, AF, ZF, SF, OF}---are the
ALU reporting on its own work. The trick that makes them powerful is
that a comparison is just a subtraction that throws away its result.
The instruction \code{cmp a, b} computes $a - b$ and sets the flags
exactly as a subtraction would, but stores nothing. The branch that
follows then reads the flags as an answer.

Consider the simplest possible overflow, on paper:

\begin{verbatim}
   add rax, rbx      rax = 0x7fff...ffff   rbx = 1

   result:           0x8000...0000
     unsigned view:  2^63 - 1 + 1 = 2^63, which fits  -> CF = 0
     signed view:    (2^63 - 1) + 1 = 2^63, which as a signed
                     64-bit number is negative -> overflow -> OF = 1
     zero?  no -> ZF = 0       sign? top bit set -> SF = 1
\end{verbatim}

One operation, two honest answers: \code{CF} reports the truth to
unsigned arithmetic, \code{OF} to signed. That is why there are two
families of conditional jumps. Unsigned comparisons (\code{jb}, \code{ja},
``below/above'') test \code{CF} and \code{ZF}; signed comparisons
(\code{jl}, \code{jg}, ``less/greater'') test \code{SF} against
\code{OF}. The same bytes, the same flags, two different readings.

Let a compiler show it. This C function

\begin{lstlisting}[caption={A branch the compiler will turn into flags.}]
int max(int a, int b) {
    if (a > b) {
        return a;
    } else {
        return b;
    }
}
\end{lstlisting}

compiles (unoptimized, so the branch is visible) into something very
close to this:

\begin{verbatim}
   max:
       pushq  %rbp
       movq   %rsp, %rbp
       movl   %edi, -8(%rbp)     a lives on the stack
       movl   %esi, -12(%rbp)    b lives on the stack
       movl   -8(%rbp), %eax     eax = a
       cmpl   -12(%rbp), %eax    compute a - b, set flags, keep nothing
       jle    .L2                if a <= b (signed) go load b
       movl   -8(%rbp), %eax     the a > b path: return a
       jmp    .L3
   .L2:
       movl   -12(%rbp), %eax    the else path: return b
   .L3:
       popq   %rbp
       ret
\end{verbatim}

(Exact output varies with compiler and version, but the shape is
universal.) The compare sets \code{CF}, \code{ZF}, \code{SF}, and
\code{OF} from $a - b$; \code{jle} then jumps if the subtraction said
``less than or equal'' in signed terms---that is, if $\text{ZF}=1$ or
$\text{SF} \ne \text{OF}$. A modern optimizer notices the branch is
short and replaces the whole dance with \code{cmov}, a conditional
\emph{move} that reads the same flags without branching. The flags
are the machine's internal vocabulary for ``how did that
comparison go?'', spoken in every branch your program ever takes.

\begin{insight}[Flags are a side effect]
Arithmetic sets the flags whether you want them or not. A stray
\code{add} in unrelated code silently rewrites \code{ZF}, which is
why the hardware convention is that branches must immediately follow
the comparison they depend on---and why compilers reorder your code
with such care. In assembly, a comparison is not data until the
branch consumes it. This pairing---set flags, consume flags---is the
lowest-level control flow your code will ever see.
\end{insight}

\begin{ruleofthumb}[Unsigned or signed: pick by the type]
When you compare two \code{size\_t} or \code{unsigned} values, the
compiler emits unsigned branches testing \code{CF} and \code{ZF};
when you compare \code{int} or \code{long}, it emits signed branches
testing \code{SF} against \code{OF}. Mixing the two is the classic
one-off bug (``\code{size\_t} is huge when it underflows'') that
Chapter~1's Loop Counter Rule exists to prevent. The flags remember
both readings; the instruction chooses which to consult.
\end{ruleofthumb}

\section{Practical Exercise: See Your Program's Register Image}
\label{sec:practical}

The opening scene of this chapter was a register image, but you have
probably never seen your \emph{own} program's registers. Let us fix
that. Two tools show you the fastest memory there is: a debugger to
dump it live, and the compiler to show how your C becomes register
moves.

\subsection{Dumping the image with gdb}

Write a tiny program, compile it with debug symbols, and stop it
inside a function:

\begin{lstlisting}[caption={max.c: the function from the previous section.}]
#include <stdio.h>

int max(int a, int b) {
    if (a > b) {
        return a;
    } else {
        return b;
    }
}

int main(void) {
    printf("max(7, 11) = %d\n", max(7, 11));
    return 0;
}
\end{lstlisting}

Compile with debug info and run it under the debugger:

\begin{verbatim}
   $ gcc -O0 -g -o max max.c
   $ gdb ./max
   (gdb) break max
   (gdb) run
   (gdb) info registers
\end{verbatim}

Because we break at the \emph{entry} of \code{max}, before any of its
own instructions run, the general-purpose registers still show the
calling convention in action: the first two integer arguments arrived
in \code{rdi} and \code{rsi}. You are looking at the register image
of a live computation---the 128 bytes of Chapter~2, mid-flight.

\begin{verbatim}
   rax            0x0                 0
   rbx            0x0                 0
   rdi            0x7                 7        <- first argument
   rsi            0xb                 11       <- second argument
   rsp            0x7fffffffe3d8      0x7fffffffe3d8
   rbp            0x7fffffffe3d0      0x7fffffffe3d0
   rip            0x401136            0x401136 <max>
   eflags         0x202               [ IF ]
\end{verbatim}

Try \code{next} a few times and watch \code{eflags} change as the
\code{cmp} and \code{jle} execute; then \code{continue} and watch the
return value land in \code{rax}. This is not a simulation. Every value
you see is held by cross-coupled loops of doped silicon, and a
\code{context switch} would freeze exactly these bytes to disk-shaped
memory and restore them later.

\subsection{Following your data through the compiler}

The debugger shows you the live image; the compiler shows you how it
was produced. Ask \code{gcc} for assembly instead of a binary:

\begin{verbatim}
   $ gcc -O2 -S max.c
   $ cat max.s
\end{verbatim}

At \code{-O2} the optimizer has made the branch disappear---the
\code{cmov} mentioned in the previous section---and the whole
function is a handful of instructions. The calling convention is
readable in the input names: \code{edi} and \code{esi} (the low 32
bits of \code{rdi} and \code{rsi}) enter, and the answer leaves in
\code{eax}. If you know where the arguments and the return value
live, you can read almost any function's skeleton without a manual.

\begin{verbatim}
   The calling convention (x86-64 System V), in one line:

   args 1-6 : rdi, rsi, rdx, rcx, r8, r9    return: rax
   callee-saved: rbx, rbp, r12-r15          caller-saved: the rest
\end{verbatim}

This convention is part of why the register image has the shape it
does. The callee-saved registers are the ones a function must restore
before returning---so a context switch only needs to save \emph{its}
working set, while the hardware and the compiler split the rest by
this rule. When Chapter~5 serializes a register image and Chapter~11
saves it across a thread switch, this table is the key to what gets
saved, and when.

\subsection{What to do when registers are not enough}

A running program needs more than sixteen words. Function arguments
beyond the sixth, large locals, and array indices all spill into
memory---onto the stack and heap of Chapter~1. The register file is
the fastest memory a program has, but it is also the smallest, and
the compiler's endless job is to decide which words deserve a
register and which must be spilled. Run \code{gcc -O0 -S} on the same
\code{max.c} and you will see the spilled version from the previous
section: everything parked on the stack, registers barely used. Run
\code{-O2} and the spills vanish. That trade---registers for speed,
memory for capacity---is the oldest trade in computing, and it is
exactly the hierarchy of Chapter~1 playing out inside a single
function.

\begin{ruleofthumb}[The register keyword is a fossil]
C has a keyword, \code{register}, that asks the compiler to keep a
variable in a register. Modern compilers ignore it almost entirely;
they have better information than you do about which words are hot,
and they reallocate freely. The reason it exists is that C was
written for machines where the answer was not obvious and the
compilers were not good. Let your compiler do its job, and reserve
your attention for the part it cannot decide: what \emph{lifetime}
your data needs, which is the stack-versus-heap question of
Chapter~1.
\end{ruleofthumb}

\section{Structures in Bytes}
\label{sec:structures}

The register file on x86-64, given byte addresses (each register is a
word, 8 bytes):

\begin{verbatim}
   general-purpose registers (16 x 8 bytes = 128 bytes):

   RAX   0     RCX   8     RDX  16     RBX  24
   RSP  32     RBP  40     RSI  48     RDI  56
   R8   64     R9   72     R10  80     R11  88
   R12  96     R13  104    R14  112    R15  120

   special registers:
   RIP      the instruction pointer: which program byte runs next
   RFLAGS   status bits: carry, zero, sign, overflow, direction...

   one D flip-flop: 1 bit
   one register:    64 flip-flops = 8 bytes
   the image:       128 bytes of the fastest storage that exists
\end{verbatim}

RSP is simultaneously the stack pointer and one of the sixteen
general registers; RIP and RFLAGS sit just outside the 128-byte file
but are saved with it on every context switch. This byte map is
exactly what Chapter~11's thread-scheduling code copies out and back
in, and it is the ``$\sim$100s B, $\sim$1 ns'' row of Chapter~1's
hierarchy table, made visible.

Each register is itself a small structure: one storage unit with
several overlapping views. The \code{rax} word can be addressed as a
full 64 bits, its low 32, its low 16, or its individual bytes:

\begin{verbatim}
   RAX, seen as overlapping windows (all of them are the same bytes):

   byte 0   1   2   3   4   5   6   7
   +---+---+---+---+---+---+---+---+
   | AL| AH|   AX (16 bits)         |
   |   EAX (32 bits)                |
   |   RAX (64 bits)                |
   +---+---+---+---+---+---+---+---+

   AL   = low 8 bits     AH  = bits 8..15
   AX   = low 16 bits    EAX = low 32 bits    RAX = all 64

   The register file's second word layer: RAX at bytes 0..7, RCX at
   8..15, and so on down the file. Writing EAX zeroes the top 32
   bits of RAX; writing AL leaves them alone.
\end{verbatim}

This partial-register trick is why the register image in the opening
scene is a clean 128 bytes while the machine inside happily reads
half-words and bytes. The file is 128 bytes of structure---sixteen
8-byte words---and every word is itself a byte map.

The flags word is a structure too, and a bit-level one: each status
flag is a single bit at a fixed position, which is why the kernel
reads and writes \code{RFLAGS} with bit tests and masks rather than
as a number. Section~\ref{sec:rflags} listed the positions; here they
are as bytes and bits:

\begin{verbatim}
   RFLAGS as a byte map (meaningful bits only):

   +-+-+-+-+-+-+-+-+-+-+-+-+-+-+-+-+-+-+-+-+-+-+-+-+-+-+-+-+
   |...|OF|DF|IF|TF|SF|ZF| AF |PF| |CF|
   +-+-+-+-+-+-+-+-+-+-+-+-+-+-+-+-+-+-+-+-+-+-+-+-+-+-+-+-+
      11 10  9  8  7  6   4  2    0

   the other bits are reserved and stay 0, which is why a clean
   flags word reads 0x202 (IF, plus the always-on reserved bits).
\end{verbatim}

Two more banks live beside the 128-byte file and are saved with it on
context switches: the \emph{vector registers} \code{XMM0}--\code{XMM15}
(sixteen 128-bit words, used by floating point and SIMD) and the
\emph{segment registers} \code{CS}, \code{DS}, \code{ES}, \code{FS},
\code{GS}, \code{SS}. Most programs touch only the general file, but
any code that does heavy float or string work will find its state
spread across all three banks. When you read a register dump, these
are the three blocks you are looking at.

\section{Summary}

\begin{verbatim}
   The register model, complete:

   register     a word of storage on the CPU die
   flip-flop    one bit: two cross-coupled inverting gates
   SR latch     set / reset / hold; the ancestor of all of it (1918)
   D flip-flop  captures D on a clock edge; 64 of them = a register
   bistability  two stable states; metastable only on the razor edge
   the clock    the shared edge that makes the whole machine agree
   setup/hold   a flop needs quiet around the edge; break it -> meta
   synchronizer two flops that absorb an async signal's uncertainty
   SRAM         6 transistors, static, no refresh
   DRAM         1 transistor + 1 capacitor, leaking, refreshed
   retention    V(t) = V0 e^(-t/RC): why DRAM refreshes every 64 ms
   datapath     RIP -> decode -> register file -> ALU -> RFLAGS
   RFLAGS       CF PF AF ZF SF OF: the ALU reporting on its work
   cmp/jcc      subtract, keep only the flags, branch on the flags
   chemistry    silicon + phosphorus/boron = the switches

   Rules that keep the model straight:
     1. a register holds its value only while power flows;
     2. writes happen on clock edges, all bits of a word together;
     3. DRAM must be refreshed; SRAM and registers must not;
     4. the register image is a byte map: 16 x 8 = 128 bytes;
     5. async inputs get synchronized; the register file never does;
     6. comparisons set flags; branches read them; never separate them.
\end{verbatim}

\begin{researchnote}[Two century-spanning inventions]
The flip-flop is older than the transistor. William Eccles and
F.~W.~Jordan patented their ``trigger relay''---two amplifying valves
cross-coupled in a positive feedback loop, the same loop as
Section~\ref{sec:flipflop}---in 1918, and a quarter-century later it
was storing the state of the Colossus codebreaking computers. The
gates-drawn diagram of the SR NOR latch in this chapter is the direct
descendant of their two vacuum tubes.

The other half of the story is the scaling rule that made registers
and DRAM cheap. R.~H.~Dennard and colleagues showed in 1974 that
shrinking transistors could keep power density constant while running
faster; the chemistry of leakage in Section~\ref{sec:chemistry} is
the reason that law eventually broke, and it is the same physics that
forces DRAM to refresh today. The metastability side of the story
has its own founding paper: Chaney and Molnar's 1973 proof, cited in
Section~\ref{sec:clock}, that the problem is inherent to synchronizing
two unrelated clocks.

\paperref{Wikipedia, ``Flip-flop (electronics)''}{https://en.wikipedia.org/wiki/Flip-flop_(electronics)}
\end{researchnote}

With the machine's fastest memory understood bit by bit, the next
chapter returns to the programmer's view: Chapter~3 combines scalars
into \emph{structures}, and shows how the machine's demand for
aligned access quietly inserts invisible padding bytes between our
fields.
}{}
\IfFileExists{chapter3.tex}{% ============================================================================
%  chapter3.tex
%  Building Abstractions with Structs
% ============================================================================

\chapter{Building Abstractions with Structs}

\section{From Scalars to Records}

Chapter~1 gave us bytes, addresses, and pointers, and Chapter~2
opened the registers that hold them while the machine runs. With only
scalar
types (\code{int}, \code{char}, \code{double}) we can already do real
work, but real systems are full of \emph{aggregates}: a connection has
a socket, a buffer, a state, a timeout; a record has a key and a
value; a node has data and links to its neighbors. In C, the
\emph{struct} is the tool that groups related data into one object
whose bytes live contiguously in memory.

A struct is a \emph{layout contract}: it tells the compiler exactly how
to arrange a set of fields one after another in memory, so that the
whole group can be passed around, pointed at, and stored as a single
unit. Everything else in this book---linked lists, hash tables, trees,
graphs, sockets---is, at bottom, structs pointing at other structs.

\section{Declaring and Using a Struct}

\begin{lstlisting}[caption={A minimal struct: a 2D point.}]
#include <stdio.h>

struct point {
    int x;
    int y;
};

int main(void) {
    struct point origin;       /* the struct is a variable like any other */

    origin.x = 0;              /* fields are accessed with the dot */
    origin.y = 0;

    printf("origin = (%d, %d)\n", origin.x, origin.y);
    return 0;
}
\end{lstlisting}

The declaration \code{struct point \{ int x; int y; \};} introduces a
new type named \code{struct point}. It does \emph{not} allocate any
memory; it merely describes the shape. Only when we declare a variable
of that type---\code{struct point origin;}---does the compiler reserve
the bytes. The memory reserved is exactly the sum of the field sizes,
plus whatever padding alignment requires, as we are about to see.

\subsection{typedef: Giving the Struct a Short Name}

Writing \code{struct point} everywhere is verbose, and in this book we
will write a great many struct types. \code{typedef} introduces an
alias for an existing type, so we can say \code{point} and mean
\code{struct point}. From here on we will declare our structs with a
\code{typedef} and refer to them by the bare name.

\begin{lstlisting}[caption={A typedef names the struct; a companion pointer typedef follows.}]
#include <stdio.h>

typedef struct point {
    int x;
    int y;
} point;

typedef struct node {
    int          value;
    struct node *next;   /* inside the definition we still use "struct node" */
} node;

int main(void) {
    point origin;        /* "point" is an alias for "struct point" */

    origin.x = 0;
    origin.y = 0;

    printf("(%d, %d)\n", origin.x, origin.y);
    return 0;
}
\end{lstlisting}

Two notes on discipline. First, the tag name \code{point} in
\code{typedef struct point \{...\} point;} is optional but useful for
self-documentation. Second, \code{typedef} does \emph{not} create a
new type; it creates an alias. The rules of argument passing, copying,
and layout apply to the alias exactly as they do to the original. And a
pointer typedef, \code{typedef struct node *node\_ptr;}, is a style
choice many codebases make; we avoid it in this course, because
hiding a pointer inside a typedef makes it easy to forget whether you
hold a pointer or a value. Pointers stay explicit.

\section{Memory Layout of a Struct}

Because a struct's fields are laid out in declaration order, we can
draw the layout of \code{struct point} with perfect precision:

\begin{verbatim}
   struct point { int x; int y; }   -- size 8 bytes

   offset:  0     1     2     3     4     5     6     7
           +-----+-----+-----+-----+-----+-----+-----+-----+
           |           x           |           y           |
           +-----+-----+-----+-----+-----+-----+-----+-----+
            ^                       ^
            address of origin       address of origin.y
            == address of origin.x
\end{verbatim}

The address of a struct is the address of its first field. The two are
the same number, which is exactly why \code{\&origin.x} and
\code{\&origin} compare equal. Every later field sits at a fixed
\emph{offset} from the start of the struct.

\begin{lstlisting}[caption={A struct is an int and another int, one after another.}]
#include <stdio.h>

struct point {
    int x;
    int y;
};

int main(void) {
    struct point origin;

    origin.x = 10;
    origin.y = 20;

    printf("sizeof(struct point) = %zu\n", sizeof(struct point));
    printf("&origin      = %p\n", (void *) &origin);
    printf("&origin.x    = %p\n", (void *) &origin.x);
    printf("&origin.y    = %p\n", (void *) &origin.y);
    printf("offset of y  = %ld\n",
           (long) ((char *) &origin.y - (char *) &origin.x));

    return 0;
}
\end{lstlisting}

We printed the offset by subtracting two \code{char*} pointers. Pointer
subtraction between two pointers into the same array of \code{char}
yields the number of \code{char} elements between them, which is the
offset in bytes. This is a legitimate, explicit technique; note that we
cast to \code{char*} precisely so that the arithmetic counts bytes and
not elements.

\section{Data Alignment: Why the CPU Demands It}

Here is the fact about modern hardware that explains everything that
looks ``wasteful'' about structs: a CPU reads memory in aligned units.
On x86-64, a 4-byte \code{int} is best placed at an address divisible
by 4; an 8-byte \code{double} at an address divisible by 8. When a
multi-byte object straddles a boundary, the hardware may need two
fetches instead of one, or---on stricter architectures such as some
RISC designs---may raise a fault outright.

\begin{verbatim}
   Correctly aligned (good):

   address: 1000  1001  1002  1003  1004  1005  1006  1007
           +-----+-----+-----+-----+-----+-----+-----+-----+
           |           an int at 1000         |  (next obj) |
           +-----+-----+-----+-----+-----+-----+-----+-----+
            ^ int starts at a multiple of 4

   Misaligned (bad):

   address: 1002  1003  1004  1005  1006  1007  1008  1009
           +-----+-----+-----+-----+-----+-----+-----+-----+
           | (prev) |      an int starting at 1002      |   |
           +-----+-----+-----+-----+-----+-----+-----+-----+
            ^ 1002 is not a multiple of 4: the int straddles a
              boundary and may need two fetches to be read.
\end{verbatim}

To honor alignment, the compiler may insert \emph{padding} bytes
between struct fields. These bytes exist, are addressable, and are
completely invisible to you unless you look at the layout directly.
They cost memory but buy speed and, on many machines, correctness.

\begin{ruleofthumb}[Never assume a field's offset]
The C standard guarantees fields are laid out in declaration order, but
\emph{not} their exact offsets, because padding depends on the machine
and the compiler. Any code that assumes \code{int} sits at offset 0, or
that a struct is exactly the sum of its parts, is broken by definition.
Use \code{sizeof} and \code{offsetof} if you must know the numbers.
\end{ruleofthumb}

\section{Padding: The Hidden Bytes}

Consider a structure whose fields have different alignments:

\begin{lstlisting}[caption={A struct with mixed field sizes.}]
struct entry {
    char   tag;     /* 1 byte,  alignment 1 */
    int    count;   /* 4 bytes, alignment 4 */
    short  level;   /* 2 bytes, alignment 2 */
};
\end{lstlisting}

The compiler must place each field at an offset divisible by its own
alignment. The resulting layout, with padding marked \code{//}:

\begin{verbatim}
   struct entry -- size 12, with 5 padding bytes hidden inside

   offset:  0     1     2     3     4     5     6     7
           +-----+-----+-----+-----+-----+-----+-----+-----+
           | tag | //  | //  | //  |        count         |
           +-----+-----+-----+-----+-----+-----+-----+-----+
   offset:  8     9    10    11
           +-----+-----+-----+-----+
           | level |     //    | // |
           +-----+-----+-----+-----+

   - tag    at offset 0  (alignment 1: any address)
   - count  at offset 4  (must be a multiple of 4; bytes 1-3 are padding)
   - level  at offset 8  (must be a multiple of 2)
   - total raw payload:  1 + 4 + 2 = 7 bytes
   - struct size: 12     (rounded up to a multiple of 4, the largest
                           alignment in the struct, so that arrays of
                           these structs stay aligned)
\end{verbatim}

The struct ended up at 12 bytes for 7 bytes of real data. Five of the
twelve bytes are silent. This is not a bug; it is the machine telling
you how it wants to be spoken to. But notice the waste is not forced:
nothing about the fields themselves demands that layout. If we
redeclare the fields in the right order, the compiler can satisfy every
alignment without a single padding byte between fields.

\section{Field Ordering to Shrink Structs}

Place fields in decreasing order of their alignment. The largest
alignment dominates; once the big fields are placed, the small ones
slot into the gaps.

\begin{lstlisting}[caption={The same data, reordered.}]
struct entry {
    int    count;   /* 4 bytes, alignment 4 */
    short  level;   /* 2 bytes, alignment 2 */
    char   tag;     /* 1 byte,  alignment 1 */
};
\end{lstlisting}

\begin{verbatim}
   struct entry, reordered -- size 8

   offset:  0     1     2     3     4     5     6     7
           +-----+-----+-----+-----+-----+-----+-----+-----+
           |        count         | level | tag | //  |
           +-----+-----+-----+-----+-----+-----+-----+-----+

   - count  at offset 0
   - level  at offset 4
   - tag    at offset 6
   - one padding byte at offset 7 (to round the size up to 8)
   - raw payload still 7 bytes, but the struct is now 8, not 12.
\end{verbatim}

The same information shrank from 12 bytes to 8 bytes, a third smaller,
for free. In a table of a million entries that is 4~MB saved, and on a
cache-bound hot path it is substantially more than a memory saving.
Field ordering is the cheapest optimization in C and one that beginners
never think about.

\begin{insight}[The struct is a memory-layout contract, not just a type]
Reading a struct definition is reading a memory diagram. Before you
write a struct for a hot path or a wire format, order its fields from
largest alignment to smallest, and get into the habit of sketching the
layout by hand. The machine will reward you.
\end{insight}

\section{offsetof and sizeof}

C gives us two operators for interrogating layouts at compile time.
\code{sizeof(type)} returns the total size including trailing padding;
\code{offsetof(type, field)} returns the byte offset of a field.

\begin{lstlisting}[caption={Reading the layout without guessing.}]
#include <stdio.h>
#include <stddef.h>

struct entry {
    int    count;
    short  level;
    char   tag;
};

int main(void) {
    printf("sizeof(struct entry)      = %zu\n", sizeof(struct entry));
    printf("offsetof count            = %zu\n", offsetof(struct entry, count));
    printf("offsetof level            = %zu\n", offsetof(struct entry, level));
    printf("offsetof tag              = %zu\n", offsetof(struct entry, tag));
    return 0;
}
\end{lstlisting}

\code{offsetof} is what lets a generic serializer walk a struct without
special-casing each field: given a pointer to the struct, the address
of a field is simply
\code{(char *) base + offsetof(struct entry, level)}. We will use
exactly this idea when we build a data serialization layer in the
chapters on file I/O and networking.

\begin{ruleofthumb}[Arrays of structs inherit the alignment]
The compiler rounds a struct's size up to a multiple of its strictest
member so that an array of such structs has every element aligned. That
trailing padding is why \code{sizeof} can exceed the sum of the
\code{offsetof} values plus the last field's size. Treat
\code{sizeof(struct)} as authoritative; never compute it by hand.
\end{ruleofthumb}

\section{Structs and Pointers: The Arrow Operator}

Pointers to structs are everywhere, because structs are the units we
link together. To reach a field through a pointer, C offers the
\emph{arrow operator}: \code{p->field} means \code{(*p).field}.

\begin{lstlisting}[caption={Accessing fields through a pointer.}]
#include <stdio.h>

struct point {
    int x;
    int y;
};

/* Move the point that p points at by (dx, dy). */
void translate(struct point *p, int dx, int dy) {
    p->x = p->x + dx;   /* equivalent to (*p).x = (*p).x + dx; */
    p->y = p->y + dy;
}

int main(void) {
    struct point here;
    struct point *cursor;

    here.x = 3;
    here.y = 4;

    cursor = &here;               /* cursor now points at here */
    translate(cursor, 5, 6);      /* here becomes (8, 10)     */

    printf("here = (%d, %d)\n", here.x, here.y);
    printf("via pointer = (%d, %d)\n", cursor->x, cursor->y);
    return 0;
}
\end{lstlisting}

\begin{verbatim}
   Memory state after cursor = &here:

        +--------------------+
        |      here          |
        |  x = 3  y = 4      |
        +--------------------+
            ^
            |
        cursor  (a struct point* holding &here)
\end{verbatim}

Why pointers instead of passing the struct by value? A pointer passes 8
bytes no matter how large the struct is; copying a big struct by value
copies every byte. More importantly, a function that receives a pointer
can \emph{modify} the caller's struct, exactly as \code{translate} did
above. If you want a function to change a struct, pass its address.

\section{const Correctness with Structs}

\code{const} is the compiler's way of recording your intent, and in
struct-heavy code it earns its keep. The keyword has two different
places to attach: to the pointed-to struct (\code{const struct person
*p}, read-only object) or to the pointer itself (\code{struct person
*const p}, a pointer that cannot be reassigned). The first is the one
we use constantly in function signatures.

\begin{lstlisting}[caption={const expresses ``read-only'' on function boundaries.}]
#include <stdio.h>

typedef struct point {
    int x;
    int y;
} point;

/* "p" points at a struct the caller promises not to modify through us. */
int distance_squared(const point *p) {
    return p->x * p->x + p->y * p->y;
}

/* The pointer itself is const: inside this function, q cannot be reassigned. */
void print_point(point *const q) {
    printf("(%d, %d)\n", q->x, q->y);
}

int main(void) {
    point origin = {0, 0};

    print_point(&origin);
    printf("dist^2 = %d\n", distance_squared(&origin));
    return 0;
}
\end{lstlisting}

Read the signatures as promises. \code{distance\_squared} promises it
will not modify the point; a caller can therefore hand it any point,
confident nothing is corrupted. \code{const} does not make memory
physically read-only---it is a compile-time contract, and a
\code{const} pointer to a struct is still an ordinary 8-byte address.
But the contract is enforced: writing through a \code{const} pointer is
a compile error, which turns a whole class of lifetime and aliasing
bugs into loud, immediate failures at build time.

\begin{ruleofthumb}[Declare read-only parameters const]
Any function that receives a pointer to a struct and does not modify
it should declare that parameter \code{const}. It costs nothing, it
documents the interface, and it lets the compiler catch accidental
writes. In Chapter~4 every ``print'' and ``lookup'' function takes
\code{const} pointers; every ``insert'' and ``remove'' does not.
\end{ruleofthumb}

\section{Structs on the Heap}

So far our structs have lived on the stack. But the stack frame dies
with the function, so any struct that must outlive its creator, or
whose count is known only at runtime, is allocated on the heap with
\code{malloc}. This is the moment when structs and pointers become
genuinely powerful: \code{malloc} returns the address of a fresh,
uninitialized region big enough for the struct, and we then fill its
fields through that address.

\begin{lstlisting}[caption={A struct living on the heap, with ownership handled.}]
#include <stdio.h>
#include <stdlib.h>

struct account {
    int   id;
    long  balance;
};

struct account *account_new(int id) {
    struct account *a;

    a = malloc(sizeof(struct account));   /* ask for exactly one struct */
    if (a == NULL) {
        return NULL;
    }

    a->id = id;
    a->balance = 0;

    return a;
}

int main(void) {
    struct account *acc;

    acc = account_new(7);
    if (acc == NULL) {
        fprintf(stderr, "allocation failed\n");
        exit(1);
    }

    printf("account %d has balance %ld\n", acc->id, acc->balance);

    free(acc);       /* we own it, so we free it; no more acc->id after this */
    return 0;
}
\end{lstlisting}

Three details deserve emphasis, because they are the core discipline of
every data structure we build in this book:

\begin{itemize}
  \item \code{malloc(sizeof(struct account))} asks for exactly the
    struct's size---padding included---so the layout contract carries
    over to the heap without us ever computing byte counts.
  \item The \code{if (a == NULL)} check is the contract of every
    allocation. Heap exhaustion is a real, reportable outcome.
  \item Ownership is explicit: \code{account\_new} allocates, and the
    caller frees. If \code{acc} is later embedded in a linked list or a
    hash table, the data structure becomes the owner, and it is the
    data structure's \emph{destroy} function that frees it.
\end{itemize}

\section{Nested Structs}

A struct may contain another struct \emph{by value}. The inner struct's
bytes are laid out inline, within the outer struct, at the offset its
alignment requires. The memory diagram makes the word ``nested'' literal:

\begin{lstlisting}[caption={A struct containing a struct.}]
#include <stdio.h>

typedef struct point {
    int x;
    int y;
} point;

typedef struct rect {
    point  top_left;     /* a whole point, embedded */
    point  bottom_right; /* another whole point, embedded */
} rect;
\end{lstlisting}

\begin{verbatim}
   struct rect -- size 16

   offset:  0     1     2     3     4     5     6     7
           +-----+-----+-----+-----+-----+-----+-----+-----+
           |      top_left.x       |      top_left.y       |
           +-----+-----+-----+-----+-----+-----+-----+-----+
   offset:  8     9    10    11    12    13    14    15
           +-----+-----+-----+-----+-----+-----+-----+-----+
           |   bottom_right.x     |   bottom_right.y       |
           +-----+-----+-----+-----+-----+-----+-----+-----+

   top_left   occupies bytes 0..7
   bottom_right occupies bytes 8..15
\end{verbatim}

Access uses the dot twice: \code{box.top\_left.x}. There is no arrow
involved unless you reach the nested field through a pointer to the
outer struct, in which case it is \code{box->top\_left.y}. The nested
struct has no separate identity; it is a set of fields at fixed offsets
inside its parent. This is how we build hierarchy from flat memory.

\section{Flexible Array Members: The Variable-Length Tail}

C99 introduced the \emph{flexible array member}: a struct whose last
field is an unsized array \code{data[]}. The struct itself occupies
only the fixed fields; the array bytes are tacked on at the end by
whichever allocation you make. This is exactly how real kernels and
libraries carry a length plus a payload in one allocation, and it is
the pattern behind messages, packets, and buffers throughout this book.

\begin{lstlisting}[caption={One allocation holds header and payload together.}]
#include <stdio.h>
#include <stdlib.h>
#include <string.h>

typedef struct message {
    unsigned int len;      /* number of valid bytes in data */
    char data[];           /* flexible array member: no size here */
} message;

int main(void) {
    message *msg;
    size_t payload;

    payload = 5;                        /* we want to carry 5 bytes */
    msg = malloc(sizeof(message) + payload);
    if (msg == NULL) {
        exit(1);
    }

    msg->len = payload;
    memcpy(msg->data, "hello", payload);

    printf("message: len=%u data=%.*s\n", msg->len, (int) msg->len, msg->data);

    free(msg);                          /* one allocation, one free */
    return 0;
}
\end{lstlisting}

\begin{verbatim}
   The allocation from malloc(sizeof(message) + 5):

   +----------------+------------------------------+
   |   len  (4B)    |   data[0..4]  ("hello")      |
   +----------------+------------------------------+
   ^  sizeof(message) = 4            ^  exactly payload bytes

   The whole thing is ONE contiguous allocation, freed with one free().
   No pointer chase, no leak, no fragmentation from a second malloc.
\end{verbatim}

The rules are strict: the flexible array must be the last member, there
must be at least one named member before it, and you cannot
\emph{store} an array of these structs---you must allocate each one
individually, because its size is not fixed at compile time. Used
honestly, it removes an entire class of double-allocation and
free-order bugs from protocols, virtual machines, and caches.

\section{Self-Referential Structures}

A struct may contain a pointer to an object of its own type. The type
is incomplete during its own definition, but a \emph{pointer} to it is
always allowed, because a pointer is a fixed-size address regardless of
what it targets. This one trick---structs pointing at themselves---is
the seed from which linked lists, trees, and graphs all grow.

\begin{lstlisting}[caption={The birth of every linked structure.}]
struct node {
    int value;
    struct node *next;   /* pointer to another node, or NULL */
};
\end{lstlisting}

\begin{verbatim}
   Three nodes linked into a chain:

   head
    |
    v
   +---------+      +---------+      +---------+
   | value=10|      | value=20|      | value=30|
   | next ---|----->| next ---|----->| next ---|-----> NULL
   +---------+      +---------+      +---------+
      ^                 ^                 ^
      node A            node B            node C
      (heap)            (heap)            (heap)

   Each box is a struct node. The last "next" points at NULL, marking
   the end of the chain. The whole list is named by the single pointer
   "head".
\end{verbatim}

The list is a web of arrows, exactly the diagrams we drew in
Chapter~1. A linked list is not a special language feature; it is a
pointer stored inside a struct pointing at another instance of that
struct. That realization, once felt, unlocks every container in the
standard toolbox.

\section{Copying and Returning Structs}

Unlike arrays, structs can be copied and returned by value. The
compiler copies the whole layout, padding included. A struct returned
by value is a complete, independent object.

\begin{lstlisting}[caption={Structs copy and return by value; arrays do not.}]
#include <stdio.h>

struct point {
    int x;
    int y;
};

struct point make_point(int x, int y) {
    struct point p;

    p.x = x;
    p.y = y;
    return p;          /* a copy of p is handed back; safe, p is local */
}

int main(void) {
    struct point a = make_point(2, 3);
    struct point b;

    b = a;             /* b is now an independent copy of a */
    b.x = 100;

    printf("a = (%d, %d)\n", a.x, a.y);   /* still (2, 3) */
    printf("b = (%d, %d)\n", b.x, b.y);   /* (100, 3)    */
    return 0;
}
\end{lstlisting}

Returning a \emph{struct} by value is safe; returning a pointer to a
local struct is the dangling-pointer bug of Chapter~1. The rule of
thumb: small structs (a few words) are fine by value; large ones are
cheaper as pointers, but then you own the lifetime problem that comes
with the heap.

\begin{insight}[Shallow vs.\ deep copy]
\code{b = a} copies the struct's bytes verbatim. If the struct contains
a pointer (like our \code{node} above), the \emph{pointer} is copied,
not the node it points at: \code{b.next} and \code{a.next} name the
same memory. That is a shallow copy, and it is usually what we want in
containers. Deep copying---duplicating everything reachable through
pointers---is a deliberate, manual operation we will write by hand when
we build graphs in Chapter~4.
\end{insight}

\section{Array of Structs vs.\ Struct of Arrays}

When you hold many records, there are two ways to arrange them in
memory, and the choice is a measurable performance decision that
systems programmers make on every hot path.

\textbf{Array of Structs} (AoS) interleaves the fields: element 0's all
fields, then element 1's all fields. \textbf{Struct of Arrays} (SoA)
groups like fields together: all \code{x} values, then all \code{y}
values.

\begin{verbatim}
   Array of Structs (AoS):              Struct of Arrays (SoA):

   point points[4];                     struct points {
                                            int xs[4];
       [x0 y0 x1 y1 x2 y2 x3 y3]            int ys[4];
                                        };

   each element is contiguous:          all xs live together:
   points[i].x and points[i].y         points.xs[i]
   are adjacent bytes                  points.ys[i] are far apart

   good when you touch both x and y     good when a pass touches only x
   together (all fields of a record)    (a single field across records)
\end{verbatim}

The modern memory hierarchy makes this matter. The CPU loads memory in
cache lines (typically 64 bytes). AoS is ideal when each record is used
wholesale---the whole record arrives in one line. SoA is ideal when an
algorithm scans one field of every record, because consecutive accesses
land in consecutive addresses, and the fields you never touch are never
pulled into cache. Scientific code, game engines, and database column
stores all exploit SoA.

\begin{insight}[Which layout does your access pattern want?]
Ask: do my loops walk \emph{records} or \emph{fields}? Walking records
wants AoS; walking one field at a time wants SoA. There is no
universal winner, and getting this right is often worth more than any
micro-optimization inside a loop.
\end{insight}

\section{Bit Fields: Squeezing Flags into Bytes}

For control registers, file formats, and network headers, C offers
\emph{bit fields}: members whose width is specified in bits. The layout
of bit fields is implementation-defined in painful detail, so we treat
them as a specialized tool, not a general one. They are invaluable
when a hardware or protocol spec forces bits to be packed.

\begin{lstlisting}[caption={Bit fields: several flags in one storage unit.}]
#include <stdio.h>

struct flags {
    unsigned int valid  : 1;   /* one bit  */
    unsigned int dirty  : 1;   /* one bit  */
    unsigned int kind   : 3;   /* three bits (0..7) */
};

int main(void) {
    struct flags f;

    f.valid = 1;
    f.dirty = 0;
    f.kind  = 5;

    printf("valid=%u dirty=%u kind=%u\n", f.valid, f.dirty, f.kind);
    return 0;
}
\end{lstlisting}

Notice we declared them \code{unsigned int}. A signed one-bit field can
only hold $-1$ or $0$, which is almost never what anyone means.

\begin{ruleofthumb}[Bit fields are for hardware and wire formats]
If your struct is destined for a control register or a protocol header
whose bit layout is dictated by a spec, bit fields are the right tool.
If your struct merely wants ``small flags'', prefer ordinary
\code{int} fields (or bitwise \code{\&} / \code{|} on a single
integer); the exact layout of bit fields across compilers is not
portable, and portability is a promise we keep in this course.
\end{ruleofthumb}

\section{Unions: Overlapping Storage}

A \code{union} is a struct whose members all occupy the \emph{same}
bytes. Its size is the size of its largest member, not the sum. Writing
one member overwrites the others, so a union is only as meaningful as
the convention that says which member is live. Pair a union with a
``kind'' tag and you have a \emph{tagged union}, the C ancestor of the
enum variant that powers trees, graphs, and protocol parsers everywhere.

That tag is an \code{enum}---short for \emph{enumeration}, a type that
names a set of integer constants. \code{enum kind \{ KIND\_INT,
KIND\_DOUBLE, KIND\_TEXT \};} declares three constants
(\code{KIND\_INT} is 0, the next 1, the next 2, unless you say
otherwise) and a type \code{kind} that holds one of them. It exists so
your code can say \code{v.tag = KIND\_INT} instead of the
incomprehensible \code{v.tag = 0}. In C an enum is just a set of named
integer constants, so the compiler does not police which values you
assign or compare---the names exist for humans, not for enforcement.
This is the C version of a named choice, and you will see it in every
parser and protocol handler later in the book.

\begin{verbatim}
   struct or union?  Same members, very different size.

   struct thing {                   union thing {
       int n;      // 4 bytes            int n;      // 4 bytes
       double d;   // 8 bytes            double d;   // 8 bytes
       char c[16]; // 16 bytes           char c[16]; // 16 bytes
   };               // 32 bytes       };              // 16 bytes

   A struct keeps every member:               A union reuses the space:
   +----+----+----+----+----+----+----+      +----+----+----+----+----+
   | n  | pad |      d      |    c[16]   |   |   n, or d, or c[16]    |
   +----+----+----+----+----+----+----+----+  +----+----+----+----+----+
   0    4   8                16     32        0                        16
\end{verbatim}

The standard trick is to record \emph{which} member is live in a
sibling field, and to treat writes to the union as a single operation.

\begin{lstlisting}[caption={A tagged union: kind says which member is live.}]
#include <stdio.h>

typedef enum kind {
    KIND_INT,
    KIND_DOUBLE,
    KIND_TEXT
} kind;

typedef struct value {
    kind tag;              /* which member of u is valid */
    union {
        int    as_int;
        double as_double;
        char   as_text[16];
    } u;
} value;

void value_print(const value *v) {
    /* We promise to read only the member named by v->tag. */
    if (v->tag == KIND_INT) {
        printf("%d\n", v->u.as_int);
    } else if (v->tag == KIND_DOUBLE) {
        printf("%f\n", v->u.as_double);
    } else {
        printf("%s\n", v->u.as_text);
    }
}

int main(void) {
    value v;

    v.tag = KIND_INT;
    v.u.as_int = 42;
    value_print(&v);

    v.tag = KIND_TEXT;
    v.u.as_text[0] = 'h';
    v.u.as_text[1] = 'i';
    v.u.as_text[2] = '\0';
    value_print(&v);

    return 0;
}
\end{lstlisting}

The danger is symmetrical and classic: read the wrong member and you
interpret one kind of bits as another, silently producing garbage
numbers or strings. The tag discipline is the price of memory sharing.
Used with discipline, tagged unions give us a single \code{value} type
for the key/value engines, expression trees, and configuration parsers
we build in later chapters.

\section{Practical Exercise: A Person Registry}

Let us assemble the chapter into a working registry: a linked list of
\code{struct person} records, built entirely with the techniques we
have seen. This is the skeleton on which the serious containers of
Chapter~4 will be built, so read its shape carefully: constructor,
push-front, traversal, destroy.

\begin{lstlisting}[caption={registry.c: a minimal linked list of persons.}]
#include <stdio.h>
#include <stdlib.h>
#include <string.h>

struct person {
    char name[32];
    int  age;
    struct person *next;      /* the link to the following person, or NULL */
};

/* Create a new person on the heap. Returns NULL on failure. */
struct person *person_new(const char *name, int age) {
    struct person *p;

    p = malloc(sizeof(struct person));
    if (p == NULL) {
        return NULL;
    }

    strncpy(p->name, name, sizeof(p->name) - 1);
    p->name[sizeof(p->name) - 1] = '\0';
    p->age = age;
    p->next = NULL;

    return p;
}

/* Insert p at the front of the list named by *head. */
void list_push_front(struct person **head, struct person *p) {
    p->next = *head;          /* p points at whatever was first before  */
    *head = p;                /* and *head now names p as the new first */
}

/* Print every person in the list. */
void list_print_all(const struct person *head) {
    const struct person *cur;

    cur = head;
    while (cur != NULL) {
        printf("%s, age %d\n", cur->name, cur->age);
        cur = cur->next;      /* advance to the next node */
    }
}

/* Free every node in the list. */
void list_free_all(struct person *head) {
    struct person *cur;
    struct person *next;

    cur = head;
    while (cur != NULL) {
        next = cur->next;     /* remember the successor before freeing  */
        free(cur);            /* free this node                          */
        cur = next;           /* move on                                */
    }
}

int main(void) {
    struct person *head = NULL;

    list_push_front(&head, person_new("ada", 36));
    list_push_front(&head, person_new("linus", 55));
    list_push_front(&head, person_new("ken", 81));

    list_print_all(head);

    list_free_all(head);
    head = NULL;   /* head is now dangling; clearing it is good hygiene */

    return 0;
}
\end{lstlisting}

Study \code{list\_push\_front} closely. It takes \code{struct person
**head} because it must \emph{rewrite} the caller's \code{head}
variable---the exact pointer-to-pointer pattern from Chapter~1.
And \code{list\_free\_all} saves \code{next} \emph{before} calling
\code{free}, because after \code{free(cur)} the node is gone and
reading \code{cur->next} would be a use-after-free. Every line here
exists to keep the arrows honest.

\section{Structures in Bytes}

The layouts you computed by hand earlier in this chapter, collected on
one page. On a 64-bit machine with 4-byte \code{int} and 8-byte
\code{double}:

\begin{verbatim}
   struct point { int x; int y; }        size  8   x:0  y:4
   struct entry (tag, count, level)      size 12   tag:0 count:4 level:8
   struct entry (count, level, tag)      size  8   count:0 level:4 tag:6
   struct rect  { point tl; point br; }  size 16   tl:0  br:8
   union thing  { int n; double d; char c[16]; }   size 16
   struct flags (three bit fields)       size  4   one storage unit
   struct message { uint len; char data[]; }       size  4  (FAM adds nothing)

   "field:off" is the byte offset a debugger reports; size is what
   sizeof() returns. Both come from the alignment rules, never from
   hand arithmetic.
\end{verbatim}

The three rules in byte terms: every field starts at a byte offset
divisible by its own alignment; the struct's size rounds up to a
multiple of its strictest member's alignment; and a union occupies
exactly its largest member. Bit fields are the exception to ``every
field has an address''---a 1-bit field has no byte address, only a bit
position inside its storage unit. And \code{offsetof(struct entry,
level)} being 8 while \code{sizeof(struct entry)} is 12 is the
padding: 4 bytes the compiler reserved for you. Those padding bytes
are exactly why Chapter~16 will refuse to \code{fwrite} a struct to
disk and will write a wire format instead.

\section{Summary}

We now possess the two great tools of C: addresses/pointers
(Chapter~1) and structured, contiguous records (this chapter). The
combination---structs that point at each other---is the raw material of
every data structure that follows.

\begin{verbatim}
   The toolbox so far:

   +--------------------------+-----------------------------------------+
   | idea                     | the tool                                |
   +--------------------------+-----------------------------------------+
   | an address               | pointer: int *p;                        |
   | a grouped record         | struct point { int x; int y; };         |
   | a short alias for a type | typedef struct point {...} point;       |
   | record on the heap       | malloc(sizeof(point));                  |
   | field through ptr        | p->x                                    |
   | read-only contract       | const point *p                          |
   | a struct inside a struct | point top_left;  (nested by value)      |
   | a chain of records       | struct node { ...; struct node *next; } |
   | header + payload         | char data[];   (flexible array member)  |
   | silent extra bytes       | alignment padding (order fields!)       |
   | one storage, many views  | union { ... }  + a tag                  |
   | rows or columns          | AoS vs SoA (choose by access pattern)   |
   +--------------------------+-----------------------------------------+
\end{verbatim}

In Chapter~4 we go to work. The Mega Chapter builds and dissects the
entire menagerie---linked lists in every flavor, skip lists, hash
tables, trees from $N$-ary to B+ trees, and graphs with edges and
weights---each one drawn with memory diagrams and implemented in this
explicit, pointer-honest style.
}{}
\IfFileExists{chapter4.tex}{% ============================================================================
%  chapter4.tex
%  Data Structures and Algorithms in C (The Mega Chapter)
% ============================================================================

\chapter{Data Structures and Algorithms in C}

\section{Welcome to the Menagerie}

Everything in the previous chapters was preparation for a single
realization: \emph{a data structure is a disciplined web of pointers}. A
linked list is a struct pointing at a struct. A tree is a struct with
two or more pointers to structs. A graph is an array of lists of
structs. A hash table is an array of pointers to structs, chosen by
hashing. There is nothing else under the hood---only structs, pointers,
and the ownership rules we insisted upon.

This chapter is the biggest in the book, and deliberately so. We
implement, draw, and dissect the entire family of structures a systems
programmer reaches for daily. For each one we answer the same four
questions, the ones that settle every design:

\begin{verbatim}
   1. SHAPE  : what does the struct look like?       (draw it)
   2. COST   : what do search/insert/delete cost?    (O(...))
   3. OWNER  : who allocates, who frees?             (write it down)
   4. CACHE  : how do the bytes sit in memory?       (think about it)
\end{verbatim}

The code follows the house style without exception: pointers move on
their own lines, \code{size\_t} is reserved for loop counters, every
\code{malloc} is checked, and every diagram is drawn before any code is
trusted.

\section{Linked Lists}

\subsection{The Singly Linked List}

The singly linked list is the first structure, the seed from which the
others grow. A node holds a value and one pointer: the address of the
next node.

\begin{lstlisting}[caption={The singly linked node.}]
typedef struct node {
    int value;
    struct node *next;
} node;
\end{lstlisting}

\begin{verbatim}
   head
    |
    v
   +-------+      +-------+      +-------+
   | val=3 |      | val=7 |      | val=1 |
   | next -----> | next -----> | next -----> NULL
   +-------+      +-------+      +-------+

   - The list is named by the single pointer "head".
   - Every "next" is an arrow to the following node.
   - The final "next" is NULL: the arrow that points at nothing.
   - You may only walk FORWARD; to reach node k you must pass through
     nodes 0..k-1, one dereference at a time.
\end{verbatim}

Push at the front is the operation we already know from Chapter~3: it
rewrites the caller's \code{head} through a pointer-to-pointer.

\begin{lstlisting}[caption={Push at the front. Returns 1 on success.}]
int list_push_front(node **head, int value) {
    node *new_node;

    new_node = malloc(sizeof(node));
    if (new_node == NULL) {
        return 0;
    }
    new_node->value = value;
    new_node->next = *head;
    *head = new_node;
    return 1;
}
\end{lstlisting}

Singly lists reward one trick above all others: keeping the address of
the \emph{link} you are interested in, rather than the node. Inserting
in sorted order and deleting both collapse to the same pattern---walk a
\code{node **} and splice.

\begin{lstlisting}[caption={Insert keeping the list sorted ascending.}]
int list_insert_sorted(node **head, int value) {
    node *new_node;
    node **cur;

    new_node = malloc(sizeof(node));
    if (new_node == NULL) {
        return 0;
    }
    new_node->value = value;
    new_node->next = NULL;

    /* cur is the ADDRESS of a "next" link, not the address of a node.
       It walks until the first node whose value is >= value.        */
    cur = head;
    while (*cur != NULL && (*cur)->value < value) {
        cur = &(*cur)->next;
    }

    /* Splice: the new node points at the node cur names, and the
       link cur names now points at the new node.                  */
    new_node->next = *cur;
    *cur = new_node;

    return 1;
}
\end{lstlisting}

\begin{verbatim}
   Insert 5 into  3 -> 7 -> NULL

   head
    |
    v
   +-------+      +-------+
   | val=3 |      | val=7 |
   | next -----> | next -----> NULL
   +-------+      +-------+

   cur begins as &head (an address OF A POINTER).
     step 1: *cur = 3-node, 3 < 5, so cur = &(3-node.next)
     step 2: *cur = 7-node, 7 not < 5, stop

   Now:
     new_node->next = *cur;   (new node -> 7-node)
     *cur = new_node;         (3-node.next -> new node)

   head
    |
    v
   +-------+      +-------+      +-------+
   | val=3 |      | val=5 |      | val=7 |
   | next -----> | next -----> | next -----> NULL
   +-------+      +-------+      +-------+

   Removing a node uses the same trick, and never needs a special case
   for "removing the head".
\end{verbatim}

\begin{lstlisting}[caption={Remove the first node holding value.}]
int list_remove(node **head, int value) {
    node **cur;
    node *victim;

    cur = head;
    while (*cur != NULL && (*cur)->value != value) {
        cur = &(*cur)->next;
    }

    if (*cur == NULL) {
        return 0;              /* not present */
    }

    victim = *cur;
    *cur = victim->next;       /* the predecessor now skips victim */
    free(victim);
    return 1;
}
\end{lstlisting}

Reading \code{list\_remove}: because \code{cur} is an address of a
link, unlinking is one assignment. If the victim was the head,
\code{cur} is \code{\&head} and \code{*cur} updates the caller's head
directly. No branch, no special case, no bug.

\begin{verbatim}
   Removing the 7-node from  3 -> 7 -> 9 -> NULL:

   BEFORE                                  AFTER
   +-------+   +-------+   +-------+       +-------+   +-------+
   | val=3 |   | val=7 |   | val=9 |       | val=3 |   | val=9 |
   | next ----->| next ----->| next --->   | next ----->| next --->
   +-------+   +-------+   +-------+       +-------+   +-------+
                 ^ victim

   *cur (which is 3-node.next) is set to victim->next (the 9-node).
   Then free(victim). The 7-node is gone; the chain is unbroken.
\end{verbatim}

Traversal and destruction are the loops we saw in Chapter~3: walk with
\code{cur = cur->next}, and in \code{free\_all} save the successor
before freeing. Search is the same walk returning the matching node, or
\code{NULL}.

\begin{insight}[Why we keep the link address, not the node]
The singly linked list has only forward arrows. To delete a node you
must be able to reach the arrow that points at it---the \emph{previous}
node's \code{next}. Rather than tracking the previous node in a second
variable, we walk a pointer \emph{to} the link itself. \code{node **}
is not advanced; it is exactly the right precision: it is the address
of the arrow.
\end{insight}

\subsection{The Doubly Linked List}

Add a second pointer to each node and you can walk both ways, delete a
node in \code{O(1)} given its address, and insert before a node without
searching. The cost is one extra pointer per node.

\begin{lstlisting}[caption={The doubly linked node.}]
typedef struct dnode {
    int value;
    struct dnode *prev;
    struct dnode *next;
} dnode;
\end{lstlisting}

\begin{verbatim}
   head
    |
    v
   +--------+      +--------+      +--------+
   | prev=NULL |   | prev=3 |       | prev=7 |
   |  value=3 |   | value=7 |       | value=9 |
   | next=7 ----+  | next=9 ----+   | next=NULL |
   +--------+  |  +--------+  |    +--------+
        ^       +----->|       +----->|
        |               |               |
        +----(the prev arrows)----------+

   Two chains of arrows: "next" runs forward, "prev" runs backward.
   head->next is the first node; head->prev is NULL (or the tail, in
   the circular variant below).
\end{verbatim}

Inserting after a given node touches four arrows. Do them in an order
that never reads a node you are about to unlink.

\begin{lstlisting}[caption={Insert a new node right after node at.}]
int dlist_insert_after(dnode *at, int value) {
    dnode *new_node;

    new_node = malloc(sizeof(dnode));
    if (new_node == NULL) {
        return 0;
    }
    new_node->value = value;

    /* Wire the new node into the gap. */
    new_node->prev = at;
    new_node->next = at->next;

    /* Repair the node that used to follow "at". */
    if (at->next != NULL) {
        at->next->prev = new_node;
    }
    at->next = new_node;

    return 1;
}
\end{lstlisting}

\begin{verbatim}
   Inserting value=5 right after the 7-node.  BEFORE and AFTER:

   BEFORE:                          AFTER:
   +------+   +------+   +------+  +------+   +------+   +------+   +------+
   | val=3|   | val=7 |   | val=9|  | val=3|   | val=7 |   | val=5|   | val=9|
   | prev |<->| prev  |<->| prev |  | prev |<->| prev  |<->| prev |<->| prev |
   | next |   | next  |   | next |  | next |   | next  |   | next |   | next |
   +------+   +------+   +------+  +------+   +------+   +------+   +------+
                  ^ at                            ^ at

   The new node (value=5) is spliced in between 7 and 9.
\end{verbatim}

Deleting a node given only its address is the doubly linked list's
signature advantage:

\begin{lstlisting}[caption={Delete a node in O(1), given its address.}]
void dlist_remove(dnode *victim) {
    /* Repair the left neighbor, if any. */
    if (victim->prev != NULL) {
        victim->prev->next = victim->next;
    }
    /* Repair the right neighbor, if any. */
    if (victim->next != NULL) {
        victim->next->prev = victim->prev;
    }
    free(victim);
}
\end{lstlisting}

No search is needed: the victim itself carries the arrows to its two
neighbors, and each neighbor points back at the victim. Two assignments
unlink it. This is why hash-table buckets, LRU caches, and memory
allocators all use doubly linked lists internally.

\subsection{The Circular Linked List}

Make the tail point back at the head instead of at \code{NULL}, and the
list becomes a ring. With a single \code{tail} pointer, both push-front
and push-back are \code{O(1)}, and a walk never ends---which is exactly
what a round-robin scheduler wants.

\begin{lstlisting}[caption={The circular node is just a node; the ring is a convention.}]
typedef struct cnode {
    int value;
    struct cnode *next;
} cnode;
\end{lstlisting}

\begin{verbatim}
   A circular list. The last node's next points back at the first.

   tail
    |
    v
   +-------+      +-------+      +-------+
   | val=1 |      | val=3 |      | val=7 |
   | next -----> | next -----> | next --+
   +-------+      +-------+      +-------+
       ^                             |
       +-----------------------------+

   - tail->next is the first node (O(1) "front").
   - tail itself is the last node (O(1) "back").
   - Walking: cur = cur->next revisits 1 after 7, forever.

   Round-robin across 3 processes:
     process 1 runs, then 2, then 3, then back to 1...
     cur = cur->next;  -- a single pointer, a single step.
\end{verbatim}

Inserting at the front of a circular list is the singly linked push
plus one extra fix-up to keep the ring closed.

\begin{lstlisting}[caption={Push at the front of a circular list.}]
int clist_push_front(cnode **tail, int value) {
    cnode *new_node;

    new_node = malloc(sizeof(cnode));
    if (new_node == NULL) {
        return 0;
    }
    new_node->value = value;

    if (*tail == NULL) {
        /* First node: the ring is just itself. */
        new_node->next = new_node;
        *tail = new_node;
        return 1;
    }

    new_node->next = (*tail)->next;   /* new node points at the old first */
    (*tail)->next = new_node;         /* tail still points at the last    */
    return 1;
}
\end{lstlisting}

\begin{verbatim}
   Why "tail" and not "head"? With a tail pointer:

     front = tail->next     O(1)
     back  = tail           O(1)

   With only a head pointer, finding the last node to push at the back
   would require walking the whole ring. The tail pointer buys O(1) on
   both ends at the cost of one field.
\end{verbatim}

\begin{insight}[The sentinel head]
The empty ring is the only case that needs care. Our push handles it
with an explicit \code{if (*tail == NULL)}. An alternative that
eliminates the branch entirely is a permanent \emph{sentinel} node that
is never data---the ring always exists, and ``empty'' means the
sentinel points at itself. Linux kernel lists (next subsection) are
built exactly this way, so that the empty list and the non-empty list
are the same code path.
\end{insight}

\subsection{The Intrusive Linked List}

Every list so far owned its data: the node \emph{was} the value. The
\emph{intrusive} list inverts this: your data structure contains the
links. Instead of \code{struct node \{ int value; node *next; \}}, you
embed a generic \code{list\_head} inside your own struct, and the list
machinery manipulates only the embedded links. The data structure never
knows what its elements are---it owns no \code{malloc}, it frees
nothing, it just arranges arrows.

\begin{lstlisting}[caption={The generic link: no data, just arrows.}]
typedef struct list_head {
    struct list_head *next;
    struct list_head *prev;
} list_head;
\end{lstlisting}

\begin{lstlisting}[caption={A struct that carries a link inside it.}]
typedef struct process {
    int pid;
    int priority;
    list_head links;     /* the intrusive hook: the list arranges this */
} process;
\end{lstlisting}

\begin{verbatim}
   Intrusive list: the chain is made of list_head fields EMBEDDED in
   the processes.

   +----------------+---------------+
   | process { pid = 1000, prio 3 } |
   |   links.next --+               |
   |   links.prev <-- NULL          |
   +----------------+---------------+
                    |   (next: down, prev: up)
                    v
   +----------------+---------------+
   | process { pid = 1001, prio 1 } |
   |   links.next --+               |
   |   links.prev <-+               |
   +----------------+---------------+
                    |
                    v
   +----------------+---------------+
   | process { pid = 1002, prio 5 } |
   |   links.next --> NULL          |
   |   links.prev <-+               |
   +----------------+---------------+

   The list code sees only the links; it never knows pid exists.
   Given a pointer to a link, the process is recovered by arithmetic:

       process *p = (process *) ((char *) link_ptr - offsetof(process, links));
\end{verbatim}

That last step, going from a member pointer back to the containing
struct, is the entire trick, and the Linux kernel gives it a name:
\code{container\_of}. It is just pointer subtraction with
\code{offsetof}.

\begin{lstlisting}[caption={container\_of: from a member back to its container.}]
#include <stddef.h>

/* Given a pointer to a member inside a struct, return the struct. */
#define container_of(ptr, type, member) \
    ((type *) ((char *) (ptr) - offsetof(type, member)))
\end{lstlisting}

The list operations are uniform and tiny. We use a permanent sentinel
\code{head} that forms a ring with itself when empty, which removes
every ``empty list'' special case.

\begin{lstlisting}[caption={Circular intrusive list with a sentinel head.}]
/* A head is a link too; it never carries data. */
void list_init(list_head *head) {
    head->next = head;
    head->prev = head;
}

/* Insert a new link between prev and next. */
static void list_add_between(list_head *new_node, list_head *prev,
                             list_head *next) {
    next->prev = new_node;
    new_node->next = next;
    new_node->prev = prev;
    prev->next = new_node;
}

/* Insert at the front of the list headed by head. */
void list_push_front(list_head *head, list_head *new_node) {
    list_add_between(new_node, head, head->next);
}

/* Unlink a link from its list. */
void list_remove(list_head *entry) {
    entry->prev->next = entry->next;
    entry->next->prev = entry->prev;
    entry->next = entry;   /* leave the link self-contained (safe) */
    entry->prev = entry;
}

/* Walk the list; on each pass you recover the container with container_of. */
#define list_for_each(pos, head) \
    for (pos = (head)->next; pos != (head); pos = pos->next)
\end{lstlisting}

\begin{verbatim}
   The sentinel head forms the ring:

   empty:  head.next = head, head.prev = head

            +-------+
            | head  |<-+
            | next--|--+
            | prev--|--+
            +-------+

   after pushing two processes: head sits between the front and the
   back, and the ring is closed through it. Walking stops when you
   return to the sentinel, not at NULL.

   Why intrusive?  Because the list is FREE: no allocation per element,
   no node ownership, and an element can belong to several lists at
   once (e.g. one "all processes" list and one "ready queue"). The cost
   is that container_of is your responsibility, and that an element
   cannot outlive its containers.
\end{verbatim}

\begin{insight}[Why Linux uses intrusive lists everywhere]
A kernel may have millions of tasks, and each task can sit in several
lists simultaneously. An intrusive list adds zero allocations and zero
per-element nodes, and the link is wherever you put it. This is the
same philosophy behind \code{container\_of}: ownership and layout
belong to the user, the list only arranges arrows.
\end{insight}

\subsection{The Unrolled Linked List}

A plain linked list scatters one small node per element across the
heap, and each \code{next} dereference may miss the cache. An
\emph{unrolled} list trades a little insertion work for locality: each
node holds a small \emph{array} of values plus one \code{next} pointer,
so a walk visits many elements per node and per cache line.

\begin{lstlisting}[caption={The unrolled node: a block plus one arrow.}]
#define BLOCK_CAP 4

typedef struct unrolled_node {
    int count;                 /* how many slots are currently live */
    int data[BLOCK_CAP];       /* the fixed-size block              */
    struct unrolled_node *next;
} unrolled_node;
\end{lstlisting}

\begin{verbatim}
   An unrolled list holding 1,2,3,4,5,6:

   +-----------------------------------------------+
   | count=4 | 1 | 2 | 3 | 4 | next ----------+    |
   +-----------------------------------------------+ |
   +-----------------------------------------------+ |
   | count=2 | 5 | 6 | _ | _ | next -> NULL     |<--+
   +-----------------------------------------------+

   - Four elements arrive in ONE node: one allocation, one cache line
     (or two), four dereferences avoided.
   - A block that fills to capacity is SPLIT into two half-full blocks.
   - Deletion can MERGE adjacent half-empty blocks to stay dense.
\end{verbatim}

Insertion walks to the right block, then either places the value in a
free slot or splits the block first.

\begin{lstlisting}[caption={Insert at a 0-based index, splitting when full.}]
int ulist_insert(unrolled_node **head, int index, int value) {
    unrolled_node *cur;
    unrolled_node *block;
    size_t pos;
    size_t i;

    if (*head == NULL) {
        block = malloc(sizeof(unrolled_node));
        if (block == NULL) {
            return 0;
        }
        block->count = 0;
        block->next = NULL;
        *head = block;
    }

    /* Advance to the block containing position "index". */
    cur = *head;
    pos = (size_t) index;
    while (cur->next != NULL && pos > (size_t) cur->count) {
        pos = pos - (size_t) cur->count;
        cur = cur->next;
    }

    /* A full block is split: move the upper half to a fresh block. */
    if (cur->count == BLOCK_CAP) {
        block = malloc(sizeof(unrolled_node));
        if (block == NULL) {
            return 0;
        }
        block->count = BLOCK_CAP / 2;
        for (i = 0; i < (size_t) block->count; i++) {
            block->data[i] = cur->data[BLOCK_CAP / 2 + i];
        }
        block->next = cur->next;
        cur->next = block;
        cur->count = BLOCK_CAP / 2;

        /* The insert belongs to the new block if it passed the middle. */
        if (pos > (size_t) cur->count) {
            pos = pos - (size_t) cur->count;
            cur = block;
        }
    }

    /* Shift the block's tail right by one, then store the value. */
    for (i = cur->count; i > pos; i--) {
        cur->data[i] = cur->data[i - 1];
    }
    cur->data[pos] = value;
    cur->count = cur->count + 1;

    return 1;
}
\end{lstlisting}

\begin{verbatim}
   Inserting 9 into the middle of a full block:

   BEFORE                                AFTER SPLIT + INSERT
   +---------------------------+         +---------------------------+
   | count=4 | 1 | 2 | 3 | 4 |          | count=2 | 1 | 2 |   |   | ---> block
   +---------------------------+         +---------------------------+
                                           count=3 | 9 | 3 | 4 |   |
                                         (new block holds the tail)

   The block no longer has to move when the list grows a bit more;
   locality is preserved at the cost of O(block) shifting per insert.
\end{verbatim}

The unrolled list shines for queues and work-stealing stacks where
traversal dominates, and it is the ancestor of the \emph{block} idea
inside B-trees (Section~4.6). It also demonstrates the Loop Counter
Rule in a shift loop where the counter must count \emph{down} without
underflowing: \code{for (i = cur->count; i > pos; i--)} stops at
\code{i == pos}, never below.

\subsection{The XOR Linked List}

A doubly linked list stores two pointers per node. The \emph{XOR
list} stores one integer: the XOR of the addresses of its two
neighbors. From one neighbor and the stored value you always recover
the other, because \code{(a \^{} b) \^{} a == b}.

\begin{lstlisting}[caption={The XOR node: one integer, two neighbors recovered.}]
#include <stdint.h>

typedef struct xor_node {
    int value;
    uintptr_t link;   /* prev XOR next, stored as an unsigned address */
} xor_node;
\end{lstlisting}

\begin{verbatim}
   head ----> [3] ----> [7] ----> [9] ----> NULL

   stored in each node:
     node(3).link = NULL  ^ node(7)     = node(7)
     node(7).link = node(3) ^ node(9)
     node(9).link = node(7) ^ NULL      = node(7)

   walking forward you know the previous node, so:
     next = node.link  XOR  prev

   walking backward you know the next node, so:
     prev = node.link  XOR  next

   It is a doubly linked list with a single pointer per node.
\end{verbatim}

Insertion at the front must repair the old head's link, and every walk
must carry the predecessor explicitly.

\begin{lstlisting}[caption={Push at the front of an XOR list.}]
xor_node *xlist_push_front(xor_node *head, int value) {
    xor_node *new_node;

    new_node = malloc(sizeof(xor_node));
    if (new_node == NULL) {
        return NULL;
    }
    new_node->value = value;

    /* new_node->link must be: prev(NULL) XOR next(head) = head. */
    new_node->link = (uintptr_t) head;

    /* The old head's link was (0 XOR its_next); it becomes
       (its_next XOR new_node), so that its_next is still recoverable. */
    if (head != NULL) {
        head->link = head->link ^ (uintptr_t) new_node;
    }

    return new_node;
}
\end{lstlisting}

\begin{lstlisting}[caption={Walk forward, remembering the predecessor.}]
void xlist_print(const xor_node *head) {
    const xor_node *cur;
    const xor_node *prev;
    const xor_node *next;
    uintptr_t stored;

    prev = NULL;
    cur = head;
    while (cur != NULL) {
        printf("%d ", cur->value);

        stored = cur->link;
        next = (const xor_node *) (stored ^ (uintptr_t) prev);

        prev = cur;
        cur = next;
    }
    printf("\n");
}
\end{lstlisting}

\begin{verbatim}
   The walk keeps TWO pointers alive: the current node and the one
   before it. From those two, the next node is arithmetic, not a
   dereference:

     state:   prev        cur         (stored link)      next
              NULL        [3]         node(7)     -->    [7]
              [3]         [7]         3^9         -->    [9]
              [7]         [9]         node(7)     -->    NULL
\end{verbatim}

\begin{ruleofthumb}[XOR lists are a curiosity, use them sparingly]
The XOR list saves one pointer per node, and pays for it in three ways:
the walk always needs the neighbor (you cannot jump into the middle of
a list and walk from there); casting pointers to \code{uintptr\_t} and
back is implementation-defined, so the trick is not portable to every
platform; and it is incompatible with garbage collectors and with any
form of bounds checking that inspects pointer provenance. Admire it,
understand the arithmetic, then prefer a conventional doubly linked
list in production code.
\end{ruleofthumb}

\section{The Skip List}

The skip list is the structure that made Salvatore Sanfilippo's
Redis famous as a sorted-set implementation, and it deserves its
reputation. It is a stack of linked lists: the bottom level holds every
element, and each higher level holds a random subset, spaced widely
enough that searches leap over large ranges and land near their target
in \code{O(log n)} expected time.

\begin{verbatim}
   A skip list over 3,9,12,17,21,25. Towers run vertically.

   level 4:  head ----------------------------------------> NULL
   level 3:  head ----------------------> 17 ------------> NULL
   level 2:  head ---------> 9 --------> 17 ------------> NULL
   level 1:  head --> 3 ---> 9 --------> 17 ---> 21 ---> NULL
   level 0:  head --> 3 --> 9 --> 12 --> 17 --> 21 --> 25 -> NULL

   Searching for 21, starting at the head's highest level:
     level 4:  next is NULL,            drop down
     level 3:  17 < 21  move; NULL,     drop down
     level 2:  17 < 21  move; NULL,     drop down
     level 1:  21 found in 2 hops here; overall ~4 hops instead of 6.

   The expected search length is O(log n): each level cuts the space
   by a constant factor, exactly like the levels of a balanced tree--
   but a skip list is a SORTED SET built from linked lists, with no
   rotations and no rebalancing passes.
\end{verbatim}

A node carries an array of forward pointers; its height is chosen
randomly, so tall nodes are rare. We build the node with the flexible
array member from Chapter~3.

\begin{lstlisting}[caption={The skip list structures.}]
#include <stdio.h>
#include <stdlib.h>
#include <string.h>

#define SKIP_MAX_LEVEL 16
#define SKIP_P          4   /* 1/P probability of one more level */

typedef struct skiplist_node {
    int value;
    int level;                        /* how many forward links this node has */
    struct skiplist_node *forward[];  /* forward[0..level-1]                 */
} skiplist_node;

typedef struct skiplist {
    skiplist_node *head;   /* tower of SKIP_MAX_LEVEL forward links */
    int level;             /* highest level currently in use        */
    int size;              /* number of elements                    */
} skiplist;

/* Random tower height: geometric with expectation 1/(1-1/P). */
static int random_level(void) {
    int level;

    level = 1;
    while ((rand() % SKIP_P) == 0 && level < SKIP_MAX_LEVEL) {
        level = level + 1;
    }
    return level;
}

void skiplist_init(skiplist *sl) {
    int i;

    sl->head = malloc(sizeof(skiplist_node) +
                      (size_t) SKIP_MAX_LEVEL * sizeof(skiplist_node *));
    if (sl->head == NULL) {
        exit(1);
    }
    for (i = 0; i < SKIP_MAX_LEVEL; i++) {
        sl->head->forward[i] = NULL;
    }
    sl->head->level = SKIP_MAX_LEVEL;
    sl->level = 1;
    sl->size = 0;
}
\end{lstlisting}

Search first, because insert and delete reuse its shape: descend level
by level, each time advancing as far right as possible without
overshooting.

\begin{lstlisting}[caption={Search: leap down and to the right.}]
/* Returns 1 if value is present. */
int skiplist_search(const skiplist *sl, int value) {
    const skiplist_node *cur;
    size_t i;

    cur = sl->head;
    for (i = sl->level; i > 0; i--) {
        while (cur->forward[i - 1] != NULL &&
               cur->forward[i - 1]->value < value) {
            cur = cur->forward[i - 1];
        }
    }

    /* After the top-down descent, the answer is the next node at level 0. */
    cur = cur->forward[0];
    if (cur != NULL && cur->value == value) {
        return 1;
    }
    return 0;
}
\end{lstlisting}

Notice the down-counting loop written safely: the \code{size\_t} counter
\code{i} counts \emph{down to 1}, never to 0, and the actual level
index is \code{i - 1}. The loop counter never underflows.

Insertion records, for each level, the last node visited before the
insertion point; those nodes' forward links are exactly where the new
tower must be spliced.

\begin{lstlisting}[caption={Insert: record the update path, then splice the tower.}]
/* Insert value. Returns 1 on success, 0 on allocation failure. */
int skiplist_insert(skiplist *sl, int value) {
    skiplist_node *update[SKIP_MAX_LEVEL];
    skiplist_node *cur;
    skiplist_node *new_node;
    size_t i;
    int new_level;

    /* Descend, remembering the last node at each level. */
    cur = sl->head;
    for (i = sl->level; i > 0; i--) {
        while (cur->forward[i - 1] != NULL &&
               cur->forward[i - 1]->value < value) {
            cur = cur->forward[i - 1];
        }
        update[i - 1] = cur;
    }

    /* Choose the tower height; extend the head if we got taller. */
    new_level = random_level();
    if (new_level > sl->level) {
        for (i = sl->level; i < (size_t) new_level; i++) {
            update[i] = sl->head;
        }
        sl->level = new_level;
    }

    new_node = malloc(sizeof(skiplist_node) +
                      (size_t) new_level * sizeof(skiplist_node *));
    if (new_node == NULL) {
        return 0;
    }
    new_node->value = value;
    new_node->level = new_level;

    /* Splice the new node into each level of the update path. */
    for (i = 0; i < (size_t) new_level; i++) {
        new_node->forward[i] = update[i]->forward[i];
        update[i]->forward[i] = new_node;
    }

    sl->size = sl->size + 1;
    return 1;
}
\end{lstlisting}

\begin{verbatim}
   Inserting 10 into the list drawn above, with a 2-tall tower:

   level 2:  update[2] = head      level 2:  head -------------------> 17 -> ...
                                             ^
                                             |new
   level 1:  update[1] = 9         level 1:  head --> 9 ---> [10] ---> 17 -> ...
   level 0:  update[0] = 9         level 0:  head -> 3 -> 9 -> [10] -> 12 -> ...

   update[k] is the node whose forward[k] must point at the new node.
   After splicing, 10 is reachable at every level of its own tower.
\end{verbatim}

Deletion is the mirror image: record the update path, and for each
level of the victim's tower, route \code{update[i]->forward[i]} around
the victim; then free it.

\begin{verbatim}
   The skip list completes the "all the lists" tour:

   structure       search     insert   delete    memory/node
   singly          O(n)       O(1)+    O(n)      1 pointer
   doubly          O(n)       O(1)+    O(1)+     2 pointers
   circular        O(n)       O(1)+    O(1)+     ring
   intrusive       O(n)       O(1)+    O(1)+     embedded, free
   unrolled        O(n)       O(n)     O(n)      block of values
   XOR             O(n)       O(1)+    O(1)+     1 pointer, needs neighbor
   skip list       O(log n)*  O(log n)* O(log n)* towers (probabilistic)
   * expected, and with the same code for sorted ranges.
\end{verbatim}

\begin{insight}[Lists beat arrays for one thing: stable pointers]
An array moves its elements on insert and delete; a list does not. The
skip list keeps that property and adds logarithmic search. That is why
Redis implements sorted sets with skip lists rather than balanced
trees: the code is dramatically simpler, the pointers stay stable, and
the expected performance matches. When you must keep addresses valid
while sorting and searching, reach for a skip list.
\end{insight}

\begin{researchnote}[The skip list's origin]
The skip list was invented by William Pugh in 1990, and its original
paper is worth reading because it makes the central claim of this
section explicit: a sorted structure built from linked lists can match
a balanced tree's guarantees without the tree's machinery.

\begin{verbatim}
   An extract from the paper's abstract:

       "Skip lists are a probabilistic data structure that seem likely
        to supplant balanced trees as the implementation method of
        choice for many applications, since skip lists are
        substantially simpler to implement than balanced trees, and
        provide expected search time proportional to log n."

   In plain terms: "probabilistic" means the towers are chosen by
   coin flips rather than by enforcing a strict balance invariant;
   "expected search time proportional to log n" is the same bound a
   balanced tree guarantees, but the skip list achieves it with far
   less bookkeeping.
\end{verbatim}

The practical payoff is visible in this book's own example: Redis, a
production database serving billions of operations, uses exactly this
structure for sorted sets. A data structure invented in a research
paper is not a classroom curiosity; it is running in production all
around you.

\paperref{W.~Pugh, ``Skip Lists: A Probabilistic Alternative to Balanced
Trees,'' Communications of the ACM 33(6), 1990}{https://15721.courses.cs.cmu.edu/spring2018/papers/08-oltpindexes1/pugh-skiplists-cacm1990.pdf}
\end{researchnote}

\section{Hash Tables}

A hash table trades a function for a search. Instead of comparing a key
against every stored key, we compute a number from the key---the
\emph{hash}---reduce it modulo the table size to get a \emph{bucket
index}, and go straight there. Perfect hashing would give
\code{O(1)}. Reality gives \code{O(1)} expected, provided we handle
collisions (two different keys hashing to the same bucket) and keep the
table from filling up.

\begin{verbatim}
   The core idea:

      key "abc" --hash--> 281872  --mod 8-->  bucket 0
      key "xyz" --hash--> 450817  --mod 8-->  bucket 1
      key "ghi" --hash--> 226514  --mod 8-->  bucket 2

      buckets
        +------+------+------+------+------+------+------+------+
        | [0]  | [1]  | [2]  | [3]  | [4]  | [5]  | [6]  | [7]  |
        +------+------+------+------+------+------+------+------+
         key     key    key
         "abc"   "xyz"  "ghi"

   Two keys with the same bucket collide; the two major collision
   strategies are the two subsections below.
\end{verbatim}

The hash function must be fast, must spread keys evenly, and must be
\emph{deterministic}. The classic djb2 reads a string and mixes it with
\code{h = h * 33 + c}.

\begin{lstlisting}[caption={djb2: a simple, proven string hash.}]
static unsigned int hash_str(const char *s) {
    unsigned int h;
    int c;

    h = 5381;
    c = (unsigned char) *s;
    while (c != 0) {
        h = (h * 33) + (unsigned int) c;
        s = s + 1;
        c = (unsigned char) *s;
    }
    return h;
}
\end{lstlisting}

Note the cast of each character to \code{unsigned char}: \code{char}
may be signed, and negative bytes would otherwise corrupt the mix.
The multiplication by 33 and the accumulation spread the bits across
the whole word, which is what turns similar keys into dissimilar
bucket indexes.

\subsection{Separate Chaining}

Separate chaining makes each bucket a small linked list. A collision is
not an error; it just makes the bucket's chain one node longer. The
whole table is an array of chain heads.

\begin{lstlisting}[caption={A chained hash table.}]
#include <stdio.h>
#include <stdlib.h>
#include <string.h>

typedef struct ht_node {
    char *key;                /* owned copy of the key */
    int   value;
    struct ht_node *next;     /* next node in this bucket */
} ht_node;

typedef struct hashtable {
    ht_node **buckets;        /* array of chain heads */
    int nbuckets;             /* current number of buckets */
    int used;                 /* number of live entries   */
} hashtable;
\end{lstlisting}

\begin{verbatim}
   buckets is an array of pointers; each pointer is the head of a chain.

        +------+      +------+      +------+      +------+
        | [0]  |      | [1]  |      | [2]  |      | [3]  |
        +------+      +------+      +------+      +------+
          |              |             |             |
          v              v             v             v
        +-------+      +-------+     +-------+     +-------+
        | "abc" |      | "ghi" |     | "def" |     | "xyz" |
        | 42    |      | 7     |     | 99    |     | 13    |
        | next -->NULL | next -->NULL | next -->NULL | next -->NULL
        +-------+      +-------+     +-------+     +-------+
        +-------+
        | "zzz" |      <-- "zzz" collided with "ghi": the chain grows
        | 1     |
        | next -->NULL
        +-------+

   A collision simply lengthens a chain. Load factor = used/nbuckets;
   chains average exactly the load factor.
\end{verbatim}

Lookup walks the bucket's chain comparing with \code{strcmp}; update
replaces a value when the key is found; otherwise a new node is pushed
onto the chain.

\begin{lstlisting}[caption={Insert or update. Returns 1 on success.}]
int ht_set(hashtable *ht, const char *key, int value) {
    unsigned int idx;
    ht_node *cur;
    ht_node *new_node;

    idx = hash_str(key) % (unsigned int) ht->nbuckets;

    /* Update in place if the key already exists in its bucket. */
    cur = ht->buckets[idx];
    while (cur != NULL) {
        if (strcmp(cur->key, key) == 0) {
            cur->value = value;
            return 1;
        }
        cur = cur->next;
    }

    /* Not present: allocate a node and its owned key copy. */
    new_node = malloc(sizeof(ht_node));
    if (new_node == NULL) {
        return 0;
    }
    new_node->key = strdup(key);
    if (new_node->key == NULL) {
        free(new_node);
        return 0;
    }
    new_node->value = value;

    new_node->next = ht->buckets[idx];
    ht->buckets[idx] = new_node;

    ht->used = ht->used + 1;
    return 1;
}
\end{lstlisting}

\begin{lstlisting}[caption={Lookup. Returns 1 and sets *out when found.}]
int ht_get(const hashtable *ht, const char *key, int *out) {
    unsigned int idx;
    const ht_node *cur;

    idx = hash_str(key) % (unsigned int) ht->nbuckets;

    cur = ht->buckets[idx];
    while (cur != NULL) {
        if (strcmp(cur->key, key) == 0) {
            *out = cur->value;
            return 1;
        }
        cur = cur->next;
    }
    return 0;
}
\end{lstlisting}

The killer detail in \code{ht\_set}: the key is \emph{copied} with
\code{strdup}, so the table owns the key for its whole life and the
caller may reuse or free the original buffer. Ownership is explicit on
both sides. The destroy function then frees every node and every owned
key copy.

When the load factor climbs (say past 0.75), we \emph{grow} the table:
allocate a new, larger bucket array and rehash every live node into it.
Rehashing matters---\code{idx} depends on \code{nbuckets}, so the old
bucket number is meaningless in the new array.

\begin{lstlisting}[caption={Grow the table and rehash every node.}]
static int ht_resize(hashtable *ht, int new_count) {
    ht_node **old_buckets;
    ht_node **new_buckets;
    ht_node *cur;
    ht_node *next;
    unsigned int idx;
    size_t i;

    new_buckets = calloc((size_t) new_count, sizeof(ht_node *));
    if (new_buckets == NULL) {
        return 0;
    }

    old_buckets = ht->buckets;
    for (i = 0; i < (size_t) ht->nbuckets; i++) {
        cur = old_buckets[i];
        while (cur != NULL) {
            next = cur->next;                /* save before relinking */

            idx = hash_str(cur->key) % (unsigned int) new_count;
            cur->next = new_buckets[idx];    /* rehang into the new table */
            new_buckets[idx] = cur;

            cur = next;
        }
    }

    free(old_buckets);
    ht->buckets = new_buckets;
    ht->nbuckets = new_count;
    return 1;
}
\end{lstlisting}

Because the nodes themselves are reused---only their \code{next}
pointers are rewired---the resize is a cheap pointer shuffle, not a
reallocation of every key.

\begin{ruleofthumb}[Keep the load factor bounded]
Chaining degrades gracefully, but long chains defeat the \code{O(1)}
promise. Grow when \code{used / nbuckets} exceeds roughly 0.75. The
resize doubles the array, so the amortized cost of growth is a
constant per insert---the same argument that pays for dynamic arrays.
\end{ruleofthumb}

\subsection{Open Addressing}

Open addressing keeps all entries \emph{inside} the array itself:
no chains, no per-node allocations, better cache locality. A collision
is resolved by \emph{probing} forward to the next free slot. The
simplest probe is \emph{linear}: try slot \code{i}, then \code{i+1},
then \code{i+2}, wrapping around the end of the array.

\begin{lstlisting}[caption={An open-addressing table with linear probing.}]
typedef enum slot_state {
    SLOT_EMPTY,       /* never used: the search can stop here */
    SLOT_USED,        /* holds a live key */
    SLOT_TOMBSTONE    /* was used, then deleted: search continues */
} slot_state;

typedef struct oa_table {
    char       **keys;
    int        *values;
    slot_state *states;
    int        capacity;
    int        used;
} oa_table;
\end{lstlisting}

\begin{verbatim}
   capacity = 8.  hash("abc")%8 = 3, hash("xyz")%8 = 4, hash("ghi")%8 = 5,
   hash("def")%8 = 3  (collides with "abc").

   Insert in order: "abc", "xyz", "ghi", "def".

   slot:  0        1        2        3        4        5        6        7
        +--------+--------+--------+--------+--------+--------+--------+--------+
        | EMPTY  | EMPTY  | EMPTY  | "abc"  | "xyz"  | "ghi"  | "def"  | EMPTY  |
        |        |        |        |  42    |   7    |   1    |  99    |        |
        +--------+--------+--------+--------+--------+--------+--------+--------+

   "def" wanted slot 3, found it full, probed 4, full, 5, full, 6, empty.
   Linear probing turns collisions into a short run of nearby probes.
   Deletion leaves a TOMBSTONE so later searches keep walking past the
   hole instead of stopping at a false "not found".
\end{verbatim}

Insertion probes forward, stopping at the first empty or tombstone
slot; if it finds the key already present, it updates instead.

\begin{lstlisting}[caption={Insert with linear probing. Returns 1 on success.}]
int oa_set(oa_table *t, const char *key, int value) {
    unsigned int idx;
    size_t i;
    int slot;

    idx = hash_str(key) % (unsigned int) t->capacity;

    for (i = 0; i < (size_t) t->capacity; i++) {
        slot = (int) ((idx + (unsigned int) i) % (unsigned int) t->capacity);

        if (t->states[slot] == SLOT_EMPTY ||
            t->states[slot] == SLOT_TOMBSTONE) {
            /* Free slot: the probe stops here. */
            if (t->keys[slot] != NULL) {
                free(t->keys[slot]);
            }
            t->keys[slot] = strdup(key);
            if (t->keys[slot] == NULL) {
                return 0;
            }
            t->values[slot] = value;
            t->states[slot] = SLOT_USED;
            t->used = t->used + 1;
            return 1;
        }

        if (t->states[slot] == SLOT_USED && strcmp(t->keys[slot], key) == 0) {
            t->values[slot] = value;      /* update an existing key */
            return 1;
        }
    }

    return 0;   /* table full */
}
\end{lstlisting}

Lookup must probe exactly the same way, and it must not stop at a
tombstone---the key it seeks may be past it. It stops only at an empty
slot, which is the proof the key is absent.

\begin{lstlisting}[caption={Lookup with the same probe sequence.}]
int oa_get(const oa_table *t, const char *key, int *out) {
    unsigned int idx;
    size_t i;
    int slot;

    idx = hash_str(key) % (unsigned int) t->capacity;

    for (i = 0; i < (size_t) t->capacity; i++) {
        slot = (int) ((idx + (unsigned int) i) % (unsigned int) t->capacity);

        if (t->states[slot] == SLOT_EMPTY) {
            return 0;        /* the key is definitely not present */
        }
        if (t->states[slot] == SLOT_USED && strcmp(t->keys[slot], key) == 0) {
            *out = t->values[slot];
            return 1;
        }
    }

    return 0;
}
\end{lstlisting}

\begin{verbatim}
   Looking up "def" (hash %8 = 3):

     slot 3: USED "abc"  -- not a match, keep probing
     slot 4: USED "xyz"  -- not a match
     slot 5: USED "ghi"  -- not a match
     slot 6: USED "def"  -- match! return 99

   Looking up "qrs" (hash %8 = 3):

     slot 3..5: used, no match
     slot 6:    USED "def", no match
     slot 7:    EMPTY  -- stop. "qrs" is absent. We cannot skip slot 7;
     if we stopped at the first non-EMPTY we would report "absent" for
     keys that live further right. That is why tombstones exist.
\end{verbatim}

\begin{ruleofthumb}[Open addressing must stay loose]
Probing chains merge into clusters as the table fills, and long
clusters make every lookup slower. Keep the load factor below about
0.7, grow by reinserting all entries into a doubled array, and clear
tombstones on resize so they stop accumulating. Clustering is the
reason variants such as quadratic probing and double hashing exist:
they spread probes so clusters do not form.
\end{ruleofthumb}

\begin{verbatim}
   chaining vs open addressing:

   concern            separate chaining        open addressing
   ----------------   --------------------     ----------------------
   memory             nodes + bucket array     one flat array
   allocation         one per insert           none during insert
   locality           poor (chains scatter)    excellent (same array)
   deletion           free the node            tombstone + cleanup
   safe load factor   up to ~1.0               must stay < 0.7
   worst case         long chain               cluster / full table
   best for           generic, unknown keys    hot, cache-bound tables
\end{verbatim}

\begin{researchnote}[Collision handling has been studied since 1956]
The choice between collision strategies---our chaining versus our
linear probing---is not folklore. It was studied quantitatively in one
of the very first papers on hash tables, by E.~J.~Isaac and
R.~C.~Singleton, published in 1956, and the core trade-off they
identified is still the one you just saw in code.

\begin{verbatim}
   What the 1956 study established (paraphrased, not quoted):

       "The number of probes required to find an item depends on how
        the table is filled and on the collision technique used;
        comparing the techniques requires computing this cost."

   In our terms: chaining keeps clusters short by making the bucket a
   chain (cost grows with load factor), while open addressing pays a
   probe penalty that worsens as clusters form. That distinction is the
   load-factor table above, and it has not changed in seventy years.
\end{verbatim}

Why cite a paper from 1956? Because it shows that the design questions
you are making choices about today---probe sequences, load factors,
memory vs.\ locality---are the same questions the field has measured
since its beginning. The details have evolved (double hashing, Robin
Hood hashing), but the fundamental trade-off is a research result, not
a guess.

\paperref{E.~J.~Isaac, R.~C.~Singleton, ``An Analysis of Alternative
Collision Detection Techniques,'' Communications of the ACM, 1956}{https://dl.acm.org/doi/10.1145/320604.320609}
\end{researchnote}

\section{Trees}

\subsection{The General (N-ary) Tree}

A tree is a set of nodes, each with a value and a set of \emph{children}
that are themselves trees. When a node may have any number of children,
the simplest encoding carries two pointers: one to the first child, one
to the next sibling. The child list is a linked list, so the tree is a
list of lists of trees.

\begin{lstlisting}[caption={The N-ary node: first child, next sibling.}]
typedef struct tree_node {
    int value;
    struct tree_node *first_child;   /* head of the children list */
    struct tree_node *next_sibling;  /* next sibling in that list */
} tree_node;
\end{lstlisting}

\begin{verbatim}
   The tree:

                    (1)
                   / | \
                 (2)(3)(4)
                 / \    |
               (5) (6) (7)

   Encoded with first_child / next_sibling:

   (1) -- first_child --> (2) -- next_sibling --> (3) -- next_sibling --> (4)
                              |                      |                    |
                              v                      v                    v
                      +---------+               (NULL)                  (7)
                      | (5)     |
                      | next_sib --+            (3) has no children:
                      +------------+            first_child = NULL
                                   v
                              +---------+
                              | (6)     |
                              | next_sib --> NULL
                              +---------+

   Reading the picture: every node's children hang under it in a
   singly linked chain; the chain's links are "next_sibling".
\end{verbatim}

A preorder traversal visits a node, then each of its children in turn;
because the children form a list, the traversal is a loop over siblings
recursing into each child.

\begin{lstlisting}[caption={Preorder traversal of an N-ary tree.}]
void tree_print_preorder(const tree_node *root) {
    const tree_node *child;

    if (root == NULL) {
        return;
    }

    printf("%d ", root->value);

    child = root->first_child;
    while (child != NULL) {
        tree_print_preorder(child);
        child = child->next_sibling;
    }
}
\end{lstlisting}

This encoding is how compilers hold syntax trees (an operator with an
arbitrary number of operands), how file systems hold directory
hierarchies, and how HTML/XML documents are represented in memory.
Its cost is that a node needs a \emph{list} to reach its children; its
gift is that children of one node never share storage with anyone.

\subsection{Binary Trees and Binary Search Trees}

A binary tree is the special case where every node has at most two
children, conventionally named \code{left} and \code{right}. It is the
most studied shape in computing, because a \emph{binary search tree}
(BST) with the ordering property---left subtree holds smaller values,
right subtree holds larger---gives \code{O(h)} search, where
\code{h} is the tree height.

\begin{lstlisting}[caption={The BST node.}]
typedef struct bst_node {
    int value;
    struct bst_node *left;
    struct bst_node *right;
} bst_node;
\end{lstlisting}

\begin{verbatim}
   Inserting 3,1,4,2,5 in that order:

                3
               / \
              1   4
               \   \
                2   5

   The ordering property, stated at every node:
     everything in the left subtree  is smaller than the node;
     everything in the right subtree is larger than the node.
   Because of it, search is a descent: at each node, go left, go right,
   or stop. The value 2 is found in three steps, not five.
\end{verbatim}

Insertion reuses the pointer-to-pointer technique from the linked
lists: walk a \code{bst\_node **} down until it names an empty link,
then hang the new leaf there.

\begin{lstlisting}[caption={Insert: descend to an empty link and hang a leaf.}]
int bst_insert(bst_node **root, int value) {
    bst_node *new_node;
    bst_node **cur;

    cur = root;
    while (*cur != NULL) {
        if (value < (*cur)->value) {
            cur = &(*cur)->left;
        } else if (value > (*cur)->value) {
            cur = &(*cur)->right;
        } else {
            return 1;       /* already present: nothing to do */
        }
    }

    new_node = malloc(sizeof(bst_node));
    if (new_node == NULL) {
        return 0;
    }
    new_node->value = value;
    new_node->left = NULL;
    new_node->right = NULL;

    *cur = new_node;        /* the empty link now points at the leaf */
    return 1;
}
\end{lstlisting}

\begin{verbatim}
   Inserting 2 into the tree above. The descent:

     cur = &root        (names the 3-node)
     2 < 3  -> cur = &(3.left)     (names the 1-node)
     2 > 1  -> cur = &(1.right)    (names NULL -- an empty link)

   *cur = new_node hangs the leaf; no parent variable, no special case.

                3
               / \
              1   4
             / \   \
            ... 2   5
\end{verbatim}

Search is the same descent, returning true when the value is found.

\begin{lstlisting}[caption={Search: the descent in miniature.}]
int bst_search(const bst_node *root, int value) {
    const bst_node *cur;

    cur = root;
    while (cur != NULL) {
        if (value < cur->value) {
            cur = cur->left;
        } else if (value > cur->value) {
            cur = cur->right;
        } else {
            return 1;
        }
    }
    return 0;
}
\end{lstlisting}

A BST walked in-order---left, node, right---visits every value in
sorted order, which is exactly what a sorted ``print the whole
collection'' needs.

\begin{lstlisting}[caption={In-order traversal yields sorted output.}]
void bst_print_inorder(const bst_node *root) {
    if (root == NULL) {
        return;
    }
    bst_print_inorder(root->left);
    printf("%d ", root->value);
    bst_print_inorder(root->right);
}
\end{lstlisting}

Deletion is the one operation with three cases, and it is where BSTs
earn their reputation. The three cases are a leaf, a node with one
child, and a node with two children.

\begin{lstlisting}[caption={Remove, handling the three structural cases.}]
/* Remove value. Returns 1 if it was present. */
int bst_remove(bst_node **root, int value) {
    bst_node *victim;
    bst_node *successor;

    if (*root == NULL) {
        return 0;
    }

    if (value < (*root)->value) {
        return bst_remove(&(*root)->left, value);
    }
    if (value > (*root)->value) {
        return bst_remove(&(*root)->right, value);
    }

    /* Found it: *root is the node to remove. */
    victim = *root;

    if (victim->left == NULL && victim->right == NULL) {
        /* Case 1: a leaf. Drop it. */
        free(victim);
        *root = NULL;
        return 1;
    }

    if (victim->left == NULL) {
        /* Case 2: only a right child; it replaces the node. */
        *root = victim->right;
        free(victim);
        return 1;
    }
    if (victim->right == NULL) {
        /* Case 2: only a left child; it replaces the node. */
        *root = victim->left;
        free(victim);
        return 1;
    }

    /* Case 3: two children. Steal the value of the in-order successor
       (the leftmost node of the right subtree), then delete that
       successor node, which has at most one child.               */
    successor = victim->right;
    while (successor->left != NULL) {
        successor = successor->left;
    }
    victim->value = successor->value;
    return bst_remove(&victim->right, successor->value);
}
\end{lstlisting}

\begin{verbatim}
   Case 1 (leaf):   delete 9               6
                                             / \
        6            6                    3   9      ->  3's right = empty
       / \          / \                  / \   \
      3   9   ==>  3   (empty)          1   5   (gone)
     / \          / \
    1   5        1   5

   Case 2 (one child):   delete 3         6
                                             / \
        6            6                    1   9     ->  3's right promoted
       / \          / \                  / \   \
      3   9   ==>  1   9                1   5   7
     / \   \      / \   \
    1   5   7    1   5   7

   Case 3 (two children):   delete 6
        6                   7
       / \                 / \
      3   9       ==>     3   9
     / \   \             / \   \
    1   5   7           1   5   (the 7-node is gone)
\end{verbatim}

\begin{insight}[The height is the whole story]
A BST is only fast if it stays shallow. Inserting already-sorted
values yields a right-only chain of height \code{n}, and every
operation becomes \code{O(n)}---a linked list wearing a tree costume.
Balanced variants (AVL, red--black) guarantee \code{O(log n)} by
rotating on insertion. We do not implement them here because the
B-trees below give the same guarantee with a structure the disk and
cache actually like.
\end{insight}

\subsection{B-Trees}

A B-tree is the tree that databases and file systems actually use. It
is a balanced search tree with \emph{wide} nodes: each node holds many
keys and many children, so a tree of a million keys is only a few
levels tall, and each visit to a node can read an entire block in one
I/O. The balance is absolute---every leaf sits at the same depth---and
it is maintained by splitting nodes that overflow, not by rotations.

\begin{lstlisting}[caption={A B-tree node of order 4 (max 3 keys, 4 children).}]
#define BTREE_ORDER 4

typedef struct btree_node {
    int nkeys;                     /* keys currently stored */
    int keys[BTREE_ORDER - 1];     /* up to 3 keys          */
    struct btree_node *child[BTREE_ORDER];  /* up to 4 children */
    int leaf;                      /* 1 if this node has no children */
} btree_node;
\end{lstlisting}

\begin{verbatim}
   A B-tree of order 4 ("2-3-4 tree").  Keys in [ ] , children below.

                    [  20  |  40  ]
                   /       |       \
          [ 5 | 10 ]  [ 25 | 30 | 35 ]   [ 60 | 70 ]
          /    |    \    ...              ...
       (1,3)  (7,8) (11,15)

   Invariants (order m = 4):
     - a node with k keys has k+1 children;
     - every node (except the root) holds between m/2 - 1 and m - 1 keys;
     - all leaves sit at the same depth;
     - keys inside a node are ordered, and children bracket them:
         child[i]   holds keys  <  keys[i]
         child[i+1] holds keys  >  keys[i]

   Search is one node visit per level, and a node visit is a linear
   scan of a handful of keys. With wide nodes, height stays tiny:
   order 200 and a 2-level tree already covers 40,000 keys.
\end{verbatim}

Search first, because insert descends the same path.

\begin{lstlisting}[caption={Search a B-tree, one node visit per level.}]
/* Returns 1 if key is present. */
int btree_search(const btree_node *root, int key) {
    const btree_node *cur;
    size_t i;

    cur = root;
    while (cur != NULL) {
        /* Find the first key >= key in this node. */
        i = 0;
        while (i < (size_t) cur->nkeys && key > cur->keys[i]) {
            i = i + 1;
        }

        if (i < (size_t) cur->nkeys && key == cur->keys[i]) {
            return 1;              /* found in this node */
        }
        if (cur->leaf) {
            return 0;              /* reached a leaf without finding it */
        }
        cur = cur->child[i];       /* descend into the correct child */
    }
    return 0;
}
\end{lstlisting}

Insertion never grows a node beyond \code{m - 1} keys, because we
\emph{split} any node that is full on the way down. The middle key
moves up into the parent, and the two halves become children of the
parent. This is the single idea behind all B-tree variants.

\begin{lstlisting}[caption={Split a full child of cur at child index i.}]
/* Precondition: cur->child[i] is full (holds BTREE_ORDER-1 keys). */
void btree_split_child(btree_node *cur, size_t i) {
    btree_node *child;
    btree_node *right;
    size_t mid;
    size_t j;

    child = cur->child[i];
    right = malloc(sizeof(btree_node));
    if (right == NULL) {
        exit(1);
    }

    mid = (BTREE_ORDER - 1) / 2;
    right->leaf = child->leaf;
    right->nkeys = child->nkeys - (int) mid - 1;

    /* Move the upper half of child's keys into right. */
    for (j = 0; j < (size_t) right->nkeys; j++) {
        right->keys[j] = child->keys[mid + 1 + j];
    }
    /* Move the upper half of child's children into right. */
    if (child->leaf == 0) {
        for (j = 0; j <= (size_t) right->nkeys; j++) {
            right->child[j] = child->child[mid + 1 + j];
        }
    }

    /* Make room in the parent, promote the middle key, link right in. */
    for (j = cur->nkeys; j > i; j--) {
        cur->keys[j] = cur->keys[j - 1];
        cur->child[j + 1] = cur->child[j];
    }
    cur->keys[i] = child->keys[mid];
    cur->child[i + 1] = right;

    child->nkeys = (int) mid;
    cur->nkeys = cur->nkeys + 1;
}
\end{lstlisting}

\begin{verbatim}
   Splitting a full child during insert.

   BEFORE (child at index 1 is full):          AFTER split:
        [ 20 ]                                     [ 20 | 30 ]
        /      \                                 /    |      \
   [5|10]  [25|30|35]                     [5|10]  [25]  [35]
   (full: 3 keys)

   - the middle key 30 is PROMOTED into the parent;
   - the child splits into two nodes of 1 key each (for order 4,
     a node is at minimum half full after a split);
   - the parent grew by one key, and may itself need a split next.
\end{verbatim}

Insert descends from the root, splitting every full node it meets, so
that when the new key reaches a leaf there is always room for it. The
root is the only node that may ever be split alone, and splitting the
root is the only way a B-tree gains a level---balance is never
temporarily broken.

\begin{verbatim}
   The growth of a B-tree is the height growing UPWARD:

   [1|2|3]  insert 4  ->  [1]  [2]  [3|4]   (2 moves up, root splits)
                              |
                          [ 2 ]
                         /     \
                      [1]     [3|4]   ->   insert more... height grows at
                                          the root, never in the middle.
\end{verbatim}

\begin{ruleofthumb}[B-trees trade memory for I/O]
Each node is sized to a block (a disk sector or a cache line), so one
node read fetches many keys at once. Search becomes ``read a block,
scan a few keys, repeat,'' with a path length that is the logarithm of
the key count in base \emph{order}. That is why database indexes and
filesystem metadata use B-trees and not BSTs.
\end{ruleofthumb}

\subsection{B+ Trees}

The B+ tree takes the B-tree and separates two jobs: \emph{internal
nodes hold only keys, for routing}, and \emph{all data lives in the
leaves}, which are linked together in sorted order. A lookup or range
scan descends once and then walks the leaf chain, which is sequential
reads---exactly what rotating disks and SSDs love.

\begin{lstlisting}[caption={The B+ leaf adds a sibling pointer.}]
typedef struct bplus_node {
    int leaf;                 /* 1 for a leaf, 0 for an internal node */
    int nkeys;
    int keys[BTREE_ORDER - 1];
    struct bplus_node *child[BTREE_ORDER];
    struct bplus_node *next;  /* used ONLY by leaves: the sibling link */
} bplus_node;
\end{lstlisting}

\begin{verbatim}
   A B+ tree. Internal nodes hold routing keys only; the data rows
   live in the leaves, and the leaves are a sorted linked list.

            [ 20 | 40 ]              <- internal: pure routing
           /     |      \
    [ 5|10 ]  [ 25|30 ]  [ 60|70 ]   <- internal: still routing
       |          |          |
       v          v          v
   +--------------------------------------------------------------+
   | leaf (1,3,5,7) -> leaf (20,25,30) -> leaf (40,60,70) -> NULL |
   +--------------------------------------------------------------+

   "SELECT * FROM t WHERE key BETWEEN 20 AND 60":
     1. descend to the leaf holding 20;
     2. walk the leaf chain, reading leaves sequentially.
     No backtracking, no internal-node data: a straight line of reads.
\end{verbatim}

\begin{verbatim}
   B vs B+ at a glance:

                 B-tree               B+ tree
   data          in every node        only in leaves
   internal nodes        keys + data + children   keys + children
   range scan    scatter, no chain    leaf chain = sorted sequence
   space reuse   duplicate keys in internal nodes   keys duplicated only
                               between levels (routing copies)
   best for      general purpose     databases, ordered range queries
\end{verbatim}

\subsection{B* Trees}

The B* tree tightens the B-tree's occupancy rule: nodes are kept at
least \emph{two-thirds} full instead of half-full, by redistributing
keys with a sibling before splitting, and by splitting one full node
into \emph{three} when both it and its sibling are full. The result is
fewer, denser nodes---fewer I/Os for the same data---at the cost of a
more complex insertion.

\begin{verbatim}
   B* tree: before splitting, borrow from a neighbor.

   order 5, node X full, sibling Y half-full:

        BEFORE                              AFTER (redistribute)
        [ 15 ]                              [ 17 ]
       /      \                            /       \
   [1|5|10|14] [17|20]              [1|5|10|14|15]   [17|20]
   (full)      (room)

   Only when BOTH are full does a B* split produce THREE nodes,
   each about 2/3 full, pushing two keys into the parent.

   cost/benefit:
     B-tree   min 1/2 full   simpler
     B* tree  min 2/3 full   33% fewer nodes, more complex inserts
\end{verbatim}

\begin{ruleofthumb}[When to use which tree]
Use a BST for in-memory data where simplicity wins and the data is
unpredictable. Use a B-tree when nodes map to blocks (disk, flash,
caches). Use a B+ tree when ordered range scans matter, as in nearly
every real database index. Use a B* tree when storage density
dominates and the extra insertion complexity is justified---they are
rare outside research and niche storage engines.
\end{ruleofthumb}

\begin{researchnote}[The B-tree family in the research record]
The B-tree has one of the longest and best-documented research
histories in computer science, and two papers in particular define
what you learned in this section. The first is the paper that
\emph{invented} the B-tree (Bayer and McCreight, 1972). The second is
the survey that catalogued its influence on databases (Comer, 1979).

\begin{verbatim}
   Extracts (condensed for the book):

       Bayer & McCreight (1972): the B-tree was designed so that
       searching a large ordered file requires only a few page reads,
       by keeping the tree balanced and each node the size of one
       storage page.

       Comer (1979), on the family's reach: the B-tree has become the
       standard method for organizing indexes in file systems and
       databases, and variants -- B+, B*, prefix B-trees -- were
       introduced to address specific needs.
\end{verbatim}

These extracts are paraphrases, not exact sentences; what matters is
the claim they establish: the wide, shallow, block-sized B-tree is not
a niche idea but \emph{the} industry-standard index structure, and the
B+ variant we implemented exists because ordered range scans are what
real queries need. When you meet \code{btree} indexes in every
production database, you are looking at a research lineage from 1972.

\paperref{D.~Comer, ``The Ubiquitous B-Tree,'' ACM Computing Surveys
11(2), 1979}{https://infolab.usc.edu/csci585/Spring2010/den_ar/indexing.pdf}

\paperref{R.~Bayer, E.~McCreight, ``Organization and Maintenance of
Large Ordered Indexes,'' Acta Informatica 1(3), 1972}{https://dl.acm.org/doi/10.1007/BF00288683}
\end{researchnote}

\section{Graphs}

\subsection{Graphs as Structs and Arrows}

A graph is a set of \emph{vertices} connected by \emph{edges}. Where a
tree is a graph with exactly one path between any two nodes and a
designated root, a graph drops every such restriction: cycles, multiple
edges, and unreachable regions are all allowed. Graphs model networks of
every kind---the internet, social graphs, road maps, dependency graphs,
circuit nets---and they are the most general structure in this chapter.

A graph is \emph{directed} if its edges have a direction (a one-way
street), \emph{undirected} if they do not (a two-way street), and
\emph{weighted} if each edge carries a cost. These are three
independent properties; all eight combinations exist.

\begin{verbatim}
   Undirected graph:                      Directed graph:

         0 -- 1                              0 --> 1
         |    |                              ^     |
         |    |                              |     v
         2 -- 3                              2 <-- 3

   edges are unordered pairs           edges are ordered pairs
   0-1 == 1-0                          0->1 is NOT 1->0

   Weighted (undirected):

            5
       0 -------- 1
       |          |
       3          2
       |          |
       4 -------- 3
            1

   each edge carries a number:         DAG (directed acyclic graph):
   the cost of traveling it.
                                          A --> B --> D
                                          |     ^
                                          v     |
                                          C --> E
                                          \
                                           v
                                           F
                                          (no cycle exists; topo-sortable)
\end{verbatim}

Two representations dominate, and they are the same trade-off we have
seen all chapter: pointer-honest locality versus constant-time
directness.

\subsection{Adjacency List}

Each vertex owns a linked list of its edges. The list is exactly the
linked list of this chapter; the graph is an array of such lists.

\begin{lstlisting}[caption={Adjacency list representation.}]
typedef struct graph_edge {
    int to;                  /* the vertex at the other end */
    int weight;              /* cost of the edge            */
    struct graph_edge *next; /* next edge leaving this vertex */
} graph_edge;

typedef struct graph {
    int n;                   /* number of vertices 0..n-1  */
    graph_edge **adj;        /* adj[v] = head of v's edge list */
} graph;
\end{lstlisting}

\begin{verbatim}
   The directed graph  0 -> 1 -> 2
                        ^    |
                        |    v
                        3 <- 4

   as an adjacency list (each vertex's outgoing edges):

     adj[0] -> (to=1) -----------> NULL
     adj[1] -> (to=2) -----------> NULL
     adj[2] -> NULL
     adj[3] -> (to=0) -----------> NULL
     adj[4] -> (to=1) -> (to=3) -> NULL

   Memory is O(V + E): one node per edge, one head per vertex.
   Asking "who are the neighbors of v?" is walking v's list.
   Asking "is u adjacent to v?" is walking u's list (or v's).
\end{verbatim}

Adding an edge pushes onto a vertex's list. For an undirected graph,
every edge must be added twice---once in each direction---and a
weighted edge stores its cost in the node.

\begin{lstlisting}[caption={Add a directed edge, and its undirected twin.}]
/* Add edge u -> v with the given weight. Returns 1 on success. */
int graph_add_edge(graph *g, int u, int v, int weight) {
    graph_edge *e;

    e = malloc(sizeof(graph_edge));
    if (e == NULL) {
        return 0;
    }
    e->to = v;
    e->weight = weight;

    e->next = g->adj[u];
    g->adj[u] = e;
    return 1;
}

/* Add an undirected edge: it is two directed edges. */
int graph_add_undirected(graph *g, int u, int v, int weight) {
    int r1;
    int r2;

    r1 = graph_add_edge(g, u, v, weight);
    r2 = graph_add_edge(g, v, u, weight);
    if (r1 == 0 || r2 == 0) {
        return 0;
    }
    return 1;
}
\end{lstlisting}

\begin{verbatim}
   Memory for one undirected edge u-v:

     adj[u] -> (to=v, w=w) -> ...
     adj[v] -> (to=u, w=w) -> ...

   Two directed edges model one undirected edge. The weight is stored
   on both sides, so both directions cost the same--as an undirected
   edge must.
\end{verbatim}

\subsection{Adjacency Matrix}

The alternative is a 2D matrix where entry \code{m[u][v]} is 1 (or the
weight) when the edge \code{u -> v} exists, and 0 otherwise. Any graph
problem that asks ``is this edge present?'' or needs all-pairs
information becomes a constant-time lookup.

\begin{lstlisting}[caption={Adjacency matrix representation.}]
typedef struct matrix_graph {
    int n;
    int **m;   /* m[u][v] = weight, or 0 when no edge u->v exists */
} matrix_graph;
\end{lstlisting}

\begin{verbatim}
   Directed graph  0 -> 1 -> 2
                    ^    |
                    |    v
                    3 <- 4

   as a matrix:

         0  1  2  3  4
       +----------------+
     0 | 0  1  0  0  0 |     m[0][1] = 1 : the edge 0->1 exists
     1 | 0  0  1  0  0 |     m[4][3] = 1 : the edge 4->3 exists
     2 | 0  0  0  0  0 |     every diagonal 0 : no self-loops
     3 | 1  0  0  0  0 |
     4 | 0  1  0  1  0 |
       +----------------+

   Memory is O(V^2) even for a sparse graph; each edge check is O(1).
   For a WEIGHTED matrix, a 0 is "no edge" and the value is the cost.
\end{verbatim}

\begin{verbatim}
   adjacency list vs adjacency matrix:

   concern              adjacency list      adjacency matrix
   ----------------     -----------------   ------------------
   memory               O(V + E)            O(V^2)
   edge check (u,v)?    O(degree) walk      O(1)
   iterate neighbors    O(degree), cached   O(V), must scan a row
   add edge             O(1)                O(1)
   best for             sparse, real nets   dense, small, or
                                              all-pairs algorithms
\end{verbatim}

\subsection{DAGs and Topological Order}

A DAG---directed acyclic graph---is a graph with no cycles. DAGs are
the shape of every dependency system: build systems, package managers,
compilation order, and spreadsheet cell evaluation all order a set of
jobs so that every job runs after its prerequisites. A \emph{topological
order} is an ordering of the vertices in which every edge points
forward. A graph admits one if and only if it is a DAG.

\begin{verbatim}
   A build system as a DAG (an arrow u -> v means "u must build
   before v"):

        lex -----------+
                       v
        yacc -----> make ---> main.o ---> myprog
                   ^                            ^
                   |                            |
                   +----------------------------+   (yacc also feeds myprog)

   Topological order must list each target after all of its inputs.
   Here: lex, yacc, make, main.o, myprog is valid; yacc before make
   is valid; make before lex is NOT (lex is a dependency of make).

   If a dependency graph has a cycle, no order exists: A needs B,
   B needs A, someone has to come first and cannot.
\end{verbatim}

Kahn's algorithm produces a topological order without recursion: count
every vertex's \emph{indegree} (incoming edges), then repeatedly emit a
vertex with zero indegree and delete it, decrementing the indegrees of
its successors. If we run out of zero-indegree vertices before emitting
everything, a cycle exists.

\begin{lstlisting}[caption={Topological sort by Kahn's algorithm.}]
#include <stdlib.h>

/* Write a topological order of g into order[]. Returns 1 on success,
   0 if g has a cycle. */
int graph_toposort(const graph *g, int *order) {
    int *indegree;
    int *queue;
    int head;
    int tail;
    int v;
    int emitted;
    const graph_edge *e;
    size_t i;

    indegree = calloc((size_t) g->n, sizeof(int));
    queue = malloc(sizeof(int) * (size_t) g->n);
    if (indegree == NULL || queue == NULL) {
        free(indegree);
        free(queue);
        return 0;
    }

    /* Pass 1: count incoming edges for every vertex. */
    for (i = 0; i < (size_t) g->n; i++) {
        e = g->adj[i];
        while (e != NULL) {
            indegree[e->to] = indegree[e->to] + 1;
            e = e->next;
        }
    }

    /* Seed the queue with the vertices that depend on nothing. */
    head = 0;
    tail = 0;
    for (i = 0; i < (size_t) g->n; i++) {
        if (indegree[i] == 0) {
            queue[tail] = (int) i;
            tail = tail + 1;
        }
    }

    /* Pass 2: emit zero-indegree vertices, one at a time. */
    emitted = 0;
    while (head < tail) {
        v = queue[head];
        head = head + 1;
        order[emitted] = v;
        emitted = emitted + 1;

        /* Removing v lowers the indegree of each successor. */
        e = g->adj[v];
        while (e != NULL) {
            indegree[e->to] = indegree[e->to] - 1;
            if (indegree[e->to] == 0) {
                queue[tail] = e->to;
                tail = tail + 1;
            }
            e = e->next;
        }
    }

    free(indegree);
    free(queue);

    if (emitted != g->n) {
        return 0;   /* some vertices never reached zero indegree: cycle */
    }
    return 1;
}
\end{lstlisting}

\begin{verbatim}
   Kahn's algorithm on the DAG above, step by step:

   indegrees:  lex=0 yacc=0 make=2 main.o=1 myprog=2
   (lex and yacc depend on nothing, so they are the "sources")

   queue: [lex, yacc]
   emit lex    ->  make indegree 2-1 = 1
   emit yacc   ->  make 1-1 = 0, enqueue make; myprog 2-1 = 1
   emit make   ->  main.o 1-1 = 0, enqueue main.o
   emit main.o ->  myprog 1-1 = 0, enqueue myprog
   emit myprog

   order: lex, yacc, make, main.o, myprog    -- every edge points forward.
\end{verbatim}

\subsection{Traversal: DFS and BFS}

Graphs are explored by two walks, and every serious graph algorithm
rides on one of them. \emph{Depth-first search} (DFS) dives down a
path as far as it goes before backtracking; \emph{breadth-first search}
(BFS) explores in expanding waves. Both mark visited vertices so they
never revisit, and both visit each vertex and edge once:
\code{O(V + E)}.

\begin{lstlisting}[caption={DFS, written with an explicit recursion.}]
/* Print the graph's vertices reachable from start, depth-first. */
void graph_dfs(const graph *g, int start, int *visited) {
    const graph_edge *e;

    visited[start] = 1;
    printf("%d ", start);

    e = g->adj[start];
    while (e != NULL) {
        if (visited[e->to] == 0) {
            graph_dfs(g, e->to, visited);
        }
        e = e->next;
    }
}
\end{lstlisting}

\begin{verbatim}
   DFS order from 0 on:                graph:
   visit 0, then 0's first neighbor 1,      0 - 1
   then 1's first unvisited neighbor 3,      | / |
   then 3's neighbor 2, done; backtrack.    2 - 3
                                           (undirected)

   order: 0, 1, 3, 2
   DFS goes as deep as it CAN before trying another neighbor of 0.
\end{verbatim}

BFS needs a queue; we reuse the array queue pattern from Kahn's
algorithm.

\begin{lstlisting}[caption={BFS, exploring in waves with a queue.}]
#include <stdlib.h>

/* Print the vertices reachable from start, breadth-first. */
void graph_bfs(const graph *g, int start, int *visited) {
    int *queue;
    int head;
    int tail;
    int v;
    const graph_edge *e;

    queue = malloc(sizeof(int) * (size_t) g->n);
    if (queue == NULL) {
        exit(1);
    }
    head = 0;
    tail = 0;

    visited[start] = 1;
    queue[tail] = start;
    tail = tail + 1;

    while (head < tail) {
        v = queue[head];
        head = head + 1;
        printf("%d ", v);

        e = g->adj[v];
        while (e != NULL) {
            if (visited[e->to] == 0) {
                visited[e->to] = 1;
                queue[tail] = e->to;
                tail = tail + 1;
            }
            e = e->next;
        }
    }

    free(queue);
}
\end{lstlisting}

\begin{verbatim}
   BFS order from 0 on the same graph:

   queue: [0]
   pop 0, enqueue 1, 2          queue: [1, 2]
   pop 1, enqueue 3             queue: [2, 3]
   pop 2 (3 already visited)    queue: [3]
   pop 3                        order: 0, 1, 2, 3

   BFS visits in distance order: everything one hop away, then two,
   then three. That is why BFS finds SHORTEST paths in an unweighted
   graph, while DFS is the natural tool for reachability and for
   detecting cycles.
\end{verbatim}

\begin{verbatim}
   DFS vs BFS at a glance:

                DFS                         BFS
   tool         recursion (or explicit      queue
                stack)
   order        deep before wide            distance order
   best for     reachability, cycles,       shortest paths, routing,
                topological, mazes          level order, social hops
   memory       O(depth)                    O(width)
   both run in  O(V + E), each vertex/edge visited once
\end{verbatim}

\section{Structures in Bytes}

Every container in this chapter is a struct, so every container has a
byte size. On a 64-bit machine (\code{int}=4, pointer=8):

\begin{verbatim}
   structure                node bytes   layout (offsets)
   ----------------------   ----------   ---------------------------
   singly linked list       16           value:0  next:8
   doubly linked list       24           value:0  prev:8  next:16
   circular list            16           value:0  next:8
   intrusive list_head      16           next:0   prev:8  (embedded)
   unrolled node            4 + 4B + 8   count:0 data:4 next:4B+8
   XOR list                 16           value:0  link:8
   skip list node           variable     value:0 level:4 forward[] at 8
   hash-table node          24           key:0 value:8 next:16
   B-tree node (order 4)    56           nkeys:0 keys:4 child:16 leaf:52
   graph edge               16           to:0  weight:4  next:8
\end{verbatim}

Two rows are worth reading twice. The B-tree node rounds up to 56
bytes because the \code{child[4]} pointer array must start at a
multiple of 8 (offset 16) and the trailing \code{leaf} int plus
padding fill the last 8. The skip-list node is variable because its
\code{forward[]} is a flexible array member: its true size is
\code{sizeof(skiplist\_node) + level * sizeof(pointer)}. These numbers
are what decide how many nodes fit in one 64-byte cache line---the
locality argument every data structure in this chapter is built on.

\section{Summary of the Menagerie}

Every structure in this chapter is the same two tools in a different
uniform: structs for shape, pointers for arrows, and a rule for who
owns each allocation. The cost table below is the one page you should
be able to redraw from memory.

\begin{verbatim}
   structure           search       insert     delete     notes
   ------------------  -----------  ---------  ---------  --------------------
   singly linked       O(n)         O(1)+      O(n)       1 ptr/node
   doubly linked       O(n)         O(1)+      O(1)+      2 ptrs/node
   circular            O(n)         O(1)+      O(1)+      ring, tail ptr
   intrusive           O(n)         O(1)+      O(1)+      links embedded
   unrolled            O(n)         O(n)       O(n)       blocks, cache
   XOR                 O(n)         O(1)+      O(1)+      1 ptr, needs prev
   skip list           O(log n)*    O(log n)*  O(log n)*  towers, sorted
   hash (chaining)     O(1)*        O(1)*      O(1)*      buckets+chains
   hash (open addr)    O(1)*        O(1)*      O(1)*      probing
   BST                 O(h)         O(h)       O(h)       h worst-case = n
   B-tree / B+/B*      O(log n)     O(log n)   O(log n)   wide, shallow
   graph adj list      O(deg)       O(1)       O(1)       O(V+E) memory
   graph adj matrix    O(1)         O(1)       O(1)       O(V^2) memory

   +  given a position/neighbor already found
   *  expected / probabilistic
\end{verbatim}

The four questions that opened the chapter are the checklist you apply
to every one of these---and to every structure you invent:

\begin{verbatim}
   SHAPE   : draw the struct.  How many pointers?  What points where?
   COST    : what does search/insert/delete actually touch?
   OWNER   : who mallocs, who frees, and what happens on the error path?
   CACHE   : are the bytes the algorithm wants in the same cache line?
\end{verbatim}

Answer those four honestly and you have designed a data structure;
code is the last, easy part. With this toolbox in hand we turn from the
shape of data to the services of the operating system: files, syscalls,
processes, signals, and the machinery every real C program sits on.





}{}
\IfFileExists{chapter5.tex}{% ============================================================================
%  chapter5.tex
%  System Calls, File I/O, and Process Creation
% ============================================================================

\chapter{System Calls, File I/O, and Process Creation}

\section{The Operating System Is a Service Provider}

Everything before this chapter happened entirely inside a single
process: bytes in memory, structs on the heap, arrows between them.
That world is a comfortable one, but it is also a closed one. Real
software must reach \emph{outside} itself---to read a file, to spawn a
worker, to send a packet, to ask what time it is. Every one of those
acts crosses the boundary between your program and the operating
system.

An operating system is, at its core, a \emph{service provider with a
monopoly on the hardware}. Three jobs define it:

\begin{description}
  \item[Multiplexing] Many programs run at once on one CPU, one RAM,
    and a handful of disks. The OS interleaves their use of every
    device so that each program believes it has the machine to itself.
  \item[Protection] A program that misbehaves must not corrupt another
    program, the OS, or the hardware. The OS enforces this with
    hardware-enforced privilege levels.
  \item[Abstraction] The OS hides the terrifying variety of hardware---a
    dozen disk controllers, a hundred network cards---behind a small
    number of clean interfaces: files, processes, sockets, clocks.
\end{description}

\begin{verbatim}
   Without the OS, each program would have to:           With the OS:

   - drive the disk controller itself,                    read(fd, buf, n);
   - know the exact geometry of every disk,
   - police itself against corrupting other programs,
   - schedule its own CPU time,                          fork(); exec();
   - share the network card without colliding.           socket(); send();
\end{verbatim}

This chapter is the pivot of the book: we stop arranging memory inside
one process and start asking the OS for services. Every later chapter
is a deeper dive into one of these services. The chapter also teaches
the habit that makes systems programming survivable: \emph{every
OS call returns either success or an error, and a real program treats
both as normal control flow}.

\begin{osconcept}
An operating system exists precisely because hardware is shared,
dangerous, and diverse. Individual applications cannot be trusted to
multiplex the CPU or the disk among each other (they would starve and
corrupt one another), so resource management, protection, and uniform
device interfaces are problems that only a privileged, global piece of
software---the OS---can solve. That is why ``how do I read a file?''
is an \emph{OS} question and not a library question: the answer is a
contract with the kernel.
\end{osconcept}

\subsection{What the Kernel Is}

The word \emph{kernel} has been used casually here; let us pin it
down. The \emph{kernel} is the core of the operating system: the
privileged program that runs in kernel mode, owns the hardware, and
provides the system calls. When this book says ``the OS decides'' or
``the kernel enforces,'' it means this one program---the software
that controls the CPU, the memory, the disks, and the network card,
and that no other program may modify or even fully see.

\begin{verbatim}
   The kernel, drawn in the machine:

   +--------------------------------------------------------+
   |  THE KERNEL  -- one privileged program, in kernel mode  |
   |   owns: the CPU, memory pages, disks, network card      |
   |   provides: every system call you have used             |
   |   is: loaded at boot, resident for the machine's life   |
   +--------------------------------------------------------+
        ^                        ^
        | syscalls               | the hardware
        v                        v
   +--------------------------------------------------------+
   |  USER PROCESSES  (your programs, in user mode)          |
   |  many, untrusted, isolated from each other              |
   +--------------------------------------------------------+

   - "Kernel" is also the name of the actual binary: on Linux,
     /boot/vmlinuz (Chapter 15) IS the kernel, loaded at boot.
   - The kernel is the ONLY program allowed to touch the hardware.
     Everything your program does with the machine goes through it.
   - That is why every chapter's refrain--"only the kernel can do
     this"--is literally true: the CPU refuses to let anyone else.
\end{verbatim}

The kernel is not a distant concept; it is the direct addressee of
every system call in this book. \code{open}, \code{fork},
\code{mmap}, \code{poll}, and \code{connect} are all \emph{requests
to the kernel}. Everything since Chapter~5 has been a conversation
with one privileged program, and the rest of this book is learning
the terms of that conversation.

\section{User Mode and Kernel Mode}

To enforce protection, the CPU itself has at least two privilege
levels. In \emph{user mode}, your program runs. In \emph{kernel mode},
the OS runs. The difference is absolute: many instructions (turning off
interrupts, writing to device registers, remapping memory) simply
\emph{fault} if executed from user mode. The hardware forbids your
program from doing the dangerous things, so the OS does not have to
trust you.

\begin{verbatim}
   User mode (your program)              Kernel mode (the OS)

   +-----------------------------+       +--------------------------------+
   | main()                      |       | device drivers                |
   |   read(fd, buf, 4096);     |       | file systems (ext4, xfs)       |
   |   fork();                  |       | process scheduler             |
   |   write(1, msg, n);        |       | virtual memory manager        |
   |                            |       | open file table               |
   | your code can ONLY touch   |       | process table                 |
   | its OWN address space      |       |                              |
   +----------------------------+       | may touch ALL hardware        |
                |                       +--------------------------------+
                |  "please handle this for me"  (a system call)
                v
        +------------------+
        | privilege switch |
        | (CPU exception)  |
        +------------------+
\end{verbatim}

\begin{osconcept}
Protection cannot be a polite agreement between programs; it must be
\textbf{mechanically enforced}, and only the hardware can do that. By
having the CPU refuse dangerous instructions in user mode, the OS makes
misbehavior impossible-by-construction instead of just discouraged.
Deciding who gets kernel privilege, and keeping that boundary
leak-proof, is a core operating system problem.
\end{osconcept}

\section{System Calls: The Only Door}

The privileged operations are exposed to user programs through
\emph{system calls}: the tiny, fixed set of doors into the kernel.
Your C code rarely issues a system call directly; the C library
(\code{glibc}) provides wrappers that look like ordinary functions but
perform the privilege switch underneath.

\begin{verbatim}
   One system call, three layers:

   your code:            read(fd, buf, 4096);
        |
        v
   glibc wrapper:        place syscall number 0 and the arguments in
                         registers; execute the CPU instruction "syscall"
        |
        v
   kernel:               switch to kernel mode, look up syscall number 0,
                         validate the arguments, perform the operation,
                         store the result, switch back to user mode.

   A system call is NOT an ordinary function call. It changes privilege
   level, and the CPU does real work to make that safe. That is why
   system calls are comparatively expensive (hundreds of nanoseconds)
   and why programs do their best to do many file operations per call.
\end{verbatim}

\begin{verbatim}
   On x86-64 Linux, some system call numbers you are about to use:

   read      = 0        openat   = 257
   write     = 1        close    = 3
   exit      = 60       fork     = 57 (glibc wraps clone)
   lseek     = 8        execve   = 59
\end{verbatim}

Two properties of system calls matter enormously. First, they are
\emph{stable}: application code survives decades because the kernel
keeps syscall semantics constant. Second, they return
\emph{status}: almost every system call returns \code{-1} on failure
and records a reason in the variable \code{errno}. There is no
exception mechanism crossing the kernel boundary; there is a return
value, and you must read it.

The mechanical basis of the privilege switch is the \emph{interrupt}---
in this context also called a \emph{trap}. An interrupt is the CPU
abandoning its normal instruction stream to handle an exceptional
event: a key was pressed (a hardware interrupt), a page was missing
(Chapter~1), or a program executed the \code{syscall} instruction (a
software interrupt, the trap). Whatever the cause, the CPU saves its
state and jumps to a fixed handler \emph{in the kernel}. The word
``interrupt'' will return in Chapter~6 (signals interrupt your
program) and Chapter~10 (a system call interrupted by a signal): all
of them are the same CPU mechanism---``stop what you are doing, deal
with this, come back.''

\begin{osconcept}
The system call interface is the operating system's public API, and
designing it is an OS problem because it must be minimal (every entry
point is an attack surface and a maintenance contract), stable (users
depend on it for decades), and safe (the kernel must validate every
argument from an untrusted caller). Getting that boundary right is one
of the oldest and hardest problems in systems software.
\end{osconcept}

\section{Files: The Universal Abstraction}

The single most successful abstraction in computing is the \emph{file}:
a named sequence of bytes. Almost every device the OS owns is exposed
as a file---disks, terminals, pipes, even some clocks. The virtual file
system (VFS) lets the same \code{open}/\code{read}/\code{write} calls
work regardless of whether the bytes live on an ext4 disk, a network
share, or in another program's output.

\begin{verbatim}
   The file abstraction unifies everything:

     /dev/tty     -> a terminal            (read/write characters)
     /proc/cpuinfo-> data manufactured on the fly by the kernel
     fifo          -> a pipe between two programs
     socket        -> network data
     /home/you/f  -> bytes on a disk

   From your program's point of view each is "a sequence of bytes you
   can read or write." The OS sorts out where the bytes really live.
\end{verbatim}

\begin{osconcept}
Naming, layout, and storage of persistent data are an OS problem
because only the kernel can give programs a \emph{namespace} (the file
tree) that is shared, protected by permissions, and backed by arbitrary
physical devices. A library cannot create a file; it can only call the
syscall that does. Files are the OS's way of saying ``here is a uniform
interface to everything persistent or remote.''
\end{osconcept}

\section{File Descriptors: Your Handles on the Kernel}

When you \code{open} a file, the kernel creates an entry in its own
bookkeeping---tracking the file position, the access mode, the
permissions---and hands you back a small integer, the \emph{file
descriptor}. From then on you never name the path again; you name the
descriptor. The kernel keeps the path lookup and the caching behind
your back.

\begin{lstlisting}[caption={Opening a file for reading.}]
#include <stdio.h>
#include <stdlib.h>
#include <fcntl.h>
#include <unistd.h>

int main(void) {
    int fd;

    fd = open("hello.txt", O_RDONLY);
    if (fd < 0) {
        perror("open");
        exit(1);
    }

    printf("opened file, fd = %d\n", fd);

    close(fd);     /* release the kernel's entry */
    return 0;
}
\end{lstlisting}

\begin{verbatim}
   Your process                          The kernel
   +-------------------------+            +---------------------------+
   | file descriptor table   |            | open file table           |
   |  0  ->  stdin      ----+----------->|  position: 0              |
   |  1  ->  stdout     ----+----------->|  mode: O_RDONLY           |
   |  2  ->  stderr     ----+----------->|  -> inode (disk file)     |
   |  3  ->  hello.txt  ----+----------->|                           |
   |  4  ->  (free)         |            |                           |
   +-------------------------+            +---------------------------+

   - Descriptors 0, 1, 2 are ALWAYS present: stdin, stdout, stderr.
   - open() returns the LOWEST free descriptor (here 3).
   - close(3) removes the entry; the number is then reusable.
   - The path "hello.txt" is resolved ONCE at open(); everything after
     uses the number.
\end{verbatim}

The flags matter. \code{O\_RDONLY} opens read-only; \code{O\_WRONLY}
write-only; \code{O\_RDWR} both. \code{O\_CREAT} says ``create the file
if it does not exist'' and then requires a third argument, the mode
bits. \code{O\_TRUNC} empties an existing file first. \code{O\_APPEND}
forces every write to go to the end, which is how logs stay logs even
when several processes write to them.

\begin{osconcept}
A file descriptor is a \emph{capability}: a small token that says
``this process may access that kernel object.'' Managing the per-process
descriptor table, the shared open-file table, and the reference
counting that decides when a file is truly closed is pure OS work.
Your program sees an integer; the kernel sees a reference into
structures it alone can touch. That indirection is what allows
permissions to be checked once at open time and then enforced cheaply
for every subsequent call.
\end{osconcept}

\begin{researchnote}[The paper that invented all of this]
Every abstraction in this chapter---files as byte streams, descriptors,
pipes, processes---was designed in one place, in 1974, and written up
in the most influential systems paper ever published: ``The UNIX
Time-Sharing System,'' by Ken Thompson and Dennis Ritchie. If you read
one paper from this entire book, read this one.

\begin{verbatim}
   An extract (a real sentence from the paper):

       "The most important job of UNIX is to provide a file system.
        From the point of view of the user, there are three kinds of
        files: ordinary disk files, directories, and special files."

   What this emphasis on the file system means, in our chapter's terms:
   the authors chose to make nearly EVERYTHING a file -- ordinary data,
   directories, even devices -- so that the tools for reading and
   writing one thing work for all of them. That single decision is why
   our open/read/write pattern is so uniform, and why descriptors 0,
   1, and 2 exist at all.
\end{verbatim}

Thompson and Ritchie published it after years of internal use at Bell
Labs, and the paper is famous for being both short and astonishingly
clear. Reading it is a masterclass in explaining an operating system
in prose---exactly the style this book tries to imitate.

\paperref{K.~Thompson, D.~M.~Ritchie, ``The UNIX Time-Sharing System,''
Communications of the ACM 17(7), 1974}{https://www.cs.princeton.edu/courses/archive/spring14/cos217/reading/unix.pdf}
\end{researchnote}

\section{Reading and Writing}

\subsection{What Is a Buffer?}

You will meet the word \emph{buffer} in every chapter of this book, so
let us fix its meaning once and for all. A \emph{buffer} is a
contiguous region of memory used to hold data temporarily while it is
in transit between two places. The key word is \emph{temporarily}: a
buffer is a staging area, not a destination. Data flows in, sits
briefly, and flows out.

Three facts of real hardware and real kernels make buffers
unavoidable in systems programming:

\begin{itemize}
  \item Devices deliver data in \emph{chunks}, never one byte at a
    time on demand. A disk reads a whole sector; a network card
    delivers a whole packet. A buffer is where a chunk waits while
    your code catches up.
  \item Crossing the kernel boundary is expensive. Doing I/O in big
    batches through a buffer costs far less than many tiny transfers.
  \item Producers and consumers run at different speeds. A buffer is
    the shock absorber between them.
\end{itemize}

\begin{verbatim}
   The central picture: data in transit through a staging area.

      producer -----> [  B U F F E R  ] -----> consumer
                       (contiguous bytes
                        waiting their turn)

   In the copy_file() program below:
      read()  fills buf from the disk   (producer: the kernel)
      write() drains buf to the disk    (consumer: the kernel)
      your buffer sits between the two.
\end{verbatim}

In C a buffer is simply an array---or a block from \code{malloc}.
Our copy loop's \code{char buf[CHUNK]} is a buffer of 4096 bytes:
\code{read} deposits bytes into it, \code{write} drains them, and the
two operations alternate until the file is exhausted. Notice that the
buffer never grows; it is reused, which is the whole point of a
staging area.

\begin{verbatim}
   Where buffers appear throughout this book:

     file I/O      :  char buf[4096];  this chapter
     strings       :  strncpy targets, flexible array tails (ch.3)
     the kernel    :  the buffer cache -- disk blocks cached in RAM
     pipes         :  the kernel's pipe buffer  (ch.7)
     networking    :  a socket's receive/send buffers (ch.17)
     producer/consumer : a RING buffer (circular queue) (ch.11)

   They are all the same idea: contiguous bytes held in transit.
\end{verbatim}

The magnitude of the win is worth stating with numbers. A system call
costs hundreds of nanoseconds because of the privilege switch we saw
earlier in this chapter.

\begin{verbatim}
   Reading a 1 GB file:

     one byte at a time:   read(fd, &b, 1);      ~1,000,000,000 syscalls
     4096-byte chunks:     read(fd, buf, 4096);    ~250,000      syscalls

     The buffer turns a billion expensive crossings of the kernel
     boundary into a quarter-million. That is the difference between a
     program that finishes in seconds and one that takes an hour.
\end{verbatim}

\begin{osconcept}
Buffering is not just a user-mode habit; the kernel buffers too. The
\emph{buffer cache} keeps recently-read disk blocks in RAM so that a
second read of the same block never touches the disk again, and write
coalescing lets the kernel schedule big sequential writes instead of
thousands of tiny ones. Deciding how much to buffer, when to flush,
and how to share the cache fairly among processes is an operating
system problem. Your user-space buffer is the interface to that
machinery.
\end{osconcept}

\subsection{The Partial I/O Contract}

\code{read(fd, buf, n)} asks the kernel to copy up to \code{n} bytes
from the file into your buffer; it returns how many it actually copied.
\code{write(fd, buf, n)} does the reverse. The subtlety that beginners
miss, and that makes systems programming distinct from scripting, is
that these calls may transfer \emph{fewer bytes than requested}, and
that is not an error---it is the normal contract of an I/O device.

\begin{lstlisting}[caption={A correct copy loop: reads and writes in chunks.}]
#include <stdio.h>
#include <stdlib.h>
#include <fcntl.h>
#include <unistd.h>

#define CHUNK 4096

/* Copy src into dst. Returns 0 on success, -1 on failure. */
int copy_file(const char *src, const char *dst) {
    int in_fd;
    int out_fd;
    char buf[CHUNK];
    ssize_t got;
    ssize_t written;
    ssize_t r;
    int failed;

    in_fd = open(src, O_RDONLY);
    if (in_fd < 0) {
        perror("open(src)");
        return -1;
    }

    out_fd = open(dst, O_WRONLY | O_CREAT | O_TRUNC, 0644);
    if (out_fd < 0) {
        perror("open(dst)");
        close(in_fd);
        return -1;
    }

    failed = 0;
    while (1) {
        got = read(in_fd, buf, sizeof(buf));
        if (got < 0) {
            perror("read");
            failed = 1;
            break;
        }
        if (got == 0) {
            break;          /* end of file: nothing left to copy */
        }

        /* write() may transfer fewer bytes than we asked; loop until
           every byte we read has been written.                    */
        written = 0;
        while (written < got) {
            r = write(out_fd, buf + written, (size_t) (got - written));
            if (r < 0) {
                perror("write");
                failed = 1;
                break;
            }
            written = written + r;
        }
        if (failed) {
            break;
        }
    }

    close(in_fd);
    close(out_fd);
    return failed ? -1 : 0;
}

int main(int argc, char *argv[]) {
    if (argc != 3) {
        fprintf(stderr, "usage: %s src dst\n", argv[0]);
        return 1;
    }
    return copy_file(argv[1], argv[2]) == 0 ? 0 : 1;
}
\end{lstlisting}

\begin{verbatim}
   Why the inner loop?  Because the contract says "at most n bytes":

   write(out_fd, buf + written, got - written);
                                   ^^^^^^^^^^^
   may write 1, or 10, or all 4096. The loop keeps calling until the
   count is exhausted. For files on disk this rarely matters; for
   pipes, terminals, and network sockets it matters constantly. Build
   the loop once, copy it into every program you write.
\end{verbatim}

\begin{verbatim}
   The read contract:

     read(fd, buf, n)  =  n          n bytes copied, exactly as asked
     read(fd, buf, n)  =  k < n      k bytes copied (partial); try again
     read(fd, buf, n)  =  0          end of file, no more bytes
     read(fd, buf, n)  =  -1         error; inspect errno
\end{verbatim}

Two design choices in this program are deliberate. The buffer is
\code{4096} bytes---one page---because that matches the granularity the
kernel and the disk like. And we check every return value, because in
systems code the ``impossible'' failure is not impossible; it is just
waiting for the right day to happen.

\begin{osconcept}
I/O devices are asynchronous, shared, and slow, and the kernel sits
between your program and all of that: it decides when a read can
complete (the device may still be spinning), it buffers and caches so
repeated small reads do not hammer the disk, and it interleaves
requests from all processes fairly. The partial-I/O contract, the
buffer cache, and the scheduling of device requests are therefore
operating system problems; your program merely states what it wants and
pays attention to what it got.
\end{osconcept}

\section{Seeking: Moving the File Position}

Each open file carries a \emph{position}: the offset where the next
read or write will happen. Reads and writes advance it. \code{lseek}
moves it without touching any data---seeking past the end then writing
creates a \emph{hole} of zero bytes, which is how sparse files are born.

\begin{lstlisting}[caption={Finding a file's size by seeking to its end.}]
#include <stdio.h>
#include <fcntl.h>
#include <unistd.h>
#include <sys/stat.h>

/* Return the size in bytes of the file opened as fd, or -1. */
off_t file_size(int fd) {
    off_t end;

    end = lseek(fd, 0, SEEK_END);
    if (end < 0) {
        perror("lseek");
        return -1;
    }
    return end;
}
\end{lstlisting}

\begin{verbatim}
   A file with a hole:

        logical bytes:   0       1024     2048    4096
                         +--------+--------+--------+
                         | 'A'*1k | (hole) | 'B'*2k |
                         +--------+--------+--------+
   physical disk blocks: allocated  nothing  allocated
                         (the hole costs no disk space; reads of the
                          hole return zero bytes)

   lseek(fd, 2048, SEEK_SET);  then write(fd, buf, 2048);
   fills the hole only if you actually write into it.
\end{verbatim}

\begin{osconcept}
The file position is kernel state attached to an open file, not to the
file on disk or to your program. When two processes open the same file,
each gets its own position, while \code{O\_APPEND} deliberately makes
the kernel bump the position before every write so concurrent appenders
never clobber each other. Deciding where the position lives, and how
concurrent access to it is serialized, is operating system work.
\end{osconcept}

\section{errno: The Error Channel}

Every system call that can fail records the reason in a thread-local
integer named \code{errno}. Two functions turn it into human words:
\code{perror(prefix)} prints \code{prefix} plus the message, and
\code{strerror(errno)} returns the message string alone. The discipline
is to check the failure and report \emph{immediately}, while the
meaning is still fresh.

\begin{lstlisting}[caption={The error-handling pattern used everywhere in this book.}]
#include <stdio.h>
#include <stdlib.h>
#include <fcntl.h>
#include <errno.h>

int open_or_die(const char *path) {
    int fd;

    fd = open(path, O_RDONLY);
    if (fd < 0) {
        fprintf(stderr, "cannot open %s: %s\n",
                path, strerror(errno));
        exit(1);
    }
    return fd;
}
\end{lstlisting}

\begin{verbatim}
   What each failure actually means:

     open("f.txt", O_RDONLY)  ->  -1, errno = ENOENT   No such file
     open("/root/x", O_RDONLY)->  -1, errno = EACCES   Permission denied
     open("d/", O_RDONLY)     ->  -1, errno = EISDIR   Is a directory
     read(fd, ...)            ->  -1, errno = EINTR    Interrupted by signal

   perror("open")  prints:    open: No such file or directory
\end{verbatim}

\begin{ruleofthumb}[Report errors where they happen]
Do not swallow errors, and do not report them a mile from the cause.
The instant a call returns \code{-1}, print or log it with enough
context (\code{errno} and the operation) and decide the failure policy
right there: retry, degrade, or abort. Programs that defer error
handling are the ones that ship the confusing bug reports.
\end{ruleofthumb}

\section{Standard Streams and Redirection}

Descriptors 0, 1, and 2 are already open: standard input, standard
output, and standard error. \code{printf} writes to descriptor 1;
\code{fprintf(stderr, ...)} to descriptor 2. The shell's
\code{>}, \code{<}, and \code{2>} operators are nothing but the shell
rearranging what those descriptors point at before it \code{exec}s your
program.

\begin{verbatim}
   What the shell does for "myprog > out.txt 2> err.txt":

     1. open("out.txt", O_WRONLY|O_CREAT|O_TRUNC, 0644) -> fd 3
     2. open("err.txt", O_WRONLY|O_CREAT|O_TRUNC, 0644) -> fd 4
     3. dup2(3, 1);   /* make descriptor 1 ALSO name out.txt  */
     4. dup2(4, 2);   /* make descriptor 2 ALSO name err.txt  */
     5. close(3); close(4);
     6. exec("myprog");

   myprog sees stdout as out.txt and stderr as err.txt, with no
   knowledge that a redirection ever happened. That is the power of
   the descriptor indirection: the program does not care what its
   streams are connected to.
\end{verbatim}

\begin{lstlisting}[caption={Redirecting stdout inside your own program.}]
#include <stdio.h>
#include <stdlib.h>
#include <fcntl.h>
#include <unistd.h>

int main(void) {
    int fd;

    fd = open("out.txt", O_WRONLY | O_CREAT | O_TRUNC, 0644);
    if (fd < 0) {
        perror("open");
        exit(1);
    }

    /* Make descriptor 1 (stdout) name the same file as fd. */
    if (dup2(fd, 1) < 0) {
        perror("dup2");
        exit(1);
    }
    close(fd);       /* we no longer need the original descriptor */

    printf("this line lands in out.txt\n");
    return 0;
}
\end{lstlisting}

\code{dup2(fd, 1)} is the magic word: it copies the kernel object that
\code{fd} names onto descriptor 1, atomically. This single call is the
foundation of every shell redirection, every pipeline, and most of the
inter-process communication we will build in Chapter~7.

\begin{osconcept}
Redirection works because the kernel, not the program, owns what each
descriptor names. A program cannot tell whether descriptor 1 is a
terminal, a file, or a pipe---and it does not need to. Making the
connection between a descriptor number and the underlying object a
changeable piece of kernel state, mutable by \code{dup2}, is what lets
an external agent (the shell) rewire a program's world from outside.
That is a subtle, deliberate piece of OS design.
\end{osconcept}

\section{Processes: What the OS Schedules}

A \emph{process} is a running program: one virtual address space, one
set of registers, one execution thread, and a set of kernel resources
(open files, a working directory, environment variables). The OS keeps
one record per process---the process control block (PCB)---and the
\emph{scheduler} decides, thousands of times per second, which process
gets the CPU next. From the OS's point of view, a process is a unit of
accounting and isolation; from your program's point of view, it is
``everything I am right now.''

\begin{verbatim}
   The OS keeps MANY processes in flight on a few CPUs:

     CPU 0:  [ A ] [ C ] [ A ] [ D ] ...    (each column a time slice)
     CPU 1:  [ B ] [ A ] [ D ] [ B ] ...

   Each process THINKS it owns the machine; the scheduler interleaves
   them. Switching is fast but not free: registers, the memory map,
   and the page tables must be swapped ("context switch").

   Isolation: A cannot read or write B's memory. The OS gives each
   process its own private address space:

     Process A               Process B
     +---------+             +---------+
     | stack   |             | stack   |
     | heap    |             | heap    |
     | data    |             | data    |
     | text    |             | text    |
     +---------+             +---------+
       \                       /
        \  shared ONLY through explicit OS channels: files, pipes,
         sockets -- never by poking another process's memory.
\end{verbatim}

\begin{osconcept}
Multiplexing the CPU is the quintessential OS problem. The hardware has
one instruction stream per core; the OS manufactures the illusion of
concurrency by switching processes, and it must do so \emph{fairly}
(no process starves), \emph{safely} (one process cannot corrupt
another's memory), and \emph{cheaply} (switches are not too frequent).
Without the OS, cooperating programs could not share a CPU at all.
\end{osconcept}

\section{fork(): The Process Duplication Primitive}

How are new processes created? On Unix, the primitive is \code{fork}:
the kernel makes a near-exact copy of the calling process. The child is
a new process with its own copy of the parent's address space, its own
copy of the descriptor table, and its own process id. After the call,
\emph{two} processes are running the same code from the same line.

The cleverness of the implementation is that the copy is nearly free at
first. The parent's pages are shared with the child and marked
\emph{copy-on-write} (COW): the physical pages are only duplicated when
one of the two processes actually writes to them. Most programs
\code{fork} and then immediately \code{exec}, so most pages are never
copied at all.

\begin{lstlisting}[caption={fork returns twice: once in each process.}]
#include <stdio.h>
#include <stdlib.h>
#include <unistd.h>
#include <sys/types.h>

int main(void) {
    pid_t pid;

    pid = fork();
    if (pid < 0) {
        perror("fork");
        exit(1);
    }

    if (pid == 0) {
        /* This branch runs ONLY in the child. */
        printf("child:  my pid = %d, my parent's pid = %d\n",
               getpid(), getppid());
        exit(0);
    }

    /* This branch runs ONLY in the parent. */
    printf("parent: my pid = %d, my child's pid = %d\n",
           getpid(), pid);

    return 0;
}
\end{lstlisting}

\begin{verbatim}
   One call, two worlds:

                              fork()
                 parent -----------------> parent (returns child's pid)
                    |   |                     (the caller CONTINUES here)
                    |   +---- child ---------> child (returns 0)
                    |       (the child starts here, at the SAME line,
                    |        with a COPY of the parent's memory)

   The two processes are now independent: the child can change its
   copy of a variable and the parent's copy is untouched, thanks to
   copy-on-write the change happens on a page that is now private.
\end{verbatim}

Two facts about the return value drive every use of fork:

\begin{verbatim}
   In the PARENT, fork returns the child's pid (a positive number).
   In the CHILD,  fork returns 0.
   On failure,    fork returns -1 and no child is created.

   The single branch -- if (pid == 0) child else parent -- is the
   universal shape. After fork, parent and child have DIFFERENT pids,
   DIFFERENT copies of memory, but the SAME open files (the descriptor
   table is copied, so both can write to fd 1, for example).
\end{verbatim}

\begin{osconcept}
Creating a process is a privileged act: the kernel must allocate a new
process control block, a new address space, new page tables, and a new
scheduler entry, then wire the child into the process tree. A user
program cannot do any of that. The elegance of \code{fork} is that it
makes process creation almost free by deferring all copying via
copy-on-write---a decision only the OS can make, because only the OS
controls the page tables.
\end{osconcept}

\subsection{vfork: The Half-Fork}

There is an older, rawer sibling of \code{fork} that you will meet in
old code and in some fast libraries, and it deserves a name:
\code{vfork} is the \emph{half-fork}. Where \code{fork} gives the
child a copy (or a copy-on-write view) of the address space,
\code{vfork} gives the child \emph{the parent's very own pages}, and
\emph{suspends the parent} until the child either \code{exec}s or
\code{\_exit}s.

\begin{verbatim}
   fork  :  parent and child BOTH run, each with (eventually) its own
            copy of memory. The child can do anything the parent can.

   vfork :  only the child runs; the parent is FROZEN. The child is
            using the parent's memory directly -- no copy at all --
            and may ONLY call exec() or _exit(). That is the entire
            deal.

   Why it existed: before copy-on-write was invented, fork() copied
   the whole address space. For the overwhelmingly common
   "fork, then immediately exec a new program," that copy was pure
   waste. vfork skipped it: no copy, no COW, no memory at all -- the
   child borrows the parent's address space for the instant it needs
   it, then execs and leaves it behind.

   Why it is dangerous: because the child runs on the PARENT'S stack
   and memory, it must not modify anything, must not call ordinary
   functions (malloc, printf), and must never return from its own
   main -- if it does, it corrupts the suspended parent. The POSIX
   rules are absolute: in the vfork child you may call only
   async-signal-safe functions, exec, or _exit.
\end{verbatim}

On modern systems, \code{fork} is copy-on-write and costs almost
nothing, so \code{vfork} has lost most of its reason to exist. POSIX
keeps it for rare, performance-critical fork-and-exec paths, and the
standard is explicit about what the vfork child may do: call only
async-signal-safe functions, \code{exec}, or \code{\_exit}. The
half-fork survives in tight code that wants the absolute minimum
between two processes---and it is worth recognizing on sight, because
its constraints are the opposite of \code{fork}'s freedom.

\begin{ruleofthumb}[Never write new code with vfork]
Use \code{fork}. The copy-on-write \code{fork} of any modern kernel
is cheap, and it does not freeze the parent or forbid ordinary
library calls in the child. If you ever meet \code{vfork} in the
wild, read the word as a warning sign: the code around it is older
and stricter, and the child may only \code{exec} or \code{\_exit}.
\end{ruleofthumb}

\section{exec(): Replacing the Process Image}

\code{fork} duplicates the process; \code{exec} \emph{replaces} it.
The \code{exec} family (\code{execl}, \code{execv}, \code{execvp},
\code{execve}) loads a new program file into the current process,
discarding the old address space and starting the new program at its
\code{main}. The process id does not change. If \code{exec} succeeds,
it never returns---there is no old program left to return to. If it
fails, it returns \code{-1} and you are still in the old program.

\begin{lstlisting}[caption={Fork + exec: the standard recipe for "run this program".}]
#include <stdio.h>
#include <stdlib.h>
#include <unistd.h>
#include <sys/types.h>
#include <sys/wait.h>

/* Run the program named by argv[0] with arguments argv[1..].
   Returns the child's exit status, or -1 on failure. */
int run_command(char *const argv[]) {
    pid_t pid;
    int status;

    pid = fork();
    if (pid < 0) {
        perror("fork");
        return -1;
    }

    if (pid == 0) {
        /* CHILD: become the new program. execvp searches the PATH
           for argv[0]. It only returns here if it FAILED.       */
        execvp(argv[0], argv);
        perror("execvp");
        exit(127);
    }

    /* PARENT: wait for the child to finish, then report its status. */
    if (waitpid(pid, &status, 0) < 0) {
        perror("waitpid");
        return -1;
    }
    if (WIFEXITED(status)) {
        return WEXITSTATUS(status);
    }
    return -1;
}

int main(int argc, char *argv[]) {
    if (argc < 2) {
        fprintf(stderr, "usage: %s command [args...]\n", argv[0]);
        return 1;
    }
    return run_command(argv + 1);
}
\end{lstlisting}

\begin{verbatim}
   Life of a process that runs "ls -l":

   main()                       (your program)
     fork() --> [child]
     parent: waitpid()          (sits, waiting)
     child:  execvp("ls", ...)
             (the child's address space is REPLACED by /bin/ls;
              the pid is UNCHANGED)
     /bin/ls runs, prints, exits with status 0
   parent: waitpid() returns 0  (the child's exit status)
\end{verbatim}

\begin{verbatim}
   fork + exec = "build a new shell of a process, then fill it":
     fork creates a private address space with your image;
     exec destroys that image and loads a new one.
   The two steps are kept separate so the child can do setup between
   them -- close a descriptor, dup2 a pipe, change directories --
   which is exactly how shells and daemons are built.
\end{verbatim}

\begin{osconcept}
Loading an executable image is privileged work: the kernel must parse
the file format, map the code and data segments into a fresh address
space, set up the stack and the \code{argv}/\code{envp} arrays, and
honor the file's permissions and security flags. No user-mode code can
legitimately conjure a running program from a file---only the kernel
can. That \code{exec} preserves the pid and the open descriptors while
replacing the memory is a carefully chosen OS semantic that makes
processes composable.
\end{osconcept}

\section{wait() and waitpid(): Reaping Children}

When a child exits, it does not vanish. It enters the \emph{zombie}
state: all of its resources are released except a tiny record holding
its pid and its exit status. That record exists so the parent can
discover \emph{how} the child died. The parent must call
\code{wait} or \code{waitpid} to collect the record, at which point
the zombie is reaped and the pid becomes reusable.

\begin{verbatim}
   The life cycle of a child:

      fork() ---> running ---> exiting ---> ZOMBIE ---> reaped
                        |                      |
                        +------ exits --------++------ wait() collects it

   - A ZOMBIE consumes almost no resources, but its pid is taken.
   - If the parent NEVER waits, the zombie lingers until the parent
     exits; then init (pid 1) adopts it and reaps it.
   - If the parent exits FIRST, the child becomes an ORPHAN and is
     adopted by init (pid 1), which reaps it when it finally dies.

   The rule: every fork() deserves a matching wait(). Forgetting the
   wait leaks a zombie per child -- a classic "my server grew to
   thousands of processes" bug.
\end{verbatim}

\begin{lstlisting}[caption={waitpid with the exit-status macros.}]
#include <stdio.h>
#include <stdlib.h>
#include <unistd.h>
#include <sys/types.h>
#include <sys/wait.h>

int main(void) {
    pid_t pid;
    int status;

    pid = fork();
    if (pid < 0) {
        perror("fork");
        exit(1);
    }

    if (pid == 0) {
        printf("child exiting with status 3\n");
        exit(3);
    }

    if (waitpid(pid, &status, 0) < 0) {
        perror("waitpid");
        exit(1);
    }

    if (WIFEXITED(status)) {
        printf("child exited with status %d\n", WEXITSTATUS(status));
    } else if (WIFSIGNALED(status)) {
        printf("child killed by signal %d\n", WTERMSIG(status));
    }

    return 0;
}
\end{lstlisting}

\begin{verbatim}
   The status integer packs three stories:
     WIFEXITED(s)    : did it exit normally?
     WEXITSTATUS(s)  : if so, its exit code (0..255)
     WIFSIGNALED(s)  : was it killed by a signal?
     WTERMSIG(s)     : which signal?

   Without wait(), the parent has NO way to know these. And because
   a signal can kill a process at ANY moment, a parent that ignores
   wait() cannot even be sure its children succeeded.
\end{verbatim}

\begin{osconcept}
Reaping children is an invariant the OS maintains for the whole
process tree: no pid is ever reused while a parent might still ask
about it, and no orphan is ever left uncollected. Keeping the exit
status available until the parent explicitly retrieves it, adopting
orphans, and recycling pids safely are all kernel bookkeeping
decisions---your program just calls \code{waitpid}, and the OS handles
the rest. The zombie rule exists so the answer to ``what happened to my
child?'' is never lost.
\end{osconcept}

\section{Practical: A Tiny Command Runner}

Let us combine everything from this chapter into one tool: a miniature
shell for a single command. It forks, lets the child \code{exec} the
requested program, waits for it, and prints how it ended. This program
is a few dozen lines, yet it is the beating heart of every shell,
service manager, and job-control system ever written.

\begin{lstlisting}[caption={runone.c: fork, exec, wait -- the whole chapter at once.}]
#include <stdio.h>
#include <stdlib.h>
#include <unistd.h>
#include <sys/types.h>
#include <sys/wait.h>

int main(int argc, char *argv[]) {
    pid_t pid;
    int status;

    if (argc < 2) {
        fprintf(stderr, "usage: %s command [args...]\n", argv[0]);
        return 1;
    }

    pid = fork();
    if (pid < 0) {
        perror("fork");
        return 1;
    }

    if (pid == 0) {
        /* Child: become the requested program. */
        execvp(argv[1], argv + 1);
        perror("execvp");
        exit(127);          /* conventional "command not found" status */
    }

    /* Parent: wait for the child and report the outcome. */
    if (waitpid(pid, &status, 0) < 0) {
        perror("waitpid");
        return 1;
    }

    if (WIFEXITED(status)) {
        printf("command exited with status %d\n", WEXITSTATUS(status));
    } else if (WIFSIGNALED(status)) {
        printf("command killed by signal %d\n", WTERMSIG(status));
    }

    return 0;
}
\end{lstlisting}

\begin{verbatim}
   Running it:

     $ ./runone echo hello world
     hello world
     command exited with status 0

     $ ./runone ls /nonexistent
     ls: cannot access '/nonexistent': No such file or directory
     command exited with status 2

     $ ./runone /bin/sh -c 'kill -9 $$'
     command killed by signal 9

   One small program demonstrates: file descriptors survive exec
   (echo used the same fd 1 we inherited), argv is rebuilt by the
   kernel, and wait() reveals whether the child exited or was killed.
\end{verbatim}

The next time you type a command into a shell, know that you are
watching this chapter run: the shell \code{fork}s, you are the child
mid-\code{exec}, and the prompt you get back is the shell's
\code{waitpid} returning. Everything we built with structs and pointers
is what runs \emph{inside} such a process; this chapter is what creates
it.

\section{Structures in Bytes}

The byte structures this chapter touched are the tables the kernel
builds and the stack it prepares at \code{exec} time.

\begin{verbatim}
   The argv array on the new process's stack (high to low):

   +-------------------------------+
   | "two words\0"     argv[2]     |
   | "alpha\0"         argv[1]     |   the strings, each NUL-terminated
   | "./prog\0"        argv[0]     |
   +-------------------------------+
   | NULL                          |   argv[argc] == NULL
   | argv[2]  (pointer, 8B)        |
   | argv[1]  (pointer, 8B)        |   an array of 8-byte pointers
   | argv[0]  (pointer, 8B)        |
   +-------------------------------+
   | argc (8B)                     |   the count
   +-------------------------------+

   The kernel's descriptor table is an array of pointers to file
   objects; each file object carries the position, mode, and flags,
   and points at the inode. dup2() copies a pointer, and close()
   drops a reference count.
\end{verbatim}

A stack frame has the same shape as Chapter~1's: an 8-byte return
address pushed by \code{call}, the saved frame pointer, then the
locals. System calls add no byte structures to your process; they are
requests that the kernel's own records---descriptor table, open-file
table, process table---do the work.

\section{Summary}

The chapter's arc, in one picture:

\begin{verbatim}
   your program (user mode)
   +------------------------------------------------------------+
   | structs & pointers   (chapters 1-3)                        |
   | library calls: malloc, printf, strcmp                       |
   +----------------------------+-------------------------------+
                                | system calls (the only door)
                                v
   the kernel (kernel mode)  +---------------------------------+
   | open/read/write/lseek/close  -> files                     |
   | fork/exec/waitpid           -> processes                  |
   | dup2                        -> redirection & wiring      |
   | descriptor table, open file table, scheduler, zombies    |
   +-----------------------------------------------------------+

   discipline: every call is checked, every error is reported with
   errno, every fork is matched with a wait, every fd is closed.
\end{verbatim}

\begin{osconcept}
This entire chapter is ``the OS problem'' in miniature: a safe,
stable, shared interface between untrusted user code and privileged
hardware. The abstractions you met---files, descriptors, processes,
wait--exist so that many programs can coexist safely on one machine.
Understanding that these are \emph{services with contracts}, not
helper functions, is what separates a C programmer from a systems
programmer.
\end{osconcept}

In the next chapter we meet the events that can interrupt a process at
any moment---signals---and learn how to handle, ignore, and block
them without corrupting our own state.

}{}
\IfFileExists{chapter6.tex}{% ============================================================================
%  chapter6.tex
%  Signals: Asynchronous Events and Their Handlers
% ============================================================================

\chapter{Signals: Asynchronous Events and Their Handlers}

\section{What Is a Signal?}

Everything so far has been synchronous: your code ran, made a system
call, and waited for the answer. Signals are the operating system
breaking that rule. A \emph{signal} is a small, asynchronous
notification delivered to a process. It can arrive at \emph{any}
moment---between any two of your instructions---and it is not requested
by your code. It simply happens to you.

The kernel generates signals for a dizzying variety of reasons:

\begin{verbatim}
   WHO OR WHAT FIRES   SIGNAL        WHY
   ----------------   ----------    ----------------------------------
   you press Ctrl-C    SIGINT        the terminal driver asks you to stop
   you press Ctrl-\    SIGQUIT       stop and leave a core dump
   kill command        SIGTERM       polite shutdown request
   kill -9 command     SIGKILL       unignorable kill
   your program derefs SIGSEGV       invalid memory access (ch.1's abyss)
   a timer expires     SIGALRM       alarm() went off
   a child exits       SIGCHLD       "your child is done; go collect it"
   writing a dead pipe SIGPIPE       the reader of a pipe has gone away
   stop command        SIGSTOP       freeze this process (uncatchable)
   you type a TSTP     SIGTSTP       Ctrl-Z: stop, but catchable
\end{verbatim}

\begin{verbatim}
   Time flows left to right. Signals arrive OUT OF BAND:

     your code:  ----instr--instr--instr--instr--instr--instr--instr---
                              |                  |
                              v                  v
                          SIGINT           SIGSEGV
                          (Ctrl-C)          (bug!)

     The handler runs, returns, and your code resumes at the very next
     instruction. From your code's point of view, a signal is a hidden
     function call injected by the kernel.
\end{verbatim}

A signal is really just a small integer (1--31 on Linux, with a range
of real-time signals beyond that). All the meaning lives in the
\emph{disposition} the process has registered for it and in the default
action the kernel would take if the process has not.

\begin{osconcept}
Only the kernel can deliver signals, because only the kernel can
\emph{interrupt a running process}. A process is a passive guest of the
CPU; it cannot pause itself to receive a notification. The kernel, which
controls the CPU, must notice events (keys, timers, child deaths) and
inject them into the victim process's execution. Asynchronous event
delivery is thus an operating system problem by construction: the OS is
the only actor with the power and the global view to do it.
\end{osconcept}

\section{The Default Actions}

When a signal arrives and the process has no custom handler, the kernel
applies one of a small set of \emph{default actions}. There are exactly
five of them, and every signal's default is one of these:

\begin{verbatim}
   default action     meaning                          examples
   ----------------   -------------------------------  ---------
   terminate          kill the process                SIGHUP SIGINT SIGTERM
   core dump          terminate + save core image     SIGSEGV SIGQUIT SIGABRT
   stop               suspend the process             SIGSTOP SIGTSTP
   continue           resume a stopped process        SIGCONT
   ignore             do nothing                      SIGCHLD SIGURG

   So "Ctrl-C kills my program" is not a rule of the language; it is
   the DEFAULT action for SIGINT, which is why Ctrl-C works the same
   in every program whether or not it knows about signals.
\end{verbatim}

Two signals have no handler, period. \code{SIGKILL} and
\code{SIGSTOP} are the kernel's own kill switch: they cannot be caught,
blocked, or ignored. If any process could block \code{SIGKILL}, a
runaway program could make itself immortal; the OS would lose its
ultimate authority. Their existence is the sign that signals are not a
programmer convenience but an operating system mechanism.

\begin{osconcept}
The default-action table is operating system policy. The kernel must
decide, for every possible event, what happens to a process that has
not expressed an opinion---and that decision must hold even when the
process is broken (a dereferenced \code{NULL}, an illegal instruction).
Core dumps, process reaping, and the uncatchable \code{SIGKILL} are all
designed at the OS level to keep the whole system rescuable, not merely
each individual process.
\end{osconcept}

\section{The Three Dispositions}

For every signal except \code{SIGKILL} and \code{SIGSTOP}, a process
chooses one of three \emph{dispositions}: let the kernel use the
default, ignore the signal outright, or run a function of your own.
That function is the \emph{signal handler}.

\begin{verbatim}
   disposition          what the kernel does when the signal arrives
   ------------------   ----------------------------------------------
   SIG_DFL  (default)   apply the default action from the table above
   SIG_IGN  (ignore)    discard the signal, as if it never happened
   handler  (custom)    call your function, then resume your code

   Default is the starting state of every signal. You override it.
\end{verbatim}

The modern, portable way to install a disposition is
\code{sigaction}, which tells the kernel everything about your handler
in one structure. (The older \code{signal} call behaves differently on
different Unixes; in this book we use \code{sigaction} and forget
\code{signal} exists.)

\begin{lstlisting}[caption={Installing a handler with sigaction.}]
#include <stdio.h>
#include <signal.h>
#include <unistd.h>

/* Runs when SIGINT arrives. "sig" tells you which signal fired. */
static void handle_sigint(int sig) {
    write(1, "interrupted!\n", 13);
}

int main(void) {
    struct sigaction sa;

    sa.sa_handler = handle_sigint;
    sigemptyset(&sa.sa_mask);     /* nothing extra blocked during it */
    sa.sa_flags = 0;

    if (sigaction(SIGINT, &sa, NULL) < 0) {
        perror("sigaction");
        return 1;
    }

    printf("press Ctrl-C (or send SIGINT) ...\n");
    while (1) {
        sleep(1);
    }
    return 0;
}
\end{lstlisting}

\begin{verbatim}
   What just happened:

     struct sigaction sa;      describes the desired behavior
     sa.sa_handler            the function to run
     sa.sa_mask               signals to block WHILE the handler runs
     sa.sa_flags              options (more below)

     sigaction(SIGINT, &sa, NULL)
         ^signal   ^new behavior   ^old behavior (NULL = don't care)

   When Ctrl-C arrives, the kernel saves your state, runs
   handle_sigint(2), restores your state, and your loop continues
   exactly where it was interrupted.
\end{verbatim}

Notice the handler writes with \code{write(2)}, not \code{printf}.
That is not a stylistic tic; it is a hard rule, and we will justify it
in Section~6.6.

\begin{osconcept}
Dispositions are kernel state attached to the process, not program
state. The kernel must remember, for each of dozens of signals, what to
do on delivery, and it must honor that even when the process is deep in
some library you do not control. Registering a handler, validating it,
and switching behavior at the exact moment of delivery is kernel work;
your \code{sigaction} call is just writing an entry in a table the
kernel maintains and consults.
\end{osconcept}

\section{Blocking, Pending, and Delivered}

A signal does not necessarily get delivered the instant it is
generated. The kernel runs a small state machine for each signal of
each process, with three states:

\begin{verbatim}
   GENERATED          BLOCKED?            DELIVERED
   (event happens)    process asked       kernel runs the handler
                      to delay it?        (or default action)
        |                   |
        +----> pending -----+----> if unblocked, deliver now
        |     (waiting)            if blocked, keep pending
        |
        +-----------------+-----------> delivered immediately

   - A BLOCKED signal is not lost: it sits "pending" and is delivered
     the moment the process unblocks it.
   - Standard signals COALESCE: ten SIGINTs generated while blocked
     are delivered as ONE. The kernel remembers "a SIGINT is pending",
     not "ten SIGINTs are pending".
   - A signal that arrives while its own handler is running may be
     delivered again (or blocked, per sa.sa_mask).
\end{verbatim}

Blocking is done with a \emph{signal set}, manipulated through
\code{sigset\_t} and friends:

\begin{lstlisting}[caption={Block, then unblock, a signal.}]
#include <stdio.h>
#include <stdlib.h>
#include <signal.h>
#include <unistd.h>

static void handle_sigint(int sig) {
    write(1, "SIGINT handled\n", 15);
}

int main(void) {
    struct sigaction sa;
    sigset_t block_set;
    sigset_t old_set;
    int i;

    sa.sa_handler = handle_sigint;
    sigemptyset(&sa.sa_mask);
    sa.sa_flags = 0;
    sigaction(SIGINT, &sa, NULL);

    /* 1. Block SIGINT: it will now go pending instead of being delivered. */
    sigemptyset(&block_set);
    sigaddset(&block_set, SIGINT);
    sigprocmask(SIG_BLOCK, &block_set, &old_set);

    printf("SIGINT is now blocked. Send one (Ctrl-C in another terminal,\n");
    printf("or `kill -INT %d` from a shell). Sleeping 3 seconds...\n", getpid());
    sleep(3);

    /* 2. Unblock: any SIGINT that arrived is delivered NOW. */
    printf("unblocking...\n");
    sigprocmask(SIG_SETMASK, &old_set, NULL);

    sleep(2);
    return 0;
}
\end{lstlisting}

\begin{verbatim}
   The blocked interval:

     you send SIGINT           3 seconds later you unblock
        |                             |
        v                             v
   +---------+  blocked,       +-----+  delivered now, handler runs
   | pending |  then one        |     |
   | SIGINT  |  unblock --->    |  !  |----->  write("SIGINT handled")
   +---------+                  +-----+
\end{verbatim}

\begin{ruleofthumb}[Block the signal during its own handler]
Unless you want a handler to be re-entered, add the signal you are
handling to \code{sa.sa\_mask}. Simpler still: make handlers set a
\emph{flag} and do all real work in the main loop (Section~6.7); then
re-entrancy stops being a question you have to answer.
\end{ruleofthumb}

\begin{osconcept}
The pending/blocked/delivered state machine lives entirely in the
kernel. Only the kernel can decide when a process is at a safe point to
run a handler, and only the kernel can hold a signal in limbo without
losing it. A process cannot "pause" its own execution safely---it has
no authority over the scheduler. The signal lifecycle is therefore
kernel bookkeeping, exposed to you as three tiny calls.
\end{osconcept}

\section{What You May Do in a Handler}

Here is the discipline that keeps signal programming honest: a handler
runs \emph{asynchronously}, between two arbitrary instructions of your
main program. If the handler calls a function that your main program is
midway through using---\code{malloc}, for example---the two can corrupt
each other. A handler may therefore call only \emph{async-signal-safe}
functions, a short list blessed by POSIX. \code{printf}, \code{malloc},
\code{strcpy}, and \code{free} are \emph{not} on it. \code{write},
\code{read}, \code{waitpid}, \code{kill}, \code{\_exit}, and
\code{sigaction} are.

\begin{verbatim}
   ASYNC-SIGNAL-SAFE (safe in a handler)     NOT safe in a handler
   -------------------------------          ----------------------
   write(2), read(2)                         printf, sprintf, puts
   _exit(2)                                  exit, malloc, free
   kill(2), raise(3), waitpid(2)             strcpy, strlen, memcpy
   sigaction, sigprocmask                    fopen, fclose, fflush
   alarm(2), getpid(2)                       any code that locks
\end{verbatim}

The consequence is that a handler usually does almost nothing. The
canonical pattern is a \emph{flag}: the handler sets a variable, the
main loop checks it. The variable must be \code{volatile} (so the
compiler does not cache it in a register across the async change) and
of type \code{sig\_atomic\_t} (an integer type guaranteed to be
readable and writable in one atomic step).

\begin{lstlisting}[caption={The flag pattern: handler sets, main loop acts.}]
#include <stdio.h>
#include <signal.h>
#include <unistd.h>

/* Set by the handler, read by main. volatile sig_atomic_t is the
   one type guaranteed safe for this two-way traffic.              */
static volatile sig_atomic_t got_sigint = 0;

static void handle_sigint(int sig) {
    got_sigint = 1;       /* the ONLY thing the handler does */
}

int main(void) {
    struct sigaction sa;
    int ticks;

    sa.sa_handler = handle_sigint;
    sigemptyset(&sa.sa_mask);
    sa.sa_flags = 0;
    sigaction(SIGINT, &sa, NULL);

    ticks = 0;
    while (got_sigint == 0) {
        ticks = ticks + 1;
        printf("tick %d\n", ticks);
        sleep(1);
    }

    printf("SIGINT received; the loop noticed within a second.\n");
    return 0;
}
\end{lstlisting}

\begin{verbatim}
   The division of labor:

     signal arrives ----> handler runs:   got_sigint = 1;
     (kernel thread of execution)         (safe: one atomic store)

     main loop       ----> checks:        while (got_sigint == 0)
     (your thread of execution)           ... does the real work

   The handler and the main loop are like two coroutines sharing one
   int. Everything that is complicated--printing, allocating, closing
   files--belongs in the main loop, where the whole library is safe.
\end{verbatim}

\begin{osconcept}
The async-signal-safe list exists because the kernel delivers signals
at moments it alone chooses, without coordinating with your libraries.
POSIX (the OS standard) must therefore guarantee that a specific subset
of functions never internally allocates, never locks, and never needs
your program's state to be consistent. Whether a function may be called
from a handler is an operating-system-contract decision, not a language
decision. It is the kernel's interrupt model that forces it.
\end{osconcept}

\section{SIGCHLD: Reaping Children Asynchronously}

Chapter~5 ended with a parent that calls \code{waitpid} and blocks
until its child finishes. But what if the parent is busy serving
requests and children finish at unpredictable moments? It cannot sit in
\code{waitpid}. The answer is \code{SIGCHLD}: the kernel sends it to a
parent \emph{whenever} a child changes state---exits, stops, or
continues. The parent installs a \code{SIGCHLD} handler, and the handler
reaps whatever children have died since the last sweep.

The handler must loop: several children may have exited, and
\code{SIGCHLD} \emph{coalesces}, so one signal may stand for many dead
children. The loop drains the pipe:

\begin{lstlisting}[caption={Reaping every dead child with WNOHANG.}]
#include <stdio.h>
#include <stdlib.h>
#include <signal.h>
#include <sys/types.h>
#include <sys/wait.h>
#include <unistd.h>

static volatile sig_atomic_t children_changed = 0;

static void handle_child(int sig) {
    children_changed = 1;
}

/* Called from the MAIN loop, never from the handler itself. */
void reap_children(void) {
    pid_t done;
    int status;

    while (1) {
        done = waitpid(-1, &status, WNOHANG);
        if (done == 0) {
            break;          /* no more children have exited */
        }
        if (done < 0) {
            break;          /* ECHILD: no children left at all */
        }
        if (WIFEXITED(status)) {
            printf("child %d exited with status %d\n",
                   (int) done, WEXITSTATUS(status));
        }
    }
}

int main(void) {
    struct sigaction sa;
    int i;

    sa.sa_handler = handle_child;
    sigemptyset(&sa.sa_mask);
    sa.sa_flags = 0;
    sigaction(SIGCHLD, &sa, NULL);

    /* Spawn three short-lived children. */
    for (i = 0; i < 3; i++) {
        pid_t pid;

        pid = fork();
        if (pid == 0) {
            exit(10 + i);
        }
    }

    /* Busy loop: do work, and check whether reaping is due. */
    for (i = 0; i < 10; i++) {
        sleep(1);
        if (children_changed) {
            children_changed = 0;
            reap_children();
        }
    }

    return 0;
}
\end{lstlisting}

\begin{verbatim}
   What SIGCHLD reaping prevents:

      children exit            zombie count
      -------------            ------------
      3 children finish        the handler's waitpid(-1, WNOHANG)
      (3 SIGCHLD coalesce      loop reaps ALL THREE in one call,
       into possibly 1 signal) so no zombie survives more than a
                               signal delivery.

   WITHOUT the loop (just one waitpid in the handler): two children
   stay zombies until the next SIGCHLD or the parent exits. WITH it,
   the drain loop guarantees every corpse is collected.
\end{verbatim}

\begin{ruleofthumb}[WNOHANG is non-negotiable here]
Inside a \code{SIGCHLD} handler, you must use \code{waitpid(-1,
\&status, WNOHANG)} and loop. A plain \code{waitpid(pid,...)} would
block if no child has exited, freezing the process inside the handler
forever. \code{WNOHANG} means ``don't wait; tell me right now what's
ready,'' which is exactly the contract an asynchronous reaper needs.
\end{ruleofthumb}

\begin{osconcept}
Zombie reaping, introduced in Chapter~5, becomes an event-driven
problem here, and the OS solves it by \emph{fusing} the two mechanisms:
child death generates \code{SIGCHLD} automatically, and \code{waitpid}
collects the result. The kernel chooses to deliver one signal per state
change (coalescing repeats), and it is what keeps the process tree
clean without the parent babysitting every child. Connecting process
lifecycle events to the signal mechanism is a deliberate piece of
operating system design.
\end{osconcept}

\section{Sending Signals: kill and raise}

Signals flow in both directions. \code{kill(pid, sig)} sends a signal
to another process (given permission); \code{raise(sig)} sends one to
yourself; \code{kill(0, sig)} sends to your whole process group. The
permission rule is the kernel's: you may signal a process you own, or
one owned by the same user (or root can signal anyone).

\begin{lstlisting}[caption={Sending signals, three ways.}]
#include <stdio.h>
#include <stdlib.h>
#include <signal.h>
#include <unistd.h>

int main(void) {
    pid_t pid;

    pid = fork();
    if (pid < 0) {
        perror("fork");
        exit(1);
    }

    if (pid == 0) {
        /* Child: announce itself, then idle forever. */
        printf("child %d sleeping...\n", getpid());
        while (1) {
            sleep(1);
        }
    }

    /* Parent: give the child a moment, then say goodnight. */
    sleep(1);
    printf("sending SIGTERM to child %d\n", (int) pid);
    if (kill(pid, SIGTERM) < 0) {
        perror("kill");
        exit(1);
    }

    sleep(1);
    return 0;
}
\end{lstlisting}

\begin{verbatim}
   kill(pid, 0) is a useful trick: it sends NO signal, but succeeds if
   the process exists and you may signal it. "kill -0" is a test of
   existence and permission, widely used in scripts.

   raise(SIGTERM) == kill(getpid(), SIGTERM)   -- signal yourself.
   alarm(5)      : the kernel arranges to send SIGALRM in 5 seconds.
   pause()       : block until ANY signal arrives.
\end{verbatim}

\begin{osconcept}
Permission checking on \code{kill} is kernel work: before delivering,
the kernel verifies the sender owns the target (or is root), because
letting arbitrary processes terminate each other would be a denial of
service. The \emph{capability} to signal someone is a form of
authority the kernel hands out and enforces. \code{kill(pid, 0)}
exists only because that same authority needs a non-destructive test.
\end{osconcept}

\section{alarm: A Signal in the Future}

\code{alarm(n)} asks the kernel to send \code{SIGALRM} to your process
in \code{n} seconds---exactly once. It is the classic way to enforce
timeouts on operations that might hang, because \code{SIGALRM}
interrupts the current system call. Its default action is termination,
so installing a handler (or using it with \code{sigaction}) turns a
hang into a manageable event.

\begin{lstlisting}[caption={A timeout guard around an I/O read.}]
#include <stdio.h>
#include <stdlib.h>
#include <signal.h>
#include <unistd.h>
#include <errno.h>

static volatile sig_atomic_t timed_out = 0;

static void handle_alarm(int sig) {
    timed_out = 1;
}

int main(void) {
    struct sigaction sa;
    char buf[64];
    ssize_t got;

    sa.sa_handler = handle_alarm;
    sigemptyset(&sa.sa_mask);
    sa.sa_flags = 0;
    sigaction(SIGALRM, &sa, NULL);

    printf("you have 5 seconds to type a line\n");

    alarm(5);
    got = read(0, buf, sizeof(buf));     /* may be interrupted by SIGALRM */
    alarm(0);                            /* cancel the timer */

    if (timed_out) {
        printf("\ntimeout; giving up\n");
        return 1;
    }
    if (got < 0) {
        perror("read");
        return 1;
    }

    printf("you typed %zd bytes\n", got);
    return 0;
}
\end{lstlisting}

\begin{verbatim}
   The timeline when the user types nothing:

     t=0  alarm(5)  begins
     t=5  SIGALRM delivered  -->  handle_alarm() sets timed_out
     t=5  read() returns -1, errno = EINTR   (see next section)
     main sees timed_out == 1 and gives up.

   alarm(0) cancels a pending alarm. Real applications time out
   network reads, locks, and whole requests this way--the signal is
   the kernel's ticking clock.
\end{verbatim}

\section{EINTR: When a Signal Interrupts a System Call}

Here is the subtle interaction that causes the classic bug in signal
programming. While your process is blocked inside a \emph{slow}
system call (\code{read} on a pipe or terminal, \code{waitpid},
\code{sleep}), a signal arrives. The kernel must run the handler. The
only safe way is to make the system call return \emph{prematurely}, so
the handler can run in your context. The call returns \code{-1} with
\code{errno = EINTR}, ``interrupted by signal.''

\begin{verbatim}
   blocked in read()               SIGTERM arrives
        |                                |
        |  read returns -1,             v
        |  errno = EINTR          handler runs
        |                                |
        +----- your code: is EINTR an   |
              error or a "retry"?  <-----+
              It is a RETRY: the operation did not fail, it was
              postponed. Standard practice:

                 while (got < 0 && errno == EINTR) {
                     got = read(fd, buf, n);
                 }
\end{verbatim}

Two ways to cope. The manual way wraps every slow call in an
\code{EINTR} retry loop. The automatic way sets \code{SA\_RESTART} in
\code{sa.sa\_flags}, asking the kernel to restart interrupted calls
itself---but note that \code{read} on a pipe, \code{waitpid}, and a few
others are \emph{not} restartable by it, so the manual loop is the
portable truth.

\begin{lstlisting}[caption={The universal EINTR retry loop.}]
#include <unistd.h>
#include <errno.h>

/* Read up to n bytes, restarting if a signal interrupts us. */
ssize_t read_eintr(int fd, void *buf, size_t n) {
    ssize_t got;

    while (1) {
        got = read(fd, buf, n);
        if (got < 0 && errno == EINTR) {
            continue;      /* a signal arrived; retry, don't give up */
        }
        return got;        /* real result: n, k<n, 0 (EOF), or an error */
    }
}
\end{lstlisting}

\begin{osconcept}
The kernel's decision to interrupt slow system calls is what makes
asynchronous delivery \emph{possible at all}: a handler cannot run
inside a blocked call, so the kernel must unwind the call, deliver, and
let your code decide what to do next. \code{EINTR} is therefore not a
bug in your code; it is the visible trace of the kernel's interrupt
mechanism. Designing which calls restart automatically and which do
not is an operating system judgment call, and \code{SA\_RESTART} is the
kernel exposing that choice to you.
\end{osconcept}

\section{Practical: A Graceful Shutdown Server}

Let us assemble the chapter into the pattern that real daemons use:
handle \code{SIGTERM}/\code{SIGINT} by setting a flag, check it in the
main loop, and reap children via \code{SIGCHLD}. The handlers do almost
nothing; all work happens in the main loop where the full standard
library is safe.

\begin{lstlisting}[caption={server.c: flag-based shutdown plus asynchronous reaping.}]
#include <stdio.h>
#include <stdlib.h>
#include <signal.h>
#include <unistd.h>
#include <sys/types.h>
#include <sys/wait.h>

static volatile sig_atomic_t g_running = 1;
static volatile sig_atomic_t g_children = 0;

static void handle_shutdown(int sig) {
    g_running = 0;          /* ask the main loop to finish */
}

static void handle_child(int sig) {
    g_children = 1;         /* ask the main loop to reap   */
}

void reap_children(void) {
    pid_t done;
    int status;

    while (1) {
        done = waitpid(-1, &status, WNOHANG);
        if (done <= 0) {
            break;
        }
        if (WIFEXITED(status)) {
            printf("reaped child %d (status %d)\n",
                   (int) done, WEXITSTATUS(status));
        } else if (WIFSIGNALED(status)) {
            printf("reaped child %d (signal %d)\n",
                   (int) done, WTERMSIG(status));
        }
    }
}

int main(void) {
    struct sigaction sa;

    /* One handler for both termination signals. */
    sa.sa_handler = handle_shutdown;
    sigemptyset(&sa.sa_mask);
    sa.sa_flags = 0;
    sigaction(SIGINT, &sa, NULL);
    sigaction(SIGTERM, &sa, NULL);

    sa.sa_handler = handle_child;
    sigemptyset(&sa.sa_mask);
    sa.sa_flags = 0;
    sigaction(SIGCHLD, &sa, NULL);

    printf("server running (pid %d). Try: kill -TERM %d\n",
           getpid(), getpid());

    while (g_running) {
        if (g_children) {
            g_children = 0;
            reap_children();
        }
        /* In a real server: accept a connection, fork a worker... */
        sleep(1);
    }

    printf("shutdown complete; resources cleaned up\n");
    return 0;
}
\end{lstlisting}

\begin{verbatim}
   What happens end to end:

     kill -TERM <pid>  --->  SIGTERM  --->  handle_shutdown
                                              |  g_running = 0
     the main loop's next check of g_running  |
     sees 0, exits the loop, cleans up,       |
     returns from main.  <-------------------+

     A child exits  ----->  SIGCHLD  --->  handle_child
                                              |  g_children = 1
     the main loop reaps every zombie with    |
     waitpid(-1, WNOHANG).  <----------------+

   Both handlers are two assignments. The main loop is the only place
   that prints, allocates, or calls anything complex. That is the whole
   art of signal programming.
\end{verbatim}

\begin{researchnote}[Why the flag-and-loop pattern wins]
The graceful-shutdown server you just built is \emph{event-driven}: one
loop, a few flags, and all slow work done inline. This is not merely a
style preference; it is a design conclusion with a famous champion.
John Ousterhout, the creator of Tcl/Tk and a leading systems
researcher, argued in a widely-cited essay that a single event-driven
loop is easier to make correct than a design where many concurrent
activities each block on their own.

\begin{verbatim}
   The essay's central claim (paraphrased):

       "Concurrency via threads is hard to program correctly, and for
        most applications an event-driven design -- one loop that
        checks state and handles events as they arrive -- is simpler
        and more robust."

   Translate to this chapter: our main loop checks g_running and
   g_children instead of blocking; the signal handlers only set flags.
   That is event-driven thinking applied to asynchronous events.
\end{verbatim}

Why this matters beyond signals: it is the same architecture we will
use for network servers in Chapter~17, where one loop multiplexes
thousands of connections. The claim is contested (some argue threads
scale better), but the pattern is proven and it is what our code uses.

\paperref{J.~Ousterhout, ``Why Threads Are A Bad Idea (for most
purposes),'' invited talk, USENIX 1996}{https://web.stanford.edu/~ouster/cgi-bin/papers/threads.pdf}
\end{researchnote}

\section{Structures in Bytes}

A signal is a small integer (1--31 standard, plus 34--64 real-time on
Linux), so the kernel's per-process signal state is mostly bits. The
\code{sigset\_t} you pass to \code{sigprocmask} is exactly a bit array:

\begin{verbatim}
   sigset_t on Linux/x86-64: 128 bytes = 1024 bits, one bit per signal.

   byte:   0        1        2        ...   127
          +--------+--------+--------+-- --+--------+
          |        |        |        |     |        |   bit i = signal i+1
          +--------+--------+--------+-- --+--------+
          SIGINT (2)  -> bit 1 of byte 0
          SIGKILL (9) -> bit 8: byte 1, bit 0
          SIGTERM(15) -> bit 14: byte 1, bit 6

   struct sigaction on glibc/x86-64: 152 bytes
   offset  size  field
   0       8     sa_handler      (or sa_sigaction: same storage)
   8       128   sa_mask         (a sigset_t)
   136     4     sa_flags        (SA_RESTART, SA_SIGINFO, ...)
   144     8     sa_restorer     (legacy; you never set it)
\end{verbatim}

(The exact sizes are an ABI detail; the shape is what matters: one
pointer, one bit-array, one int, one pointer.) You never touch these
bytes directly---\code{sigemptyset} and \code{sigaddset} manage the
bit array, \code{sigaction} fills the struct---but it is useful to
know that blocking a signal is just setting a bit.

\section{Summary}

\begin{verbatim}
   The signal model in one picture:

      event ------> GENERATED
                       |      blocked? yes ----> PENDING (coalesced)
                       |      blocked? no       | (delivered on unblock)
                       v                        v
                  DELIVERED ----> default action or your handler
                       |
                       +--- handler may set flags (volatile sig_atomic_t)
                       +--- handler may call only async-signal-safe fns
                       +--- slow syscalls return EINTR; retry or SA_RESTART

   The rules that keep signal code correct:
     1. install with sigaction, not signal;
     2. handlers set flags; the main loop does the work;
     3. flags are volatile sig_atomic_t;
     4. never call malloc/printf in a handler; use write/_exit;
     5. reap SIGCHLD with a WNOHANG loop;
     6. wrap slow syscalls against EINTR;
     7. SIGKILL and SIGSTOP are not yours to catch.
\end{verbatim}

\begin{osconcept}
Signals are the oldest asynchronous mechanism in Unix, and they remain
the OS's primary way of injecting events into running processes. The
whole chapter is a tour of operating system design: privileged event
delivery, per-process dispositions, a kernel-side state machine
(pending/blocked/delivered), an async-safety contract, and the
interaction of signals with blocking system calls. Master it and you
can build daemons, supervisors, and servers that behave correctly under
the chaos of the real machine.
\end{osconcept}

In the next chapter we give signals a channel: inter-process
communication through pipes, which let one process pour its output into
another---the mechanism behind every command pipeline you have ever
typed.

}{}
\IfFileExists{chapter7.tex}{% ============================================================================
%  chapter7.tex
%  Inter-Process Communication: Pipes and the Art of Wiring Processes
% ============================================================================

\chapter{Inter-Process Communication: Pipes}

\section{Why Processes Cannot Just Talk}

Chapter~5 gave each process a \emph{private} address space and told you
the kernel forbids one process from poking another's memory. That
isolation is the price of safety. But isolation would be useless if
processes could not cooperate. A compiler must feed its output to an
assembler; a web server must hand a request to a worker; a database
must talk to its replicas. The general problem is called
\emph{inter-process communication} (IPC): \emph{how do two independent
processes exchange data when they cannot touch each other's memory?}

\begin{verbatim}
   What processes CANNOT do:            What the OS provides instead:

   Process A reads Process B's heap.    channels that cross the kernel:
        A -X-> B                         pipes, files, shared memory,
                                         sockets, signals (ch.6)

   Both point at "the same" variable.   The OS copies bytes across the
   No shared addresses exist.            boundary; the kernel does the
                                         moving.
\end{verbatim}

There are five classic IPC mechanisms, each a different trade-off, and
we will meet them all by the end of the book. This chapter builds the
first and most elegant one---the \emph{pipe}---because it is the
simplest possible answer and because it is the mechanism hiding behind
every command line you have ever piped.

\begin{osconcept}
IPC is an operating system problem for a reason that flows directly
from Chapter~5: processes are isolated by the kernel, so only the
kernel can connect them. A channel must be a kernel object---something
both processes name through descriptors---because the kernel is the
only party both can trust and the only party with the power to move
bytes between private address spaces. IPC is the deliberate exception
to the isolation rule, and the OS decides exactly which exceptions are
allowed.
\end{osconcept}

\section{What a Pipe Is}

A \emph{pipe} is a unidirectional byte stream connecting two endpoints.
You create one with \code{pipe()}, which gives you a pair of
descriptors: \code{pfd[0]} is the \emph{read end}, \code{pfd[1]} is
the \emph{write end}. Bytes written to \code{pfd[1]} appear, in order,
at \code{pfd[0]}.

\begin{lstlisting}[caption={A pipe is two descriptors: read end and write end.}]
#include <stdio.h>
#include <stdlib.h>
#include <unistd.h>

int main(void) {
    int pfd[2];

    if (pipe(pfd) < 0) {
        perror("pipe");
        exit(1);
    }

    printf("read end  = fd %d\n", pfd[0]);
    printf("write end = fd %d\n", pfd[1]);

    /* Write two bytes, then close the write end to signal EOF. */
    write(pfd[1], "hi", 2);
    close(pfd[1]);

    {
        char buf[8];
        ssize_t got;

        got = read(pfd[0], buf, sizeof(buf));
        if (got < 0) {
            perror("read");
            exit(1);
        }
        printf("read %zd bytes: %.*s\n", got, (int) got, buf);
    }

    close(pfd[0]);
    return 0;
}
\end{lstlisting}

\begin{verbatim}
   The pipe is not a file. It is a KERNEL buffer with two doors:

        write end (fd[1])                read end (fd[0])
             |                                ^
             v                                |
        +--------------------------------------------+
        |  [ 'h' ][ 'i' ][ ][ ][ ][ ]  ...          |
        |         the kernel's pipe buffer           |
        +--------------------------------------------+

   - "Byte stream" means there are NO message boundaries: whatever the
     writer writes is a continuous sequence of bytes to the reader. If
     you write "ab" and "cd", the reader may see "abcd", "abc"+"d", or
     "a"+"bcd" depending on timing. Boundaries are NOT preserved.
   - The pipe is UNIDIRECTIONAL: bytes flow only from write end to
     read end. For two-way conversation you create two pipes.
   - The bytes live in the KERNEL. The processes never share memory;
     they only pass bytes across the syscall boundary.
\end{verbatim}

Two properties of this picture matter enormously. First, the pipe
synchronizes the two sides \emph{automatically}: the kernel makes one
side wait for the other. Second, the pipe has a finite capacity---on
Linux, around 64 kilobytes---and that capacity, not the file system, is
what makes blocking real.

\begin{osconcept}
A pipe is a kernel object: kernel memory for the buffer, kernel
bookkeeping for the two endpoints, and kernel rules about when a
process may proceed. Your program never sees the buffer; it sees two
integers. This is the file-descriptor model of Chapter~5 applied to
live data flow. Only the kernel can own the buffer, because only the
kernel can mediate between two isolated address spaces without copying
through either one's memory.
\end{osconcept}

\section{The Blocking Contract}

Pipes are \emph{blocking}, and that is their superpower. The kernel
makes the two ends wait for each other:

\begin{verbatim}
   EMPTY pipe:  read() blocks until bytes arrive.

        write end                 read end
             |                       v  (waiting...)
        +------------------------------------+
        |                                    |   <-- nothing to read
        +------------------------------------+
             (writer writes)
        +------------------------------------+
        | [ 'h' ][ 'i' ][ ][ ][ ][ ]        |
        +------------------------------------+
        read() wakes up and returns the bytes.

   FULL pipe:   write() blocks until there is room.

        +------------------------------------+
        | [X][X][X][X][X][X][X][X]           |   <-- no room
        +------------------------------------+
             (reader reads)
        +------------------------------------+
        | [X][X][X][X][ ][ ][ ][ ]          |
        +------------------------------------+
        write() wakes up and deposits its bytes.
\end{verbatim}

This is worth slowing down for, because it is the whole point. The
kernel is acting as a \emph{scheduler}: the reader does not busy-wait
spinning the CPU; it goes to sleep and the kernel wakes it exactly when
data is available. The writer sleeps when the buffer is full and wakes
when the reader frees space. Neither process polls, and neither needs
to know the other exists.

\begin{verbatim}
   The producer/consumer dance:

   producer:                    consumer:
     write(fd1, data, n);         got = read(fd0, buf, n);
       |                              ^
       |  if the pipe is full,        |  if the pipe is empty,
       |  the kernel parks us         |  the kernel parks us
       |  and wakes us later          |  and wakes us later
       v                              |
   bytes cross the kernel buffer, paced by whichever side is slower.
\end{verbatim}

The blocking contract is also where the classic \emph{deadlock} lives:
if two processes each write into a pipe the other is supposed to read,
and both pipes are full, each process sits blocked forever, each
waiting for the other to make room. The kernel cannot rescue you; it
only enforces the waiting.

\begin{osconcept}
Blocking I/O is the kernel's oldest concurrency solution. By parking a
process that cannot make progress, the OS wastes no CPU and needs no
coordination from the program; by waking it the instant it can proceed,
the OS delivers an efficient rendezvous. Deciding when to sleep and
when to wake---for every pipe, on every byte---is the scheduler's
problem, and it is exactly why a pipe is an OS service rather than a
library function.
\end{osconcept}

\section{fork + pipe: Wiring Two Processes Together}

A pipe inside one process is a curiosity. The pipe becomes IPC the
moment you \code{fork}: the child inherits a \emph{copy} of the
descriptor table, so both processes hold the same two ends. One process
keeps the write end, the other keeps the read end, and bytes flow
between two independent address spaces.

\begin{lstlisting}[caption={Parent writes, child reads: the canonical IPC dance.}]
#include <stdio.h>
#include <stdlib.h>
#include <unistd.h>
#include <sys/types.h>
#include <sys/wait.h>

int main(void) {
    int pfd[2];
    pid_t pid;

    if (pipe(pfd) < 0) {
        perror("pipe");
        exit(1);
    }

    pid = fork();
    if (pid < 0) {
        perror("fork");
        exit(1);
    }

    if (pid == 0) {
        /* CHILD: keep the read end, drop the write end. */
        char buf[64];
        ssize_t got;

        close(pfd[1]);

        got = read(pfd[0], buf, sizeof(buf));
        if (got < 0) {
            perror("read");
            exit(1);
        }
        printf("child read %zd bytes: %.*s\n", got, (int) got, buf);

        close(pfd[0]);
        exit(0);
    }

    /* PARENT: keep the write end, drop the read end. */
    close(pfd[0]);

    write(pfd[1], "hello from parent", 17);

    close(pfd[1]);            /* closing the last write end sends EOF */

    waitpid(pid, NULL, 0);
    return 0;
}
\end{lstlisting}

\begin{verbatim}
   The state after fork, before each side closes its unused end:

     parent                          child
   +-------------+                 +-------------+
   | fd table    |                 | fd table    |
   | 0 stdin     |                 | 0 stdin     |
   | 1 stdout    |                 | 1 stdout    |
   | 2 stderr    |                 | 2 stderr    |
   | 3 pfd[0] --+----------------->| 3 pfd[0]    |   read end
   | 4 pfd[1] --+----------------->| 4 pfd[1]    |   write end
   +-------------+                 +-------------+

   The kernel pipe object is shared: it is reference-counted, like the
   open file table. BOTH processes name it, so EITHER can read or write.

   After each side closes its unused end, the roles are clean:

     parent: [write end] =============> [read end]: child
                 \ byte stream flows through the kernel /
\end{verbatim}

The discipline here is absolute and it is the number-one pipe bug in
all of systems programming: \emph{close the end you do not use, in
every process}. If the parent keeps the read end open while the child
waits for EOF, the EOF never arrives (the kernel still counts the
parent's copy as a reader), and the child hangs forever.

\begin{ruleofthumb}[Close both ends you don't use, in both processes]
After \code{fork}, each process must close the pipe end it will not
use. Leaving a spare copy of the read end open is the classic way to
make a reader wait forever for an EOF that can never come. Draw the
descriptor table for every process in your program; there are never
extra descriptors on purpose.
\end{ruleofthumb}

\section{The EOF Rule}

When does \code{read} on a pipe return 0, signaling end of stream? The
rule is precise: \emph{when every copy of the write end is closed}.
EOF is not ``the writer has nothing more to say right now''; it is
``no writer exists anymore, anywhere.''

\begin{verbatim}
   EOF rule: read() returns 0 only when NO write end remains open.

   writer open                            all writers gone
        |                                      |
        v                                      v
   read() returns bytes...            read() returns 0  ("end of stream")

   This is why the child in the previous program could finish: the
   parent wrote its bytes AND closed pfd[1]. If the parent had left
   pfd[1] open, the child's read() would sit forever, even after the
   last byte, because a writer still existed.
\end{verbatim}

The EOF rule is what makes pipes compose into pipelines: a downstream
process knows the input is finished because the upstream process
\emph{closed its end}. The closing, not any special marker byte, is
the signal.

\begin{osconcept}
The EOF rule is a kernel decision about reference counting. The kernel
tracks how many open descriptors name each pipe endpoint, and it is the
one authority that knows when the count reaches zero. Your program
cannot ``force'' an EOF; it can only close descriptors and let the
kernel draw the conclusion. This is descriptor-table bookkeeping
exposed as a data-flow event, and it is what lets processes signal
completion without any extra protocol.
\end{osconcept}

\section{Pipelines: cmd1 | cmd2}

With the EOF rule and \code{dup2} from Chapter~5 in hand, we can build
the thing that made Unix famous: the \emph{pipeline}. ``Run the first
program with its stdout connected to the second program's stdin.''
The recipe has exactly four moves: create a pipe, fork the first
command with its stdout re-pointed at the write end, fork the second
with its stdin re-pointed at the read end, and have the parent close
both pipe ends and wait.

\begin{lstlisting}[caption={A two-stage pipeline, end to end.}]
#include <stdio.h>
#include <stdlib.h>
#include <unistd.h>
#include <sys/types.h>
#include <sys/wait.h>

/* Run "prog1..." with its stdout piped into "prog2...". */
int run_pipeline(char *const argv1[], char *const argv2[]) {
    int pfd[2];
    pid_t p1;
    pid_t p2;

    if (pipe(pfd) < 0) {
        perror("pipe");
        return -1;
    }

    p1 = fork();
    if (p1 < 0) {
        perror("fork");
        return -1;
    }
    if (p1 == 0) {
        /* First command: stdout must name the pipe's write end. */
        close(pfd[0]);
        if (dup2(pfd[1], 1) < 0) {
            perror("dup2");
            exit(1);
        }
        close(pfd[1]);
        execvp(argv1[0], argv1);
        perror("execvp");
        exit(127);
    }

    p2 = fork();
    if (p2 < 0) {
        perror("fork");
        return -1;
    }
    if (p2 == 0) {
        /* Second command: stdin must name the pipe's read end. */
        close(pfd[1]);
        if (dup2(pfd[0], 0) < 0) {
            perror("dup2");
            exit(1);
        }
        close(pfd[0]);
        execvp(argv2[0], argv2);
        perror("execvp");
        exit(127);
    }

    /* The parent is only the referee: it needs neither pipe end. */
    close(pfd[0]);
    close(pfd[1]);

    waitpid(p1, NULL, 0);
    waitpid(p2, NULL, 0);
    return 0;
}

int main(int argc, char *argv[]) {
    char *cmd1[3];
    char *cmd2[3];

    if (argc != 5) {
        fprintf(stderr, "usage: %s prog1 a1 prog2 a2\n", argv[0]);
        return 1;
    }
    cmd1[0] = argv[1];
    cmd1[1] = argv[2];
    cmd1[2] = NULL;
    cmd2[0] = argv[3];
    cmd2[1] = argv[4];
    cmd2[2] = NULL;

    return run_pipeline(cmd1, cmd2);
}
\end{lstlisting}

\begin{verbatim}
   The wiring, drawn process by process:

     p1 (child 1)                  p2 (child 2)
     fd 1 --> pipe write end       fd 0 --> pipe read end
           \                        /
            \  +----------------+ /
             \ | kernel buffer  |/
              >| pfd[1] -> pfd[0]|
              <|                |<   bytes flow p1 -> p2
               +----------------+

     parent: closes both ends; just waitpid()s both children.
     dup2() copies the pipe descriptor ONTO descriptors 1 and 0, so
     exec'd programs write to stdout / read from stdin without knowing
     a pipe exists -- Chapter 5's redirection, exactly as the shell
     does it.
\end{verbatim}

The magic ingredient is that \code{exec} preserves open descriptors:
after \code{dup2} and \code{close}, the child \code{exec}s a totally
new program whose \code{stdout} is the pipe, with no knowledge of the
plumbing. That is why every Unix tool composes with every other tool:
each reads \code{stdin} and writes \code{stdout}, and the shell---or
your program---rewires the descriptors.

\begin{verbatim}
   The full lifecycle of "cat f.txt | sort":

     1. shell pipe()s -> pfd = [3, 4]
     2. shell forks child A: dup2(4, 1), close(3), close(4), exec cat
        (cat's stdout is now the pipe)
     3. shell forks child B: dup2(3, 0), close(3), close(4), exec sort
        (sort's stdin is now the pipe)
     4. shell closes 3 and 4, then waits for both
     5. cat writes f.txt into the pipe; sort reads it, sorts, prints;
        when cat exits and its stdout closes, the last write end dies,
        sort's read() hits EOF, sort prints and exits, shell reaps both.
\end{verbatim}

\begin{osconcept}
Pipelines work because the OS treats descriptors as \emph{rewireable
connections} whose ownership can be transferred between processes and
onto the standard slots. The kernel's descriptor-copying primitive
(\code{dup2}) plus the EOF rule are the entire protocol---no message
headers, no handshaking. Teaching processes to ``just write to fd 1''
and letting the kernel route it is a design decision so good that it
became the model for every later IPC mechanism.
\end{osconcept}

\section{SIGPIPE: When the Channel Dies}

What happens if the reader is gone? Suppose the second command in a
pipeline exits early, closing the read end, while the first command
keeps writing. The kernel detects that a write has no readers and does
two things: it sends \code{SIGPIPE} to the writer, and if the signal
is ignored or handled, the \code{write} returns \code{-1} with
\code{errno = EPIPE}.

\begin{verbatim}
   reader closes its read end, exits
        |
        v
   writer calls write(fd1, ...)
        |
        +--> kernel finds no readers
        |
        +--> raises SIGPIPE on the writer
        |      default action: terminate the writer (like `yes | head`)
        |
        +--> if SIGPIPE ignored/handled:
             write() returns -1, errno = EPIPE

   This is why `yes | head -1` terminates `yes` promptly: `head` closes
   the pipe after one line, `yes` writes, the kernel SIGPIPEs it.
\end{verbatim}

For a well-behaved server or filter, the sensible choice is to ignore
\code{SIGPIPE} and treat \code{EPIPE} as ``the peer went away''---an
ordinary, expected condition, not a crash.

\begin{lstlisting}[caption={Handling a dead reader gracefully.}]
#include <stdio.h>
#include <stdlib.h>
#include <signal.h>
#include <unistd.h>
#include <errno.h>

int main(void) {
    struct sigaction sa;
    char msg[16];
    int i;

    /* The reader may vanish; treat that as EPIPE, not death. */
    sa.sa_handler = SIG_IGN;
    sigemptyset(&sa.sa_mask);
    sa.sa_flags = 0;
    sigaction(SIGPIPE, &sa, NULL);

    /* Write a line repeatedly until the reader goes away. */
    for (i = 0; i < 100000; i++) {
        ssize_t r;

        snprintf(msg, sizeof(msg), "line %d\n", i);
        r = write(1, msg, strlen(msg));
        if (r < 0) {
            if (errno == EPIPE) {
                printf("reader is gone; stopping politely\n");
                break;
            }
            perror("write");
            break;
        }
    }
    return 0;
}
\end{lstlisting}

\begin{osconcept}
\code{SIGPIPE} is the kernel's way of converting a dead channel into
an interruptible event. The OS cannot know whether a missing reader is
a bug or an expected teardown, so it gives the writer the choice:
die by default, or catch the signal and inspect \code{EPIPE}. Choosing
the default is the reason so many quick pipelines terminate cleanly;
choosing the handler is what production services do so that a client
disconnecting is not a crash.
\end{osconcept}

\section{Named Pipes (FIFOs)}

An unnamed pipe lives and dies with the processes that created it. A
\emph{named pipe}, or \emph{FIFO}, is a pipe with a name in the file
system. Any process that can open the path can connect to it, so two
unrelated processes---neither forking the other---can talk through a
FIFO. You create one with \code{mkfifo} (the program or the syscall);
you open it like a file; the data still flows through a kernel buffer,
not through the file system.

\begin{lstlisting}[caption={Creating and opening a FIFO.}]
#include <stdio.h>
#include <stdlib.h>
#include <sys/stat.h>
#include <fcntl.h>
#include <unistd.h>

int main(int argc, char *argv[]) {
    int fd;

    if (argc != 2) {
        fprintf(stderr, "usage: %s fifo_name\n", argv[0]);
        return 1;
    }

    /* Create the FIFO if it does not exist yet. */
    if (mkfifo(argv[1], 0644) < 0) {
        perror("mkfifo");   /* fine if it already exists (EEXIST) */
    }

    printf("opening %s for writing (blocks until a reader opens it)...\n",
           argv[1]);
    fd = open(argv[1], O_WRONLY);
    if (fd < 0) {
        perror("open");
        return 1;
    }

    write(fd, "hello through a named pipe", 26);
    close(fd);
    return 0;
}
\end{lstlisting}

\begin{verbatim}
   FIFO = a pipe with an address in the file namespace:

      $ mkfifo /tmp/chat
      $ ./writer /tmp/chat &     (blocks: no reader yet)
      $ cat /tmp/chat            (reader opens the other end)
      hello through a named pipe

   - The FIFO is a directory entry, but it holds NO data of its own.
     Bytes go through the kernel pipe buffer; the path is just the
     meeting point.
   - open() BLOCKS until both a reader and a writer have opened the
     FIFO. That rendezvous is the FIFO's version of pipe() + fork().
   - When all readers or writers close, the pipe state resets: a later
     open() gets a fresh buffer, so data written earlier is gone.
\end{verbatim}

\begin{osconcept}
A FIFO is the file system serving as a \emph{namespace} for IPC. The
kernel already owns the directory tree, permissions, and lookup; giving
a pipe a path reuses all of that machinery so that unrelated processes
can rendezvous by name. Deciding that the path is only an address---not
a store of data---is an OS design choice that keeps FIFOs cheap and
their semantics identical to anonymous pipes.
\end{osconcept}

\section{The IPC Menu}

Pipes are one of five classic channels. The others appear later in the
book; a map now will save you confusion when we meet them.

\begin{verbatim}
   mechanism        data model          who can talk        we meet it
   --------------   ----------------    ----------------    ----------
   pipe / FIFO      byte stream,        parent-child or     this chapter
                    kernel-buffered     any named pair
   signal           tiny events         any processes       chapter 5
                    (an integer)
   shared memory    raw bytes, no       processes that      chapter 11
   (mmap)           kernel copy         agree on a region
   message queue    typed messages,     any processes       (surveyed)
                    kernel-buffered     with permissions
   socket           byte stream or      any processes,      chapter 16
                    datagrams           even across hosts

   The trade-off in one line: pipes and sockets BUFFER and COPY (safe,
   automatic), while shared memory does NOT (fast, dangerous). You
   choose based on whether you want the kernel to babysit your data.
\end{verbatim}

\begin{verbatim}
   The rule of thumb for this book's projects:

     one writer, one reader, stream data  -> pipe
     unrelated processes, named channel   -> FIFO
     lots of small structured messages    -> message queue
     fast shared state, both are careful  -> shared memory (ch.12)
     any machine, any distance            -> sockets (ch.17)
\end{verbatim}

\begin{researchnote}[Pipes are communication as mathematics]
The pipe is not just a Unix convenience; it is the direct implementation
of a celebrated idea. In 1978, C.~A.~R.~Hoare published
``Communicating Sequential Processes,'' arguing that the cleanest way
to think about concurrency is as \emph{processes that exchange data by
message passing}, with no shared state at all. A pipe is exactly that:
two processes, one channel, bytes in order.

\begin{verbatim}
   Hoare's thesis, condensed to one line:

       "Programs are built from sequential processes that interact
        solely by input and output; shared memory is unnecessary and
        harmful."

   This is the philosophy the pipe implements: the writer and reader
   never share a variable, never lock anything, and are correct purely
   because the channel delivers bytes in order. Hoare's later book of
   the same name turned this into a formal language for reasoning
   about concurrency.
\end{verbatim}

Why it matters to you: the same idea drives most real servers
(Chapter~17) and the concurrency chapter (Chapter~11). When a design
uses pipes or sockets to connect worker processes, it is inheriting a
four-decade-old research result about how to make concurrency
\emph{provable} instead of merely probable.

\paperref{C.~A.~R.~Hoare, ``Communicating Sequential Processes,''
Communications of the ACM 21(8), 1978}{https://www.cs.cmu.edu/~crary/819-f09/Hoare78.pdf}
\end{researchnote}

\section{Practical: A Two-Stage Filter Shell}

Let us wire the chapter together into a small program that behaves like
a real shell pipeline: read a command line like
\code{prog1 arg prog2 arg}, split it, and run it through a pipe---the
exact four-move recipe, with every error checked and every unused
descriptor closed.

\begin{lstlisting}[caption={piperun.c: parse two commands and pipe them together.}]
#include <stdio.h>
#include <stdlib.h>
#include <string.h>
#include <unistd.h>
#include <sys/types.h>
#include <sys/wait.h>

#define MAX_ARG 32

/* Split "cmd arg1 arg2 | cmd arg1" into two NULL-terminated argv arrays. */
int parse_pipeline(char *line, char ***left, char ***right) {
    char *pipe_pos;
    char **a;
    char *tok;
    int n;

    pipe_pos = strchr(line, '|');
    if (pipe_pos == NULL) {
        return -1;      /* no pipe found */
    }
    *pipe_pos = '\0';

    a = malloc(MAX_ARG * sizeof(char *));
    if (a == NULL) {
        return -1;
    }
    n = 0;
    tok = strtok(line, " \t\n");
    while (tok != NULL && n < MAX_ARG - 1) {
        a[n] = tok;
        n = n + 1;
        tok = strtok(NULL, " \t\n");
    }
    a[n] = NULL;
    *left = a;

    a = malloc(MAX_ARG * sizeof(char *));
    if (a == NULL) {
        return -1;
    }
    n = 0;
    tok = strtok(pipe_pos + 1, " \t\n");
    while (tok != NULL && n < MAX_ARG - 1) {
        a[n] = tok;
        n = n + 1;
        tok = strtok(NULL, " \t\n");
    }
    a[n] = NULL;
    *right = a;

    return 0;
}

/* Run the pipeline built by parse_pipeline. */
int run_pipeline(char *const left[], char *const right[]) {
    int pfd[2];
    pid_t p1;
    pid_t p2;

    if (pipe(pfd) < 0) {
        perror("pipe");
        return -1;
    }

    p1 = fork();
    if (p1 < 0) {
        perror("fork");
        return -1;
    }
    if (p1 == 0) {
        close(pfd[0]);
        dup2(pfd[1], 1);
        close(pfd[1]);
        execvp(left[0], left);
        perror("execvp");
        exit(127);
    }

    p2 = fork();
    if (p2 < 0) {
        perror("fork");
        return -1;
    }
    if (p2 == 0) {
        close(pfd[1]);
        dup2(pfd[0], 0);
        close(pfd[0]);
        execvp(right[0], right);
        perror("execvp");
        exit(127);
    }

    close(pfd[0]);
    close(pfd[1]);

    waitpid(p1, NULL, 0);
    waitpid(p2, NULL, 0);
    return 0;
}

int main(void) {
    char line[512];
    char **left;
    char **right;

    printf("piperun> ");
    fflush(stdout);
    while (fgets(line, sizeof(line), stdin) != NULL) {
        if (parse_pipeline(line, &left, &right) == 0) {
            run_pipeline(left, right);
            free(left);
            free(right);
        } else {
            printf("need two commands separated by a pipe\n");
        }
        printf("piperun> ");
        fflush(stdout);
    }
    return 0;
}
\end{lstlisting}

\begin{verbatim}
   Try it:

     $ ./piperun
     piperun> echo hello world | tr a-z A-Z
     HELLO WORLD
     piperun> seq 10 | tail -3
     8
     9
     10
     piperun> ^D

   Every line you typed exercised the whole chapter: a kernel buffer,
   two forks, two dup2 rewirings, four closes, and two waits. That is
   a shell's pipeline in miniature.
\end{verbatim}

\section{Structures in Bytes}

A pipe is a kernel object, so its byte structure lives in the kernel;
in your process a pipe is just two small integers. The object they
name looks like this (Linux):

\begin{verbatim}
   the pipe: a ring buffer of 16 pages = 64 KB.

   +---------+---------+---------+-- ... --+---------+
   | page 0  | page 1  | page 2  |          | page 15 |
   +---------+---------+---------+-- ... --+---------+
        ^ read position               ^ write position
        (byte offsets into the ring)

   the two file objects name the same pipe:

   file object (read end)  \   +---------------------+
   file object (write end) /-->| pipe: 16 pages      |
                              | read pos            |
                              | write pos           |
                              | nreaders, nwriters  |  <- the EOF rule
                              +---------------------+
\end{verbatim}

Writes land at the write position, wrap around the ring, and advance
it; reads do the same at the read position. The reference counts are
the byte-level truth behind the EOF rule: \code{read()} returns 0 not
because a marker byte exists, but because the \code{nwriters} counter
reached zero. The ring is also why blocking is cheap---a writer with no
room sleeps until a reader advances the read position.

\section{Summary}

\begin{verbatim}
   The pipe model, complete:

   pipe()      : one call, two descriptors, one kernel buffer
   fork()      : duplicates the descriptors into both processes
   close()     : each side drops the end it does not use
   blocking    : the kernel parks readers on empty, writers on full
   EOF         : read() returns 0 when the last write end closes
   dup2()      : rewires a descriptor, building pipelines
   SIGPIPE     : the kernel tells a writer its reader is gone
   FIFO        : the same channel, given a name in the file system

   Rules that keep pipe programs correct:
     1. close the unused end in EVERY process (the EOF killer);
     2. the parent (or middleman) closes BOTH ends;
     3. remember pipes are byte streams: no message boundaries;
     4. expect blocking: it is the feature, not a bug;
     5. handle SIGPIPE (or EPIPE) when a peer may disconnect.
\end{verbatim}

\begin{osconcept}
This chapter was a complete answer to ``how do isolated processes
talk?'' The OS answers: through channels it alone owns---buffers that
synchronize, reference counts that produce EOF, signals that report a
dead peer, and names that let strangers rendezvous. Every later IPC
mechanism is a variation on these same ideas, and the discipline
(close your ends, check your returns) never changes.
\end{osconcept}

The pipes we just built are \emph{synchronous}: the writer is paced by
the reader. In the chapters ahead we make communication more powerful
and more dangerous: threads sharing one address space (Chapter~11),
shared memory with no kernel copy (Chapter~12), and finally sockets
that cross machine boundaries (Chapter~17).

}{}
\IfFileExists{chapter8.tex}{% ============================================================================
%  chapter8.tex
%  Filesystems in Depth: Metadata, Links, and Directories
% ============================================================================

\chapter{Filesystems in Depth}

\section{The Filesystem Is a Namespace}

Chapter~5 treated files as ``a sequence of bytes you can read or
write,'' and that was the right first cut. Now we go one layer deeper:
a \emph{filesystem} is the kernel's system for \emph{naming,
organizing, and persisting} those byte sequences. Understanding it
means answering three questions the earlier chapter deferred: where do
the bytes physically live, what is a \emph{name}, and what does a
\emph{directory} really contain?

The three answers, in one picture:

\begin{verbatim}
   A filesystem has three layers:

     NAMES        a tree of directory entries:   /home/you/data.txt
        |
        v
     INODES       per-file metadata + a map of where the data is:
                  { type, permissions, size, [block 4][block 9]... }
        |
        v
     BLOCKS       fixed-size chunks of disk (or memory) that hold
                  the bytes themselves.

   - A NAME is a directory entry: a string in a directory file.
   - An INODE ("index node") is the kernel's record for ONE file:
     its type, its permissions, its size, and pointers to its blocks.
   - A BLOCK is the actual storage, typically 4 KB.
\end{verbatim}

The term to keep: a filesystem is a \emph{namespace}. The tree of
paths you type is the interface; behind it sits a machinery of inodes
and blocks that the kernel maintains and you almost never see. When we
say ``the file system,'' we mean this whole apparatus---not just the
directory tree.

\begin{osconcept}
Persistence and naming are an OS problem for the same reason every
service so far has been: only the kernel can give programs a
\emph{shared, protected, persistent} namespace. A library cannot make a
name survive a reboot; the filesystem does, because the kernel owns
the storage devices and the mapping from names to blocks. Providing
that namespace---and keeping it consistent across crashes and
processes---is one of the oldest responsibilities of an operating
system.
\end{osconcept}

\section{stat: Reading the Inode's Record}

The kernel exposes the contents of an inode through the \code{stat}
system call. It fills a \code{struct stat} with every fact about the
file that is not its data: its type, its permissions, its size, how
many hard links point to it, who owns it, and when it was touched.

\begin{lstlisting}[caption={stat: dump the metadata of a path.}]
#include <stdio.h>
#include <stdlib.h>
#include <sys/stat.h>
#include <sys/sysmacros.h>

int main(int argc, char *argv[]) {
    struct stat st;
    int r;

    if (argc != 2) {
        fprintf(stderr, "usage: %s path\n", argv[0]);
        return 1;
    }

    r = stat(argv[1], &st);
    if (r < 0) {
        perror("stat");
        return 1;
    }

    printf("size        : %ld bytes\n",  (long) st.st_size);
    printf("mode        : 0%o (octal)\n", (unsigned) st.st_mode);
    printf("hard links  : %ld\n",          (long) st.st_nlink);
    printf("owner       : uid=%d gid=%d\n", (int) st.st_uid, (int) st.st_gid);
    printf("device      : major=%d minor=%d\n",
           major(st.st_dev), minor(st.st_dev));
    printf("blocks      : %ld (x512 bytes)\n", (long) st.st_blocks);

    return 0;
}
\end{lstlisting}

\begin{verbatim}
   What each field is (the inode's record, in plain words):

     st_mode    a bit field packing BOTH the file type and the
                permissions. Two kinds of info in one integer.
     st_nlink   how many directory entries name this same inode.
                A file with no name left has nlink == 0 and is
                actually deleted (or waiting for its last close).
     st_size    the length in bytes of the DATA.
     st_uid, st_gid   who owns it (the numbers behind ls's names).
     st_mtime   last time the DATA changed.  ("m" = modified)
     st_ctime   last time the INODE metadata changed. (NOT creation)
     st_atime   last time it was read.
     st_blocks  how many 512-byte units of disk it occupies. Compare
                st_blocks*512 to st_size: a sparse file (ch.5) shows
                blocks*512 < size, because holes cost nothing.

   st_mode's two halves:

     bits 12-15 : the TYPE      bits 0-11 : the PERMISSIONS
     0100000 regular (S_IFREG)   0755     = rwx r-x r-x
     0040000 directory (S_IFDIR)
     0120000 symbolic link (S_IFLNK)
\end{verbatim}

The type bits are read through the macros \code{S\_ISREG},
\code{S\_ISDIR}, \code{S\_ISLNK}, and friends---never by comparing raw
numbers:

\begin{lstlisting}[caption={The S\_IS* macros decode the type half of st\_mode.}]
#include <stdio.h>
#include <sys/stat.h>

const char *type_name(mode_t mode) {
    if (S_ISREG(mode))  return "regular file";
    if (S_ISDIR(mode))  return "directory";
    if (S_ISLNK(mode))  return "symbolic link";
    if (S_ISCHR(mode))  return "character device";
    if (S_ISBLK(mode))  return "block device";
    if (S_ISFIFO(mode)) return "fifo";
    if (S_ISSOCK(mode)) return "socket";
    return "unknown";
}
\end{lstlisting}

One warning before we go on: there are \emph{two} stat-like calls.
\code{stat(path, ...)} follows a symbolic link and reports the
\emph{target's} inode. \code{lstat(path, ...)} reports the link itself.
The difference matters the moment we walk directories, because a
directory may contain links we must not chase.

\begin{osconcept}
An inode is kernel state: one per file, owned by the kernel, updated
by the kernel (permissions, timestamps, link counts) as the world
changes around you. \code{stat} is a read-only window into that state.
Only the kernel can maintain the mapping from names to inodes to
blocks, so ``what is this file?'' is answered by the OS, not by your
program.
\end{osconcept}

\section{File Types}

The type half of \code{st\_mode} distinguishes seven kinds of object,
and every one of them is ``a file'' in the eyes of the namespace.
This is the design decision of Chapter~5 paying off: a uniform
interface over wildly different objects.

\begin{verbatim}
   the seven types, as the first character of `ls -l`:

     -  regular file     bytes of data (the kind you usually mean)
     d  directory        a file whose bytes ARE directory entries
     l  symbolic link    a file whose bytes are a target path
     c  character device a device that transfers bytes one at a time
                          (terminals, /dev/null)
     b  block device     a device that transfers in blocks (disks)
     p  fifo             a named pipe from Chapter 7
     s  socket           a network endpoint from Chapter 17

   `ls -l /dev/null` shows:  c...
   `ls -ld /home`     shows:  d...
   `ls -l /proc`      shows:  d...  (and inside it, virtual files)
\end{verbatim}

\begin{verbatim}
   The surprising consequence:

     A directory is a FILE. Its "data" is a list of names -> inode
     numbers. `read()` a directory directly is awkward and
     discouraged, which is why we will use the opendir/readdir API
     (Section 8.5). But underneath, a directory is just bytes the
     kernel happens to interpret as entries.
\end{verbatim}

\begin{osconcept}
Treating devices, pipes, and directories as file types is the
\emph{everything is a file} doctrine, and it is an OS design decision
of enormous power. Because every object in the namespace answers the
same \code{stat}/\code{open}/\code{read} interface, tools never need
to know what they are touching. The kernel's job is to translate the
generic calls into whatever the underlying object actually needs---
hardware for a device, bytes for a regular file.
\end{osconcept}

\section{Permissions}

The low 12 bits of \code{st\_mode} describe who may do what. There are
three classes---user, group, other---and three operations per class---
read, write, execute. In octal, each class is one digit: 4 for read,
2 for write, 1 for execute, summed.

\begin{verbatim}
   mode 0644, drawn bit by bit:

       octal:     0     6       4       4
             special  u w x   g r w x   o r w x   (class: rwx)
       bits:  sticky  | | |    | | |    | | |
                     user     group    other
       0644 = rw- r-- r--

   The three special bits (not shown here) are setuid, setgid, and
   the sticky bit -- we mention them only to note they exist.

   mode 0755 (a program): rwx r-x r-x   -- owner may edit, all may run
   mode 0600 (a secret):  rw- --- ---   -- only the owner
   mode 0777 (a warning): rwx rwx rwx   -- everyone everything
\end{verbatim}

The kernel enforces these on every \code{open}: it compares the
requester's uid/gid against the file's owner and group, picks the
matching class, and checks the appropriate bits. This check happens
\emph{in the kernel}, at \code{open} time, once---which is why a file
can be opened and then its permissions changed without breaking the
already-open descriptor.

\begin{verbatim}
   Where the bits are consulted, and who changes them:

     open("f", O_RDONLY)   -> kernel checks read bits for your class
     open("f", O_WRONLY)   -> kernel checks write bits
     chmod("f", 0600)      -> you change the bits (if you own it)
     umask(022)            -> your process asks the kernel to mask
                              off bits from every new file's mode

   The umask explains a common mystery: you open() with mode 0666 and
   get 0644. The kernel ANDed your mode with ~umask (022) at open
   time. The kernel, not the library, owns this masking.
\end{verbatim}

\begin{lstlisting}[caption={Decoding permissions bit by bit, like ls does.}]
#include <stdio.h>
#include <sys/stat.h>

char file_type_char(mode_t mode) {
    if (S_ISREG(mode))  return '-';
    if (S_ISDIR(mode))  return 'd';
    if (S_ISLNK(mode))  return 'l';
    if (S_ISCHR(mode))  return 'c';
    if (S_ISBLK(mode))  return 'b';
    if (S_ISFIFO(mode)) return 'p';
    if (S_ISSOCK(mode)) return 's';
    return '?';
}

/* Print the 10-character ls -l permissions field. */
void print_mode_bits(mode_t mode) {
    const char *rwx = "rwxrwxrwx";
    int i;

    putchar(file_type_char(mode));

    /* Bit 8 = user read ... bit 0 = other execute. */
    for (i = 0; i < 9; i++) {
        int bit;

        bit = 8 - i;
        if (mode & (1 << bit)) {
            putchar(rwx[i]);
        } else {
            putchar('-');
        }
    }
}
\end{lstlisting}

\begin{osconcept}
The permission model is enforced by the kernel because users cannot be
trusted to enforce it on each other, and because the check must happen
at the exact moment of every \code{open}. The kernel compares real and
effective ids, walks the three classes, and refuses---returning
\code{EACCES}---before any data moves. Making permission checks part of
the syscall boundary is what makes multi-user sharing safe at all.
\end{osconcept}

\section{Links: Two Kinds of Indirection}

The \code{st\_nlink} field in our \code{stat} dump is a count of
\emph{hard links}. Understanding the two kinds of link is the key to
the inode model, and both kinds come from a single idea: a \emph{name}
and a \emph{file} are different things, and several names can point at
one file.

\subsection{Hard Links}

A hard link is a second name for the \emph{same inode}. Creating one
increments \code{st\_nlink}; deleting one decrements it. The data is
freed only when the count reaches zero. The two names are perfectly
equal---there is no ``real'' file and ``link.'' When you create a file,
\code{st\_nlink} is 1; \code{link()} makes it 2.

\begin{lstlisting}[caption={Creating both kinds of link.}]
#include <stdio.h>
#include <unistd.h>

int main(void) {
    if (link("data.txt", "data_hard.txt") < 0) {
        perror("link");
    }
    if (symlink("data.txt", "data_soft.txt") < 0) {
        perror("symlink");
    }
    return 0;
}
\end{lstlisting}

\begin{verbatim}
   HARD LINK: two directory entries, ONE inode.

   directory                   the single inode (nlink = 2)
   +--------------------+       +------------------------+
   | data.txt ----------+------>| type, perms, size,     |
   | data_hard.txt -----+       | and pointers to the    |
   +--------------------+       | actual data blocks     |
                                +------------------------+

   - Both names are the SAME file. Renaming one leaves the other.
   - Deleting data.txt only decrements nlink to 1; data survives.
   - Hard links CANNOT cross filesystems (they are inode numbers,
     and inode numbers are only meaningful within one filesystem),
     and cannot point at a directory.
\end{verbatim}

\subsection{Symbolic Links}

A symbolic link (symlink) is a different beast: a tiny file whose
\emph{content is a path string}. It is a pointer in the namespace, not
a second name for an inode. The kernel resolves the path whenever you
open through it. That is why a symlink \emph{can} cross filesystems,
\emph{can} point at a directory, and \emph{dangles} when its target
vanishes.

\begin{verbatim}
   SYMBOLIC LINK: a name containing a path.

   directory                       inodes
   +--------------------+       +------------------------+
   | data.txt ----------+------>| inode for data.txt     |
   | data_soft.txt -----+       +------------------------+
   +--------------------+             ^
        "data_soft.txt" is a file whose bytes are the STRING
        "data.txt". Opening it makes the kernel open data.txt.

   - stat()  follows the link (reports the target).
   - lstat() reports the link itself.
   - Rename or delete data.txt, and data_soft.txt still "exists"
     but resolves to nothing: a dangling link.
\end{verbatim}

\begin{osconcept}
Links are kernel bookkeeping about the namespace. A hard link is a
reference count on an inode; a symlink is a small file containing a
path. Only the kernel can maintain the inode table, resolve paths,
and decide when a file's data may finally be reclaimed. The \emph{nlink}
count in \code{stat} is literally a kernel-owned counter that your
program can only read.
\end{osconcept}

\section{Directories and Traversal}

A directory is a file whose contents are \emph{directory entries}: a
list of (name, inode number) pairs. To read it you use the directory
stream API: \code{opendir} opens a handle, \code{readdir} returns one
entry at a time, \code{closedir} releases it.

\begin{lstlisting}[caption={Listing a directory, entry by entry.}]
#include <stdio.h>
#include <stdlib.h>
#include <dirent.h>

int list_dir(const char *path) {
    DIR *d;
    struct dirent *e;

    d = opendir(path);
    if (d == NULL) {
        perror("opendir");
        return -1;
    }

    while (1) {
        e = readdir(d);
        if (e == NULL) {
            break;
        }
        if (e->d_name[0] == '.') {
            continue;    /* skip "." and ".." (and dotfiles here) */
        }
        printf("%s\n", e->d_name);
    }

    closedir(d);
    return 0;
}
\end{lstlisting}

\begin{verbatim}
   What readdir() hands you, one struct per entry:

     struct dirent {
         char d_name[256];   the entry's name
         ...                 (the inode number, and a type hint that
                              MAY be present -- never rely on it)
     };

   - Every directory contains "." (itself) and ".." (its parent);
     real listings skip them.
   - readdir gives you only NAMES. To get size/permissions you must
     stat() each one -- that is exactly why `ls -l` reads the whole
     directory AND stats every entry.
   - The order of entries is not meaningful and not sorted.
\end{verbatim}

Recursing into subdirectories turns a lister into a tree walker, and
that is where \code{lstat} becomes essential: using \code{stat} here
would follow symlinks and could loop forever on a link that points at
an ancestor.

\begin{lstlisting}[caption={walk.c: a recursive tree walker, using lstat.}]
#include <stdio.h>
#include <stdlib.h>
#include <string.h>
#include <dirent.h>
#include <sys/stat.h>

/* Recursively print every entry under base, indented by depth. */
void walk(const char *base, int depth) {
    DIR *d;
    struct dirent *e;
    int i;

    d = opendir(base);
    if (d == NULL) {
        return;
    }

    while (1) {
        struct stat st;
        char path[1024];
        int plen;

        e = readdir(d);
        if (e == NULL) {
            break;
        }
        if (strcmp(e->d_name, ".") == 0 || strcmp(e->d_name, "..") == 0) {
            continue;
        }

        plen = snprintf(path, sizeof(path), "%s/%s", base, e->d_name);
        if (plen < 0 || plen >= (int) sizeof(path)) {
            continue;          /* name too long; skip this entry */
        }

        /* lstat, not stat: a symlink must be reported AS a link. */
        if (lstat(path, &st) < 0) {
            continue;
        }

        for (i = 0; i < depth; i++) {
            printf("  ");
        }
        printf("%s\n", e->d_name);

        if (S_ISDIR(st.st_mode)) {
            walk(path, depth + 1);
        }
    }

    closedir(d);
}

int main(int argc, char *argv[]) {
    const char *root;

    root = (argc >= 2) ? argv[1] : ".";
    walk(root, 0);
    return 0;
}
\end{lstlisting}

\begin{verbatim}
   Output shape for a small tree:

     .
       src
         main.c
         util.c
       README.md
       build.sh

   Every level is an opendir/readdir/closedir cycle; every entry is
   lstat()ed. That is precisely the loop `find` and `du` run.
\end{verbatim}

\begin{osconcept}
Directory traversal is kernel-managed iteration. The directory is a
file whose format the kernel owns, so user code reads it only through
the kernel's \code{readdir} interface---never by guessing at bytes.
The \code{lstat}-vs-\code{stat} distinction exists because the kernel
must honor your choice of whether to follow a link, and getting that
choice wrong (as \code{stat} does in a walker) is how infinite
recursion and path loops enter a program.
\end{osconcept}

\section{/proc: The Virtual Filesystem}

\textbf{Concept:} a \emph{virtual filesystem} is a filesystem whose
``files'' are not stored anywhere---their contents are \emph{generated
on the fly} by the kernel when you read them. Linux's \code{/proc} is
the canonical example. Every running process has a directory
\code{/proc/<pid>} full of files describing it, manufactured from
kernel data structures the moment you open them.

\begin{verbatim}
   /proc is not a disk. It is a window into the kernel:

     /proc/self/status     the calling process's state (parent,
                           threads, memory, ...)
     /proc/self/cmdline    the command line that started it
     /proc/self/fd/        a directory of its open descriptors
     /proc/<pid>/fd/       the same, for any other process
     /proc/meminfo         whole-system memory numbers
     /proc/cpuinfo         what the CPUs are
     /proc/net/tcp         the kernel's TCP connection table

   Reading them is just read(): the bytes are composed on demand.
   They answer the OS questions of Chapters 4-6 from inside a C
   program -- without signals, pipes, or privilege tricks.
\end{verbatim}

\begin{lstlisting}[caption={Reading the kernel through a virtual file.}]
#include <stdio.h>
#include <stdlib.h>
#include <fcntl.h>
#include <unistd.h>

int main(void) {
    int fd;
    char buf[512];
    ssize_t got;

    fd = open("/proc/self/status", O_RDONLY);
    if (fd < 0) {
        perror("open");
        return 1;
    }

    got = read(fd, buf, sizeof(buf) - 1);
    if (got < 0) {
        perror("read");
        return 1;
    }
    buf[got] = '\0';
    printf("%s\n", buf);

    close(fd);
    return 0;
}
\end{lstlisting}

The output is ordinary text (its first lines name the pid and parent,
the thread count, and the memory sizes). Reading it is nothing but the
Chapter~5 \code{open}/\code{read}/\code{close} pattern applied to a
file that the kernel creates from its own tables---which is why
\code{/proc} is a wonderful teaching tool and a beloved debugging
instrument for systems programmers.

\begin{osconcept}
\code{/proc} shows the kernel \emph{reusing the file abstraction for
introspection}. The same namespace machinery, inodes, and read path
now serve as a read-only API to kernel state. This is an OS design
choice: rather than inventing a new interface for ``ask the kernel
about itself,'' Linux made the kernel a filesystem. It only works
because the file abstraction is general enough to serve both disks and
data structures.
\end{osconcept}

\begin{researchnote}[The paper that made filesystems fast]
The metadata and layout story of this chapter has a famous ancestor.
In 1984, Marshall Kirk McKusick and colleagues at Berkeley published
``A Fast File System for UNIX,'' describing the redesign of the
original Unix filesystem. Its innovations are precisely the ideas we
have been handling all chapter: separating inode metadata from data
blocks, grouping related blocks for sequential reads, and laying out
data to minimize disk seeks.

\begin{verbatim}
   The paper's central design principle (condensed):

       "The file system is organized to keep the disk head moving
        in the right direction: data for one file is placed in
        contiguous block groups, and the inode for a file sits near
        its data."

   In our chapter's terms: st_blocks, inode-vs-data separation, and
   the reason `ls -l` can be fast are all consequences of this layout
   design. The 4 KB block size used throughout this book descends
   directly from this paper's block-group scheme.
\end{verbatim}

Why read it: it is the clearest demonstration that a filesystem is a
\emph{performance-and-durability} design as much as a namespace, and
that the questions we still answer today (block size, metadata
placement, crash consistency) were first posed in this paper.

\paperref{M.~K.~McKusick, W.~N.~Joy, S.~J.~Leffler, R.~S.~Fabry, ``A Fast
File System for UNIX,'' ACM Transactions on Computer Systems 2(3),
1984}{https://www.cs.utexas.edu/~dahlin/classes/cs375/papers/fast-ffs.pdf}
\end{researchnote}

\section{Practical: A Mini ls -l}

Let us combine \code{opendir}, \code{lstat}, and our permission
decoder into a genuine tool: a small \code{ls -l} that lists a
directory with type, permissions, link count, size, and name---the
output you have read a thousand times, now understood byte by byte.

\begin{lstlisting}[caption={minils.c: a real directory listing tool.}]
#include <stdio.h>
#include <stdlib.h>
#include <dirent.h>
#include <sys/stat.h>

char file_type_char(mode_t mode) {
    if (S_ISREG(mode))  return '-';
    if (S_ISDIR(mode))  return 'd';
    if (S_ISLNK(mode))  return 'l';
    if (S_ISCHR(mode))  return 'c';
    if (S_ISBLK(mode))  return 'b';
    if (S_ISFIFO(mode)) return 'p';
    if (S_ISSOCK(mode)) return 's';
    return '?';
}

void print_permissions(mode_t mode) {
    const char *rwx = "rwxrwxrwx";
    int i;

    putchar(file_type_char(mode));
    for (i = 0; i < 9; i++) {
        int bit;

        bit = 8 - i;
        if (mode & (1 << bit)) {
            putchar(rwx[i]);
        } else {
            putchar('-');
        }
    }
}

int main(int argc, char *argv[]) {
    const char *path;
    DIR *d;
    struct dirent *e;

    path = (argc >= 2) ? argv[1] : ".";

    d = opendir(path);
    if (d == NULL) {
        perror("opendir");
        return 1;
    }

    while (1) {
        struct stat st;
        char full[1024];
        int plen;

        e = readdir(d);
        if (e == NULL) {
            break;
        }
        if (e->d_name[0] == '.') {
            continue;
        }

        plen = snprintf(full, sizeof(full), "%s/%s", path, e->d_name);
        if (plen < 0 || plen >= (int) sizeof(full)) {
            continue;
        }

        if (lstat(full, &st) < 0) {
            continue;
        }

        print_permissions(st.st_mode);
        printf(" %2ld %8ld %s\n",
               (long) st.st_nlink,
               (long) st.st_size,
               e->d_name);
    }

    closedir(d);
    return 0;
}
\end{lstlisting}

\begin{verbatim}
   $ ./minils src
   -rw-r--r--  1     142 main.c
   -rw-r--r--  1      89 util.c

   You are now reading those ten characters and three columns the way
   the kernel writes them: type bits, permission bits, a reference
   count, and a size from the inode -- fetched by lstat per entry.
\end{verbatim}

\section{Structures in Bytes}

The inode record that \code{stat} reads is a byte structure. On
x86-64 glibc, \code{struct stat} is 144 bytes; the fields this chapter
used sit at these offsets:

\begin{verbatim}
   struct stat (x86-64 glibc), the fields you met:

   offset  size  field        meaning
   0       8     st_dev       device the file lives on
   8       8     st_ino       inode number
   16      8     st_nlink     hard-link count
   24      4     st_mode      type + permission bits
   28      4     st_uid       owner
   32      4     st_gid       group
   40      8     st_size      length in bytes
   72      8     st_blocks    x512-byte disk units
   96      16    st_mtime     data last modified
   112     16    st_atime     last access
   128     16    st_ctime     metadata changed (NOT creation)
\end{verbatim}

The on-disk ext4 inode is 256 bytes and holds the same ideas in
compact form:

\begin{verbatim}
   ext4 inode, key offsets:

   offset  size  field
   0       2     mode
   2       2     uid
   4       4     size_lo
   8       4     atime
   12      4     ctime
   16      4     mtime
   24      2     gid
   26      2     links_count      (st_nlink's source)
   28      4     blocks_lo
   40      60    block[15]        the map to data blocks
   106     2     size_high        (size above 4 GB)
\end{verbatim}

And a directory entry is bytes too: 4 bytes of inode number, 2 of name
length, 1 type byte, then the name. The kernel parses these records;
\code{readdir} hands you the pieces, and \code{ls -l} reads the rest
from the 144-byte \code{stat} record above.

\section{Summary}

\begin{verbatim}
   The filesystem in one paragraph:

   names (directory entries) -> inodes (metadata + block map) -> blocks

   Tools learned this chapter:
     stat()/lstat()    read an inode's record
     S_IS*()           decode the type half of st_mode
     mode bits         the permission model (rwx, umask, chmod)
     link()/symlink()  hard vs symbolic indirection
     opendir/readdir   iterate a directory
     /proc             the kernel as a virtual filesystem

   Rules:
     1. lstat when walking directories, stat when you mean the target;
     2. readdir gives names only -- stat each entry for details;
     3. nlink==0 means the file's data can be reclaimed;
     4. permissions are enforced at open() time by the kernel.
\end{verbatim}

\begin{osconcept}
Everything in this chapter was the OS maintaining a namespace: inodes,
reference counts, permission checks, virtual files, and directory
iteration are all kernel-managed state behind the paths you type. The
filesystem is where the OS's oldest promise lives---``your data will
survive, named, and protected''---and understanding its internals is
what lets you build databases, editors, and tools that treat storage
as a real resource instead of magic.
\end{osconcept}

We now have files and the metadata around them. In the next chapter we
turn back to processes for their other half: how they are grouped and
supervised---process groups, sessions, and the discipline of writing
daemons that run forever without a terminal.

}{}
\IfFileExists{chapter9.tex}{% ============================================================================
%  chapter9.tex
%  Process Control: Groups, Sessions, and Daemons
% ============================================================================

\chapter{Process Control and Daemons}

\section{From Single Processes to a Community}

Chapter~5 created processes; Chapter~7 connected them with pipes;
Chapter~6 taught them to react. This chapter asks the next question
up: \emph{how does a collection of processes live together?} Real
machines run hundreds of processes at once, and the operating system
must group them so that it can signal them, schedule them, and police
them as units. The two grouping concepts are the \emph{process group}
and the \emph{session}, and they are the backbone of every shell,
every terminal, and every daemon.

\begin{verbatim}
   The process tree. Every process has a parent; the root is init:

     init (pid 1)
       |
       +-- systemd (pid 1's job: supervise everything)
             |
             +-- your login shell (a session leader)
                   |
                   +-- emacs      (foreground process group)
                   +-- make       (a background group)
                   |     +-- gcc
                   |     +-- cc1
                   +-- server     (another background group)

   The tree is one fact; the groups and sessions are the labels the
   kernel uses to decide who gets which signal and who may read the
   terminal.
\end{verbatim}

\begin{osconcept}
The kernel does not manage processes one by one; it manages them as
units. Groups exist so that a signal or a job can address many
processes at once, and sessions exist so that the kernel can tie one
set of processes to one terminal. Deciding what those units are---and
how a process joins or leaves them---is operating system bookkeeping,
and it is what makes job control and daemons possible at all.
\end{osconcept}

\section{Process Groups}

A \emph{process group} is a set of processes that share a group id. The
group's purpose is to be a target: \emph{you can signal a whole group
at once}. When you press Ctrl-C, the terminal does not signal your one
process; it signals your entire foreground \emph{group}---which is why
a pipeline like \code{cat f | sort} stops as one unit.

\begin{lstlisting}[caption={Who am I, and who is my group?}]
#include <stdio.h>
#include <unistd.h>

int main(void) {
    printf("pid       = %d\n", getpid());
    printf("process group = %d\n", getpgrp());
    return 0;
}
\end{lstlisting}

A child inherits its parent's process group; a session leader, or a job
that a shell places in its own group, has a group id equal to its own
pid---it starts a new group alone. The calls that move processes between groups are
\code{setpgid} (join or create a group) and \code{getpgrp} (ask who
your group is). The shell uses \code{setpgid} to gather a pipeline's
children into one group so it can signal them together.

\begin{lstlisting}[caption={Signal a whole group: note the minus sign.}]
#include <stdio.h>
#include <stdlib.h>
#include <unistd.h>
#include <signal.h>

/* Ask every process in "pgrp" to terminate. */
void kill_group(pid_t pgrp) {
    int r;

    r = kill(-pgrp, SIGTERM);    /* -pgrp means "the group named pgrp" */
    if (r < 0) {
        perror("kill group");
    }
}
\end{lstlisting}

\begin{verbatim}
   kill(-pgrp, sig)  vs  kill(pid, sig):

     kill(pid, sig)    signal ONE process
     kill(-pgrp, sig)  signal EVERY process in the group pgrp
     kill(0, sig)      signal every process in MY group
     kill(-1, sig)     signal every process I may signal (extreme)

   This is why Ctrl-C feels like it hits "the whole job": the terminal
   driver runs kill(-foreground_pgrp, SIGINT). One call, many victims.
\end{verbatim}

\begin{osconcept}
Process groups are the kernel's unit of \emph{collective action}. The
OS lets one call address many processes because stopping or starting a
job should be atomic with respect to the whole job, not piecemeal.
The group id is kernel state; \code{setpgid} is a syscall that
reassigns it, and only the kernel can guarantee that a group's members
are tracked consistently.
\end{osconcept}

\section{Sessions and the Controlling Terminal}

A \emph{session} is a collection of process groups, headed by a
\emph{session leader} (the process that created the session). A
session may own one \emph{controlling terminal}. The session is the
unit that owns a terminal; within the session, exactly one process
group---the \emph{foreground group}---is allowed to read it.

\begin{verbatim}
   One session, one terminal, three groups:

      /dev/tty1  (the controlling terminal)
         |
         +---- session led by the login shell (sid = 100)
                |
                +---- FOREGROUND group  (your `cat bigfile`)
                |       -- the ONLY group allowed to read the tty
                |       -- receives Ctrl-C as SIGINT
                |
                +---- background group  (your `sleep 100 &`)
                |       -- NOT allowed to read the tty
                |       -- does NOT receive Ctrl-C
                |
                +---- background group  (your `make &`)

   Two kernel-enforced rules that follow from this picture:

     1. A background group that tries to READ the terminal is STOPPED
        by SIGTTIN. Try it: run `cat` in the background of a shell
        and it freezes, because the kernel refuses to let it steal
        your keystrokes.

     2. When you log out, the terminal closes; the kernel sends SIGHUP
        (hang up) to the session's leader, which the shell uses to
        clean up its jobs. Daemons, as we will see, detach from all
        of this on purpose.
\end{verbatim}

\begin{lstlisting}[caption={Forking a background job is group arithmetic.}]
#include <stdio.h>
#include <stdlib.h>
#include <unistd.h>
#include <sys/types.h>
#include <sys/wait.h>

int main(void) {
    pid_t pid;
    pid_t my_group;

    pid = fork();
    if (pid < 0) {
        perror("fork");
        exit(1);
    }

    if (pid == 0) {
        /* Child: put itself in its own new group, like a shell
           backgrounding a job. Then run a command.                 */
        setpgid(0, 0);            /* new group, id = my pid */
        execlp("sleep", "sleep", "100", NULL);
        perror("execlp");
        exit(127);
    }

    /* Parent: also record the child in its new group, then wait. */
    my_group = getpgrp();
    setpgid(pid, pid);
    printf("parent in group %d, child %d is in its own group %d\n",
           (int) my_group, (int) pid, (int) pid);

    waitpid(pid, NULL, 0);
    return 0;
}
\end{lstlisting}

\begin{verbatim}
   Why both parent AND child call setpgid? A race:

     - if the child runs first, IT calls setpgid(0,0) before the
       parent's setpgid(pid,pid);
     - if the parent runs first, the parent's call fixes the group
       before the child execs.

   Both orders are correct because setpgid is idempotent and allowed
   for a process still in the same session. Shells do exactly this to
   avoid a window where a backgrounded child could read the terminal.
\end{verbatim}

\begin{osconcept}
Sessions and the controlling terminal are the kernel's answer to a
very concrete problem: many processes share one screen, and the screen
must have exactly one reader and one interrupt target at a time. The
kernel owns the terminal device, so it can enforce the foreground rule
mechanically---stopping an intruder with \code{SIGTTIN} and routing
\code{SIGINT} only to the foreground group. This is why job control
feels ``instant and fair'': it is hardware-level policy.
\end{osconcept}

\section{setsid: Leaving the Terminal Behind}

\code{setsid} is the escape hatch. It creates a brand-new session,
makes the caller its leader, and---crucially---\emph{detaches the
caller from any controlling terminal}. After \code{setsid}, your
process is a session leader with no terminal, no foreground group, and
no route for Ctrl-C to reach it. This is the first step of becoming a
daemon.

\begin{lstlisting}[caption={setsid: new session, no terminal.}]
#include <stdio.h>
#include <unistd.h>

int main(void) {
    pid_t sid;

    sid = setsid();
    if (sid < 0) {
        perror("setsid");
        return 1;
    }

    printf("new session id = %d\n", (int) sid);
    printf("(I am now a session leader with no controlling terminal)\n");
    return 0;
}
\end{lstlisting}

The catch: \code{setsid} fails if the caller is already a process
group leader, because a group leader could accidentally make its whole
group leave together. That failure is the entire reason for the
\emph{first fork} in daemonization (next section).

\begin{osconcept}
\code{setsid} is the kernel's ``give me a clean room'' call. It
allocates a new session, reassigns the process, and disconnects the
terminal---all state only the kernel can change. The restriction that
a group leader may not do it exists to keep the group atomic: the OS
refuses to split a group mid-session, so your program must
\emph{fork} to create a child that is not a leader. The detour is the
kernel's safety rule, and daemons are built around it.
\end{osconcept}

\section{Daemons: Processes Without a User}

A \emph{daemon} is a long-running background service: a web server, a
database, a cron scheduler, an inotify watcher. It has no terminal, no
interactive user, no expectation of ever being \code{Ctrl-C}'d from a
shell. It is created to \emph{survive}: it must outlive the login
session that started it, ignore the terminal's death, and keep running
until told to stop.

To become a daemon, a process must undo every link to the user world
it inherited. The standard recipe, in order, is:

\begin{verbatim}
   The daemonization ceremony:

   1. fork()          so the child is NOT a process group leader,
                      and the parent may exit (the child becomes an
                      orphan; init adopts it).

   2. setsid()        the child becomes a new session leader with
                      NO controlling terminal.

   3. fork() AGAIN    now the child is not a session leader either,
                      so it can never reacquire a controlling
                      terminal. The first fork's child exits.

   4. chdir("/")      leave whatever directory we were started in
                      (that mount point may be unmounted later).

   5. umask(022)      set a sane default for files the daemon creates.

   6. close 0,1,2 and re-point all three at /dev/null, so any
      accidental stdout write cannot spam a terminal that no longer
      exists, and stdin reads return EOF instead of blocking.

   The double fork is the part beginners get wrong. Explain it in one
   line: the FIRST fork makes setsid() legal (child is not a group
   leader); the SECOND fork makes the daemon unable to grab a
   controlling terminal again (not a session leader).
\end{verbatim}

The double fork deserves a picture, because it is the most ritualistic
step in all of Unix programming:

\begin{verbatim}
   Why fork twice?

     process A (original)
       |  fork()
       v
     process B   (not a group leader: setsid() is now ALLOWED)
       |  setsid()  -> B becomes session leader
       |  fork()
       v
     process C   (NOT a session leader: can never reacquire a tty)
       |          this is the daemon; B exits immediately
       v
     the daemon lives on, adopted by init, immune to terminals.

   A session leader that later opens a terminal device can be made
   its controlling terminal again. A NON-leader cannot. The second
   fork buys permanent detach.
\end{verbatim}

\begin{lstlisting}[caption={daemonize.c: the whole ceremony in one function.}]
#include <stdio.h>
#include <stdlib.h>
#include <unistd.h>
#include <fcntl.h>
#include <sys/stat.h>
#include <sys/types.h>

/* Re-point descriptor fd at /dev/null. */
static void point_at_devnull(int fd) {
    int nfd;

    nfd = open("/dev/null", O_RDWR);
    if (nfd >= 0) {
        dup2(nfd, fd);
        close(nfd);
    }
}

/* Become a daemon. Returns 0 on success, -1 on failure. */
int daemonize(void) {
    pid_t pid;

    /* 1. First fork: the child may now call setsid(). */
    pid = fork();
    if (pid < 0) {
        return -1;
    }
    if (pid > 0) {
        exit(0);       /* parent's job is done */
    }

    /* 2. New session, no controlling terminal. */
    if (setsid() < 0) {
        return -1;
    }

    /* 3. Second fork: ensure we are not a session leader. */
    pid = fork();
    if (pid < 0) {
        return -1;
    }
    if (pid > 0) {
        exit(0);
    }

    /* 4. Work somewhere safe. */
    if (chdir("/") < 0) {
        return -1;
    }

    /* 5. Default permissions for files we create. */
    umask(022);

    /* 6. Detach all standard descriptors from any terminal. */
    close(0);
    close(1);
    close(2);
    point_at_devnull(0);
    point_at_devnull(1);
    point_at_devnull(2);

    return 0;
}
\end{lstlisting}

After \code{daemonize()} returns, the process has no terminal, no
interactive parent that matters, and no user pressing Ctrl-C. That is
both the point and the danger: nothing will stop it except a signal,
so a well-behaved daemon installs signal handlers (Chapter~6) for
\code{SIGTERM} and, very commonly, treats \code{SIGHUP} as
``reload your configuration'' rather than ``die.''

\begin{osconcept}
Detaching from the user session is entirely an OS ceremony. Forking,
creating sessions, unmounting-proof directories, and re-pointing
descriptors are all kernel services; a user-space process cannot make
itself terminal-independent any other way. (A \emph{mount point} is a
directory where a filesystem is attached: the kernel hangs a disk's
tree at that path. A daemon that keeps working inside a directory
whose filesystem gets unmounted holds a broken view of the world,
which is why the ceremony moves it to \code{/}.) The double-fork rule
exists because the kernel defines who may own a controlling terminal
(only session leaders), so the daemon must manipulate process/session
state
through syscalls to become permanently un-ownable.
\end{osconcept}

\section{Resource Limits}

A daemon runs forever, so it must not grow forever. The kernel enforces
per-process \emph{resource limits}---max open files, max memory, max
CPU time---and your program can read and set them with
\code{getrlimit}/\code{setrlimit}. Each limit has two values: the
\emph{soft} limit (what the process normally gets) and the \emph{hard}
limit (the ceiling the process may raise the soft limit to).

\begin{lstlisting}[caption={Reading and raising the open-file limit.}]
#include <stdio.h>
#include <stdlib.h>
#include <sys/resource.h>

int main(void) {
    struct rlimit rl;
    int r;

    r = getrlimit(RLIMIT_NOFILE, &rl);
    if (r < 0) {
        perror("getrlimit");
        return 1;
    }

    printf("open-file limit: soft = %lu, hard = %lu\n",
           (unsigned long) rl.rlim_cur,
           (unsigned long) rl.rlim_max);

    /* A server wanting more concurrent sockets raises the soft limit
       toward the hard limit. */
    if (rl.rlim_cur < rl.rlim_max) {
        rl.rlim_cur = rl.rlim_max;
        if (setrlimit(RLIMIT_NOFILE, &rl) < 0) {
            perror("setrlimit");
            return 1;
        }
        printf("raised soft limit to %lu\n", (unsigned long) rl.rlim_cur);
    }

    return 0;
}
\end{lstlisting}

\begin{verbatim}
   The limits you will actually need, in one table:

     RLIMIT_NOFILE  how many descriptors you may hold open
                    (a server with 10,000 connections needs a big one)
     RLIMIT_AS      total address space a process may map
     RLIMIT_CORE    size of core dumps (0 = disable them)
     RLIMIT_CPU     seconds of CPU before the kernel kills you

   `ulimit -n` in a shell is the same knob, set by the shell before
   your program even starts. getrlimit() reads it from inside.
\end{verbatim}

\begin{osconcept}
Resource limits are kernel-enforced policy: the kernel counts every
descriptor and every byte of address space against the limit at the
moment it is granted, and refuses to exceed it. A user-space program
cannot ``borrow'' a bit more RAM when the kernel says no. Limits are
the OS's way of containing runaway processes---which is exactly why
daemons both read them and raise them, in full knowledge of what the
kernel will and will not allow.
\end{osconcept}

\section{Practical: A Production-Shaped Daemon}

Let us assemble the chapter---and the signal discipline of Chapter~6---
into a complete daemon: double-forked, detached, signal-handled, with
a \emph{pidfile} (a file holding its pid so \code{stop} scripts can
find it). This is the skeleton that real services are built from.

\begin{lstlisting}[caption={mydaemon.c: the full daemon life cycle.}]
#include <stdio.h>
#include <stdlib.h>
#include <unistd.h>
#include <signal.h>
#include <fcntl.h>
#include <sys/stat.h>
#include <sys/types.h>

static volatile sig_atomic_t g_running = 1;

static void handle_term(int sig) {
    g_running = 0;             /* the only work the handler does */
}

/* Record our pid in a file so `kill $(cat pidfile)` finds us. */
static void write_pidfile(const char *path) {
    int fd;
    char buf[64];
    int n;

    fd = open(path, O_WRONLY | O_CREAT | O_TRUNC, 0644);
    if (fd < 0) {
        return;                /* failure to write the pidfile is not
                                  fatal for the daemon itself */
    }
    n = snprintf(buf, sizeof(buf), "%d\n", getpid());
    write(fd, buf, (size_t) n);
    close(fd);
}

int daemonize(void);           /* defined in Section 9.5 */

int main(void) {
    struct sigaction sa;

    if (daemonize() < 0) {
        perror("daemonize");
        return 1;
    }

    /* Handle stop requests; SIGHUP would normally mean "reload". */
    sa.sa_handler = handle_term;
    sigemptyset(&sa.sa_mask);
    sa.sa_flags = 0;
    sigaction(SIGTERM, &sa, NULL);
    sigaction(SIGINT, &sa, NULL);

    write_pidfile("/tmp/mydaemon.pid");

    printf("daemon started, pid %d\n", getpid());

    while (g_running) {
        /* Real work would go here: accept, read, process, log.
           printf() now goes to /dev/null, so logs must go to a
           file or syslog instead. */
        sleep(1);
    }

    unlink("/tmp/mydaemon.pid");
    printf("daemon stopping\n");
    return 0;
}
\end{lstlisting}

\begin{verbatim}
   Try the full cycle from a shell:

     $ ./mydaemon
     daemon started, pid 12345        (the prompt returns instantly:
                                        the parent exited at fork 1)

     $ ps -ef | grep mydaemon         (it is alive, no tty column)
     $ cat /tmp/mydaemon.pid
     12345
     $ kill -TERM $(cat /tmp/mydaemon.pid)
     $ ps -ef | grep mydaemon         (gone; pidfile unlinked)

   Notice the pidfile made the whole lifecycle scriptable: anyone can
   find the daemon by name, without knowing where it went.
\end{verbatim}

\begin{researchnote}[The evolution of process control]
The session/daemon machinery of this chapter is old, and its design
history is documented by the person who designed much of it. In 1979,
Dennis Ritchie published ``The Evolution of the UNIX Time-sharing
System,'' a retrospective on how the operating system grew: the fork
model, signals, and the shell-based process control we take for
granted were choices made and revised over a decade of use.

\begin{verbatim}
   Ritchie's retrospective, in one line:

       "The first (shell) programs were simple, but the underlying
        process-control mechanisms were general enough to support
        everything that came later."

   Read that as this chapter's thesis: fork, signals, and groups were
   designed as GENERAL primitives, and the shell (and later job
   control and daemons) built enormous behavior on top of a very small
   set of ideas.
\end{verbatim}

It is a short paper and a charming read---part technical, part
biography. It will make the fork/session machinery you just wrote feel
less like arcane ritual and more like the deliberate engineering it is.

\paperref{D.~M.~Ritchie, ``The Evolution of the UNIX Time-sharing
System,'' Language Design and Programming Methodology, 1979}{https://www.bell-labs.com/usr/dmr/www/hist.pdf}
\end{researchnote}

\section{Structures in Bytes}

Most process-control objects live in the kernel and have no stable
byte layout, but two of them have shapes you will actually touch.

\begin{verbatim}
   struct rlimit on 64-bit Linux: 16 bytes
   offset  size  field
   0       8     rlim_cur     the soft limit
   8       8     rlim_max     the hard limit

   a pidfile is plain text:
   "/tmp/mydaemon.pid"  ->  "12345\n"     (pid plus newline)
\end{verbatim}

The kernel's per-process record---its process control block---has a
shape worth knowing even though the offsets change with every version:

\begin{verbatim}
   a minimal view of the per-process record:

   offset  size  field
   0       4     pid
   4       4     ppid
   8       4     pgid       (the process group)
   12      4     sid        (the session)
   16      1     state      (running, stopped, zombie...)
   ...     8     pointer to the open-file table
   ...     8     pointer to the signal dispositions (Chapter 6)
\end{verbatim}

What is stable is the content, not the offsets: a process is a byte
record carrying its ids, its state, and pointers to its tables---the
same facts \code{getpid()}, \code{getpgrp()}, and \code{setsid()}
read and write.

\section{Summary}

\begin{verbatim}
   The process-control model, complete:

   process group   a set of processes signaled as one unit
   session         a set of groups owning one controlling terminal
   foreground group  the group allowed to read the terminal
   setsid          detach: new session, no terminal
   daemon          a process with no user, no tty, long life
   double fork     setsid legality + permanent terminal immunity
   rlimit          kernel-enforced caps on what a process may use
   pidfile         the daemon's way of publishing its pid

   Rules that keep process-control code correct:
     1. background readers get SIGTTIN; foreground groups get SIGINT;
     2. setpgid from BOTH parent and child to beat the race;
     3. a group leader cannot setsid(); fork first;
     4. daemons: chdir("/"), umask(022), close + redirect 0,1,2;
     5. a daemon's only master is a signal: handle SIGTERM/SIGHUP.
\end{verbatim}

\begin{osconcept}
This chapter was the OS as a \emph{sociologist}: grouping processes,
assigning them to terminals, capping their appetite, and adopting
their orphans. Every mechanism was kernel state---groups, sessions,
limits---because only the kernel can enforce rules across processes
that do not trust each other. You are not writing these rules; you are
\emph{negotiating} them through syscalls, and knowing the shape of the
negotiation is what systems programming is.
\end{osconcept}

With process control in hand, we have the full single-machine toolkit:
files, processes, signals, pipes, and daemons. The next chapter turns
to the fast lane of communication---time, timers, and the polling
mechanisms that let one process watch many channels at once without a
thread per channel.

}{}
\IfFileExists{chapter10.tex}{% ============================================================================
%  chapter10.tex
%  Time, Timers, and I/O Multiplexing
% ============================================================================

\chapter{Time, Timers, and I/O Multiplexing}

\section{The Problem: One Process, Many Channels}

Everything so far treated I/O as a rendezvous: a pipe blocks a reader
until bytes arrive, a signal blocks a \code{waitpid} until a child
finishes. That model has a hidden cost. A process doing blocking I/O
on one channel is \emph{frozen} for everyone else. If a server reads a
socket and no data comes, it cannot notice that a second socket is
ready, or that a timer has fired, or that a child has died.

\begin{verbatim}
   A single-threaded process with two blocking reads:

        process
       +---------------------+
       |  read(fd1, ...)     |  <-- blocks HERE until fd1 has data
       |                     |       ...fd2 could be bursting with
       |  read(fd2, ...)     |       data the whole time, untouched
       |  (never reached)    |
       +---------------------+

   The fix is I/O MULTIPLEXING: watch many descriptors at once, and be
   told which ones are ready. The process waits for ANY of them:

        process
       +---------------------+
       |  poll([fd1, fd2,     |
       |         timer], ...) |  <-- returns when ANY is ready
       |  -> fd2 is ready     |
       |  read(fd2, ...)      |  <-- now safe: it won't block
       +---------------------+
\end{verbatim}

This chapter builds the machinery of the \emph{event loop}: one
process, many channels, no thread per channel. It is the architecture
of every network server (Chapter~17), every window system, and every
game loop. To build it we need three ingredients: a trustworthy clock,
timers that arrive as file descriptors, and a way to ask the kernel
``which of my descriptors are ready right now?''

\begin{osconcept}
Multiplexing is an OS problem because only the kernel knows the true
readiness state of every descriptor at every instant. Your process
cannot cheaply know whether a pipe has data without the kernel telling
it. So the OS offers a \emph{readiness query}: block me, but block me
on \emph{many} channels at once, and wake me the moment any of them
can make progress. That single primitive is what lets one process serve
ten thousand connections.
\end{osconcept}

\section{Time: The Two Clocks}

Every multiplexed program needs time: timeouts, heartbeats, rate
limits. C reads the clock with \code{clock\_gettime}, which fills a
\code{struct timespec}---two fields: seconds and nanoseconds. The
first design decision is \emph{which clock}. There are two you will
use, and confusing them is a classic bug:

\begin{verbatim}
   CLOCK_REALTIME    the wall clock: 2026-08-03 14:00:00
                     - matches `date`, matches the calendar
                     - CAN JUMP: NTP syncs, an admin sets the time,
                       a leap second happens, the value goes
                       backward or forward in one step

   CLOCK_MONOTONIC   time since boot, ever increasing
                     - never jumps, never goes backward
                     - useless for "what day is it?"
                     - perfect for "how long did this take?"

   Rule: measure DURATIONS with MONOTONIC, show TIMES with REALTIME.
   If you time a request with REALTIME and an admin changes the clock
   mid-request, your "elapsed" time can be negative.
\end{verbatim}

\begin{lstlisting}[caption={Timing an interval with the monotonic clock.}]
#include <stdio.h>
#include <time.h>

/* Return the monotonic clock as a double number of seconds. */
double monotonic_seconds(void) {
    struct timespec ts;

    clock_gettime(CLOCK_MONOTONIC, &ts);
    return (double) ts.tv_sec + (double) ts.tv_nsec / 1e9;
}

int main(void) {
    double start;
    double end;
    double elapsed;

    start = monotonic_seconds();

    /* Do some work (a busy pause). */
    {
        volatile unsigned long x = 0;
        unsigned long i;

        for (i = 0; i < 200000000UL; i++) {
            x = x + 1;
        }
    }

    end = monotonic_seconds();
    elapsed = end - start;

    printf("elapsed: %.6f seconds\n", elapsed);
    return 0;
}
\end{lstlisting}

\begin{verbatim}
   struct timespec, drawn:

     struct timespec { long tv_sec;  /* whole seconds     */
                       long tv_nsec; /* nanoseconds 0..999999999 */ };

     tv_sec = 10, tv_nsec = 500000000   ->  10.5 seconds exactly.

   tv_nsec counts nanoseconds, not milliseconds. 1e9 ns = 1 s.
\end{verbatim}

\begin{osconcept}
The kernel owns the clocks. It maintains them, advances the monotonic
clock on every timer tick, and is the only party that can guarantee
``this clock never jumps.'' The realtime/monotonic split is a kernel
design decision: applications measuring intervals get a stable clock,
while the calendar is free to be corrected by an administrator without
breaking every timeout in the system.
\end{osconcept}

\section{Timers}

A \emph{timer} is the kernel's promise to wake you at a future time.
There are two modern styles, and the difference is instructive:
\code{setitimer} delivers its wakeup as a \emph{signal}
(\code{SIGALRM}, Chapter~6), while \code{timerfd} delivers it as a
\emph{file descriptor} that becomes readable---which means a timer can
be watched by \code{poll} alongside real data. The second is the one
this chapter uses, because it makes timers first-class citizens of the
event loop.

\begin{lstlisting}[caption={A periodic timer as a file descriptor.}]
#include <stdio.h>
#include <stdlib.h>
#include <unistd.h>
#include <sys/timerfd.h>
#include <time.h>

/* Create a timer that expires every interval_sec seconds. */
int make_timerfd(int interval_sec) {
    int tfd;
    struct itimerspec its;

    tfd = timerfd_create(CLOCK_MONOTONIC, 0);
    if (tfd < 0) {
        perror("timerfd_create");
        return -1;
    }

    /* it_value: first expiration; it_interval: how often after that. */
    its.it_value.tv_sec = interval_sec;
    its.it_value.tv_nsec = 0;
    its.it_interval.tv_sec = interval_sec;
    its.it_interval.tv_nsec = 0;

    if (timerfd_settime(tfd, 0, &its, NULL) < 0) {
        perror("timerfd_settime");
        close(tfd);
        return -1;
    }
    return tfd;
}
\end{lstlisting}

When a timerfd expires, it becomes \emph{readable}; reading it returns
the number of expirations that have accumulated since the last read.
The word ``accumulated'' matters: if the event loop was busy and four
periods passed, the read returns 4. Timers never silently drop ticks;
they are counted.

\begin{lstlisting}[caption={Drain the timer's accumulated expirations.}]
#include <stdint.h>

/* Read and discard the timer's expiration count. */
void drain_timer(int tfd) {
    uint64_t expirations;
    ssize_t r;

    r = read(tfd, &expirations, sizeof(expirations));
    if (r < 0 && errno != EAGAIN) {
        perror("read timerfd");
    }
}
\end{lstlisting}

\begin{verbatim}
   signal timer vs file-descriptor timer:

     setitimer()        fires SIGALRM; handler runs (ch.6 rules apply:
                        async-signal-safe, flags only)
     timerfd()          becomes readable; poll() notices it; you read
                        a counter -- no signal, no handler, ordinary
                        I/O rules apply

   For an event loop, timerfd wins: the loop already knows how to
   watch descriptors, so a timer is "just another fd."
\end{verbatim}

\begin{osconcept}
A timer is the kernel scheduling a wakeup on your behalf. The kernel
maintains the timer wheel, checks it on every tick, and makes your
descriptor readable (or raises your signal) at the right moment---all
without your process consuming any CPU while it waits. Timers are
kernel state; \code{timerfd\_settime} merely writes the alarm into
that state.
\end{osconcept}

\section{Non-Blocking I/O}

The event loop has one more requirement. When \code{poll} says ``fd2
is readable,'' the subsequent \code{read} should return immediately---
there is data---so blocking is not a problem there. But some situations
defy that: the kernel reports a descriptor ready for reading, you read,
and nothing arrives; or you want to \emph{write} without waiting for
the buffer to clear. For those, a descriptor can be made
\emph{non-blocking}: any operation that would block instead returns
immediately with \code{-1} and \code{errno = EAGAIN}, ``try again
later.''

\begin{verbatim}
   The O_NONBLOCK contract:

     blocking fd:     read() parks the process until data arrives.
     non-blocking fd: read() returns EAGAIN if nothing is there NOW.

   The deal: your code calls poll(), waits for "readable", then reads
   and is guaranteed not to block. Non-blocking mode is the safety net
   that makes that guarantee hold even when the kernel's guess is
   imperfect.

     write(fd, ...) on a full pipe, blocking: parks you.
     write(fd, ...) on a full pipe, non-blocking: returns EAGAIN.
\end{verbatim}

\begin{lstlisting}[caption={Turning a descriptor non-blocking.}]
#include <stdio.h>
#include <stdlib.h>
#include <fcntl.h>
#include <unistd.h>
#include <errno.h>

/* Make fd non-blocking. Returns 0 on success, -1 on failure. */
int set_nonblocking(int fd) {
    int flags;

    flags = fcntl(fd, F_GETFL, 0);
    if (flags < 0) {
        return -1;
    }
    if (fcntl(fd, F_SETFL, flags | O_NONBLOCK) < 0) {
        return -1;
    }
    return 0;
}
\end{lstlisting}

The read-then-check pattern becomes the idiom of every event loop:

\begin{verbatim}
   got = read(fd, buf, sizeof(buf));
   if (got < 0) {
       if (errno == EAGAIN) {
           /* no data right now: go back to poll() and wait */
           continue;
       }
       perror("read");   /* real error */
       break;
   }
   if (got == 0) {
       /* EOF: the writer closed its end */
       break;
   }
   handle(buf, got);
\end{verbatim}

\begin{osconcept}
Blocking is the kernel parking a process; non-blocking is the kernel
\emph{refusing to park it}. The decision is per-descriptor state
(\code{O\_NONBLOCK}) that the kernel honors in every syscall. EAGAIN
is the kernel's answer to ``I won't wait for you''---it gives the
program the choice of what to do with its time instead of deciding
for it. The OS provides both modes because both are essential: blocking
for simple programs, non-blocking for multiplexed ones.
\end{osconcept}

\section{select: The Classic Multiplexer}

\code{select} is the original way to watch many descriptors. You build
a set of descriptors with \code{fd\_set}, tell \code{select} how long
to wait, and it blocks until any of them is ready---then it tells you,
by modifying your set in place, which ones.

\begin{lstlisting}[caption={Watching stdin and a pipe with select.}]
#include <stdio.h>
#include <stdlib.h>
#include <unistd.h>
#include <sys/select.h>

int main(void) {
    int pfd[2];
    fd_set readfds;
    int r;
    int i;

    if (pipe(pfd) < 0) {
        perror("pipe");
        return 1;
    }

    /* fd 0 = stdin, pfd[0] = pipe read end. */
    FD_ZERO(&readfds);
    FD_SET(0, &readfds);
    FD_SET(pfd[0], &readfds);

    /* Watch the highest fd + 1, with no timeout. */
    r = select(pfd[0] + 1, &readfds, NULL, NULL, NULL);
    if (r < 0) {
        perror("select");
        return 1;
    }
    if (r == 0) {
        printf("timeout (unreachable here: we passed no timeout)\n");
        return 1;
    }

    /* select rewrote readfds: only the ready fds remain set. */
    for (i = 0; i <= pfd[0]; i++) {
        if (FD_ISSET(i, &readfds)) {
            printf("fd %d is ready to read\n", i);
        }
    }

    return 0;
}
\end{lstlisting}

\begin{verbatim}
   select()'s shape, and its two famous limits:

     select(max_fd + 1, &readfds, &writefds, &exceptfds, timeout)
            ^ the HIGHEST fd + 1: select scans 0..max_fd every call
            ^ fd_set is a fixed-size bit array, historically 1024
              bits (FD_SETSIZE): fds beyond that are undetectable
            ^ it MODIFIES your sets: you must rebuild them each loop

   select was a breakthrough and remains everywhere, but its fixed
   size and linear scan make it the wrong tool for thousands of fds.
\end{verbatim}

\begin{osconcept}
\code{select} is the kernel performing a \emph{scan}: it checks every
descriptor from 0 to \code{max\_fd}, parks you if none are ready, and
wakes you the moment any is. The \code{fd\_set} is kernel-shared
state that the syscall both reads (what you want) and writes (what is
ready). The design is simple, and its limits---fixed bit width, linear
work per call---are the kernel's data structures showing through.
\end{osconcept}

\section{poll: An Array of Interest}

\code{poll} fixes select's two headaches. Instead of a fixed-size bit
set, you pass an \emph{array of structs}, one per descriptor, each
carrying the events you care about (\code{events}) and receiving the
events that actually happened (\code{revents}). There is no
\code{FD\_SETSIZE} ceiling, no rebuilding of sets, and the array does
not get destroyed by the call.

\begin{lstlisting}[caption={poll: an array of descriptors and interests.}]
#include <stdio.h>
#include <stdlib.h>
#include <unistd.h>
#include <poll.h>

int main(void) {
    int pfd[2];
    struct pollfd fds[2];
    int r;
    int i;

    if (pipe(pfd) < 0) {
        perror("pipe");
        return 1;
    }

    fds[0].fd = 0;              /* stdin */
    fds[0].events = POLLIN;     /* I want to know when it's readable */
    fds[1].fd = pfd[0];         /* the pipe read end */
    fds[1].events = POLLIN;

    /* Block until any watched descriptor is ready, no timeout. */
    r = poll(fds, 2, -1);
    if (r < 0) {
        perror("poll");
        return 1;
    }

    /* r is the number of fds with revents != 0. */
    printf("%d descriptor(s) ready:\n", r);
    for (i = 0; i < 2; i++) {
        if (fds[i].revents & POLLIN) {
            printf("  fd %d is readable\n", fds[i].fd);
        }
        if (fds[i].revents & POLLERR) {
            printf("  fd %d had an error\n", fds[i].fd);
        }
        if (fds[i].revents & POLLHUP) {
            printf("  fd %d hung up (peer closed)\n", fds[i].fd);
        }
    }

    return 0;
}
\end{lstlisting}

\begin{verbatim}
   struct pollfd, drawn:

     struct pollfd {
         int fd;          /* which descriptor to watch */
         short events;    /* what I'm interested in: POLLIN, POLLOUT */
         short revents;   /* what actually happened (kernel writes)  */
     };

   - POLLIN   data can be read without blocking
   - POLLOUT  space to write without blocking
   - POLLHUP  the peer closed (a hangup)
   - POLLERR  an error condition

   events stays untouched; revents is filled by the kernel. The array
   is rebuilt only when your interests change, not every loop.
\end{verbatim}

poll still tells the kernel your full interest list on every call, and
the kernel scans it each time. For a few dozen descriptors that is
irrelevant; for ten thousand it dominates. That is the regime
\code{epoll} was built for.

\section{epoll: The Scalable Multiplexer}

Linux's answer to scaling is \code{epoll}. The core idea is a
\emph{persistent interest table in the kernel}. You register each
descriptor once, with \code{epoll\_ctl}, and the kernel watches it
forever---no need to resend your list on every wait. When descriptors
become ready, \code{epoll\_wait} hands back only the ones that are,
so the work per call is proportional to the number of \emph{ready}
descriptors, not the total watched.

\begin{lstlisting}[caption={epoll: register once, wait many times.}]
#include <stdio.h>
#include <stdlib.h>
#include <unistd.h>
#include <sys/epoll.h>

int main(void) {
    int pfd[2];
    int epfd;
    struct epoll_event ev;
    struct epoll_event ready[8];
    int r;
    int i;

    if (pipe(pfd) < 0) {
        perror("pipe");
        return 1;
    }

    epfd = epoll_create1(0);
    if (epfd < 0) {
        perror("epoll_create1");
        return 1;
    }

    /* Register the pipe read end, and stdin, for readability. */
    ev.events = EPOLLIN;
    ev.data.fd = pfd[0];
    epoll_ctl(epfd, EPOLL_CTL_ADD, pfd[0], &ev);

    ev.events = EPOLLIN;
    ev.data.fd = 0;
    epoll_ctl(epfd, EPOLL_CTL_ADD, 0, &ev);

    /* Wait. epoll_wait returns only the fds that are actually ready. */
    r = epoll_wait(epfd, ready, 8, -1);
    if (r < 0) {
        perror("epoll_wait");
        return 1;
    }

    for (i = 0; i < r; i++) {
        if (ready[i].events & EPOLLIN) {
            printf("fd %d is readable\n", ready[i].data.fd);
        }
    }

    close(epfd);
    return 0;
}
\end{lstlisting}

\begin{verbatim}
   The epoll lifecycle:

     epoll_create1()      make an empty interest table, get its fd
     epoll_ctl(ADD)       "start watching fd X"   (register once)
     epoll_wait()         "block until some are ready, tell me which"
     epoll_ctl(DEL)       "stop watching fd X"

   Why it scales: the kernel keeps your list; epoll_wait() walks only
   the ready queue. With 10,000 idle connections and 10 active ones,
   each wait returns 10 events -- not a 10,000-entry scan.

   select/poll:  O(total fds)  per call   (rescan everything)
   epoll:        O(ready fds)  per call   (hand back only the active)
\end{verbatim}

\begin{osconcept}
epoll moves the \emph{state}---your interest list---from user memory
into the kernel, and that move is the whole optimization. The kernel
can associate each descriptor with the process watching it, wake the
process only when something in \emph{its} table changes, and hand back
a short list instead of re-scanning the world. Persistent, kernel-owned
interest tables are how the OS lets one process scale to tens of
thousands of simultaneous channels.
\end{osconcept}

\section{The Event Loop: A Reactor}

We now have every ingredient. The \emph{event loop} (also called a
\emph{reactor}) is the pattern that ties them together, and it is the
heart of every serious server:

\begin{verbatim}
   The reactor pattern, one diagram:

        +----------------------------------------------------+
        |                THE EVENT LOOP                      |
        |                                                    |
        |   +--------+   wait on ALL fds with a timeout      |
        |   | poll/  |   (including timerfds)                |
        |   | epoll  |                                       |
        |   | wait   |                                       |
        |   +--------+                                       |
        |       |                                            |
        |       v                                            |
        |   for each READY fd:                               |
        |       if timer:     drain it, do periodic work     |
        |       if data:      read (non-blocking), handle    |
        |       if EOF/hup:   close, forget this fd          |
        |       ^                                            |
        |       +--------- loop forever --------------------+
        +----------------------------------------------------+

   Invariant: no blocking call exists inside the loop except the wait
   itself. Every read/write is attempted only when the descriptor is
   known ready, and even then non-blockingly.
\end{verbatim}

The rules that keep a reactor healthy are few and absolute:

\begin{itemize}
  \item \emph{Never block in the handler.} If handling one fd could
    block, the whole world stops. Work that might stall is deferred,
    queued, or done elsewhere.
  \item \emph{Non-blocking everywhere.} A read that returns
    \code{EAGAIN} is not an error; it is the loop telling you to go
    wait again.
  \item \emph{Timers are fds.} Heartbeats and timeouts ride the same
    wait as data, so the loop wakes at deadlines even when no I/O
    arrives.
\end{itemize}

\section{Practical: A Multi-Channel Event Loop}

Let us build the real thing: one process watches \emph{two pipes and a
timer} with \code{poll}, and reacts to whichever is ready---exactly
the shape a network server's main loop will take in Chapter~17. Two
children write at different times; the loop notices each one when it
arrives, and ticks on the timer meanwhile.

\begin{lstlisting}[caption={mux.c: one loop, many channels, no thread.}]
#include <stdio.h>
#include <stdlib.h>
#include <unistd.h>
#include <poll.h>
#include <sys/timerfd.h>
#include <time.h>
#include <errno.h>
#include <stdint.h>

int make_timerfd(int interval_sec) {
    int tfd;
    struct itimerspec its;

    tfd = timerfd_create(CLOCK_MONOTONIC, 0);
    if (tfd < 0) {
        return -1;
    }
    its.it_value.tv_sec = interval_sec;
    its.it_value.tv_nsec = 0;
    its.it_interval.tv_sec = interval_sec;
    its.it_interval.tv_nsec = 0;
    if (timerfd_settime(tfd, 0, &its, NULL) < 0) {
        close(tfd);
        return -1;
    }
    return tfd;
}

void drain_timer(int tfd) {
    uint64_t expirations;
    ssize_t r;

    r = read(tfd, &expirations, sizeof(expirations));
    if (r < 0 && errno != EAGAIN) {
        perror("read timerfd");
    }
}

int main(void) {
    int pfd_a[2];
    int pfd_b[2];
    int tfd;
    struct pollfd fds[3];
    int remaining;
    int r;
    int i;

    if (pipe(pfd_a) < 0 || pipe(pfd_b) < 0) {
        perror("pipe");
        return 1;
    }

    /* Two writer children, each producing one line after a delay. */
    if (fork() == 0) {
        sleep(2);
        write(pfd_a[1], "from A\n", 7);
        _exit(0);
    }
    if (fork() == 0) {
        sleep(4);
        write(pfd_b[1], "from B\n", 7);
        _exit(0);
    }

    /* The parent never writes; drop both write ends. */
    close(pfd_a[1]);
    close(pfd_b[1]);

    tfd = make_timerfd(1);
    if (tfd < 0) {
        return 1;
    }

    /* Register all three channels. */
    fds[0].fd = pfd_a[0];
    fds[0].events = POLLIN;
    fds[1].fd = pfd_b[0];
    fds[1].events = POLLIN;
    fds[2].fd = tfd;
    fds[2].events = POLLIN;

    remaining = 2;          /* the two pipes we must hear from */

    while (remaining > 0) {
        r = poll(fds, 3, -1);
        if (r < 0) {
            if (errno == EINTR) {
                continue;              /* retry after a signal */
            }
            perror("poll");
            return 1;
        }

        for (i = 0; i < 3; i++) {
            char buf[64];
            ssize_t got;

            if ((fds[i].revents & POLLIN) == 0) {
                continue;
            }

            if (fds[i].fd == tfd) {
                drain_timer(tfd);
                printf("tick\n");
                continue;
            }

            /* A pipe is readable: read without blocking. */
            got = read(fds[i].fd, buf, sizeof(buf));
            if (got > 0) {
                printf("fd %d says: %.*s", fds[i].fd, (int) got, buf);
            } else {
                /* EOF: the writer closed. Stop watching this pipe. */
                printf("fd %d closed\n", fds[i].fd);
                fds[i].fd = -1;        /* poll ignores negative fds */
                remaining = remaining - 1;
            }
        }
    }

    printf("both channels done; exiting\n");
    return 0;
}
\end{lstlisting}

\begin{verbatim}
   Running it:

     $ ./mux
     tick                 (t=1: timer fires first)
     tick                 (t=2: timer again)
     fd 4 says: from A    (t=2: child A's line arrives the same second)
     tick                 (t=3)
     tick                 (t=4: timer)
     fd 5 says: from B    (t=4: child B's line)
     fd 5 closed          (writer B exited)
     fd 4 closed          (writer A's EOF was noticed too)
     both channels done; exiting

   One process, three channels, zero threads, zero busy-waiting. The
   kernel did all the waiting; the loop just dispatched what was ready.
\end{verbatim}

\begin{researchnote}[Event-driven servers at scale]
The event loop you just built is not a textbook toy; it is the exact
architecture that made early web servers fast. In 1999, researchers at
Rice University published ``Flash,'' a web server built around
event-driven I/O---a single loop multiplexing thousands of
connections---and showed it outperforming the threaded servers of the
day by wide margins.

\begin{verbatim}
   What Flash demonstrated (from the paper's own claims):

       "An event-driven server can handle tens of thousands of
        concurrent connections on commodity hardware, and its
        scalability comes from avoiding one-thread-per-connection
        and from careful handling of each I/O operation."

   In this chapter's terms: Flash used exactly the pattern of Section
   9.7 -- one poll-style wait, non-blocking handlers, no thread per
   connection. The lesson has aged well; today's event loops
   (epoll, kqueue, io_uring) are its descendants.
\end{verbatim}

When you build the scalable server of Chapter~17, you will be
repeating this paper's architecture on modern primitives---and now you
know where the idea came from.

\paperref{V.~S.~Pai, P.~Druschel, W.~Zwaenepoel, ``Flash: An Efficient
and Portable Web Server,'' USENIX Annual Technical Conference,
1999}{https://www.usenix.org/legacy/events/usenix99/full_papers/pai/pai.pdf}
\end{researchnote}

\section{Structures in Bytes}

Everything the event loop touches has a small, fixed byte shape:

\begin{verbatim}
   struct timespec        16 bytes   offset size field
                                     0      8    tv_sec
                                     8      8    tv_nsec
   struct itimerspec      32 bytes   0..15  it_interval
                                     16..31 it_value
   struct pollfd           8 bytes   0      4    fd
                                     4      2    events
                                     6      2    revents
   struct epoll_event     12 bytes   0      4    events (EPOLLIN, ...)
                                     4      8    data   (a 64-bit cookie)
   fd_set (select)       128 bytes   one bit per descriptor (1024 bits)

   a timerfd's readable payload: 8 bytes = a uint64_t expiration count
\end{verbatim}

Read the shapes and the API is predictable: \code{poll} takes an array
of 8-byte \code{pollfd} records and the kernel writes 2 bytes of
\code{revents} into each; \code{epoll\_wait} returns an array of
12-byte \code{epoll\_event} records; \code{clock\_gettime} fills 16
bytes. These layouts are fixed by the ABI---which is exactly why you
should never read or write them as raw bytes yourself; the helper
calls exist to keep you from doing so.

\section{Summary}

\begin{verbatim}
   The multiplexing toolbox, complete:

   clock_gettime(CLOCK_MONOTONIC)  interval timing that never jumps
   timerfd_create()                a timer as a readable descriptor
   fcntl(F_SETFL, O_NONBLOCK)      read without ever parking
   EAGAIN                          "nothing right now, wait again"
   select()                        the classic bit-set multiplexer
   poll()                          array of pollfd (events/revents)
   epoll()                         kernel-held interest table
   the event loop                  one wait, dispatch what's ready

   Rules that keep an event loop correct:
     1. no blocking calls inside handlers;
     2. reads may return EAGAIN -- go back to the wait;
     3. EOF (got == 0) and POLLHUP mean close and forget;
     4. timers are fds: wait on them like data;
     5. use MONOTONIC for durations, REALTIME for display.
\end{verbatim}

\begin{osconcept}
Multiplexing is the OS delivering on a hard promise: one process, many
channels, and the kernel does the waiting. It is an OS problem because
readiness is kernel knowledge and parking is kernel power. The
evolution from select to epoll is the OS learning to hold your
interest lists, so that scale no longer means more threads---just a
better wait.
\end{osconcept}

We have seen one process do the work of many. The next chapter goes
the other way: instead of multiplexing one process across channels, we
give a process \emph{many threads of execution} inside one address
space---sharing memory, racing on every variable, and forcing us to
learn the rules of concurrency.

}{}
\IfFileExists{chapter11.tex}{% ============================================================================
%  chapter11.tex
%  Concurrency with Threads
% ============================================================================

\chapter{Concurrency with Threads}

\section{A Process with Many Simultaneous Selves}

Chapter~10 taught one process to serve many channels by never blocking.
This chapter teaches the complementary approach: give one process
\emph{many threads of execution}. A \emph{thread} is a single sequence
of instructions running inside a process. A process with three threads
is, from the CPU's point of view, three independent instruction
streams---each with its own stack and registers---that happen to share
one address space.

\begin{verbatim}
   Two views of the same program:

   PROCESSES (ch.5-9):               THREADS (this chapter):

   +----------+  +----------+       +-----------------------+
   | stack    |  | stack    |       | stack T1  (private)   |
   | heap     |  | heap     |       | stack T2  (private)   |
   | data     |  | data     |       | stack T3  (private)   |
   | text     |  | text     |       |   THE SHARED HEAP     |
   +----------+  +----------+       |   THE SHARED GLOBALS  |
        separate address             |   text                |
        spaces, no shared             +-----------------------+
        memory                         ONE address space,
                                       many instruction streams

   Threads share: the heap, the globals, the open descriptors.
   Threads do NOT share: their stacks, their registers.

   Sharing is the whole point -- and the whole danger.
\end{verbatim}

The trade-off with processes is now concrete. Processes are isolated
but must use pipes or memory mapping to share anything; threads share
everything trivially---a thread can read and write a global variable
the way any function can---but that sharing is exactly where races
come from. The rest of this chapter is the discipline that makes
sharing safe.

\begin{osconcept}
Threads are the kernel (or the thread library) multiplexing several
instruction streams onto one process. Because they share one address
space, the \emph{context switch} between them is far cheaper than a
process switch---no page tables to swap. But the kernel can schedule
threads at any instruction boundary, so two threads can interleave
their operations arbitrarily. That interleaving is the concurrency
problem, and only disciplined coordination makes it safe.
\end{osconcept}

\section{pthread\_create and pthread\_join}

The POSIX threads library (\code{pthreads}) is the C standard for
threads. A thread is created with \code{pthread\_create}, which starts
a function in a new thread, and \emph{reaped} with
\code{pthread\_join}, which waits for that thread to finish and
collects its result.

\begin{lstlisting}[caption={Creating two threads and waiting for them.}]
#include <stdio.h>
#include <stdlib.h>
#include <pthread.h>

/* A thread runs one function to completion. */
void *work(void *arg) {
    int id;

    id = *((int *) arg);
    printf("worker %d on thread %lu\n", id, pthread_self());
    return NULL;
}

int main(void) {
    pthread_t t1;
    pthread_t t2;
    int id1;
    int id2;
    int r;

    id1 = 1;
    id2 = 2;

    r = pthread_create(&t1, NULL, work, &id1);
    if (r != 0) {
        fprintf(stderr, "pthread_create: %d\n", r);
        return 1;
    }
    r = pthread_create(&t2, NULL, work, &id2);
    if (r != 0) {
        fprintf(stderr, "pthread_create: %d\n", r);
        return 1;
    }

    /* The main thread waits for both to finish. */
    pthread_join(t1, NULL);
    pthread_join(t2, NULL);

    printf("both workers are done\n");
    return 0;
}
\end{lstlisting}

The signature is the first thing to master:

\begin{verbatim}
   pthread_create( &thread_id,     /* where to store the new thread  */
                   NULL,           /* attributes: NULL = defaults    */
                   work,           /* the function to run            */
                   &id1 );         /* ONE pointer-sized argument     */

   - The function is void *work(void *arg): it takes ONE void* and
     returns ONE void*. That single pointer carries all input and
     output. This is the same "one value in, one value out" contract
     as thread arguments everywhere.

   - pthread_create returns an ERROR NUMBER, not -1, and it does NOT
     set errno. Check it against 0 and print the number itself.

   - pthread_join(thread, &result) blocks until that thread returns;
     the thread's return value is written through the second argument.

   - The argument we pass must OUTLIVE the thread. We pass the
     address of id1, a local of main that is still alive when the
     thread runs. Passing the address of a loop variable that
     changes would be a bug -- every thread would see the last value.
\end{verbatim}

\begin{verbatim}
   The lifetime contract, drawn:

     main
       |
       +-- pthread_create(&t1, work, &id1)  ->  thread t1 starts
       |                                        running work(id1)
       +-- pthread_create(&t2, work, &id2)  ->  thread t2 starts
       |                                        running work(id2)
       |
       +-- pthread_join(t1)  <-- main blocks here until t1 returns
       +-- pthread_join(t2)  <-- ...and here until t2 returns

   A thread is like a process child, but inside the same memory.
   create = fork, join = waitpid -- with sharing instead of copying.
\end{verbatim}

\begin{osconcept}
Thread creation and reaping are, again, OS services: the kernel must
allocate a new stack, new register set, and a new scheduler entry for
the thread, then run it on whatever CPU is free. The shared address
space is what makes threads cheap and also what makes them dangerous:
the kernel gives threads \emph{no isolation}, so the safety that
processes had by construction must now come from your discipline.
\end{osconcept}

\section{Data Races: The Shared-Variable Problem}

Here is the problem that defines all of concurrency. Two threads read
and write the same variable with no coordination. The C standard calls
this a \emph{data race}, and a data race is \emph{undefined behavior}--
the compiler may assume it never happens, and can optimize your
program in ways that make the bug appear far away from its cause.

\begin{lstlisting}[caption={The classic race: two threads, one counter.}]
#include <stdio.h>
#include <pthread.h>

#define N 1000000L

static long counter = 0;

void *increment(void *arg) {
    long i;

    for (i = 0; i < N; i++) {
        counter = counter + 1;
    }
    return NULL;
}

int main(void) {
    pthread_t t1;
    pthread_t t2;

    pthread_create(&t1, NULL, increment, NULL);
    pthread_create(&t2, NULL, increment, NULL);
    pthread_join(t1, NULL);
    pthread_join(t2, NULL);

    printf("counter = %ld (expected %ld)\n", counter, 2 * N);
    return 0;
}
\end{lstlisting}

Every run prints a different number, almost never \code{2000000}.
Why? Because \code{counter = counter + 1} is not one operation. It is
three: read, add, write. The kernel can interrupt a thread between any
two of them.

\begin{verbatim}
   The three steps of "counter = counter + 1", interleaved badly:

   Thread 1                    Thread 2
   -------                     -------
   read counter  (0)
   compute 0+1   (1)
   write counter (1)
                               read counter  (0)   <-- STALE: sees 0
                               compute 0+1   (1)
                               write counter (1)

   Both read 0, both write 1. One increment was LOST.
   With a million increments per thread, thousands are lost.

   A DATA RACE: two threads access the same memory, at least one
   write, with no synchronization in between. Result: undefined
   behavior. The number can be anything.
\end{verbatim}

The fix is not ``make the threads less unlucky.'' The fix is to make
the read-modify-write a single, indivisible unit---what we call a
\emph{critical section}---protected by a lock. That is the next
section.

\begin{ruleofthumb}[Any shared, written memory needs a lock]
The moment two threads can read a memory location that another thread
writes, you have a data race, and you have undefined behavior. The
only fix is synchronization: a mutex, a condition variable, or an
atomic operation. "It probably works" is not a fix; a race is a time
bomb with a random fuse.
\end{ruleofthumb}

\section{Mutexes: Making a Section Critical}

A \emph{mutex} (from ``mutual exclusion'') is a lock. Before entering
a critical section, a thread calls \code{pthread\_mutex\_lock}; while
one thread holds the lock, any other thread trying to lock is
\emph{blocked} by the kernel until the owner unlocks. The lock turns
``read, add, write'' into a section that only one thread may be inside
at a time.

\begin{lstlisting}[caption={The race, fixed with a mutex.}]
#include <stdio.h>
#include <pthread.h>

#define N 1000000L

static long counter = 0;
static pthread_mutex_t lock = PTHREAD_MUTEX_INITIALIZER;

void *increment(void *arg) {
    long i;

    for (i = 0; i < N; i++) {
        pthread_mutex_lock(&lock);       /* enter critical section */
        counter = counter + 1;           /* now safe: exclusive    */
        pthread_mutex_unlock(&lock);     /* leave critical section */
    }
    return NULL;
}

int main(void) {
    pthread_t t1;
    pthread_t t2;

    pthread_create(&t1, NULL, increment, NULL);
    pthread_create(&t2, NULL, increment, NULL);
    pthread_join(t1, NULL);
    pthread_join(t2, NULL);

    printf("counter = %ld (expected %ld)\n", counter, 2 * N);
    return 0;
}
\end{lstlisting}

\begin{verbatim}
   The critical section, drawn:

   Thread 1                          Thread 2
   lock()   -> owns the mutex
   counter++                         lock()   -> BLOCKS (kernel parks it)
   unlock() -> releases
                                     lock()   -> now owns it
                                     counter++
                                     unlock()

   The lock makes the read-modify-write atomic WITH RESPECT TO OTHER
   THREADS USING THE SAME LOCK. Two threads using DIFFERENT locks to
   protect the SAME variable are still racing -- the lock only works
   if everyone agrees on it.
\end{verbatim}

The discipline has four parts, and each is a rule:

\begin{itemize}
  \item \emph{One lock per unit of shared state.} The lock is a
    contract: every access to \code{counter} goes through it, or the
    contract is broken.
  \item \emph{Lock and unlock always pair.} Every path that locks must
    unlock---including error paths. (This is where C's lack of
    ``finally'' makes discipline essential.)
  \item \emph{Hold the lock as briefly as possible.} The critical
    section should contain the minimum work, so threads wait the
    minimum time.
  \item \emph{Never block while holding a lock if you can avoid it.}
    Holding a lock while waiting for I/O can stall every other thread
    that wants the lock.
\end{itemize}

\begin{osconcept}
A mutex is a synchronization primitive whose guarantee comes from the
kernel. On Linux a \code{pthread\_mutex\_t} lives in user space, and
the uncontended lock and unlock path never touches the kernel; only
when a thread must actually wait does the kernel (through a futex) park
it instead of letting it spin. Either way the OS guarantees that
\emph{at most one} thread holds a given mutex, and that guarantee is
what turns a sequence of instructions into a section others cannot
enter.
\end{osconcept}

\section{Condition Variables: Waiting for a Condition}

A mutex protects a variable; a \emph{condition variable} lets a thread
\emph{wait for a state of the world}---``buffer is not empty,''
``worker is done''---and be woken when that state arrives. The
condition variable is always used \emph{together with} a mutex, and
its central primitive does something that looks contradictory:

\code{pthread\_cond\_wait(cond, mutex)} \emph{atomically releases the
mutex and sleeps}. Another thread, seeing the state change, calls
\code{pthread\_cond\_signal} to wake one waiter. The awakened thread
re-acquires the mutex and re-checks the condition.

The canonical example is the bounded buffer: a fixed-size ring of
slots shared by producers and consumers.

\begin{lstlisting}[caption={Producer/consumer with a mutex and condition variables.}]
#include <stdio.h>
#include <pthread.h>

#define CAPACITY 8

static pthread_mutex_t mutex = PTHREAD_MUTEX_INITIALIZER;
static pthread_cond_t  not_full = PTHREAD_COND_INITIALIZER;
static pthread_cond_t  not_empty = PTHREAD_COND_INITIALIZER;

static int buffer[CAPACITY];
static int count = 0;

/* Produce one value into the shared buffer. */
void producer_put(int value) {
    pthread_mutex_lock(&mutex);

    /* Full? Then sleep until a consumer frees a slot.
       The loop, not an if, is mandatory: cond_wait can wake
       spuriously, and the state can change again before we
       actually run. */
    while (count == CAPACITY) {
        pthread_cond_wait(&not_full, &mutex);
    }

    buffer[count] = value;
    count = count + 1;

    pthread_cond_signal(&not_empty);   /* wake one consumer */
    pthread_mutex_unlock(&mutex);
}

/* Consume one value from the shared buffer. */
int consumer_get(void) {
    int value;

    pthread_mutex_lock(&mutex);

    /* Empty? Then sleep until a producer adds a value. */
    while (count == 0) {
        pthread_cond_wait(&not_empty, &mutex);
    }

    count = count - 1;
    value = buffer[count];

    pthread_cond_signal(&not_full);    /* wake one producer */
    pthread_mutex_unlock(&mutex);

    return value;
}
\end{lstlisting}

\begin{verbatim}
   The wait dance, step by step (producer side):

     producer locks the mutex
        |
        |  buffer full?
        |     yes --> cond_wait(not_full, mutex):
        |              1. atomically UNLOCKS the mutex
        |              2. sleeps until signaled
        |              3. on wake, RE-LOCKS the mutex
        |                     (then the while re-checks: full again?
        |                      wait again -- this is the loop's job)
        v
     slot is free: write, signal, unlock.

   The atomic unlock-and-sleep in cond_wait is the whole trick. If it
   released the lock and THEN slept, a producer could slip in between
   and the wait would race. The kernel makes the two a single step.
\end{verbatim}

\begin{ruleofthumb}[Always wait in a while loop, never an if]
\code{pthread\_cond\_wait} may return \emph{spuriously} (the kernel
does not guarantee one signal per wakeup), and by the time the waking
thread runs, the condition may have changed again. The correct shape
is always \code{while (condition) cond\_wait(...)}---re-check, don't
assume. This one habit eliminates an entire class of concurrency bugs.
\end{ruleofthumb}

\begin{osconcept}
A condition variable is the kernel's ``park me until the world
changes'' primitive, and its atomic release-and-sleep is the 
linchpin. Only the kernel can guarantee that the mutex is released
and the thread put to sleep as one indivisible step, so that no
signal is lost between the two. This is textbook operating-system
synchronization---semaphores and monitors are the same idea under
other names---and it is exactly why condition variables are kernel
primitives rather than library conveniences.
\end{osconcept}

\section{Deadlock: The Lock You Will Never Get}

When two threads each hold a lock the other needs, neither can
proceed, and both wait forever. That is \emph{deadlock}. It is the
second great concurrency bug, and it does not show up in tests---it
shows up in production, at 3 a.m., when two requests hit the locks in
the wrong order.

\begin{verbatim}
   Thread 1                      Thread 2
   -------                       -------
   lock(A)   ok                  lock(B)   ok
   lock(B)   BLOCKS              lock(A)   BLOCKS
             (B is held           (A is held
              by T2)               by T1)
   ...forever...                 ...forever...

   The deadly embrace, drawn as a cycle:

       T1 holds A, wants B
       T2 holds B, wants A
       --> a cycle of waiters. No one can make progress.

   The cure is ORDER: give every lock a global rank, and acquire
   locks in ascending order only. If everyone locks A before B,
   the cycle above cannot form.
\end{verbatim}

The four necessary conditions for deadlock are worth memorizing,
because removing \emph{any one} of them prevents it:

\begin{verbatim}
   deadlock needs all four:
     1. mutual exclusion     -- the locks exist
     2. hold and wait        -- a thread holds one lock, wants another
     3. no preemption        -- you cannot steal a held lock
     4. circular wait        -- the waiters form a cycle

   We cannot remove 1 (locks must exclude). We control 2 and 4 by
   design: acquire all locks in a fixed order (kills the cycle), and
   hold nothing while waiting for I/O (kills hold-and-wait).
\end{verbatim}

\begin{lstlisting}[caption={The deadlock-free pattern: lock in one order, always.}]
#include <pthread.h>

void transfer(pthread_mutex_t *from_lock, pthread_mutex_t *to_lock,
              long *from, long *to, long amount) {
    /* Both threads lock in the SAME order (lowest address first).
       Because the order is fixed, no cycle of waiters can form. */
    if (from_lock < to_lock) {
        pthread_mutex_lock(from_lock);
        pthread_mutex_lock(to_lock);
    } else {
        pthread_mutex_lock(to_lock);
        pthread_mutex_lock(from_lock);
    }

    *from = *from - amount;
    *to = *to + amount;

    pthread_mutex_unlock(to_lock);
    pthread_mutex_unlock(from_lock);
}
\end{lstlisting}

Comparing lock \emph{addresses} is a real, used technique: it imposes
a total order on locks without naming them. Whatever the two threads
pass in, they acquire them in the same sequence, and the circular wait
becomes impossible.

\begin{osconcept}
Deadlock is the kernel being \emph{too} faithful: it parks a thread
forever because that is exactly what the lock primitive promises. The
OS cannot detect and fix your lock ordering for you---it only enforces
mutual exclusion. Recognizing the four conditions and designing them
away is a programming discipline, not a runtime feature, which is why
deadlock prevention appears in operating-system textbooks and in
every serious concurrency course.
\end{osconcept}

\section{Threads vs Processes: Which One}

With both toolkits in hand, the choice is a design decision, not a
religious one. The table is the decision procedure:

\begin{verbatim}
   concern            threads                      processes
   ----------------   ---------------------------  -------------------
   address space      ONE, shared                  private each
   sharing data       variables (with locks)       pipes, shared memory
   creation cost      cheap (a stack + TCB)        heavier (page tables)
   crash isolation    NONE: one thread's segfault  strong: siblings
                      kills the whole process      survive
   communication      implicit, via memory         explicit, via kernel
   best for           parallel compute, latency-   robustness, multi-
                      critical I/O workers, GUI    tenant, "one bug
                                                    must not spread"

   The rule of thumb:
     Need to move lots of data between two workers? THREADS.
     Need one worker to be immune to another's crash? PROCESSES.
     Need both? The classic pattern: a process per service, threads
     inside each service.
\end{verbatim}

\section{Practical: A Parallel Sum}

Let us put the chapter to work with a real parallel algorithm: sum a
million numbers with four threads, each summing its own slice, then
combine. The trick is that each thread writes only to its \emph{own}
range struct---no shared counters, no locks needed, because the
sharing is done once, up front, by careful partition.

\begin{lstlisting}[caption={parsum.c: four threads, four disjoint slices.}]
#include <stdio.h>
#include <stdlib.h>
#include <pthread.h>

#define NTHREADS 4
#define ELEMENTS 1000000

static int data[ELEMENTS];

typedef struct range {
    int start;          /* first index, inclusive  */
    int end;            /* last index, exclusive   */
    long total;         /* this slice's sum        */
} range;

void *sum_range(void *arg) {
    range *r;
    long acc;
    int i;

    r = (range *) arg;
    acc = 0;
    for (i = r->start; i < r->end; i++) {
        acc = acc + data[i];
    }
    r->total = acc;          /* only this thread touches r->total */
    return NULL;
}

int main(void) {
    pthread_t threads[NTHREADS];
    range ranges[NTHREADS];
    long grand;
    size_t i;

    /* Fill the array. */
    for (i = 0; i < (size_t) ELEMENTS; i++) {
        data[i] = 1;
    }

    /* Split the array into NTHREADS disjoint slices. */
    for (i = 0; i < (size_t) NTHREADS; i++) {
        ranges[i].start = (int) (i * ((size_t) ELEMENTS / NTHREADS));
        ranges[i].end = (int) ((i + 1) * ((size_t) ELEMENTS / NTHREADS));
        ranges[i].total = 0;
        pthread_create(&threads[i], NULL, sum_range, &ranges[i]);
    }

    /* Collect: join each thread, add its slice. */
    grand = 0;
    for (i = 0; i < (size_t) NTHREADS; i++) {
        pthread_join(threads[i], NULL);
        grand = grand + ranges[i].total;
    }

    printf("sum = %ld (expected %d)\n", grand, ELEMENTS);
    return 0;
}
\end{lstlisting}

\begin{verbatim}
   Why this program is race-free WITHOUT a single lock:

     thread 0  owns data[0 .. 249999]   and ranges[0].total
     thread 1  owns data[250000 .. 499999]  and ranges[1].total
     ...
     the main thread writes ranges[].start/.end BEFORE creating
     the threads, and reads ranges[].total AFTER joining them.

   Every memory location is written by exactly ONE thread at a time.
   The only synchronization needed is the create/join itself, which
   pthreads provides. Shared nothing until the combine--the design
   discipline that makes concurrency easy.
\end{verbatim}

\begin{researchnote}[The origin of mutual exclusion]
The mutex and condition variable you just used descend directly from
work done in 1965 by Edsger Dijkstra, in a technical report called
``Cooperating Sequential Processes.'' In it, Dijkstra posed the
\emph{mutual exclusion problem}---two processes sharing variables
must not interleave their critical sections---and introduced the
\emph{semaphore}, the ancestor of the mutex, along with the very idea
of a critical section.

\begin{verbatim}
   Dijkstra's formulation, condensed:

       "If two processes share variables, each must protect its
        critical section so that at most one is inside at any time;
        the semaphore is the primitive that guarantees it."

   Our pthread_mutex_lock/unlock and cond_wait are the modern,
   clarified descendants of those first semaphores. When you write a
   critical section, you are using a concept that is sixty years old
   and still state of the art.
\end{verbatim}

It is a fascinating document---early concurrency worked out by hand,
before anyone had seen it fail in production. Reading it is the
fastest way to understand why the rules of this chapter exist.

\paperref{E.~W.~Dijkstra, ``Cooperating Sequential Processes,''
Technical Report EWD-123, Technological University Eindhoven,
1965}{https://www.cs.utexas.edu/users/EWD/transcriptions/EWD01xx/EWD123.html}
\end{researchnote}

\section{Structures in Bytes}

The pthread objects are user-space byte structures whose exact layout
is an ABI detail you should treat as opaque---but their shapes are
instructive. On glibc/x86-64:

\begin{verbatim}
   pthread_t       8 bytes    an opaque handle
   pthread_mutex_t 40 bytes   owner + kind + count + wait list
   pthread_cond_t  48 bytes   wait list + clock + sequencing

   pthread_mutex_t (glibc x86-64), approximate offsets:
   offset  size  field
   0       4     lock
   4       4     count        (recursion depth)
   8       4     owner        (which thread holds it)
   12      4     nusers
   16      4     kind         (normal, recursive, errorcheck)
   20      2     spins
   24      8     wait list    (futex waiters)
\end{verbatim}

When the kernel switches threads, the byte structure it saves is the
register image: 15 general-purpose registers at 8 bytes each (120
bytes) plus the instruction and stack pointers, stored on the thread's
kernel stack. That saved image, swapped on every context switch, is the
mechanical meaning of ``two threads, two sets of registers, one
address space.''

\section{Summary}

\begin{verbatim}
   The concurrency toolbox, complete:

   pthread_create / pthread_join   spawn and reap threads
   shared address space            threads share heap and globals
   data race                       unsynchronized shared writes = UB
   pthread_mutex_t                 the critical section (mutual excl.)
   pthread_cond_t                  wait for a condition, with the lock
   deadlock                        lock cycles; fix by lock ordering

   Rules that keep threaded code correct:
     1. every shared, written variable is protected by a lock;
     2. lock and unlock always pair, on every path;
     3. hold locks briefly; never block for I/O while holding one;
     4. cond_wait in a while loop, never an if;
     5. acquire multiple locks in one global order;
     6. prefer disjoint ownership (like the parallel sum) so locks
        are rarely needed at all.
\end{verbatim}

\begin{osconcept}
Threads are the OS's cheapest form of parallelism, purchased by
surrendering the isolation that processes enjoyed. The kernel
multiplexes instruction streams and enforces the primitives (locks,
condvars) that make sharing safe; everything else---partitioning,
ordering, and the discipline of the critical section---is the
programmer's. Master that bargain and you can write programs that use
every core; break it and you get bugs that reproduce only in
production.
\end{osconcept}

The next chapter takes the sharing bargain to its logical extreme:
shared memory, where processes---not threads---map the same bytes into
their address spaces and synchronize over raw memory, with no kernel
copying and no safety net.

}{}
\IfFileExists{chapter12.tex}{% ============================================================================
%  chapter12.tex
%  Shared Memory and Memory Mapping
% ============================================================================

\chapter{Shared Memory and Memory Mapping}

\section{The Boldest Abstraction: Memory That Is a File}

Every chapter since the first has treated the address space as a
container of bytes you own. This chapter changes the deal with one
call: \code{mmap}, which \emph{maps} a file---or a region of anonymous
memory---\emph{into your address space}. After the call, reading a
file is just reading memory: \code{p[i]}, not \code{read(fd, ...)}.
And when several processes map the same object as \emph{shared}, they
see the \emph{same physical pages}---shared memory, with no pipe, no
copy, and no kernel involvement in each transfer.

\begin{verbatim}
   The two faces of mmap:

   FILE MAPPING:                    SHARED MEMORY:
   a file's bytes appear in         two (or more) processes map the
   your address space.              same region; they share pages.

        address space               process A         process B
   +----------------------+         +--------+        +--------+
   |  [ p[0] ][ p[1] ]... |         | shared |<======>| shared |
   |    (the file)        |         | pages  |        | pages  |
   +----------------------+         +--------+        +--------+
        ^                             ^ same physical RAM
        | mmap
   +----------------------+
   |   the file on disk   |
   +----------------------+

   Reading the file = dereferencing the pointer.
   Writing the file = assigning through the pointer.
   Sharing between processes = both mapping the same thing SHARED.
\end{verbatim}

To understand why this works---and why it is fast---you need the
virtual-memory model of Chapter~1 in its full form: the address space
is divided into \emph{pages} (4~KB on our machines), and the CPU
translates page-by-page through the \emph{page table}. \code{mmap}
does not copy the file into memory. It merely \emph{points the page
table} at the file's pages. The first access to a page triggers a
\emph{page fault}; the kernel then reads that one page from the file
into RAM, lazily, on demand.

\begin{osconcept}
mmap unifies the two great abstractions of this book---virtual memory
and the file system---into one: \emph{the address space becomes a view
of an object}. The OS maps files lazily, page by page, and shares
physical pages between processes by pointing multiple page tables at
them. Only the kernel can do this, because only the kernel owns the
page tables and the files. Your program just gets a pointer.
\end{osconcept}

\section{mmap: The Call}

\code{mmap} takes six arguments and returns a pointer into your address
space (or \code{MAP\_FAILED}, which is \code{(void *) -1}, not
\code{NULL}). The size must be in bytes, and the kernel works in
pages---a region is rounded up to whole pages whether you like it or
not.

\begin{lstlisting}[caption={Reading a whole file by mapping it.}]
#include <stdio.h>
#include <stdlib.h>
#include <fcntl.h>
#include <sys/mman.h>
#include <sys/stat.h>
#include <unistd.h>

int main(int argc, char *argv[]) {
    int fd;
    struct stat st;
    const unsigned char *p;
    size_t i;

    if (argc != 2) {
        fprintf(stderr, "usage: %s file\n", argv[0]);
        return 1;
    }

    fd = open(argv[1], O_RDONLY);
    if (fd < 0) {
        perror("open");
        return 1;
    }
    if (fstat(fd, &st) < 0) {
        perror("fstat");
        return 1;
    }

    p = mmap(NULL, (size_t) st.st_size, PROT_READ,
             MAP_PRIVATE, fd, 0);
    if (p == MAP_FAILED) {
        perror("mmap");
        return 1;
    }

    /* The whole file is in our address space: just print it. */
    for (i = 0; i < (size_t) st.st_size; i++) {
        putchar(p[i]);
    }

    munmap((void *) p, (size_t) st.st_size);
    close(fd);
    return 0;
}
\end{lstlisting}

The arguments, each doing one job:

\begin{verbatim}
   mmap( addr, length, prot, flags, fd, offset )
        |     |       |     |       |    |
        |     |       |     |       |    +-- start mapping at this file
        |     |       |     |       |         offset (usually 0)
        |     |       |     |       +-- the descriptor, or -1 for
        |     |       |     |            an anonymous mapping
        |     |       |     +-- MAP_SHARED  (visible to others)
        |     |       |          MAP_PRIVATE (private copy, COW)
        |     |       +-- PROT_READ / PROT_WRITE / PROT_EXEC
        |     +-- length in BYTES (rounded up to pages)
        +-- NULL = let the kernel choose the address (almost always)

   munmap(ptr, length)  undoes it, exactly once, when you are done.
\end{verbatim}

\begin{verbatim}
   The lazy mapping, drawn:

      time 0: mmap() succeeds -- NOTHING is loaded yet.
              file:     [p0][p1][p2][p3][p4]
              address:  [    mapped, but empty pages    ]

      time 1: you touch p[0]. The CPU page-faults; the kernel loads
              page 0 from the file into RAM; your read completes.

              file:     [p0][p1][p2][p3][p4]
              address:  [p0][    empty    ]   <- only p0 loaded

      Each first touch of a page is a fault and one disk read.
      Unused pages cost nothing -- the mapping is a promise, not a
      copy. This is why mapping a 1 GB file can succeed instantly.
\end{verbatim}

\begin{osconcept}
The laziness of mmap is demand paging, and it is pure OS machinery.
The page table starts pointing at ``not present,'' the CPU raises a
fault on first touch, and the kernel satisfies that one page from the
file. Your program cannot do this itself; it has no page table. The
OS is betting that you will not touch every byte, and when you do, it
reads only what you actually use.
\end{osconcept}

\section{Anonymous Mappings}

An \emph{anonymous} mapping has no file behind it: the kernel provides
zeroed memory that behaves like an array, and---with
\code{MAP\_SHARED}---can be shared between processes. It is the
modern way to allocate shared memory, and the honest way to get a
big private buffer too (the \code{malloc} of Chapter~13 is built on
it).

\begin{lstlisting}[caption={A shared counter visible to a forked child.}]
#include <stdio.h>
#include <stdlib.h>
#include <sys/mman.h>
#include <unistd.h>
#include <sys/wait.h>

int main(void) {
    long *counter;
    pid_t pid;

    /* One shared page of anonymous memory, writable by all. */
    counter = mmap(NULL, sizeof(long),
                   PROT_READ | PROT_WRITE,
                   MAP_SHARED | MAP_ANONYMOUS, -1, 0);
    if (counter == MAP_FAILED) {
        perror("mmap");
        return 1;
    }
    *counter = 0;

    pid = fork();
    if (pid < 0) {
        perror("fork");
        return 1;
    }

    if (pid == 0) {
        /* Child: add 100 to the SHARED counter. */
        int i;

        for (i = 0; i < 100; i++) {
            *counter = *counter + 1;
        }
        _exit(0);
    }

    /* Parent: wait, then read what the child wrote. */
    waitpid(pid, NULL, 0);
    printf("counter = %ld (child added 100)\n", *counter);

    munmap(counter, sizeof(long));
    return 0;
}
\end{lstlisting}

\begin{verbatim}
   What sharing means here, page by page:

     before fork:   ONE mapping, ONE physical page.
     after fork:    the child inherits the mapping. Because it is
                    MAP_SHARED, BOTH processes' page tables point at
                    the SAME physical page.

     parent: [p] --->  shared page: [ 100 ]  <--- [p] : child
                            ^ same physical RAM

   The child's writes are visible to the parent -- not by copying,
   but because they write the very same bytes. No pipe, no kernel
   buffer, no copy: the fastest possible IPC.
\end{verbatim}

The price of that speed is exactly the theme of this book: shared
writes race. Two processes incrementing the same counter without a
lock lose increments, precisely as two threads did in Chapter~11. The
difference is that now the lock itself must live \emph{in the shared
region}, so both processes can reach it.

\begin{osconcept}
Sharing physical pages is the kernel's greatest privilege: multiple
processes, one piece of RAM, no copying, no serialization. It is also
a deliberate risk, because the kernel cannot police how you use the
bytes---a shared pointer is a shared pointer. The OS provides the
sharing and the synchronization primitives; using them safely is, as
ever, the programmer's contract.
\end{osconcept}

\section{MAP\_SHARED vs MAP\_PRIVATE}

The \code{flags} argument is the one that decides what a \emph{write}
does. This is the crux of the whole mechanism:

\begin{verbatim}
   MAP_PRIVATE   -- the mapping is private to this process.
                     Writes are copy-on-write: the first time this
                     process writes a page, the kernel gives it a
                     PRIVATE COPY, and all further writes touch only
                     the copy. The file (or the other processes) are
                     NOT changed.

   MAP_SHARED    -- the mapping is shared.
                     Writes go straight to the underlying object:
                     the file on disk for a file mapping, or the
                     shared physical pages for everyone else mapping
                     the same region.

   The rule of thumb: PRIVATE to read a file without changing it;
   SHARED to edit a file in place or to communicate between
   processes.
\end{verbatim}

\begin{verbatim}
   COPY-ON-WRITE for MAP_PRIVATE, drawn:

     before any write:                 after the first write to page 0:

     [p0][p1][p2][p3]   (all shared     [p0'](copy)[p1][p2][p3]
        |      with the file)             |        still shared
        +-- the file's pages             +-- this page is now private

   The cost model: reads are free until the pages change; writes cost
   one copy per page, only for pages you actually modify. That is why
   MAP_PRIVATE is the safe default for reading.
\end{verbatim}

For a \emph{shared file} mapping, your writes modify the file's
contents in place---exactly what an editor or a database wants. The
\emph{msync} call pushes the modified pages out to the storage medium
now, rather than waiting for the kernel's background flush. Think of
it as \code{fsync} for a mapping.

\begin{lstlisting}[caption={msync: flush a shared mapping to disk.}]
#include <sys/mman.h>

/* Write the mapped bytes back to the file immediately. */
int flush_mapping(void *addr, size_t length) {
    int r;

    r = msync(addr, length, MS_SYNC);   /* MS_ASYNC = don't wait */
    if (r < 0) {
        return -1;
    }
    return 0;
}
\end{lstlisting}

\begin{osconcept}
The shared/private distinction is kernel policy with teeth. PRIVATE
asks the OS for copy-on-write: the kernel re-points one page at a
fresh copy at the moment of the first write, invisibly and
atomically. SHARED asks the OS to make every write land on the same
physical pages (and, for files, eventually on disk). Choosing between
them is choosing between isolation and communication, and the kernel
enforces that choice per page.
\end{osconcept}

\section{Synchronizing Shared Memory}

Shared memory is fast because there is no kernel in the data path---
which is precisely why there is no kernel in the \emph{coordination}
path either. Two processes that share a counter must synchronize
themselves, with a lock that lives in the shared region. Pthreads
mutexes can do this, but only if told they are \emph{process-shared}:
by default a mutex is only usable by threads of one process.

\begin{lstlisting}[caption={A process-shared mutex guarding shared memory.}]
#include <stdio.h>
#include <stdlib.h>
#include <pthread.h>
#include <sys/mman.h>
#include <unistd.h>
#include <sys/wait.h>

typedef struct shared {
    pthread_mutex_t lock;    /* must live IN the shared region */
    long counter;
} shared;

int main(void) {
    shared *s;
    pthread_mutexattr_t attr;
    pid_t pid;
    int i;

    s = mmap(NULL, sizeof(shared),
             PROT_READ | PROT_WRITE,
             MAP_SHARED | MAP_ANONYMOUS, -1, 0);
    if (s == MAP_FAILED) {
        perror("mmap");
        return 1;
    }

    /* Declare the mutex usable by ANY process (PTHREAD_PROCESS_SHARED). */
    pthread_mutexattr_init(&attr);
    pthread_mutexattr_setpshared(&attr, PTHREAD_PROCESS_SHARED);
    pthread_mutex_init(&s->lock, &attr);
    pthread_mutexattr_destroy(&attr);

    s->counter = 0;

    pid = fork();
    if (pid < 0) {
        perror("fork");
        return 1;
    }

    if (pid == 0) {
        /* Child increments 1000 times, safely. */
        for (i = 0; i < 1000; i++) {
            pthread_mutex_lock(&s->lock);
            s->counter = s->counter + 1;
            pthread_mutex_unlock(&s->lock);
        }
        _exit(0);
    }

    /* Parent increments 1000 times, safely. */
    for (i = 0; i < 1000; i++) {
        pthread_mutex_lock(&s->lock);
        s->counter = s->counter + 1;
        pthread_mutex_unlock(&s->lock);
    }

    waitpid(pid, NULL, 0);
    printf("counter = %ld (expected 2000)\n", s->counter);

    pthread_mutex_destroy(&s->lock);
    munmap(s, sizeof(shared));
    return 0;
}
\end{lstlisting}

\begin{verbatim}
   The layout that makes it work:

     shared region (one anonymous mapping):
     +---------------------------------------------+
     | pthread_mutex_t lock   |  long counter      |
     +---------------------------------------------+
            ^ mapped into BOTH processes' address spaces

     process A:  s->lock, s->counter   (the very same bytes)
     process B:  s->lock, s->counter   (the very same bytes)

   The mutex's memory is shared, so the kernel can see both
   processes contending for it. WITHOUT setpshared(...PROCESS_SHARED),
   the mutex would still be in shared memory but declared private, and
   its behavior would be undefined -- two processes, one lock, no
   contract.
\end{verbatim}

\begin{osconcept}
A mutex is kernel state, and making it process-shared is an explicit
kernel decision (\code{PTHREAD\_PROCESS\_SHARED}) so the OS knows the
lock may be contended across processes, not just threads. The
primitive is the same; only its scope changes. Shared memory without
such a lock is a data race across a process boundary---undefined
behavior with an even bigger blast radius than the thread races of
Chapter~11.
\end{osconcept}

\section{Practical: Editing a File Through a Mapping}

Let us close with the file-backed side of mmap: a program that treats
a file as an array of integers, edits it through the mapping, and
flushes the result to disk. This is the exact technique used by
databases, key-value stores, and editors to modify large files in
place without a single \code{lseek}.

\begin{lstlisting}[caption={mmapfile.c: the file as an array of ints.}]
#include <stdio.h>
#include <stdlib.h>
#include <fcntl.h>
#include <sys/mman.h>
#include <unistd.h>

#define NINTS 8

int main(int argc, char *argv[]) {
    int fd;
    int *vals;
    size_t bytes;
    size_t i;

    if (argc != 2) {
        fprintf(stderr, "usage: %s file\n", argv[0]);
        return 1;
    }

    /* Open (or create) the file read-write. */
    fd = open(argv[1], O_RDWR | O_CREAT, 0644);
    if (fd < 0) {
        perror("open");
        return 1;
    }

    /* Size the file to hold NINTS ints before mapping it. */
    bytes = NINTS * sizeof(int);
    if (ftruncate(fd, (off_t) bytes) < 0) {
        perror("ftruncate");
        return 1;
    }

    /* Map the whole thing, writable and shared. */
    vals = mmap(NULL, bytes, PROT_READ | PROT_WRITE, MAP_SHARED, fd, 0);
    if (vals == MAP_FAILED) {
        perror("mmap");
        return 1;
    }

    /* Write through the mapping: no write(), no lseek. */
    for (i = 0; i < NINTS; i++) {
        vals[i] = (int) (i * i);
    }

    /* Flush now (optional: the kernel would flush eventually). */
    msync(vals, bytes, MS_SYNC);

    printf("wrote %zu ints through the mapping:\n", NINTS);
    for (i = 0; i < NINTS; i++) {
        printf("  vals[%zu] = %d\n", i, vals[i]);
    }

    munmap(vals, bytes);
    close(fd);
    return 0;
}
\end{lstlisting}

\begin{verbatim}
   Why this is a different way to do I/O:

     classic path:   write(fd, buf, n)  -- a syscall, kernel copies
     mmap path:      vals[i] = v;       -- a store instruction; the
                        kernel's page-fault machinery moves the bytes
                        to disk when it must

   The trade-off in one line: mmap removes the copy into and out of
   the kernel for random-access edits, at the price of page-fault
   costs on first touch and giving the kernel control over when the
   data actually reaches disk.
\end{verbatim}

\begin{researchnote}[The theory behind memory-mapped I/O]
mmap is not a trick; it is virtual memory applied to files, and the
theory behind virtual memory is one of computing's foundational
results. In 1970, Peter Denning published ``Virtual Memory''---the
survey that consolidated paging, page tables, and demand loading into
the coherent model this chapter relies on.

\begin{verbatim}
   Denning's framing, condensed:

       "Virtual memory allows programs to run as if they occupied
        far more memory than is physically present, by moving pages
        between main memory and secondary storage on demand."

   Read that with mmap in mind: mapping a 1 GB file with 4 GB of RAM
   is exactly this promise. The file IS the backing store; the page
   fault is the demand; and your pointer is the program's illusion of
   having the whole file in memory.
\end{verbatim}

Denning's survey is a beautifully written summary of why the page
model works, and it explains the behavior of every mmap program you
will ever write---the lazy first touch, the copy-on-write, the
page-sized granularity. When you wonder ``why is my mapping so
fast?,'' the answer was published in 1970.

\paperref{P.~J.~Denning, ``Virtual Memory,'' ACM Computing Surveys
2(3), 1970}{https://dl.acm.org/doi/10.1145/356638.356640}
\end{researchnote}

\section{Structures in Bytes}

The unit of shared memory is the page, and the structure that makes it
work is the page-table entry (PTE)---8 bytes on x86-64.

\begin{verbatim}
   a page: 4096 bytes = 512 eight-byte words.

   a page-table entry, the 8 bytes that name a page:

   bit   0   present           (0 here means "page fault")
   bit   1   writable
   bit   2   user-accessible
   bits  3-11  more flags      (dirty, accessed, ...)
   bits 12-51  physical frame number   <- where the page really lives
   bits 52-62  available to the kernel
   bit  63   NX                (no-execute)

   the four-level walk for a 48-bit address:
   CR3 -> PML4 -> PDPT -> PD -> PT -> 4 KB page
   (each level is a page of 512 eight-byte entries)
\end{verbatim}

\code{mmap} mostly writes entries into these tables: an anonymous
mapping installs ``not present'' entries that fault in zeroed pages; a
\code{MAP\_SHARED} mapping points several processes' PTEs at the same
physical frame number. The ``structure in bytes'' of shared memory is
ultimately this: one 8-byte entry with the same frame number sitting in
two page tables.

\section{Summary}

\begin{verbatim}
   The mapping toolbox, complete:

   mmap(0, len, prot, flags, fd, off)   map a file or anonymous memory
   munmap(ptr, len)                     unmap it exactly once
   MAP_PRIVATE                          copy-on-write: private edits
   MAP_SHARED                           shared pages / edit the file
   MAP_ANONYMOUS                        no file: zeroed shared memory
   MAP_FAILED                           the error sentinel (never NULL)
   page fault                           lazy per-page load on first touch
   msync(ptr, len, MS_SYNC)             flush a shared mapping to disk
   PTHREAD_PROCESS_SHARED               a mutex shared across processes

   Rules that keep mapping code correct:
     1. size a file with ftruncate BEFORE mapping it;
     2. check against MAP_FAILED, not NULL;
     3. pair every mmap with exactly one munmap;
     4. MAP_PRIVATE to read, MAP_SHARED to edit or communicate;
     5. protect shared memory with a process-shared mutex.
\end{verbatim}

\begin{osconcept}
This chapter closed a loop opened in Chapter~1. The address space,
we said, is the process's view of memory; mmap shows it is a view of
\emph{anything the kernel can page}---files, devices, or zeroed
anonymous regions. Sharing physical pages is the fastest and most
dangerous IPC there is, and the kernel provides both the mapping and
the locks, leaving the discipline to you. Nothing in this book
rewards understanding the machine more.
\end{osconcept}

We have built data structures, processes, servers, and shared memory.
The next chapter brings it all into one library: the memory
allocator---the code that hands out and reclaims every byte of the
heap, and the performance cliffs that hide inside \code{malloc}.

}{}
\IfFileExists{chapter13.tex}{% ============================================================================
%  chapter13.tex
%  Advanced Memory Management and Custom Allocators
% ============================================================================

\chapter{Advanced Memory Management and Custom Allocators}

\section{The Problem malloc Solves}

Every data structure in Chapter~4 asked \code{malloc} for bytes, and
Chapter~1 taught us to check the result. But \code{malloc} is not a
language feature; it is a \emph{program}---an allocator---that runs
inside your process, and its behavior decides the performance and
fragmentation of everything you build. This chapter opens that black
box: how allocators work, why memory fragments, and how to write your
own allocators when the defaults are the wrong shape for your workload.

The two extremes of allocation, drawn before we go anywhere:

\begin{verbatim}
   From the kernel's point of view (ch.12):          From your program's view:

   the heap is pages the kernel hands you.           the heap is a pool of
                                                     variable-size chunks you
     heap:  [ page ][ page ][ page ][ page ]        request and release:
                                                     malloc(12) -> a 12-byte
   malloc() gets a page when it needs more,          chunk, then free() it.
   free() returns memory to the allocator,           The allocator sits in
   only rarely back to the kernel.                   between, juggling the
                                                     chunks inside the pages.
\end{verbatim}

The allocator's contract is simple and brutal: \emph{given an
arbitrary sequence of malloc and free calls, satisfy every request,
reuse freed memory when possible, waste as little as possible, and do
it fast}. Getting all four right simultaneously is impossible in the
worst case; every allocator is a set of bets about how programs
behave.

\begin{osconcept}
The kernel provides memory only in pages, and only at the granularity
of whole pages. Individual small allocations would be absurdly wasteful
if each were its own page, so the OS delegates the juggling to a
user-space allocator inside your process. That is why \code{malloc}
is a library, not a syscall---it is an OS-adjacent program built on
the kernel's page grants, and its design is where the OS and the
application meet.
\end{osconcept}

\section{Free Lists and the Split/Merge Dance}

The heart of every general allocator is the \emph{free list}: a linked
list of unused chunks, threaded through the freed memory itself. When
you \code{free} a chunk, it is added to the list. When you
\code{malloc}, the allocator searches the list for a chunk big enough,
\emph{splits} it if it is larger than needed, and removes the piece
you get. When a chunk's neighbors are also free, the allocator
\emph{merges} them back into one bigger chunk---the only defense
against fragmentation.

\begin{verbatim}
   The free list, drawn as linked memory:

     heap:  [ A ][ F1 ][ B ][ F2 ][ C ]
                ^      ^
                |      +-- free chunks threaded through the heap

     free list:  F1 --> F2 --> NULL
                 (each free chunk stores the next free chunk's address
                  INSIDE its own space -- free memory is exactly the
                  space you can afford to use for bookkeeping)

   malloc(20) from F1 (size 48):
     - is F1 big enough? yes
     - split:  F1.used = 20 bytes for you; F1.tail = 28 bytes, freed
     - return the 20-byte chunk; the 28-byte remainder rejoins the list

   free(your chunk): add it to the list; if its neighbors are free,
     MERGE them, because 28 + 20 + 16 adjacent beats three scattered
     chunks of any smaller size.
\end{verbatim}

The split/merge cycle is the allocator's entire life. Its difficulty
is that the two operations fight each other: splitting creates small
fragments, and only merging heals them, but merging only works on
\emph{adjacent} freed chunks, and adjacency is not something you can
control after the fact.

\begin{lstlisting}[caption={A minimal free-list allocator, on a static arena.}]
#include <stdio.h>
#include <stdlib.h>
#include <stdint.h>
#include <string.h>

/* The arena: one big byte buffer the allocator hands out from. */
#define ARENA_SIZE 65536
static unsigned char arena[ARENA_SIZE];

/* Each chunk starts with a header describing it. The header lives in
   the arena, right before the user's bytes.                        */
typedef struct chunk {
    size_t size;          /* user-visible size in bytes */
    int free;             /* 1 if free, 0 if allocated */
    struct chunk *next;   /* next chunk in the arena (not the free list) */
} chunk;

#define HEADER_SIZE (sizeof(chunk))

/* The first chunk in the arena. Built lazily on first use. */
static chunk *arena_head = NULL;

static void arena_init(void) {
    chunk *c;

    /* One chunk covering the whole arena, free. */
    c = (chunk *) arena;
    c->size = ARENA_SIZE - HEADER_SIZE;
    c->free = 1;
    c->next = NULL;
    arena_head = c;
}

void *my_malloc(size_t size) {
    chunk *cur;

    if (arena_head == NULL) {
        arena_init();
    }

    /* Round the request up to at least a header, and to alignment. */
    if (size < 16) {
        size = 16;
    }
    size = (size + 7) & ~(size_t) 7;

    /* First-fit search: the first free chunk big enough wins. */
    cur = arena_head;
    while (cur != NULL) {
        if (cur->free && cur->size >= size) {
            break;
        }
        cur = cur->next;
    }
    if (cur == NULL) {
        return NULL;        /* no room */
    }

    /* Split: if the leftover is big enough to be useful, do it. */
    if (cur->size >= size + HEADER_SIZE + 8) {
        chunk *tail;

        tail = (chunk *) ((unsigned char *) (cur + 1) + size);
        tail->size = cur->size - size - HEADER_SIZE;
        tail->free = 1;
        tail->next = cur->next;
        cur->next = tail;

        cur->size = size;
    }

    cur->free = 0;
    return (void *) (cur + 1);
}

void my_free(void *p) {
    chunk *c;

    if (p == NULL) {
        return;
    }
    c = (chunk *) p - 1;     /* the header lives just before the user bytes */
    c->free = 1;

    /* Coalesce with the NEXT chunk if it is free. */
    if (c->next != NULL && c->next->free) {
        c->size = c->size + HEADER_SIZE + c->next->size;
        c->next = c->next->next;
    }
}

void *my_calloc(size_t n, size_t size) {
    void *p;

    p = my_malloc(n * size);
    if (p != NULL) {
        memset(p, 0, n * size);
    }
    return p;
}
\end{lstlisting}

Study the shape of this code, because it is the template for every
allocator in this chapter: a header per chunk, first-fit search,
split on allocation, coalesce-on-next on free. Two details are
deliberate: the header lives \emph{immediately before} the user's
bytes, so \code{my\_free} finds it with \code{(chunk *) p - 1}; and
the coalescing only merges with the \emph{next} chunk, because walking
backward would need a doubly linked structure or a separate metadata
area.

\begin{verbatim}
   The memory layout my_malloc produces:

     +-------+--------------------------+-------+-------------+
     | hdr A |  your 20 bytes           | hdr B | 28 free     |
     +-------+--------------------------+-------+-------------+
       ^ cur+1 returned                  ^ header of the remainder

   free(p) finds header A (p - 1), marks it free, and merges it with
   B if B is free -- so a later malloc(40) can use A+B together.
\end{verbatim}

\begin{osconcept}
The free list is a user-space data structure, but it rests on a kernel
fact: the heap is pages you were granted, and the allocator decides
how to divide them. The kernel cannot merge your fragments, because it
sees only pages, not your chunks. Splitting, merging, and searching
are therefore application-level problems---the OS hands out the clay,
and the allocator sculpts it.
\end{osconcept}

\section{Fragmentation: The Silent Tax}

Why do long-running programs grow forever even when they free
everything? The answer is \emph{fragmentation}, and it comes in two
kinds:

\begin{verbatim}
   INTERNAL fragmentation: you asked for bytes, you got a bigger
   chunk, and the difference is wasted.

     malloc(12)  ->  chunk of 16 (the allocator rounds up):
                     +--------+-------------+
                     | header | 16 bytes    |  4 bytes never used
                     +--------+-------------+
                     internal waste = 4 bytes, invisible to you.

   EXTERNAL fragmentation: free bytes exist, but they are in
   scattered pieces, none big enough for a request.

     heap after a long life:
     [ X ][ F 8 ][ X ][ F 8 ][ X ][ F 8 ][ X ][ F 8 ]
       ^   ^----^   ^      ^      ^      ^
       used  free   used   free   used   free   32 bytes free TOTAL,
       chunks of 8: no single 16-byte request can be satisfied.

   External fragmentation is the killer. It is why a process can have
   100 MB free and still fail a 1 MB malloc.
\end{verbatim}

The classic demonstration is the \emph{allocator stone soup}: allocate
many chunks, free every other one, then try to allocate one big chunk.
The bytes are all there; they are just not adjacent.

\begin{verbatim}
   A, B, C, D allocated.   Free B and D.   malloc(big):

     [A][  F  ][C][  F  ]        free = 2 x 10 bytes, separated
       10 free    10 free        big request of 22: FAILS even though
                                 20 bytes are free, because no single
                                 stretch reaches 22.

   The allocator's only weapon is COALESCING on free (merge adjacent
   neighbors), and even that cannot heal a checkerboard pattern where
   no two free chunks touch.
\end{verbatim}

\textbf{Concept} -- the \emph{binary buddy} allocator attacks external
fragmentation structurally. It divides memory into powers of two:
a request is rounded up to the next power of two, the arena is split
in half recursively to make that size, and frees merge only with their
\emph{buddy}---the adjacent half of the same parent. Because every
chunk has exactly one buddy of its own size, merging is guaranteed to
recover full parent blocks, and the free list becomes a set of
power-of-two lists with \code{O(1)} insertion.

\begin{lstlisting}[caption={The buddy allocator: powers of two, buddy merge.}]
#include <stdio.h>
#include <stdlib.h>

/* Sizes tracked: 2^min .. 2^max, stored as lists of free blocks. */
#define MIN_ORDER 4     /* smallest block: 16 bytes  */
#define MAX_ORDER 12    /* largest block:  4096 bytes */
#define MIN_SIZE  (1 << MIN_ORDER)

typedef struct block {
    struct block *next;
} block;

static block *free_lists[MAX_ORDER + 1];

/* A block's buddy is its neighbor of the same size within the parent. */
static void *buddy_of(void *addr, int order) {
    size_t size;

    size = (size_t) 1 << order;
    return (void *) (((size_t) addr) ^ size);
}

/* Return the largest power-of-two order that fits size. */
static int order_for(size_t size) {
    int order;

    size = size - 1;
    order = MIN_ORDER;
    while ((size >> order) != 0) {
        order = order + 1;
    }
    return order;
}

void *buddy_malloc(size_t size) {
    int order;
    int i;
    block *b;
    void *p;

    order = order_for(size);
    if (order > MAX_ORDER) {
        return NULL;
    }

    /* Find the smallest free list at or above this order. */
    for (i = order; i <= MAX_ORDER; i++) {
        if (free_lists[i] != NULL) {
            break;
        }
    }
    if (i > MAX_ORDER) {
        return NULL;     /* nothing big enough exists */
    }

    /* Take the block and split it downward to the right size. */
    b = free_lists[i];
    free_lists[i] = b->next;

    while (i > order) {
        size_t half;

        half = (size_t) 1 << (i - 1);
        p = (void *) ((size_t) b + half);      /* the buddy half */
        ((block *) p)->next = free_lists[i - 1];
        free_lists[i - 1] = (block *) p;       /* free the buddy */
        i = i - 1;
    }

    return b;
}

void buddy_free(void *addr, int order) {
    void *buddy;
    block *b;

    b = (block *) addr;
    while (order < MAX_ORDER) {
        buddy = buddy_of(addr, order);
        /* If the buddy is on our free list, merge: remove it, double. */
        if (free_lists[order] == buddy) {
            free_lists[order] = free_lists[order]->next;
            if (addr < buddy) {
                b = (block *) addr;
            } else {
                b = (block *) buddy;
            }
            addr = b;
            order = order + 1;
        } else {
            break;
        }
    }

    b->next = free_lists[order];
    free_lists[order] = b;
}
\end{lstlisting}

\begin{verbatim}
   The buddy split, drawn:

     arena of 4096 (order 12):
     [             4096            ]          free_lists[12]
            |  malloc(1500) -> order_for(1500) = 11 (2048)
            v
     [     2048       ][   2048 free   ]      free_lists[11]
       ^ you get this ^ the buddy is freed, can be re-split later

     when you free the 2048, its buddy is checked: if both free, the
     pair merges back into one 4096 -- automatically, because the
     buddy is always exactly adjacent and exactly the same size.
\end{verbatim}

The beauty is the property that makes buddy systems attractive in
operating systems themselves: \emph{merging never has to search}. A
block either has its buddy free (merge, guaranteed adjacent) or not
(split again later). No coalescing pass, no walking. The cost is
internal fragmentation---rounding to powers of two can waste nearly
half a block---which is why real kernels tune the order range and use
other allocators for small objects.

\begin{osconcept}
Fragmentation is the allocator's fundamental adversary, and the kernel
feels its effects at a larger scale: a process that fragments its heap
into holes may still need the OS to grant \emph{new} pages, inflating
its footprint. Buddy systems exist because the kernel itself must
allocate physical pages and merge them back, and it wants merging to
be cheap and structural---powers of two make ``is my neighbor free''
a single XOR. Fragmentation is an OS problem because it is a physical
reality of paged memory, not just a library concern.
\end{osconcept}

\section{Arenas and Object Pools}

The general allocator is flexible and fragmented. The \emph{arena}
(also called a \emph{region} or \emph{pool}) is the opposite bet:
allocate a big block, hand out pieces by bumping a pointer, and free
\emph{everything at once}. It is the right answer when allocations
share a lifetime---a request's scratch space, a connection's
buffers, a frame of a game.

\begin{verbatim}
   Arena allocator: a pointer that only moves forward.

     arena:  [ used ][ used ][ used ][             free          ]
                    ^                          ^
               bump pointer             arena + arena_size

     malloc(n): align n, copy the bump pointer out, advance it.
     free:      NOTHING. There is no individual free.
     reset:     set the bump pointer back to the start; the whole
                arena is free again in one instruction.

   The deal: individual frees are impossible; whole-life resets are
   free. Perfect when every allocation in a phase lives as long as
   the phase.
\end{verbatim}

\begin{lstlisting}[caption={arena.c: an arena allocator.}]
#include <stdio.h>
#include <stdlib.h>
#include <stdint.h>

typedef struct arena {
    unsigned char *mem;      /* the whole block              */
    size_t cap;              /* total capacity in bytes      */
    size_t off;              /* bytes handed out so far      */
} arena;

void arena_init(arena *a, size_t cap) {
    a->mem = malloc(cap);
    if (a->mem == NULL) {
        exit(1);
    }
    a->cap = cap;
    a->off = 0;
}

void *arena_alloc(arena *a, size_t size) {
    void *p;

    /* Align to 8 bytes so any type can live there. */
    size = (size + 7) & ~(size_t) 7;

    if (a->off + size > a->cap) {
        return NULL;         /* arena exhausted */
    }
    p = a->mem + a->off;
    a->off = a->off + size;
    return p;
}

/* Free the ENTIRE arena in one step: rewind the bump pointer. */
void arena_reset(arena *a) {
    a->off = 0;
}

void arena_destroy(arena *a) {
    free(a->mem);
    a->mem = NULL;
    a->cap = 0;
    a->off = 0;
}

int main(void) {
    arena a;
    char *s1;
    char *s2;

    arena_init(&a, 4096);

    s1 = arena_alloc(&a, 32);
    s2 = arena_alloc(&a, 64);
    if (s1 == NULL || s2 == NULL) {
        fprintf(stderr, "arena exhausted\n");
        return 1;
    }

    arena_reset(&a);          /* both s1 and s2 are gone at once */

    arena_destroy(&a);
    return 0;
}
\end{lstlisting}

\begin{verbatim}
   Arena vs malloc at a glance:

                     arena              malloc
   alloc cost        a few ops          search free lists
   free individual   not supported      supported
   free everything   one reset          one call per chunk
   fragmentation     none (linear)      external fragmentation
   best for          same-lifetime      unknown lifetimes,
                     batches            general use
\end{verbatim}

\textbf{Concept} -- an \emph{object pool} is the third classic: instead
of variable sizes, a pool handles objects of \emph{one fixed size}
(threads, connections, nodes). Allocation pops from a free list;
free pushes back. Nothing is ever handed to the OS or fragmented,
because every object is interchangeable and exactly one size.

\begin{lstlisting}[caption={An object pool: fixed-size objects, no fragmentation.}]
#include <stdio.h>
#include <stdlib.h>
#include <stdint.h>

#define POOL_OBJS 64
#define OBJ_SIZE  32

typedef struct pool_obj {
    struct pool_obj *next;   /* free list threaded through free objs */
    unsigned char pad[OBJ_SIZE - sizeof(struct pool_obj *)];
} pool_obj;

typedef struct pool {
    pool_obj *free_head;     /* next free object */
    pool_obj storage[POOL_OBJS];
} pool;

void pool_init(pool *p) {
    size_t i;

    for (i = 0; i < POOL_OBJS; i++) {
        p->storage[i].next = p->free_head;
        p->free_head = &p->storage[i];
    }
}

void *pool_alloc(pool *p) {
    pool_obj *o;

    o = p->free_head;
    if (o == NULL) {
        return NULL;         /* pool exhausted */
    }
    p->free_head = o->next;
    return (void *) o;
}

void pool_free(pool *p, void *ptr) {
    pool_obj *o;

    o = (pool_obj *) ptr;
    o->next = p->free_head;      /* push back on the free list */
    p->free_head = o;
}
\end{lstlisting}

\begin{verbatim}
   The pool's memory, drawn:

     pool.storage:  [o0][o1][o2]...[o63]    one fixed-size array

     free list threads through FREE objects only:
       free_head -> o5 -> o2 -> o31 -> NULL

     alloc: pop head; free: push head. O(1), zero searches, zero
     fragmentation, zero system calls. The pool is the data
     structure of the data structure.
\end{verbatim}

\begin{osconcept}
Arena and pool allocators exploit the kernel's own strategy: hand out
large regions rarely, subdivide them locally, and reclaim wholesale.
The kernel grants pages; the arena bumps a pointer through them; the
pool never asks the kernel again. Choosing the allocator is choosing
how often your program crosses the kernel boundary---and good
allocator design is simply the discipline of crossing it as rarely as
possible.
\end{osconcept}

\section{Practical: A Connection Pool for a Server}

Let us apply the chapter to the book's own recurring hero---the
server. A connection pool hands out fixed-size connection objects,
reuses them across clients, and never fragments and never calls
\code{malloc} after startup. This is the object pool of the previous
section wearing its production costume.

\begin{lstlisting}[caption={connpool.c: fixed-size connection objects, recycled.}]
#include <stdio.h>
#include <stdlib.h>

typedef struct conn {
    int fd;                /* the socket (0 = free slot)   */
    int in_use;            /* 1 while a client owns it     */
    struct conn *next;     /* free-list threading           */
} conn;

typedef struct conn_pool {
    conn *slots;           /* the fixed array of conns      */
    int capacity;
    conn *free_head;       /* head of the free list         */
} conn_pool;

int conn_pool_init(conn_pool *p, int capacity) {
    int i;

    p->slots = malloc(sizeof(conn) * (size_t) capacity);
    if (p->slots == NULL) {
        return -1;
    }
    p->capacity = capacity;
    p->free_head = NULL;

    /* Thread every slot onto the free list. */
    for (i = 0; i < capacity; i++) {
        p->slots[i].fd = 0;
        p->slots[i].in_use = 0;
        p->slots[i].next = p->free_head;
        p->free_head = &p->slots[i];
    }
    return 0;
}

/* Hand out a free connection object, or NULL if the pool is full. */
conn *conn_pool_acquire(conn_pool *p) {
    conn *c;

    c = p->free_head;
    if (c == NULL) {
        return NULL;
    }
    p->free_head = c->next;
    c->in_use = 1;
    return c;
}

/* Return a connection to the pool. */
void conn_pool_release(conn_pool *p, conn *c) {
    c->in_use = 0;
    c->fd = 0;
    c->next = p->free_head;
    p->free_head = c;
}

int main(void) {
    conn_pool pool;
    conn *c1;
    conn *c2;

    if (conn_pool_init(&pool, 8) < 0) {
        return 1;
    }

    c1 = conn_pool_acquire(&pool);
    c2 = conn_pool_acquire(&pool);
    if (c1 == NULL || c2 == NULL) {
        fprintf(stderr, "pool exhausted\n");
        return 1;
    }
    printf("acquired two connections\n");

    conn_pool_release(&pool, c1);
    conn_pool_release(&pool, c2);

    /* Reacquire the same objects -- same addresses, no malloc. */
    c1 = conn_pool_acquire(&pool);
    printf("reacquired (pool recycles storage)\n");
    conn_pool_release(&pool, c1);

    free(pool.slots);
    return 0;
}
\end{lstlisting}

\begin{verbatim}
   Why a pool and not malloc per connection:

     malloc/free per connection   pool
     -------------------------    ----
     searches free lists          O(1) pop/push
     can fragment over a day      zero fragmentation (fixed size)
     may cross the kernel         no syscalls after init
     allocator lock contention    single-threaded pop/push (or one
                                   lock for the pool)
     allocator can't recycle      addresses recycle: caches stay warm

   The price: a fixed maximum (capacity) and objects that cannot
   outlive the pool. Both are features when you know your limits.
\end{verbatim}

\begin{researchnote}[The measured case for custom allocators]
Chapter~1 introduced the empirical study by Berger, Zorn, and
McKinley that most custom allocators do not beat \code{malloc}. Now
that you can read their argument in context, the nuance matters: they
found that allocators win when they \emph{specialize to a workload's
lifetime patterns}---exactly what arenas and pools do here---and lose
when they merely reimplement a general allocator slightly differently.

\begin{verbatim}
   The paper's real conclusion, restated honestly:

       "Generic custom allocators rarely help; allocators that match
        a program's allocation AND deallocation patterns -- regions,
        pools, per-thread caches -- are the ones that show gains."

   This chapter is that sentence in practice: the arena (reset in
   O(1)), the pool (fixed size, no fragmentation), and the buddy
   (structural merge) are lifetime-aware designs. A hand-rolled
   malloc-lookalike is the one approach the research says to skip.
\end{verbatim}

So the research note from Chapter~1 and this chapter agree: write
\emph{specialized} allocators when your workload has structure, not
generic ones on a whim. Measure first; the paper's methodology (link
in many allocators, time them) is the correct one to steal.

\paperref{E.~D.~Berger, B.~G.~Zorn, K.~S.~McKinley, ``Reconsidering
Custom Memory Allocation,'' ACM OOPSLA 2002}{https://people.cs.umass.edu/~emery/pubs/berger-oopsla2002.pdf}
\end{researchnote}

\section{Structures in Bytes}

Every allocator in this chapter hides a header in front of the bytes
you receive. On a 64-bit machine the structures are:

\begin{verbatim}
   my_malloc chunk header  24 bytes   offset size field
                                      0      8    size   (user size)
                                      8      4    free
                                      12     4    pad
                                      16     8    next   (next in arena)
   arena                   24 bytes   0      8    mem
                                      8      8    cap
                                      16     8    off
   pool_obj                32 bytes   0      8    next   (free-list link)
                                      8      24   pad[24]
   buddy free block         8 bytes   0      8    next   (threaded inside)
\end{verbatim}

The header sits immediately before the user's bytes, so
\code{my\_free} recovers it with \code{(chunk *) p - 1}; because the
header is 24 bytes and the sizes are rounded to 8, every user pointer
is 8-byte aligned. The arena is a bump pointer plus a capacity, 24
bytes total; the pool is one array of fixed-size objects whose free
list is threaded through the objects themselves. These byte counts are
the whole point: an allocator is a small set of byte records and the
rules for moving them.

\section{Summary}

\begin{verbatim}
   The allocator toolbox, complete:

   free list        freed chunks threaded through memory
   first-fit        the simplest search; split on alloc
   coalescing       merge adjacent frees (the anti-fragmentation cure)
   internal frag    waste inside your own chunk (rounding)
   external frag    free bytes in pieces too small to use
   buddy system     powers of two; buddy merge is O(1)
   arena            bump pointer, reset everything at once
   object pool      one fixed size, O(1) pop/push, no frag

   Rules that keep allocator code correct:
     1. every alloc returns NULL or a valid pointer -- check it;
     2. a header is reachable from the user pointer (p - 1);
     3. coalesce with the NEXT chunk on free, at minimum;
     4. match the allocator to the LIFETIME pattern, not to a mood;
     5. measure before you write one; the default is usually fine.
\end{verbatim}

\begin{osconcept}
Allocators are where the application and the OS shake hands: the
kernel grants pages, the allocator divides them, and fragmentation
wanders back up the chain as wasted pages. Every allocator in this
chapter is a different answer to the same kernel-shaped question--
how rarely can I ask for more?--and the best ones ask almost never.
This is the deepest layer of memory management before we leave the
heap for good.
\end{osconcept}

With memory mastered, the next chapter switches from bytes to
\emph{symbols}: parsing text with state machines, the technique that
powers every config parser, every protocol decoder, and every
compiler's front end.

}{}
\IfFileExists{chapter14.tex}{% ============================================================================
%  chapter14.tex
%  Text Parsing and State Machines
% ============================================================================

\chapter{Text Parsing and State Machines}

\section{Why Parsing Is a Systems Skill}

Every program that reads text must \emph{parse} it: take raw bytes and
extract the structure the text encodes. Configuration files, HTTP
headers, JSON, CSV, command lines, and every wire protocol are all
text with rules, and the code that reads them is a \emph{parser}.
Parsers are everywhere, they are written by hand more often than you
would expect, and they are the classic home of two bugs: crashes on
malformed input, and silent misinterpretation of ``almost valid''
input.

This chapter builds a real parser end to end. The example is the
command language of \emph{MiniDB}, a small key-value database written
in C that lives beside this book (in the workspace's
\code{MiniDB-deepseek-implementation} directory). MiniDB accepts
commands like \code{INSERT 42 100} and \code{SELECT 7} from a prompt;
before any data is touched, every line must be \emph{lexed} into
tokens and \emph{parsed} into a command. We will study its lexer and
parser as the running example, and adapt its code to this book's
style.

\begin{verbatim}
   The layers of parsing:

     raw text:   "INSERT 42 100\n"
                     |
                     v   LEX (tokenize, one char at a time)
     tokens:     (KEYWORD: INSERT) (INTEGER: 42) (INTEGER: 100) (END)
                     |
                     v   PARSE (validate shape, fill the struct)
     command:    { type = COMMAND_INSERT, key = 42, value = 100 }

   - A TOKEN is the smallest meaningful unit (a word, a number, a
     symbol). The LEXER (tokenizer) splits text into tokens and
     discards what is insignificant (whitespace).
   - The PARSER checks the token stream against the grammar and fills
     a struct (here: a Command). The database then executes it.
\end{verbatim}

\begin{osconcept}
Parsing is a systems skill because text is the universal interface
between programs: every configuration file, log line, and protocol
message your daemon reads arrives as bytes and must be understood
safely. The OS delivers you bytes; it will not protect you from
malformed input. Turning unreliable text into reliable structure, and
failing cleanly when the text is wrong, is the programmer's own
problem---which is exactly why parsers are where so much care lives.
\end{osconcept}

\section{Tokens and the Token Stream}

A \emph{token} is a run of characters that counts as one unit: an
identifier, a number, a quoted string, a symbol. The tokenizer's only
job is splitting; deciding what tokens \emph{mean} belongs to the
parser. MiniDB has just three token kinds, and that simplicity is
deliberate:

\begin{verbatim}
   MiniDB's language, in one screen:

     INSERT <key> <value>     store a key/value pair
     SELECT <key>             fetch one key
     DELETE <key>             remove one key
     RANGE  <lo> <hi>         fetch all keys in [lo, hi]
     PRINT                    dump the whole database
     EXIT                     leave the REPL
     HELP                     list the commands

   Every command is one verb plus a FIXED number of integers.
   That tiny grammar keeps the parser flat (no tree needed yet).
\end{verbatim}

\begin{lstlisting}[caption={The token types and the token struct.}]
typedef enum {
    TOKEN_KEYWORD,     /* names a command verb; read "keyword" */
    TOKEN_INTEGER,     /* a bare decimal literal; read "int_value" */
    TOKEN_END          /* sentinel: the stream ends here        */
} TokenType;

typedef enum {
    KEYWORD_INSERT, KEYWORD_SELECT, KEYWORD_DELETE, KEYWORD_RANGE,
    KEYWORD_PRINT, KEYWORD_EXIT, KEYWORD_HELP, KEYWORD_UNKNOWN
} Keyword;

typedef struct {
    TokenType type;
    Keyword   keyword;    /* valid when type == TOKEN_KEYWORD */
    int       int_value;  /* valid when type == TOKEN_INTEGER */
} Token;
\end{lstlisting}

The field comments make the \emph{discriminator} explicit: which member
of the struct is meaningful depends on \code{type}. That is the tagged
union pattern of Chapter~3, back for its first serious workout.

Two design decisions in these few lines will matter through the whole
chapter. First, every stream \emph{ends with a sentinel}
(\code{TOKEN\_END}): the parser can check for ``the stream is done''
without knowing the length by other means, and can never read past the
last real token. Second, the lexer \emph{allocates nothing}: the
caller owns the token array, and the lexer just fills it up to a fixed
bound. A command line has at most a handful of tokens, so a fixed
array is not a limitation---it is a simplification that removes an
entire class of memory bugs from the parser.

\begin{verbatim}
   "  exit  "  ====>  (KEYWORD: EXIT) (END)      whitespace vanishes
   "INSERT 42" ====>  (KEYWORD: INSERT) (INTEGER: 42) (END)
   "select 7"  ====>  (KEYWORD: SELECT) (INTEGER: 7)  (END)
   (keywords are case-insensitive; MINIDB matches them uppercased)
\end{verbatim}

\section{The Lexer as a State Machine}

The lexer is a tiny character-level \emph{state machine}: it walks the
input once, left to right, and each character moves it into exactly one
state. A \emph{state machine} (in its simplest form, a \emph{finite
automaton}) has a finite set of states, one start state, and
\emph{transitions}---rules saying what happens when a character
arrives in each state. Some states are \emph{accepting}: if the input
ends there, the input was valid.

\begin{verbatim}
   MiniDB's lexer, drawn as a state machine:

                              isspace (skip, repeat)
                              |
        START --------------->+<-- SKIP_WS
         |
         | isdigit                     | any non-whitespace
         v                             v
       INTEGER -----------------> (back to START for the next char)
         |  overflow                       ^
         v                                 |
       ERROR <--isalpha? no, no match-- WORD (matched a keyword -> token)
         |
         +-- "any character outside the grammar" also lands here

   The accepting state is END OF INPUT: when the scan runs out of
   characters, it appends TOKEN_END. Everything else either emits a
   token and loops, or enters ERROR and aborts the whole scan.

   One character, one decision, never a lookahead, never a backup:
   that is exactly what a finite automaton is.
\end{verbatim}

The lexer below is the MiniDB lexer, faithful in structure and
condensed in one place: the original spells out each of the seven
keywords with its own comparison block for teaching clarity; here we
match against a table in a loop. The states---skip whitespace,
integer, word, stray symbol, end---are unchanged.

\begin{lstlisting}[caption={The MiniDB lexer (condensed from lexer.c).}]
#include <ctype.h>
#include <limits.h>
#include <string.h>

#define LEXER_MAX_TOKENS 16

typedef enum {
    LEX_OK,
    LEX_ERROR_UNKNOWN_WORD,     /* unrecognized word or stray symbol  */
    LEX_ERROR_BAD_NUMBER,       /* digit run that would overflow int  */
    LEX_ERROR_TOO_MANY_TOKENS   /* more tokens than the array holds   */
} LexResult;

static const char *const KEYWORD_TABLE[KEYWORD_UNKNOWN] = {
    "INSERT", "SELECT", "DELETE", "RANGE",
    "PRINT", "EXIT", "HELP"
};

/* The ONLY place tokens enter the output array. Bounds-checked. */
static LexResult append_token(Token *tokens, size_t *token_count, Token t) {
    if (*token_count >= LEXER_MAX_TOKENS) {
        return LEX_ERROR_TOO_MANY_TOKENS;
    }
    tokens[*token_count] = t;
    *token_count = *token_count + 1;
    return LEX_OK;
}

LexResult lex(const char *input, Token *tokens, size_t *token_count) {
    const char *p;
    size_t count;

    p = input;
    count = 0;

    while (*p != '\0') {

        /* ---- State: skip whitespace (never significant) ---- */
        if (isspace((unsigned char) *p)) {
            p = p + 1;
            continue;
        }

        /* ---- State: integer literal ---- */
        if (isdigit((unsigned char) *p)) {
            Token t;
            int value;

            value = 0;
            while (isdigit((unsigned char) *p)) {
                int digit;

                digit = *p - '0';
                /* Overflow guard: stop before value * 10 overflows. */
                if (value > (INT_MAX - digit) / 10) {
                    return LEX_ERROR_BAD_NUMBER;
                }
                value = value * 10 + digit;
                p = p + 1;
            }

            t.type = TOKEN_INTEGER;
            t.int_value = value;
            if (append_token(tokens, &count, t) != LEX_OK) {
                return LEX_ERROR_TOO_MANY_TOKENS;
            }
            continue;
        }

        /* ---- State: word (a keyword, upcased for matching) ---- */
        if (isalpha((unsigned char) *p)) {
            char word[8];     /* longest keyword + NUL */
            size_t wlen;
            int matched;
            size_t i;

            wlen = 0;
            while (isalpha((unsigned char) *p)) {
                if (wlen >= sizeof(word) - 1) {
                    return LEX_ERROR_UNKNOWN_WORD;   /* too long */
                }
                word[wlen] = (char) toupper((unsigned char) *p);
                wlen = wlen + 1;
                p = p + 1;
            }
            word[wlen] = '\0';

            matched = 0;
            for (i = 0; i < (size_t) KEYWORD_UNKNOWN; i++) {
                if (wlen == strlen(KEYWORD_TABLE[i]) &&
                    memcmp(word, KEYWORD_TABLE[i], wlen) == 0) {
                    Token t;

                    t.type = TOKEN_KEYWORD;
                    t.keyword = (Keyword) i;
                    if (append_token(tokens, &count, t) != LEX_OK) {
                        return LEX_ERROR_TOO_MANY_TOKENS;
                    }
                    matched = 1;
                    break;
                }
            }
            if (matched == 0) {
                return LEX_ERROR_UNKNOWN_WORD;
            }
            continue;
        }

        /* ---- State: stray symbol (not part of the language) ---- */
        return LEX_ERROR_UNKNOWN_WORD;
    }

    /* ---- State: end of input -- emit the sentinel ---- */
    {
        Token end;

        end.type = TOKEN_END;
        end.keyword = KEYWORD_UNKNOWN;
        end.int_value = 0;
        if (append_token(tokens, &count, end) != LEX_OK) {
            return LEX_ERROR_TOO_MANY_TOKENS;
        }
    }

    *token_count = count;
    return LEX_OK;
}
\end{lstlisting}

Read the code against the diagram and three facts become visible. The
overflow guard is a genuine state-machine error transition: it aborts
mid-number rather than emit a wrong token. The word buffer is bounded
(\code{word[8]}, and the capacity check before every write), so a
100-character ``word'' is rejected, never overflowed. And on \emph{any}
error the function returns immediately; the caller must treat the
partially-filled stream as garbage. There is no recovery-and-continue;
MiniDB chooses fail-fast, which is the right choice for a command
line.

\begin{verbatim}
   The states, recapped in one line each:

     SKIP_WS    whitespace is never significant
     INTEGER    consume digits, guard overflow, emit INTEGER
     WORD       upcase and match against the keyword table
     STRAY      any other byte is a one-character unknown word
     END        append the TOKEN_END sentinel and finish

   No dynamic allocation anywhere. The caller owns the array; the
   lexer owns nothing. That one decision is why the whole parser can
   be reasoned about with plain C arrays and explicit bounds.
\end{verbatim}

\begin{ruleofthumb}[Draw the machine before writing the loop]
Any parser with more than one kind of token deserves a state diagram
first. Sketch the states, walk the happy path, walk two malformed
inputs, and only then write the \code{while}. The diagram is the
specification; the code is the transcription. Parsers written the
other way around are how ``almost valid'' input silently becomes
wrong data.
\end{ruleofthumb}

\section{The Parser: Validating the Token Stream}

Where the lexer is a \emph{scanner}, the parser is a \emph{validator}.
It answers two questions: \emph{is this stream a legal command, and
what does it mean?} It does not look at characters again---the lexer
already did that work. It walks the token array with a cursor and
checks the stream's \emph{shape} and \emph{arity} against the grammar.

MiniDB's grammar is the smallest possible: one verb, plus a fixed
count of integer arguments. There is nothing to nest, so the parser
barely needs recursion at all. It works in three moves:

\begin{verbatim}
   The parser's three moves:

   1. DISPATCH on tokens[0].  It must be a KEYWORD; a switch on its
      value selects the rule. A bare integer first is not a verb:
      PARSE_ERROR_UNKNOWN_COMMAND.

   2. EXACT-ARITY CHECK.  Each command calls parse_fixed_arguments()
      with the number of integers it expects. One helper consumes
      exactly that many integers, then REQUIRES the very next token
      to be END. Three failures in one place:
         - stream ends early          -> MISSING_ARGUMENT
         - a non-integer where an int belongs -> BAD_ARGUMENT
         - anything left over         -> EXTRA_ARGUMENT

   3. BUILD THE Command.  On success, the integers are copied into the
      right fields and command->type is set.
\end{verbatim}

\begin{lstlisting}[caption={The exact-arity engine shared by every command.}]
typedef enum {
    PARSE_OK,
    PARSE_ERROR_EMPTY,
    PARSE_ERROR_UNKNOWN_COMMAND,
    PARSE_ERROR_MISSING_ARGUMENT,
    PARSE_ERROR_EXTRA_ARGUMENT,
    PARSE_ERROR_BAD_ARGUMENT
} ParseResult;

/* Consume exactly argc integers starting at tokens[start]; then the
   next token MUST be TOKEN_END. Results land in args[].          */
static ParseResult parse_fixed_arguments(const Token *tokens,
                                         size_t token_count,
                                         size_t start, size_t argc,
                                         int *args) {
    size_t i;

    for (i = 0; i < argc; i++) {
        if (start >= token_count) {
            return PARSE_ERROR_MISSING_ARGUMENT;   /* no sentinel at all */
        }
        if (tokens[start].type == TOKEN_END) {
            return PARSE_ERROR_MISSING_ARGUMENT;   /* stream ran short */
        }
        if (tokens[start].type != TOKEN_INTEGER) {
            return PARSE_ERROR_BAD_ARGUMENT;       /* wrong kind of token */
        }
        args[i] = tokens[start].int_value;
        start = start + 1;
    }

    if (start >= token_count) {
        return PARSE_ERROR_MISSING_ARGUMENT;
    }
    if (tokens[start].type != TOKEN_END) {
        return PARSE_ERROR_EXTRA_ARGUMENT;         /* too many arguments */
    }
    return PARSE_OK;
}
\end{lstlisting}

\begin{verbatim}
   The same walker, three failing inputs:

     INSERT 42          ->  expects 2 ints, finds 1, then END
                            -> MISSING_ARGUMENT
     INSERT 42 100 5    ->  two ints consumed, then an INTEGER
                            where END must be -> EXTRA_ARGUMENT
     INSERT x 100       ->  an integer demanded, KEYWORD found
                            -> BAD_ARGUMENT

   One helper, three errors: the arity is checked by construction,
   because the "then it must end" line is part of the same function.
\end{verbatim}

The dispatch is a plain switch over the keyword, each case stating its
arity in one line---that number \emph{is} the grammar rule:

\begin{lstlisting}[caption={parse\_command: dispatch on the keyword, fill the Command.}]
typedef enum {
    COMMAND_INSERT, COMMAND_SELECT, COMMAND_DELETE, COMMAND_RANGE,
    COMMAND_PRINT, COMMAND_EXIT, COMMAND_HELP, COMMAND_NONE
} CommandType;

/* The proto-AST: one flat node, because every statement is one verb
   plus at most two integers. */
typedef struct {
    CommandType type;
    int key;      /* INSERT/SELECT/DELETE/RANGE: the primary key   */
    int value;    /* INSERT only: the value to store               */
    int key_end;  /* RANGE only: the inclusive upper bound         */
} Command;

ParseResult parse_command(const Token *tokens, size_t token_count,
                          Command *command) {
    int args[2];               /* largest arity in the grammar is two */
    ParseResult r;

    /* Initialize every field first: unused fields stay defined. */
    command->type = COMMAND_NONE;
    command->key = 0;
    command->value = 0;
    command->key_end = 0;

    if (token_count == 0 || tokens[0].type == TOKEN_END) {
        return PARSE_ERROR_EMPTY;          /* nothing there at all */
    }
    if (tokens[0].type != TOKEN_KEYWORD) {
        return PARSE_ERROR_UNKNOWN_COMMAND;/* an integer is not a verb */
    }

    switch (tokens[0].keyword) {
    case KEYWORD_INSERT:
        r = parse_fixed_arguments(tokens, token_count, 1, 2, args);
        if (r != PARSE_OK) return r;
        command->type = COMMAND_INSERT;
        command->key = args[0];
        command->value = args[1];
        return PARSE_OK;

    case KEYWORD_SELECT:
        r = parse_fixed_arguments(tokens, token_count, 1, 1, args);
        if (r != PARSE_OK) return r;
        command->type = COMMAND_SELECT;
        command->key = args[0];
        return PARSE_OK;

    case KEYWORD_DELETE:
        r = parse_fixed_arguments(tokens, token_count, 1, 1, args);
        if (r != PARSE_OK) return r;
        command->type = COMMAND_DELETE;
        command->key = args[0];
        return PARSE_OK;

    case KEYWORD_RANGE:
        r = parse_fixed_arguments(tokens, token_count, 1, 2, args);
        if (r != PARSE_OK) return r;
        command->type = COMMAND_RANGE;
        command->key = args[0];
        command->key_end = args[1];
        return PARSE_OK;

    /* The three zero-argument commands all reduce to: END must follow. */
    case KEYWORD_PRINT:
    case KEYWORD_EXIT:
    case KEYWORD_HELP:
        r = parse_fixed_arguments(tokens, token_count, 1, 0, args);
        if (r != PARSE_OK) return r;
        command->type = (CommandType) (COMMAND_PRINT +
            (tokens[0].keyword - KEYWORD_PRINT));
        return PARSE_OK;

    default:
        return PARSE_ERROR_UNKNOWN_COMMAND;
    }
}
\end{lstlisting}

\begin{verbatim}
   The token-to-Command mapping, drawn:

     tokens:  (KEYWORD: INSERT) (INTEGER: 42) (INTEGER: 100) (END)
                                     |              |
                                     v              v
     Command: { type = COMMAND_INSERT, key = 42, value = 100, key_end = 0 }

     tokens:  (KEYWORD: RANGE) (INTEGER: 10) (INTEGER: 20) (END)
                                     |              |
                                     v              v
     Command: { type = COMMAND_RANGE, key = 10, value = 0, key_end = 20 }
\end{verbatim}

Calling the output a \emph{proto-}AST is a hint about the future. A
real query such as \code{SELECT price FROM items WHERE id = 42} would
need nested expressions, so \code{Command} would grow pointers to
sub-structures and the parser would recurse. MiniDB's statements are
all one verb plus at most two integers, so the ``tree'' collapses to
one node---but the seam where tree-ness would return is identified,
and that is exactly the place our Chapter~4 trees would slot in.

\section{The REPL: Read, Gate, Execute}

\subsection{What a REPL Is}

A \emph{REPL}---\textbf{R}ead-\textbf{E}val-\textbf{P}rint
\textbf{L}oop---is an interactive top-level loop. It is the oldest user
interface in programming, invented for \emph{Lisp} and later adopted
by early BASIC and every language console since (Python, Ruby,
JavaScript, SQL: all of them drop you into a REPL). It means exactly
what its name says, three phases repeated forever:

\begin{verbatim}
   the three phases of every REPL:

     READ    get one complete unit of input   (one line, here)
     EVAL    do the work on it                (lex + parse + execute)
     PRINT   show the outcome to the user     (the reply message)
                ^
                +---- and loop, forever ----------+

   Contrast with a BATCH program, which reads all input, computes, and
   prints once: a REPL is *interactive* -- it must answer each input
   before waiting for the next. That single requirement shapes
   everything: prompt, flush, per-line errors, loop.
\end{verbatim}

\textbf{What Lisp was.} Lisp (for \emph{LISt Processing}) is a family
of programming languages designed by John McCarthy starting in 1958---
the second-oldest high-level language after Fortran. Its founding idea
was that the \emph{list} is the universal structure: a Lisp program is
itself a list of symbols, so a program can read another program, alter
it, and run the result. That property (\emph{homoiconicity}, ``code is
data'') is why Lisp became the language of artificial-intelligence
research for three decades. Lisp also gave computing two gifts that
every language now takes for granted: \emph{garbage collection}
(automatic reclamation of memory, which we build pieces of in
Chapter~13) and the \emph{REPL} itself. A Lisp programmer did not
write, compile, and run---they talked to the machine, one expression
at a time, and the machine answered immediately. Interactive
development was Lisp's invention, and the MiniDB console is its
descendant.

Two consequences follow from ``interactive.'' First, the loop must
\emph{survive} bad input: one bad line cannot kill the session, so
every error becomes ``print a message, continue.'' Second, the loop
must announce itself with a \emph{prompt} (\code{minidb> }), and that
prompt must be \emph{flushed}: when stdout is a pipe rather than a
terminal, it is fully buffered, and without \code{fflush(stdout)} the
prompt would sit in a buffer while the program waited for input it
never visibly requested. Interactive software must flush.

\begin{verbatim}
   The REPL is the front end of every interactive system -- a shell,
   a database console, a debugger, a calculator. MiniDB's is small,
   but it is the same shape as all of them.
\end{verbatim}

\subsection{REPL Gating: Every Stage Guards the Next}

Now the word \emph{gating}. The REPL we are about to read does not
pipe raw input to the database. It passes each line through a sequence
of validation stages, and each stage is a \emph{gate}: it either
converts the input into a more structured form and passes it through,
or it rejects the line and the whole pipeline stops there. The
database never sees a bad line, because a bad line never gets through
two gates to reach it.

\begin{verbatim}
   Gating, drawn as a funnel:

      raw line
        |
        v
   +----------+   pass?   +----------+   pass?   +-----------+
   | LEX gate |---------->| PARSE    |---------->| EXECUTE   |
   | chars -> | fail:     | tokens ->| fail:     | command - |
   | tokens   | print,    | Command  | print,    | > storage |
   +----------+ continue  +----------+ continue  +-----------+
      |                       |                       |
      v                       v                       v
   "Lex error."           "Parse error."           the reply

   Two walls protect the storage engine from the user: a line must
   survive the lexer AND the parser before any data is touched.

   Why "gate"? A gate lets only well-formed traffic through. Each
   stage raises the well-formedness bar one level: raw bytes -> tokens
   -> a validated Command. Whatever fails a bar never travels further.
\end{verbatim}

Gating is the pattern that makes layered systems safe. The lexer can
reject garbage cheaply, so the parser only ever sees sensible tokens;
the parser rejects wrong-shaped commands, so the storage engine only
ever sees validated operations. Each stage trusts its input, but the
trust is earned one gate back. MiniDB drives the whole pipeline this
way: read a line, lex it, parse it, route it, echo the reply, repeat.

\begin{lstlisting}[caption={The REPL pipeline (adapted from main.c).}]
#include <stdio.h>
#include <string.h>

/* Returns 1 if the line contains nothing but whitespace. */
static int is_blank_line(const char *line) {
    int i;

    for (i = 0; line[i] != '\0'; i++) {
        if (line[i] != ' ' && line[i] != '\t' &&
            line[i] != '\n' && line[i] != '\r') {
            return 0;
        }
    }
    return 1;
}

int main(void) {
    Token tokens[LEXER_MAX_TOKENS];   /* the caller owns the array */
    size_t token_count;
    Command cmd;
    char line[256];

    for (;;) {
        /* Prompt, and FLUSH: without this, a piped stdout would hold
           the prompt in a buffer while we wait for input it never
           shows. Interactive programs must flush.               */
        printf("minidb> ");
        fflush(stdout);

        /* EOF (Ctrl-D or a closed pipe) ends the session. */
        if (fgets(line, sizeof(line), stdin) == NULL) {
            break;
        }
        if (is_blank_line(line)) {
            continue;                  /* noise: skip silently */
        }

        /* Stage 1: LEX. A failed scan invalidates the stream. */
        token_count = 0;
        if (lex(line, tokens, &token_count) != LEX_OK) {
            printf("Lex error.\n");
            continue;
        }

        /* Stage 2: PARSE. The command struct is only filled on OK. */
        if (parse_command(tokens, token_count, &cmd) != PARSE_OK) {
            printf("Parse error. Type HELP for commands.\n");
            continue;
        }

        /* EXIT is REPL control flow, not a database operation. */
        if (cmd.type == COMMAND_EXIT) {
            printf("Goodbye!\n");
            break;
        }

        /* Stage 3: EXECUTE. MiniDB routes the Command to its B+ tree
           (Chapter 4) and prints the returned message. */
        execute_command(&cmd);
    }

    return 0;
}
\end{lstlisting}

\begin{verbatim}
   The loop, drawn as a state machine of its own:

        +----> print "minidb> " ; fflush
        |           |
        |           v
        |   fgets() -- EOF? ------> break
        |           |
        |           v
        |   blank? --------------> continue (skip)
        |           |
        |           v
        |   LEX -- error? --------> print, continue
        |           |
        |           v
        |   PARSE -- error? ------> print, continue
        |           |
        |           v
        |   EXIT? ----------------> break
        |           |
        |           v
        |   execute ; print reply  --> back to the top
        +-----------+
\end{verbatim}

The three details worth internalizing:

\begin{itemize}
  \item \emph{The stack owns the buffers.} \code{Token tokens[16]},
    \code{Command cmd}, and \code{char line[256]} are all locals. The
    lexer and parser write into caller-owned objects and allocate
    nothing. Memory reasoning in the front end is trivial: stack in,
    stack out, nothing to free.
  \item \emph{Every stage gates the next.} A lex error prints and
    skips parsing; a parse error prints and skips executing. A bad
    line can never reach the storage engine. The database is protected
    from the user by two walls of validation, and neither wall can be
    jumped.
  \item \emph{\code{EXIT} is REPL control flow, not a data command.}
    The REPL checks \code{cmd.type == COMMAND\_EXIT} before routing.
    Where the loop lives, the exit decision lives.
\end{itemize}

\section{The Complete Journey: ``INSERT 42 100''}

Let us follow one real line end to end, because the pipeline is the
point and the layers are where the meaning becomes structure.

\begin{verbatim}
   Step 1 -- the string.
   fgets() puts  'I','N','S','E','R','T',' ','4','2',' ','1','0','0',
                 '\n','\0'  into char line[256].

   Step 2 -- lex.
   'I' is a letter -> WORD mode: absorb INSERT, uppercase it, match
       the keyword table, emit { TOKEN_KEYWORD, KEYWORD_INSERT }.
   ' ' -> skipped.  '4' is a digit -> INTEGER mode: value = 4*10+2 = 42.
   ' ' -> skipped.  '1' is a digit -> INTEGER mode: value = 100.
   '\n' -> skipped; '\0' ends the loop; the END sentinel is appended.

   index:   0          1          2           3
          (KEYWORD)   (INTEGER)   (INTEGER)   (END)
           INSERT        42          100
   token_count == 4

   Step 3 -- parse.
   tokens[0] is a KEYWORD named INSERT -> arity 2:
   consume 42 and 100, require END, fill the struct:
       Command { type = COMMAND_INSERT, key = 42, value = 100 }

   Step 4 -- execute.
   dispatch on COMMAND_INSERT: the pair (42, 100) is placed in the
   storage engine -- in MiniDB, a B+ tree (our Chapter 4 structure).

   Step 5 -- reply.
   "inserted key 42 = 100" travels back to the REPL and is printed.

   One typed line, five stages, three structs, zero heap allocations
   in the front end, and the data lives in a balanced sorted tree.
\end{verbatim}

\begin{verbatim}
   A real session:

     minidb> INSERT 42 100
     inserted key 42 = 100
     minidb> SELECT 42
     key 42 = 100
     minidb> RANGE 1 50
     2 keys in [1, 50]
     minidb> DELETE 7
     deleted key 7
     minidb> SELECT 7
     key 7 not found
     minidb> EXIT
     Goodbye!
\end{verbatim}

\section{Layering: The Tree Never Learns the Language}

The deepest lesson of MiniDB is not any single function; it is the
\emph{direction of the dependencies}. The lexer knows nothing about
the parser. The parser knows nothing about the storage engine. The
tree stores \code{(int, int)} pairs and cannot tell \code{SELECT} from
\code{DELETE}---and does not need to. Each layer only understands the
value handed to it:

\begin{verbatim}
   The include graph is a strict, acyclic stack:

        main.c  (the REPL)
           |
           v
      database.h   <-- the seam: Command -> tree call -> reply
        /      \
       v        v
   parser.h   btree.h        (btree = our Chapter 4 B+ tree)
       |
       v
    lexer.h

   Read the arrows as "depends on":
     - lexer depends on nothing; it could feed any grammar.
     - parser validates streams; it has never heard of leaves.
     - tree stores pairs; it has never heard of tokens.

   Consequence, the sentence worth keeping: GROW THE LANGUAGE, AND THE
   DATABASE STAYS PUT. Add a new command and the tree does not change;
   rebalance the tree and the grammar does not change.
\end{verbatim}

That decoupling is the same front-end/back-end seam every real database
uses---the place where a ``planner'' would grow---and it is the payoff
of doing the parsing work before any storage work. The bytes become
tokens, tokens become a \code{Command}, a \code{Command} becomes tree
mutation, tree mutation becomes a reply. At each boundary, the
representation is \emph{more structured and less ambiguous}.

\begin{researchnote}[The theory of the state machine]
The lexer you just read is a \emph{finite automaton}, and the theory
behind recognizing text with one was laid down in a classic 1968 paper
by Ken Thompson: ``Regular Expression Search Algorithm.'' Thompson
showed how to compile a regular expression into a small machine that
walks the text one character at a time---exactly the discipline of
MiniDB's \code{lex()}.

\begin{verbatim}
   Thompson's insight, condensed:

       "A pattern can be compiled into a finite automaton that scans
        the input character by character, keeping only the state it
        is in; the state is all the history the machine needs."

   That is precisely MiniDB's lexer: four states, one current
   character, no lookahead, no backtracking. "The state is all the
   history" is why the lexer is a loop and not a pile of ifs -- and
   why it can reject bad input without ever having seen the future.
\end{verbatim}

Thompson's algorithm became the ancestor of every regular expression
engine and every lexer generator (lex, flex). Reading it is the
fastest way to understand why the state machine is the right shape
for character-level scanning---it is the shape with a sixty-year
pedigree.

\paperref{K.~Thompson, ``Regular Expression Search Algorithm,''
Communications of the ACM 11(6), 1968}{https://dl.acm.org/doi/10.1145/363347.363387}
\end{researchnote}

\section{Structures in Bytes}

The parser's structures are tiny, and their byte sizes are part of the
design---they are what make the fixed arrays safe.

\begin{verbatim}
   struct Token (lexer output)     12 bytes   offset size field
                                              0      4    type
                                              4      4    keyword
                                              8      4    int_value

   struct Command (parser output)  16 bytes   offset size field
                                              0      4    type
                                              4      4    key
                                              8      4    value
                                              12     4    key_end

   the REPL's buffers, in bytes:
     Token  tokens[16]     ->  16 * 12 = 192 bytes
     char   line[256]      ->  256 bytes
     Command cmd           ->  16 bytes

   LexResult and ParseResult enums: 4 bytes each (an int).
\end{verbatim}

Two consequences fall out of the sizes. First, a 256-byte line can
never overflow the 16-token array---each token needs at least one
byte, and 16 tokens is the enforced ceiling. Second, because every
field is a plain 4-byte int, a \code{Command} can be copied, compared,
and passed around byte-for-byte---exactly what the REPL does when it
routes it to the storage engine.

\section{Summary}

\begin{verbatim}
   The parsing toolbox, complete:

   token          the smallest meaningful unit
   lexer          characters -> tokens (a state machine)
   TOKEN_END      the sentinel that makes streams self-terminating
   parser         tokens -> Command (shape and arity validation)
   exact arity    consume N integers, then REQUIRE the sentinel
   proto-AST      a Command: one flat node for a flat grammar
   REPL           read, gate, execute, print -- one loop
   every stage gates the next   bad lines never reach storage

   Rules that keep parsers correct:
     1. draw the state machine before writing the loop;
     2. never read past a bound: sentinel, capacity check, buffer
        size all agree;
     3. overflow is a state transition to ERROR, not a silent wrap;
     4. on any lex/parse error, treat the stream as garbage;
     5. keep layers one-directional: the tree never learns the
        language.
\end{verbatim}

\begin{osconcept}
Parsing is where the untrusted bytes of the outside world first meet
your invariants, and the state machine is the pattern that makes the
meeting safe: every input, good or bad, ends in exactly one state. The
front-end/back-end seam---language on one side, storage on the
other---is the same seam every database draws, and it is the reason a
system can grow without its layers learning about each other.
\end{osconcept}

With a parser in hand we have the final major tool of the book. The
next chapter takes a step back from writing code to map the territory
every program we have built lives in: the Linux filesystem hierarchy---
what every directory is for, how the tree grew, and the operating
system concepts behind it.

}{}
\IfFileExists{chapter15.tex}{% ============================================================================
%  chapter15.tex
%  The Linux Filesystem Hierarchy
% ============================================================================

\chapter{The Linux Filesystem Hierarchy}

\section{The One Tree}

Every file on a Linux system lives in one tree, hanging from the single
root \code{/}. There is no C:\, no drive letters, no per-device roots.
A USB stick, a network share, and the hard disk all appear as
\emph{parts of the same tree}, attached at directories called
\emph{mount points}. This chapter is a map of that tree: what lives in
each directory, why it lives there, and---more interestingly---how the
layout came to be and what operating-system ideas it encodes.

\begin{verbatim}
   The whole machine, as seen from one path:

        /
        |-- bin/      essential user programs
        |-- sbin/     essential system administration programs
        |-- lib/      shared libraries for /bin and /sbin
        |-- etc/      host-specific configuration
        |-- home/     user home directories
        |-- root/     the superuser's home
        |-- boot/     the kernel and boot images
        |-- dev/      the hardware, as files
        |-- tmp/      world-writable scratch
        |-- var/      variable data: logs, spools, caches
        |-- run/      runtime state since boot
        |-- usr/      the bulk of the system
        |-- opt/      third-party software
        |-- srv/      service data
        |-- mnt/      temporary mount points
        |-- media/    removable media
        |-- proc/     kernel and process info, generated on the fly
        `-- sys/      kernel devices and parameters, on the fly
\end{verbatim}

We have met every idea in this tree already. Chapter~8 gave us the
\emph{namespace}: names pointing to inodes pointing to blocks. This
chapter maps that namespace, adds the history that shaped it, and
explains the OS machinery---mount points, virtual filesystems, device
nodes, tmpfs---that makes one tree possible over many kinds of storage.

\begin{osconcept}
A single tree is an operating system decision with enormous payoff. By
folding every storage device into one namespace, the OS lets programs
refer to anything---a disk, a terminal, a kernel statistic, a network
stream---with one syntax and one set of tools. The cost is that the
tree's \emph{meaning} is not self-evident; a standard layout is needed
so every program can find what it needs. That standard, and the OS
mechanisms behind it, are this chapter's subject.
\end{osconcept}

\section{The Story of the Tree}

The layout was not designed on a whiteboard. It \emph{grew}, one
survival decision at a time, and its quirks are fossils of that
history.

\begin{verbatim}
   1971 -- Version 1 Unix, Bell Labs. The entire operating system runs
   in roughly 32 KB of memory on a single disk. The tree is already
   there, and it is tiny:

        /bin     the handful of programs (cat, ls, sh, cc...)
        /etc     "et cetera" -- configuration and random admin files
        /dev     the devices (a magnetic tape, a disk, a console)
        /usr     the user files -- this is where HOME DIRECTORIES were

   That is the whole story's seed: the first Unix had about five
   directories, and /usr really did mean "user."

   The RK05 disk pack.  The story of /usr's split:
   The root filesystem lived on a small fixed disk. The /usr filesystem
   lived on a larger removable RK05 pack. To boot and repair a broken
   system you could only use the small disk, so the binaries needed to
   recover sat in /bin; everything else -- and the users' files -- went
   to /usr. When the repair disk could not mount /usr, you still had a
   working, minimal system. The split is a backup-and-repair strategy
   fossilized into the layout.
\end{verbatim}

From that seed, two evolutionary pressures shaped the rest. First,
\emph{hardware}: as disks grew, the idea of ``the tiny essential set
vs the rest'' persisted as \code{/bin} vs \code{/usr/bin}, and the
early \code{/usr} directory---which had held home directories---was
colonized by system software, while user files moved to
\code{/home}. Second, \emph{diversity}: BSD and System V rearranged
and disagreed (BSD favored \code{/usr/local} for add-ons, System V
introduced \code{/opt}), and by the early 1990s every Linux
distribution had its own slightly different tree.

\begin{verbatim}
   The tree diverged, and Linux inherited the mess:
     1980s: BSD and System V each reorganize; /home becomes the user
            area; /usr/local, /opt, /var appear in different forms.
     1993:  the Filesystem Hierarchy Standard (FHS) is started to
            make the layout CONSISTENT across Linux distributions --
            so that /etc/passwd means the same thing everywhere, and
            a package built for one distribution finds its libraries
            on another. The FHS (now version 3.0, 2015) is the tree's
            constitution.

   Two modern chapters in the story:
     The "usr merge" (2010s).  Distributions began making /bin,
       /sbin, /lib symlinks to /usr/bin, /usr/sbin, /usr/lib -- one
       copy of the system software, not two. The old split existed for
       repair-on-a-tiny-disk reasons that a 2015 disk no longer has.
     The arrival of /run (2011).  Runtime state (pids, sockets) moved
       out of /var/run into /run, a tmpfs mounted before /var is --
       because early boot needs the state before the big disk is up.
\end{verbatim}

\begin{osconcept}
The tree's history is a history of operating-system constraints. The
root-vs-usr split encodes repair strategy; the FHS encodes
\emph{portability and coordination} (a standard layout is an agreement
that package managers, bootloaders, and administrators all rely on);
the usr-merge and \code{/run} encode the modern boot sequence. Reading
the tree is reading forty years of OS decisions, each one a solution
to a problem of its era.
\end{osconcept}

\section{The OS Concept: Mount Points and the Root Filesystem}

How does one tree hold many disks? The answer is \emph{mounting}. A
\emph{mount point} is a directory in the tree where a filesystem is
\emph{attached}. The kernel hangs the root of one filesystem onto a
directory of another, and from then on the two are one tree. The
\emph{root filesystem} is the one mounted at \code{/}; everything else
is grafted onto it.

\begin{verbatim}
   Two filesystems, one tree:

     before mounting:                      after mounting:
     (root fs on disk A)                   (disk B visible at /home)
       /                                    /
       |-- bin/                             |-- bin/
       |-- etc/                             |-- etc/
       `-- home/  (empty dir,             |-- home/
                 a mount point)                 |-- alice/
             (disk B: /home)                    `-- bob/
       /home                                  (the files of disk B)
       |-- alice/
       `-- bob/

   Mounting is an OS operation (`mount`):
     1. the kernel reads disk B's superblock,
     2. it attaches disk B's root to the directory /home,
     3. from then on, "alice" is found by walking INTO /home.

   Key facts:
     - A mount point is just a directory. Its own contents are
       HIDDEN while a filesystem is mounted over it.
     - Different filesystems can sit at different mounts: ext4 at /,
       tmpfs at /tmp, proc at /proc, a USB drive at /media/usb.
     - Only the kernel can perform the graft, which is why mounting
       is a privileged operation and why /etc/fstab (the table of
       "what mounts where at boot") is a system file.
\end{verbatim}

\begin{lstlisting}[caption={Reading the mount table from inside a program.}]
#include <stdio.h>

int main(void) {
    FILE *f;
    char line[256];

    /* /proc/mounts is the kernel's own list of what is mounted where. */
    f = fopen("/proc/mounts", "r");
    if (f == NULL) {
        perror("fopen");
        return 1;
    }

    while (fgets(line, sizeof(line), f) != NULL) {
        printf("%s", line);     /* device  mount_point  type  options */
    }

    fclose(f);
    return 0;
}
\end{lstlisting}

The output (a few lines of it) looks like:

\begin{verbatim}
   /dev/sda1 / ext4 rw,relatime ...
   proc /proc proc rw,relatime ...
   sysfs /sys sysfs rw,relatime ...
   devtmpfs /dev devtmpfs rw ...
   tmpfs /tmp tmpfs rw ...
   tmpfs /run tmpfs rw ...

   Read it as: "the device (or pseudo-device) on the left is attached
   at the directory in the middle, with the filesystem type on the
   right." / is ext4 on the real disk; /proc, /sys, /dev, /tmp, /run
   are their own (often RAM-backed or virtual) filesystems.
\end{verbatim}

\begin{osconcept}
Mounting is the kernel's namespace surgery: it rewrites the tree's
resolution so that a directory names a different filesystem's root.
Only the kernel can do this, because only the kernel owns the mapping
from paths to inodes and the translation from one filesystem's
structures to another's. The mount table \emph{is} the tree's
scaffolding; the directories you see are the leaves it supports.
\end{osconcept}

\section{The Directories, One by One}

With the machinery in place, here is the tree as a reference. For each
directory: what it holds, what it does, and the problem it solves.

\subsection{The Core}

\begin{verbatim}
   /bin   essential user programs (ls, cp, mv, cat, sh...)
          -- what you need to boot, repair, and use a bare system.
          On modern distributions /bin is a symlink to /usr/bin
          (the "usr merge"): one copy, two names.
          PROBLEM SOLVED: "essential" vs "everything else" was the
          1971 repair strategy; the merge removes the duplicate.

   /sbin  essential ADMINISTRATION programs (mount, fdisk, shutdown...)
          -- the tools that repair and maintain the system itself.
          Also merged into /usr/sbin on modern systems.
          PROBLEM SOLVED: keeping the tools that can break the
          system out of the hands of ordinary PATH lookups.

   /lib   shared libraries for /bin and /sbin (libc, the dynamic
          loader) -- without these, the essential binaries cannot run.
          Merged into /usr/lib.
          PROBLEM SOLVED: one copy of each library, shared by every
          program that uses it (see Chapter 13's memory chapter).

   /etc   host-specific CONFIGURATION. /etc/passwd (users),
          /etc/fstab (what mounts where), /etc/hosts (names),
          /etc/ssh (server keys)... The name is "et cetera": the
          directory where the files that fit nowhere else went.
          PROBLEM SOLVED: separating the machine's own settings from
          the software (which lives in /usr, shared or reinstallable).
          /etc is what an administrator backs up to carry a machine's
          identity to a new one.
\end{verbatim}

\subsection{Users and Data}

\begin{verbatim}
   /home  the user home directories. /home/alice, /home/bob.
          Each user's files, dotfiles, and configuration live here.
          PROBLEM SOLVED: one place for personal data, protected by
          per-user permissions (Chapter 8).

   /root  the superuser's home. (It is not /home/root, so that root's
          home is never hidden by a /home mount or an unmount.)

   /var   VARIABLE data -- content that changes while the machine runs:
          /var/log  (logs, from Chapter 9's daemons),
          /var/spool (queues: mail, print jobs),
          /var/cache (caches that can be rebuilt),
          /var/lib   (stateful data: databases, package state).
          PROBLEM SOLVED: keeping writable, growing data off the
          root filesystem so / (or /usr) can stay read-only or be
          on a smaller, more protected device.

   /tmp   world-writable scratch space. Any process may create files
          here; the sticky bit (Chapter 8) lets the owner delete them.
          Cleared on reboot; often its own tmpfs (RAM-backed).
          PROBLEM SOLVED: a place for throwaway files without a
          per-user allocation.

   /run   runtime state since the last boot: pids of daemons, sockets,
          lock files. A tmpfs -- RAM-backed, cleared at reboot.
          PROBLEM SOLVED: early boot needs state before /var is
          mounted; /run exists from the first moment.

   /srv   service data: the files web and FTP servers actually serve
          (/srv/www). Optional, but when used it is the answer to
          "where does the website live?"
\end{verbatim}

\subsection{The System Proper}

\begin{verbatim}
   /usr   "user system resources" -- the BULK of the operating
          system, and the biggest directory on the disk:
            /usr/bin       most programs
            /usr/sbin      most admin tools
            /usr/lib       most libraries
            /usr/include   C headers
            /usr/share     architecture-independent data (docs, icons,
                           locales, man pages)
            /usr/local     software installed by the ADMIN, by hand
                           (a BSD tradition: the distribution owns
                           /usr; the admin owns /usr/local)
          PROBLEM SOLVED: /usr used to mean "user files"; when home
          directories moved to /home, the name stayed and the system
          software moved in. The merge made /bin a doorway to it.

   /opt   OPTIONAL, third-party software: vendor packages installed as
          their own trees (/opt/google/, /opt/lampp/).
          System V tradition, codified by the FHS so commercial
          software has a place that never collides with the distro's
          files.
          PROBLEM SOLVED: "my package must not fight the distribution."

   /boot  what the machine needs BEFORE the OS runs: the kernel
          (vmlinuz), the initramfs, and the bootloader's files.
          Often a separate, small, early-read partition.
          PROBLEM SOLVED: the bootloader and BIOS read this before
          any filesystem driver exists for the main disk.

   /dev   the DEVICES, as files. /dev/null, /dev/zero, /dev/tty,
          /dev/sda (your disk), /dev/nvme0n1... Writing to a device
          file talks to hardware (Chapter 5's "everything is a file").
          PROBLEM SOLVED: hardware access through the ordinary file
          interface. udev and devtmpfs create the entries at boot as
          devices are discovered.

   /mnt   a convention: temporary mount points (a filesystem you are
          about to inspect, a repair disk). Admins mount things here
          BY HAND.
   /media removable media the desktop mounts automatically (USB
          sticks, CD-ROMs) -- the automatic counterpart of /mnt.
\end{verbatim}

\subsection{Virtual Filesystems: Two Meanings of ``Virtual''}

The word \emph{virtual filesystem} is used for two related but distinct
things in Linux, and keeping them apart prevents a lot of confusion.

\begin{verbatim}
   MEANING 1 -- the VFS layer (the "Virtual File System" switch).
   This is not a filesystem at all; it is an INTERFACE that every
   filesystem must implement. The kernel presents one uniform API --
   open(), read(), write(), stat(), readdir() -- and the VFS hands
   each call to whichever filesystem the path lives on.

        open()/read()/write()/stat()   (the syscalls of ch.5, ch.8)
                        |
                  [ VFS LAYER ]
                  the one interface
             /        |       |       \        \
         ext4      tmpfs    procfs    sysfs    devtmpfs
         (disk)    (RAM)    (kernel)  (kernel)  (kernel)
            each implements the SAME operations, differently

   THIS is the mechanism behind "everything is a file": /dev/null,
   a pipe, an ext4 file, and /proc/meminfo all answer the same
   syscalls because the VFS defines one contract and every
   filesystem fulfills it.

   MEANING 2 -- filesystems with no disk behind them
   (also called PSEUDO-filesystems). Here "virtual" means: the files
   do not exist on any storage device. Their contents are either
   generated on the fly by the kernel (proc, sysfs) or are just RAM
   (tmpfs, devtmpfs). Read /proc/meminfo and the kernel formats its
   memory tables into text at that moment; there is no file to back
   it up. They live only in RAM and vanish at reboot.

   The two meanings meet: pseudo-filesystems (meaning 2) are simply
   filesystems that implement the VFS interface (meaning 1) without
   a storage device.
\end{verbatim}

When Chapter~8 called \code{/proc} a \emph{virtual filesystem}, it was
using meaning 2: files whose bytes are produced on demand. When kernel
documentation talks about \emph{the VFS}, it means meaning 1: the
abstraction layer that makes every kind of filesystem interchangeable.
Both are operating system inventions, and both are why your C program
can treat a kernel statistic and a disk file with the same
\code{open}/\code{read}/\code{close} code.

\begin{osconcept}
The VFS is the OS's plug-in architecture for storage. By defining one
interface and making every filesystem implement it, the kernel lets a
single set of system calls serve disks, RAM, and the kernel's own
internals. Pseudo-filesystems push the idea to its limit: there is no
``storage'' at all, just kernel state rendered as paths. Both exist
because only the kernel can route a syscall to the right implementation
and manufacture ``files'' from its own data structures.
\end{osconcept}

\begin{verbatim}
   /proc  the PROCESS and kernel view, as files -- generated on the
          fly when you read them (Chapter 8). /proc/self/status,
          /proc/cpuinfo, /proc/mounts, /proc/<pid>/... Nothing here
          is stored; it is the kernel's data structures rendered as
          a filesystem.

   /sys   the DEVICE MODEL and kernel parameters, as files. /sys/class,
          /sys/devices, /sys/module... the modern, structured sibling
          of /proc. Writing certain files changes kernel behavior
          ("echo 1 > /sys/.../unload").

   /dev   (listed above) is the third kernel interface: device nodes,
          maintained by udev. All three are VIRTUAL: no disk, no
          persistence, just the file abstraction over kernel state.
\end{verbatim}

\begin{verbatim}
   The five virtual / kernel filesystems, in one table:

     mount     type       backed by        holds
     ------    ---------  ---------------  --------------------------
     /         ext4       the real disk    the tree's home
     /proc     proc       kernel RAM       processes + kernel info
     /sys      sysfs      kernel RAM       the device model
     /dev      devtmpfs   kernel RAM       device nodes
     /tmp, /run, /dev/shm tmpfs           RAM (page-backed)

   "tmpfs" = a filesystem whose storage is RAM pages; it exists while
   the machine runs and vanishes at reboot. That is why /run and /tmp
   are always empty after a clean boot.
\end{verbatim}

\begin{osconcept}
These directories are not decoration; they are the OS's own
bookkeeping made visible. \code{/proc}, \code{/sys}, and \code{/dev}
are the kernel exposing its internals through the file abstraction---
the same design choice as Chapter~8's virtual filesystem, taken to the
whole system. \code{tmpfs} shows the kernel serving RAM as a
filesystem. Reading the tree is reading the kernel.
\end{osconcept}

\section{The Problem the Tree Solves}

With every directory explained, the question ``why does this layout
exist at all?'' has a four-part answer:

\begin{verbatim}
   1. COORDINATION.  A standard layout is an agreement. A program can
      assume /etc/passwd, /usr/include/stdio.h, /var/log exist and
      mean the same thing on every distribution. Without the FHS,
      every package would need a per-distro install script.

   2. SEPARATION OF CONCERNS.  Read-only vs writable (/usr vs /var),
      essential vs optional (/bin vs /opt), static vs stateful
      (/usr/share vs /var/lib), software vs configuration (/usr vs
      /etc). Each split is a different backup, mount, and permission
      policy.

   3. REPAIRABILITY.  The root-vs-usr split's fossil: a minimal,
      bootable, recoverable core exists even when the big filesystem
      is broken. /boot preserves the same idea for the bootloader.

   4. SHARED MACHINERY.  One namespace means one set of tools: find,
      ls, tar, rsync -- and every program in this book -- work on
      disks, devices, and kernel interfaces alike, because they are
      all files.
\end{verbatim}

\begin{ruleofthumb}[Let the layout answer the question for you]
When you wonder where a file belongs, ask \emph{what kind of thing it
is}: configuration goes in \code{/etc}, variable data in
\code{/var}, software in \code{/usr} or \code{/opt}, scratch in
\code{/tmp}. The tree is a decision procedure, and following it keeps
your programs portable and your system repairable.
\end{ruleofthumb}

\section{Practical: A Filesystem Tourist}

Let us build a small tool that walks the root with the directory API
of Chapter~8 and labels each top-level entry with its type---a
machine-readable version of the tree at the top of this chapter.
Everything here is code you already know: \code{opendir},
\code{readdir}, \code{lstat}, and the \code{S\_IS*} macros.

\begin{lstlisting}[caption={fstree.c: walk / and label each top-level entry.}]
#include <stdio.h>
#include <stdlib.h>
#include <string.h>
#include <dirent.h>
#include <sys/stat.h>

/* One-letter label for the file type, as in ls -l (Chapter 8). */
static char type_char(mode_t mode) {
    if (S_ISDIR(mode))  return 'd';
    if (S_ISLNK(mode))  return 'l';
    if (S_ISREG(mode))  return '-';
    if (S_ISCHR(mode))  return 'c';
    if (S_ISBLK(mode))  return 'b';
    if (S_ISFIFO(mode)) return 'p';
    if (S_ISSOCK(mode)) return 's';
    return '?';
}

/* Print the type and name of every entry in one directory. */
static void tour(const char *path) {
    DIR *d;
    struct dirent *e;

    d = opendir(path);
    if (d == NULL) {
        perror("opendir");
        return;
    }

    while (1) {
        struct stat st;
        char full[1024];
        int plen;

        e = readdir(d);
        if (e == NULL) {
            break;
        }
        if (e->d_name[0] == '.') {
            continue;          /* skip . and .. */
        }

        plen = snprintf(full, sizeof(full), "%s/%s", path, e->d_name);
        if (plen < 0 || plen >= (int) sizeof(full)) {
            continue;
        }

        if (lstat(full, &st) < 0) {
            continue;
        }

        printf("%c %s\n", type_char(st.st_mode), e->d_name);
    }

    closedir(d);
}

int main(int argc, char *argv[]) {
    const char *root;

    root = (argc >= 2) ? argv[1] : "/";
    printf("type  entry  (under %s)\n", root);
    tour(root);
    return 0;
}
\end{lstlisting}

\begin{verbatim}
   $ ./fstree
   type  entry  (under /)
   d bin
   d boot
   d dev
   d etc
   d home
   d lib
   d media
   d mnt
   d opt
   d proc
   d root
   d run
   d sbin
   d srv
   d sys
   d tmp
   d usr
   d var

   On a system with the usr-merge, ls -l /bin reveals the history at a
   glance: l bin -> usr/bin. The 1971 split has become a symlink.
\end{verbatim}

\begin{researchnote}[The tree's constitution]
The Filesystem Hierarchy Standard is the document this chapter has
been quoting, and it is worth reading from the source. The current
edition, FHS 3.0, published by the Linux Foundation in 2015, specifies
the layout precisely enough that a distribution can be checked against
it mechanically---which is exactly what makes cross-distribution
software possible.

\begin{verbatim}
   The FHS's stated purpose, condensed from its own preface:

       "The filesystem hierarchy standard describes a conventional
        layout that allows software to predict where its files and
        data reside, so that software can be written to work on any
        conforming system."

   Every table in this chapter--/bin vs /usr/bin, /etc for
   configuration, /var for variable data--is the FHS made flesh. It is
   a rare thing: a standard that is simultaneously bureaucratic and
   load-bearing.
\end{verbatim}

Reading the FHS straight through is the fastest way to cement this
chapter: it is short, unambiguous, and each entry is the "problem
solved" text of Section 15.4 in its most careful form.

\paperref{Filesystem Hierarchy Standard, Version 3.0, Linux Foundation,
2015}{https://refspecs.linuxfoundation.org/FHS_3.0/fhs-3.0.pdf}
\end{researchnote}

\section{Structures in Bytes}

The hierarchy's byte structures are mostly on disk. A symbolic link is
the simplest: its data bytes are the target path.

\begin{verbatim}
   symlink "data_soft.txt" -> "data.txt"

   inode (mode = symlink)      data
   +-----------------+         +-------------+
   | size = 8        |-------->| "data.txt\0" |   8 bytes + NUL
   +-----------------+         +-------------+
\end{verbatim}

The ext4 superblock (1024 bytes, at offset 1024 of the partition)
holds the fields the kernel reads to mount a filesystem:

\begin{verbatim}
   ext4 superblock, key offsets:

   offset  size  field
   0       4     s_inodes_count
   4       4     s_blocks_count
   20      4     s_log_block_size    (block size = 1024 << this)
   56      2     s_magic             (0xEF53 marks a real ext4)
   88      2     s_inode_size        (256 on ext4)
\end{verbatim}

And \code{/proc/mounts} renders the kernel's mount table as text---one
line per mount, each field separated by spaces:

\begin{verbatim}
   /dev/sda1 / ext4 rw,relatime 0 0
   ^device   ^mount point ^fstype ^options ^dump ^pass
\end{verbatim}

Everything the tree is built from---inodes, superblocks, mount
records---is bytes; the hierarchy is the names that point at them.

\section{Summary}

\begin{verbatim}
   The filesystem hierarchy, in one paragraph:

   one tree from / ; every device, device-node, and kernel interface
   hangs on it at a MOUNT POINT; the layout is a 1971 fossil (bin vs
   usr, /etc) standardized by the FHS (coordination, separation,
   repairability) and modernized by the usr-merge and /run.

   The directories as a decision procedure:
     configuration  /etc      variable data  /var
     software       /usr,/opt scratch        /tmp
     users          /home     admin          /sbin,/root
     boot           /boot     services       /srv
     hardware       /dev      kernel         /proc,/sys
     mounting       /mnt,/media

   Rules that keep you oriented:
     1. ls -l /bin is worth doing once on your own machine: the
        history is visible in the symlinks;
     2. /proc/mounts tells you what the tree is built from;
     3. when in doubt about where a file belongs, follow the FHS's
        "what kind of thing is it?" question.
\end{verbatim}

\begin{osconcept}
The hierarchy is the OS's address book made public. It works only
because the kernel owns mounting, only because virtual filesystems let
the kernel render itself as paths, and only because a standard keeps
the agreement stable across distributions. This chapter sat between
the code you have written and the machine it runs on: the tree is
where every path in every earlier chapter actually lived.
\end{osconcept}

The tour is complete. The next chapter returns to writing code and
builds the machine's other great promise: durable storage. We design a
wire format, append and replay a log, and make data survive a
crash---before the final chapter carries it across the network.

}{}
\IfFileExists{chapter16.tex}{% ============================================================================
%  chapter16.tex
%  Durable Storage: An Append-Only Log Database
% ============================================================================

\chapter{Durable Storage: An Append-Only Log Database}

\section{Why Persistence Is Hard}

\textbf{What persistence is.} A piece of data is \emph{persistent} when
it survives longer than the program that created it---across a program
restart, across a reboot, across a power loss. Data in RAM is
\emph{volatile}: it lives exactly as long as the machine's power and
the process's lifetime, and both can end at any instant. Data in a
file is \emph{persistent}: it lives on storage, independent of any
program and of the power state of the moment it was written. The whole
discipline of this chapter is about moving data across that line---RAM
to storage---without ever letting the transition lose the truth.

\begin{verbatim}
   Persistent vs volatile, the single idea:

     volatile   dies when the power dies          (RAM, your heap)
     persistent outlives the program and the      (files on disk)
                power cycle

   Your data is only as durable as its least persistent home:
   a "saved" file still being held in the page cache is, until
   fsync, one power loss away from vanishing (Section 16.3).
\end{verbatim}

Everything so far has stored data in \emph{RAM}, and RAM has a fatal
flaw for a database: it forgets. Turn the machine off, and the heap,
the hash tables, and the B-trees of earlier chapters are gone. A real
store must be \emph{durable}: the data must survive a power loss and
a crash, in a form that can be read back tomorrow.

\begin{verbatim}
   RAM vs storage, the two worlds:

     RAM          fast, volatile: fine while the machine runs
     file         slower, persistent: survives power loss

   The database's job is to shuttle between the two worlds:
   keep a fast copy in memory for answering queries NOW, and keep a
   safe copy in a file for answering queries AFTER the next boot.

   The hard part is the crash. A program can check every return
   value and still vanish at any instant. The database must be
   designed so that "vanished at any instant" leaves the file in a
   state from which the truth can be recovered.
\end{verbatim}

This chapter builds a real, durable key-value store: an
\emph{append-only log}. It is the smallest honest database---and the
same idea, under the names write-ahead log (WAL) and journal, sits
inside every production database you have ever used. When you finish
you will understand how data survives crashes, why you cannot write a
struct to a file and be done, and what ``durable'' actually means.

\begin{osconcept}
Durability is an operating system promise, and it is not free. When
your program calls \code{write}, the bytes do not reach the disk; they
reach the kernel's \emph{page cache}. Only \code{fsync} forces them
out. The kernel exists between your program and the hardware, and it
buffers for performance---so ``the file contains my data'' is a claim
the OS, not your program, ultimately guarantees. A durable database is
a careful conversation with the page cache.
\end{osconcept}

\section{The Append-Only Log}

The simplest durable structure is also the best one. Instead of
modifying a file in place, a database appends every change as a
\emph{record} to the end of a log. The log is the source of truth; the
in-memory structures are merely a cache of it. To rebuild after a
crash, you \emph{replay} the log: read every record and re-apply it.

\begin{verbatim}
   The append-only log, drawn:

   time --->  file "mydb.log"
              +--------+--------+--------+--------+
              | SET a=1 | SET b=2 | DEL a  | SET b=5|
              +--------+--------+--------+--------+
              records appended strictly in order, never overwritten

   Startup replay: read every record, apply it in order.
     SET a=1  ->  table[a] = 1
     SET b=2  ->  table[b] = 2
     DEL a    ->  remove a
     SET b=5  ->  table[b] = 5

   Result: { b = 5 }. The log reconstructs the exact final state.

   Why appending is the right shape:
     - appends are SEQUENTIAL writes, the fastest kind the disk does;
     - an append never destroys existing data, so a crash can only
       leave a PARTIAL RECORD at the tail -- never a corrupted middle;
     - replay is deterministic: same log, same state, every time.
\end{verbatim}

The property to notice is what a crash can and cannot do. Because we
only ever write at the tail, a crash mid-write can leave a torn record
at the end---but it cannot corrupt records that were already
completed. The database's recovery rule is therefore brutally simple:
\emph{replay from the start, and stop at the first damaged record.}
Everything before the damage is true.

\begin{verbatim}
   A crash at time X leaves:

        +--------+--------+--------+------XXXXX
        | SET a=1 | SET b=2 | SET a=1|  torn tail
        +--------+--------+--------+------XXXXX
                            ^ replay stops here: "a" and "b" are safe,
                              the torn record is discarded

   That is the entire crash story: replay the clean prefix, ignore the
   torn tail. No recovery magic, no repair pass -- the design makes
   recovery a read.
\end{verbatim}

\begin{osconcept}
Append-only is crash-friendly because the kernel's behavior matches
the structure: writes to a file are ordered, and once a write is
acknowledged and flushed, those bytes are durable. By never modifying
old bytes, the database ensures the past is always intact. Deciding
``the log is truth, memory is cache'' is a systems-design decision that
converts crash recovery from a ceremony into a read loop.
\end{osconcept}

\section{The Page Cache and fsync}

Before the record format, the durability mechanic deserves its own
section, because it decides what ``written'' means. When you call
\code{write(fd, buf, n)}, the kernel copies your bytes into the
\emph{page cache}---RAM that holds the file's pages---and returns. The
disk has not necessarily seen them. The kernel flushes cached pages
later, batched for efficiency. \code{fsync(fd)} is the syscall that
says \emph{now}: do not return until the file's data has actually
reached the storage device.

\begin{verbatim}
   Where a write() goes, layer by layer:

     your buffer  --write()-->  kernel page cache  --fsync()-->  disk
        ^                                ^
        "acknowledged"               "really there"
        (fast, not durable)          (slow, durable)

   The three-layer model:
     write()   copies bytes into the kernel's cache. Fast. Volatile:
               a power loss may drop it.
     fsync()   pushes the file's dirty pages to disk. Slow. Durable:
               after it returns, the data survives a power loss.

   Consequence for a database: a record is only durable AFTER fsync.
   A database that writes and never fsyncs is a database that lies
   to itself -- it has told the user "stored" while the truth still
   sits in volatile RAM.
\end{verbatim}

\begin{ruleofthumb}[What durable means: fsync before you claim success]
If you tell the user ``stored,'' the record must have been \code{fsync}'d.
For a group of changes that must survive together, append them all,
then \code{fsync} once. The ordering---write the record, then force it
out---is the entire discipline of crash-consistent storage.
\end{ruleofthumb}

\begin{osconcept}
The page cache is the kernel's buffer between your program and the
hardware, and it is exactly the Chapter~5 ``buffer'' idea at system
scale. The OS batches writes because disk I/O is millions of times
slower than RAM. \code{fsync} is the escape hatch that buys
durability when you need it. Deciding \emph{when} to pay that price is
the difference between a toy that loses data and a database that does
not.
\end{osconcept}

\section{The Record Format: Bytes, Not Structs}

A record is a sequence of bytes on disk. Here is the first subtlety:
you cannot simply \code{fwrite} a C struct. A struct contains
\emph{padding} (Chapter~3) whose size differs between machines and
compilers, its integers are stored in the machine's native
\emph{endianness} (Chapter~1), and reading it back on another machine
or with another compiler would misread every field. A durable format
is a \emph{wire format}: a precise, byte-by-byte description written
and read by hand.

\begin{verbatim}
   Our record, byte by byte:

   +---------+------+--------+--------+-----+-------+----------+
   | magic   | type | keylen | vallen | key | value | checksum |
   | 4 bytes | 1 B  | 1 B    | 4 B    | ... | ...   | 4 bytes  |
   +---------+------+--------+--------+-----+-------+----------+

   magic     "!DB1" = 0x31424421, detects a wrong file or a torn
            read at the record boundary
   type      1 = SET (store key = value), 2 = DEL (remove key)
   keylen    1 byte: key length 0..255
   vallen    4 bytes: value length
   key       keylen bytes
   value     vallen bytes
   checksum  a hash over everything before it, to detect corruption

   The total record is 10 + keylen + vallen + 4 bytes. Every field's
   position is fixed by the format, never by the struct layout.
\end{verbatim}

Two helpers define the wire format's integers. We choose
\emph{little-endian}: byte 0 is the least significant. On an x86 or
ARM machine that matches the hardware, and writing the bytes out
explicitly makes the format portable by construction.

\begin{lstlisting}[caption={Serialization helpers: ints as bytes, on purpose.}]
#include <stdio.h>
#include <stdlib.h>
#include <string.h>
#include <fcntl.h>
#include <unistd.h>
#include <errno.h>

#define MAGIC 0x31424421u    /* "!DB1" */

#define REC_SET 1
#define REC_DEL 2

/* Write a 32-bit unsigned integer as 4 little-endian bytes. */
static void put_u32(unsigned char *b, unsigned int v) {
    b[0] = (unsigned char) (v & 0xff);
    b[1] = (unsigned char) ((v >> 8) & 0xff);
    b[2] = (unsigned char) ((v >> 16) & 0xff);
    b[3] = (unsigned char) ((v >> 24) & 0xff);
}

/* Read a 32-bit unsigned integer from 4 little-endian bytes. */
static unsigned int get_u32(const unsigned char *b) {
    unsigned int v;

    v = (unsigned int) b[0];
    v |= ((unsigned int) b[1]) << 8;
    v |= ((unsigned int) b[2]) << 16;
    v |= ((unsigned int) b[3]) << 24;
    return v;
}

/* FNV-1a: a fast, simple 32-bit hash. We seed it so a checksum can
   be computed over several buffers as one rolling value.        */
static unsigned int fnv1a_seed(unsigned int h,
                               const unsigned char *p, size_t n) {
    size_t i;

    for (i = 0; i < n; i++) {
        h = h ^ p[i];
        h = h * 16777619u;
    }
    return h;
}

static unsigned int fnv1a(const unsigned char *p, size_t n) {
    return fnv1a_seed(2166136261u, p, n);
}
\end{lstlisting}

The checksum is the record's honesty check. We hash the header, the
key, and the value; a torn write, a flipped bit, or a truncated tail
changes the hash, and replay detects it. A checksum does not repair
data---it makes \emph{corruption detectable}, which is what lets
replay stop safely at the first bad record instead of trusting it.

\begin{osconcept}
A wire format is the boundary between two representations of the same
data: in-memory structs (padding, native endianness, pointers) and
durable bytes (no padding, fixed byte order, no pointers). Only the
kernel delivers your bytes to the disk, but \emph{you} define what
they mean. Getting this boundary explicit is the same discipline we
will need for the network in Chapter~17---a network protocol is a wire
format that travels over a socket instead of a file.
\end{osconcept}

\section{The Write Path: Append and fsync}

The write path has two responsibilities: put a well-formed record at
the end of the log, and make it durable before reporting success. The
first needs a helper that writes \emph{all} requested bytes (a plain
\code{write} may write fewer), and the second needs \code{fsync}.

\begin{lstlisting}[caption={write\_all: loop until every byte is written.}]
/* Write all n bytes, retrying on EINTR. Returns 0 or -1. */
static int write_all(int fd, const void *buf, size_t n) {
    const unsigned char *p;
    size_t written;

    p = buf;
    written = 0;
    while (written < n) {
        ssize_t r;

        r = write(fd, p + written, n - written);
        if (r < 0) {
            if (errno == EINTR) {
                continue;        /* a signal interrupted us; retry */
            }
            return -1;
        }
        written = written + (size_t) r;
    }
    return 0;
}
\end{lstlisting}

The append function builds the record into its on-disk form (header,
key, value, checksum), writes the pieces in order, and fsyncs. The
order of the bytes on disk is fixed by the format; the order of the
fsync is fixed by the discipline: \emph{write first, claim success
only after fsync}.

\begin{lstlisting}[caption={append\_record: one durable record at the tail.}]
typedef struct dblog dblog;   /* defined in the next section */

/* Append one record. Returns 0 on success, -1 on failure. */
static int append_record(dblog *db, int type,
                         const char *key, int keylen,
                         const char *value, int vallen) {
    unsigned char header[10];
    unsigned char check[4];
    unsigned int sum;

    put_u32(header, MAGIC);
    header[4] = (unsigned char) type;
    header[5] = (unsigned char) keylen;
    put_u32(header + 6, (unsigned int) vallen);

    /* Checksum over header + key + value, rolled into one hash. */
    sum = fnv1a(header, 10);
    sum = fnv1a_seed(sum, (const unsigned char *) key, (size_t) keylen);
    sum = fnv1a_seed(sum, (const unsigned char *) value, (size_t) vallen);
    put_u32(check, sum);

    /* The four pieces of the record, in the wire order. */
    if (write_all(db->fd, header, 10) < 0) {
        return -1;
    }
    if (write_all(db->fd, key, (size_t) keylen) < 0) {
        return -1;
    }
    if (write_all(db->fd, value, (size_t) vallen) < 0) {
        return -1;
    }
    if (write_all(db->fd, check, 4) < 0) {
        return -1;
    }

    /* Durable now: do not tell the caller "stored" until it is. */
    if (fsync(db->fd) < 0) {
        return -1;
    }

    db->offset = db->offset + 10 + keylen + vallen + 4;
    return 0;
}
\end{lstlisting}

\begin{verbatim}
   The order is the contract:

     1. append the record's bytes  (fast, into the page cache)
     2. fsync()                    (slow, durable -- the record is
                                    now safe against a power loss)
     3. only then update memory    (the in-memory table is a cache;
                                    the log is the truth)

   If the process crashes between 1 and 2, replay may find a torn
   tail and stop -- but it will never invent a record that was not
   fully written. The fsync is the line between "maybe" and "stored."
\end{verbatim}

\section{The Read Path: Replay}

Reading a log back is the mirror of appending, with one difference:
appending assumes the file is well-formed, while replay must be
\emph{suspicious}. A record may be torn, truncated, or corrupted, and
replay must detect it and stop.

The reader uses \code{pread}, a system call we have not needed before:
\code{pread(fd, buf, n, offset)} reads at a given offset \emph{without
moving the file position}. That is exactly right for replay---each
record is read at a position we track ourselves, and the shared file
position stays untouched.

\begin{lstlisting}[caption={read\_record: read one record, or say why you cannot.}]
/* Read one record at *pos. Returns:
     1  -> a valid record (key/value returned, *pos advanced)
     0  -> clean end of file
    -1  -> torn or corrupted (replay must stop here)             */
static int read_record(int fd, long *pos,
                       int *type, char **key, int *keylen,
                       char **value, int *vallen) {
    unsigned char header[10];
    unsigned char check[4];
    unsigned char *key_buf;
    unsigned char *value_buf;
    unsigned int want;
    unsigned int got;
    int n;

    n = (int) pread(fd, header, 10, *pos);
    if (n == 0) {
        return 0;                       /* clean end of file */
    }
    if (n != 10) {
        return -1;                      /* torn header */
    }
    if (get_u32(header) != MAGIC) {
        return -1;                      /* wrong file or bad boundary */
    }

    *type = header[4];
    *keylen = (int) header[5];
    *vallen = (int) get_u32(header + 6);

    key_buf = malloc((size_t) *keylen + 1);
    value_buf = malloc((size_t) *vallen);
    if (key_buf == NULL || value_buf == NULL) {
        free(key_buf);
        free(value_buf);
        return -1;
    }

    n = (int) pread(fd, key_buf, (size_t) *keylen, *pos + 10);
    if (n != *keylen) {
        free(key_buf);
        free(value_buf);
        return -1;
    }
    n = (int) pread(fd, value_buf, (size_t) *vallen,
                    *pos + 10 + *keylen);
    if (n != *vallen) {
        free(key_buf);
        free(value_buf);
        return -1;
    }

    n = (int) pread(fd, check, 4, *pos + 10 + *keylen + *vallen);
    if (n != 4) {
        free(key_buf);
        free(value_buf);
        return -1;
    }

    want = fnv1a(header, 10);
    want = fnv1a_seed(want, key_buf, (size_t) *keylen);
    want = fnv1a_seed(want, value_buf, (size_t) *vallen);
    got = get_u32(check);
    if (want != got) {
        free(key_buf);
        free(value_buf);
        return -1;                      /* checksum mismatch */
    }

    key_buf[*keylen] = '\0';
    *key = (char *) key_buf;
    *value = (char *) value_buf;
    *pos = *pos + 10 + *keylen + *vallen + 4;
    return 1;
}
\end{lstlisting}

Replay is then a single loop: read a record, apply it, continue; stop
at the first \code{0} (clean end) or \code{-1} (damage). The partially
rebuilt table is the truth; the records after the damage, if any, are
discarded.

\begin{lstlisting}[caption={db\_replay: read the clean prefix, apply it.}]
/* Rebuild the in-memory index by replaying the log from the start. */
static void db_replay(dblog *db) {
    long pos;
    int r;

    pos = 0;
    while (1) {
        int type;
        int keylen;
        int vallen;
        char *key;
        char *value;

        r = read_record(db->fd, &pos, &type, &key, &keylen,
                        &value, &vallen);
        if (r <= 0) {
            break;           /* clean end, or damage: stop here */
        }

        if (type == REC_SET) {
            table_put(db, key, keylen, value, vallen);
        } else if (type == REC_DEL) {
            table_del(db, key, keylen);
        }

        free(key);
        free(value);
    }
}
\end{lstlisting}

\begin{verbatim}
   Replay is a loop over a file, and that is the whole point:
   the recovery procedure is the same code as startup. There is no
   "repair mode," no interactive recovery, no special case -- the
   log was designed so that reading it IS recovering it.
\end{verbatim}

\section{The In-Memory Index and the API}

The in-memory half is the hash table of Chapter~4, with one important
role: it is a \emph{cache}, not the authority. Writes go to the log
first and then to the table; reads come from the table. A lookup never
touches the disk, which is why the database is fast; the log never
lies, which is why it is durable.

\begin{lstlisting}[caption={The database: a log plus a hash-table cache.}]
typedef struct kv {
    char *key;
    int keylen;
    char *value;
    int vallen;
    struct kv *next;
} kv;

typedef struct dblog {
    int fd;
    kv **buckets;      /* the hash table (Chapter 4) */
    int nbuckets;
    long count;        /* number of live entries      */
    long offset;       /* current end of the log      */
} dblog;

static int hash_key(const char *key, int keylen, int nbuckets) {
    return (int) (fnv1a((const unsigned char *) key,
                        (size_t) keylen) % (unsigned int) nbuckets);
}

/* Look up a key in the table; returns the node or NULL. */
static kv *table_get(const dblog *db, const char *key, int keylen) {
    kv *cur;

    cur = db->buckets[hash_key(key, keylen, db->nbuckets)];
    while (cur != NULL) {
        if (cur->keylen == keylen &&
            memcmp(cur->key, key, (size_t) keylen) == 0) {
            return cur;
        }
        cur = cur->next;
    }
    return NULL;
}

/* Insert or update. Returns 1 if a new key was added. */
static int table_put(dblog *db, const char *key, int keylen,
                     const char *value, int vallen) {
    kv *node;
    kv *cur;
    int b;

    cur = table_get(db, key, keylen);
    if (cur != NULL) {
        /* Update the existing node in place. */
        char *new_value;

        new_value = malloc((size_t) vallen);
        if (new_value == NULL) {
            return 0;
        }
        memcpy(new_value, value, (size_t) vallen);
        free(cur->value);
        cur->value = new_value;
        cur->vallen = vallen;
        return 0;
    }

    node = malloc(sizeof(kv));
    if (node == NULL) {
        return 0;
    }
    node->key = malloc((size_t) keylen);
    node->value = malloc((size_t) vallen);
    if (node->key == NULL || node->value == NULL) {
        free(node->key);
        free(node->value);
        free(node);
        return 0;
    }
    memcpy(node->key, key, (size_t) keylen);
    memcpy(node->value, value, (size_t) vallen);
    node->keylen = keylen;
    node->vallen = vallen;

    b = hash_key(key, keylen, db->nbuckets);
    node->next = db->buckets[b];
    db->buckets[b] = node;
    db->count = db->count + 1;
    return 1;
}

/* Remove a key if present. Returns 1 if it was removed. */
static int table_del(dblog *db, const char *key, int keylen) {
    kv **cur;
    kv *victim;

    cur = &db->buckets[hash_key(key, keylen, db->nbuckets)];
    while (*cur != NULL) {
        if ((*cur)->keylen == keylen &&
            memcmp((*cur)->key, key, (size_t) keylen) == 0) {
            victim = *cur;
            *cur = victim->next;
            free(victim->key);
            free(victim->value);
            free(victim);
            db->count = db->count - 1;
            return 1;
        }
        cur = &(*cur)->next;
    }
    return 0;
}
\end{lstlisting}

The public API is three functions, each a short story about the log
and the cache:

\begin{lstlisting}[caption={The API: log first, then cache.}]
/* Store key = value durably. */
int db_set(dblog *db, const char *key, int keylen,
           const char *value, int vallen) {
    if (append_record(db, REC_SET, key, keylen, value, vallen) < 0) {
        return -1;          /* the log did not take the record */
    }
    table_put(db, key, keylen, value, vallen);
    return 0;
}

/* Remove key durably. Returns 1 if it existed, 0 if not. */
int db_del(dblog *db, const char *key, int keylen) {
    int existed;

    existed = (table_get(db, key, keylen) != NULL);
    if (existed) {
        if (append_record(db, REC_DEL, key, keylen, NULL, 0) < 0) {
            return -1;
        }
        table_del(db, key, keylen);
    }
    return existed;
}

/* Fetch a key. On success returns 1 and a malloc'd copy of the
   value (the caller frees it); returns 0 if absent.             */
int db_get(const dblog *db, const char *key, int keylen,
           char **out_value, int *out_vallen) {
    kv *node;
    char *copy;

    node = table_get(db, key, keylen);
    if (node == NULL) {
        return 0;
    }
    copy = malloc((size_t) node->vallen);
    if (copy == NULL) {
        return 0;
    }
    memcpy(copy, node->value, (size_t) node->vallen);
    *out_value = copy;
    *out_vallen = node->vallen;
    return 1;
}
\end{lstlisting}

Notice the ordering in \code{db\_set}: the record goes to the log
\emph{first}, and the table is updated only if the append succeeded.
If the append fails, memory must not pretend the write happened. If
the append succeeds and we crash before touching the table, replay
rebuilds it. \emph{Log first, cache second} is the rule, and it is
what makes the two halves consistent.

\begin{osconcept}
The log/cache split is a layering decision of the same shape as
Chapter~14's parser: an authority (the log) and a fast view (the
table), with one-directional trust. The kernel provides the durability
primitive (\code{fsync}); the database chooses when to pay for it; the
cache makes reads cheap. Every production database---SQLite's journal,
PostgreSQL's WAL---is this same split, tuned and enlarged.
\end{osconcept}

\section{Compaction: Shrinking the Log}

An append-only log only grows. Every update of a key appends another
record, so after a long life the log is full of dead records---old
versions of keys that were later changed or deleted. The cure is
\emph{compaction}: rewrite a fresh log containing only the \emph{live}
entries, then swap it in atomically. The compaction walks the in-memory
table (which already holds the truth) and writes one \code{SET} record
per live entry.

The swap uses \code{rename}, and that one call carries the whole
atomicity: \code{rename} either succeeds completely or changes nothing,
so at any instant a reader sees either the old log or the new one,
never a mixture.

\begin{lstlisting}[caption={db\_compact: live entries only, then an atomic swap.}]
/* Rewrite the log with only the live entries. Returns 0 or -1. */
int db_compact(dblog *db, const char *path) {
    char tmp[1024];
    int new_fd;
    int b;
    kv *cur;

    snprintf(tmp, sizeof(tmp), "%s.tmp", path);
    new_fd = open(tmp, O_WRONLY | O_CREAT | O_TRUNC, 0644);
    if (new_fd < 0) {
        return -1;
    }

    /* Walk every bucket; write one SET record per live entry. */
    for (b = 0; b < db->nbuckets; b++) {
        cur = db->buckets[b];
        while (cur != NULL) {
            unsigned char header[10];
            unsigned char check[4];
            unsigned int sum;

            put_u32(header, MAGIC);
            header[4] = REC_SET;
            header[5] = (unsigned char) cur->keylen;
            put_u32(header + 6, (unsigned int) cur->vallen);

            sum = fnv1a(header, 10);
            sum = fnv1a_seed(sum, (const unsigned char *) cur->key,
                             (size_t) cur->keylen);
            sum = fnv1a_seed(sum, (const unsigned char *) cur->value,
                             (size_t) cur->vallen);
            put_u32(check, sum);

            if (write_all(new_fd, header, 10) < 0) {
                goto fail;
            }
            if (write_all(new_fd, cur->key,
                          (size_t) cur->keylen) < 0) {
                goto fail;
            }
            if (write_all(new_fd, cur->value,
                          (size_t) cur->vallen) < 0) {
                goto fail;
            }
            if (write_all(new_fd, check, 4) < 0) {
                goto fail;
            }

            cur = cur->next;
        }
    }

    /* Durable, then atomic: the rename is all-or-nothing. */
    if (fsync(new_fd) < 0) {
        goto fail;
    }
    close(new_fd);

    if (rename(tmp, path) < 0) {
        return -1;
    }

    close(db->fd);
    db->fd = open(path, O_RDWR, 0644);
    if (db->fd < 0) {
        return -1;
    }
    db->offset = lseek(db->fd, 0, SEEK_END);
    if (db->offset < 0) {
        return -1;
    }
    return 0;

fail:
    close(new_fd);
    unlink(tmp);
    return -1;
}
\end{lstlisting}

The \code{goto fail} deserves an explanation, because it is the one
place this book uses \code{goto} on purpose. In C, error handling has
no \code{finally} clause; the idiomatic pattern is a single cleanup
label that every error path jumps to, so that cleanup code exists
\emph{once} instead of being repeated after every failing call.
Every \code{goto fail} means ``something went wrong; clean up the
temporary file and return failure.''

\begin{verbatim}
   Compaction in pictures:

     before:  log = SET a=1, SET b=2, SET a=1, SET a=3, DEL b
              (5 records; only {a=3} is live)
     compact: walk the table, write SET a=3 to "db.tmp", fsync
     swap:    rename("db.tmp", "db")  -- atomic
     after:   log = SET a=3  (one record; the dead history is gone)

   A crash at any point leaves either the old (bigger) log or the new
   (smaller) one. Both are valid; the rename is the boundary between
   them, and it is atomic.
\end{verbatim}

\begin{osconcept}
Atomic rename is a kernel guarantee that the database leans on: the
filesystem makes the name switch indivisible, so compaction never
exposes a half-written file. This is the same ``reference counting and
naming is kernel business'' lesson as Chapter~8, applied to the most
delicate moment in a database's life. The kernel's atomicity is the
only reason swap-in-place works at all.
\end{osconcept}

\section{Practical: The Whole Database}

Let us assemble every piece---serialization, append, replay, index,
compaction---into one program, plus the two glue functions we have not
shown (\code{db\_open} opens or creates the log and replays it;
\code{db\_close} frees the index). The pieces have all appeared above;
here is how they fit.

\begin{lstlisting}[caption={dblog.c: opening, closing, and a demo main.}]
int db_open(dblog *db, const char *path) {
    int fd;

    fd = open(path, O_RDWR | O_CREAT, 0644);
    if (fd < 0) {
        return -1;
    }
    db->fd = fd;
    db->nbuckets = 1024;
    db->count = 0;
    db->offset = lseek(fd, 0, SEEK_END);
    if (db->offset < 0) {
        close(fd);
        return -1;
    }
    db->buckets = calloc((size_t) db->nbuckets, sizeof(kv *));
    if (db->buckets == NULL) {
        close(fd);
        return -1;
    }

    db_replay(db);        /* rebuild the cache from the log */
    return 0;
}

void db_close(dblog *db) {
    int b;
    kv *cur;

    for (b = 0; b < db->nbuckets; b++) {
        cur = db->buckets[b];
        while (cur != NULL) {
            kv *next;

            next = cur->next;
            free(cur->key);
            free(cur->value);
            free(cur);
            cur = next;
        }
    }
    free(db->buckets);
    close(db->fd);
    db->buckets = NULL;
    db->fd = -1;
}

int main(int argc, char *argv[]) {
    dblog db;
    char *value;
    int vallen;

    if (argc != 2) {
        fprintf(stderr, "usage: %s dbfile\n", argv[0]);
        return 1;
    }

    if (db_open(&db, argv[1]) < 0) {
        perror("db_open");
        return 1;
    }

    db_set(&db, "answer", 6, "42", 2);
    db_set(&db, "name", 4, "ada", 3);
    db_set(&db, "answer", 6, "43", 2);     /* update: appends a record */
    db_del(&db, "old", 3);                 /* absent: no record written */

    if (db_get(&db, "answer", 6, &value, &vallen) == 1) {
        printf("answer = %.*s\n", vallen, value);
        free(value);
    }

    db_compact(&db, argv[1]);
    db_close(&db);
    return 0;
}
\end{lstlisting}

\begin{verbatim}
   Run it, then look at the file -- the wire format is right there:

     $ ./dblog /tmp/mydb.log
     answer = 43
     $ xxd /tmp/mydb.log
     00000000: 21 44 42 31 01 06 02 00 00 00   !DB1.........
                ^^ ^^ ^^ ^^ ^^    ^^^^^^^^
                magic  type keylen vallen (2, little-endian)
                "!DB1"  SET  = 6
     0000000a: 61 6e 73 77 65 72               "answer"
     00000010: 34 32                           "42"
     00000012: XX XX XX XX                      checksum
     ... and so on for each record.

   Re-run the program: it opens the same file, REPLAYS it, and
   prints answer = 43 again -- the data survived the previous run.
   That is durability, demonstrated by reading the file back.
\end{verbatim}

\begin{verbatim}
   To watch recovery from a torn tail, append garbage:

     $ echo "garbage" >> /tmp/mydb.log
     $ ./dblog /tmp/mydb.log        (the magic check fails; replay
     answer = 43                     stops before the garbage, and
                                     the clean prefix is intact)

   No crash simulator needed: the checksum and magic find the damage,
   replay stops, and the truth before the damage is preserved.
\end{verbatim}

\section{Where the Real Databases Go From Here}

This log is a real, durable, replayable store in about three hundred
lines. Production databases take the same skeleton and add four things,
each of which we have met in earlier chapters:

\begin{verbatim}
   what real systems add            where we met the idea
   ----------------------------     -------------------------------
   batched, asynchronous fsync      the page cache (this chapter)
   (group commit)                   and Chapter 5's buffering
   sorted on-disk storage           the B-tree family (Chapter 4)
   instead of replay-from-zero
   concurrency and locking          threads (Chapter 11)
   a network front end              the next chapter
\end{verbatim}

The lesson is that durability is not a library---it is a discipline
of \emph{fsync ordering}, \emph{checksums}, and \emph{atomic
renames}, exactly the ones you just wrote. Everything else is
performance and convenience built on top.

\begin{researchnote}[Log-structured storage, from the source]
The decision to make a log the heart of a storage system has a famous
pedigree. In 1991, Mendel Rosenblum and John Ousterhout published
``The Design and Implementation of a Log-Structured File System,'' the
paper behind the LFS filesystem---and, later, much of the modern
flash-storage world. Its central idea is precisely this chapter's:
\emph{treat the storage medium as one continuous log, write
sequentially, and use a compacting cleaner to reclaim dead space}.

\begin{verbatim}
   The paper's thesis, condensed:

       "By making all writes sequential appends to a log, the disk's
        expensive seeks largely disappear, and a cleaner process that
        rewrites only live data keeps the log from growing without
        bound."

   Translate to this chapter: append (sequential writes), checksum
   (detect damage), compact (the cleaner). Your dblog is a
   log-structured store in miniature, and the LFS paper is the
   original argument for why that shape is fast and crash-safe.
\end{verbatim}

Reading it is worthwhile for one reason above all: it shows that the
``simple'' log you just built was, in 1991, a research breakthrough,
and that its ideas now sit inside everything from SSDs to databases.

\paperref{M.~Rosenblum, J.~K.~Ousterhout, ``The Design and
Implementation of a Log-Structured File System,'' ACM SOSP 1991}{https://web.stanford.edu/~ouster/cgi-bin/papers/lfs.pdf}
\end{researchnote}

\section{Structures in Bytes}

The chapter's whole byte structure is the record format, so this is
the page to keep. One record, with absolute offsets:

\begin{verbatim}
   offset  size  field           example (SET answer = "42")
   0       4     magic "!DB1"    21 44 42 31
   4       1     type            01 = SET
   5       1     keylen          06
   6       4     vallen          02 00 00 00   (little-endian 2)
   10      keylen key            "answer"
   10+keylen vallen  value       "42"
   10+keylen+vallen  4  checksum  XX XX XX XX  (FNV-1a)

   total length = 10 + keylen + vallen + 4.
   the file is just records back to back:

   [SET a=1][SET b=2][DEL a][SET b=5]    append-only, never rewritten
\end{verbatim}

Little-endian is a deliberate choice, written byte by byte by
\code{put\_u32} and read back by \code{get\_u32}; the checksum covers
header + key + value, so a torn tail fails the hash and replay stops
at the last good byte. Compaction rewrites the same records into a
temp file and \code{rename}s it over the original---the on-disk
structure is unchanged, only shorter.

\section{Summary}

\begin{verbatim}
   The durable-storage toolbox, complete:

   the page cache      where write() actually puts your bytes
   fsync()             the line between "maybe" and "stored"
   append-only log     sequential writes; the past is never changed
   wire format         bytes, not structs: fixed order, explicit
                       endianness, no padding
   checksum            makes corruption detectable
   replay              reading the log IS recovering the database
   pread()             read at an offset without moving the position
   compaction          rewrite only live entries
   rename()            the atomic swap between old log and new

   Rules that keep durable code correct:
     1. write the record, THEN fsync, THEN claim success;
     2. append only at the tail; never overwrite old bytes;
     3. serialize with a fixed wire format -- never fwrite a struct;
     4. on any damage, replay stops: the clean prefix is the truth;
     5. compact to a temp file, fsync, and rename atomically.
\end{verbatim}

\begin{osconcept}
Durability is the operating system's oldest promise---\emph{your data
will still be there tomorrow}---and this chapter was the conversation
that keeps it: asking the kernel's page cache to flush, trusting its
atomic renames, and using its files as the wire format's transport.
The log is where the OS's storage services meet the data structures of
the first half of this book, and it is the same architecture every
real database runs.
\end{osconcept}

The database is durable on disk, and it has a command language. The
final chapter gives it a voice across the machine: we build sockets,
learn the TCP and UDP protocols, and turn this store---and the event
loop of Chapter~10---into a server that answers the world.


}{}
\IfFileExists{chapter17.tex}{% ============================================================================
%  chapter17.tex
%  Network Programming: Sockets and the Scalable Server
% ============================================================================

\chapter{Network Programming}

\section{Beyond One Machine}

Every channel in this book so far stayed inside one machine: pipes
between processes on the same kernel, files on the same disk, signals
to the same process tree. The network is the last boundary, and the
most important one. It lets two processes on different machines---on
different continents---exchange bytes, using a stack of protocols that
the operating system provides and mostly manages for you.

The astonishing fact that makes networking feel like the chapters
before it is this: \emph{a network connection looks like a file
descriptor}. You have already written the hard parts. After the setup,
reading from a socket is \code{read}, writing is \code{write}, and the
end of the conversation is the same EOF you met in Chapter~7. The new
machinery is only the setup: naming the machine, naming the program,
and agreeing on the rules of the conversation.

\begin{verbatim}
   The whole book, in one arrow:

     our process  --read/write-->  a FILE DESCRIPTOR
                                      |
                        +-------------+-------------+
                        |             |             |
                      a file       a pipe       a SOCKET
                      (ch.5)      (ch.7)    (this chapter)
                                              |
                                        another machine,
                                        thousands of km away
\end{verbatim}

\begin{osconcept}
The network is the operating system's biggest abstraction: the kernel
owns the network hardware, implements the protocol stacks, and
presents every connection as a descriptor your program reads and
writes exactly like a file. Only the OS can do this, because only the
OS can multiplex thousands of connections onto one network card and
run the intricate protocol machinery that makes bytes survive the
unreliable network between you and your peer.
\end{osconcept}

\section{The Three Names of a Network Conversation}

To reach a program on another machine, your operating system needs
three facts, and each is a different kind of name:

\begin{description}
  \item[The IP address] names the \emph{machine}. \code{192.168.1.10}
    is a 32-bit number (IPv4) written as four bytes. It is the
    network-layer name: ``that machine, out there.''
  \item[The port] names the \emph{program} on that machine. A machine
    runs many network programs at once; the port is how a connection
    chooses one of them. \code{80} is the conventional web port,
    \code{53} DNS, and your own servers use whatever you pick.
  \item[The protocol] names the \emph{rules of conversation}. The two
    that matter are TCP and UDP: TCP for reliable streams, UDP for
    fast datagrams. Both travel on top of the same IP addressing.
\end{description}

\begin{verbatim}
   A conversation needs all three:

       192.168.1.10 : 9000     over  TCP
        ^^^^^^^^^^^   ^^^^           ^^^
        machine      program      the rules

   The socket is exactly this triple. Create one, and you have asked
   the OS for "a network endpoint on this machine, using this
   protocol."
\end{verbatim}

\begin{osconcept}
Ports are an OS-level namespace, exactly like the process tree and the
filesystem hierarchy: the kernel keeps a table of ``which socket owns
which port'' and enforces that only one program is bound to a port at
a time. Without the kernel's arbitration, two programs could both
claim port 9000 and incoming packets would have nowhere unambiguous to
go. Naming and arbitration of ports is kernel work.
\end{osconcept}

\section{Sockets and Byte Order}

A \emph{socket} is an endpoint for network communication---the thing
your program holds as a descriptor. Creating one and giving it a name
and a port is the setup dance, and it involves one concept you have
not seen since Chapter~16: \emph{byte order}.

The network's wire format is defined as \emph{big-endian}: the most
significant byte first, regardless of what your CPU does. On a
little-endian machine (our x86-64), that means integers written into
socket addresses must be \emph{converted} before the kernel will read
them correctly. The helpers are \code{htons} (host to network, short)
and \code{htonl} (host to network, long). You met the problem in
Chapter~16; the network is where it is unavoidable.

\begin{lstlisting}[caption={Creating a socket address in network byte order.}]
#include <stdio.h>
#include <stdlib.h>
#include <string.h>
#include <unistd.h>
#include <arpa/inet.h>
#include <sys/socket.h>

#define PORT 9000

int main(void) {
    int fd;
    struct sockaddr_in addr;

    /* 1. Create the socket: IPv4 + TCP (SOCK_STREAM). */
    fd = socket(AF_INET, SOCK_STREAM, 0);
    if (fd < 0) {
        perror("socket");
        return 1;
    }

    /* 2. Build the address: family, port, interface. */
    memset(&addr, 0, sizeof(addr));
    addr.sin_family = AF_INET;                 /* IPv4 */
    addr.sin_port = htons(PORT);               /* port in NETWORK order */
    addr.sin_addr.s_addr = htonl(INADDR_ANY);  /* any interface */

    printf("socket created, fd = %d\n", fd);
    printf("port in network byte order: %d\n", addr.sin_port);
    close(fd);
    return 0;
}
\end{lstlisting}

\begin{verbatim}
   The socket address struct, drawn:

     struct sockaddr_in {
         short   sin_family;     AF_INET (IPv4) or AF_INET6 (IPv6)
         unsigned short sin_port;  the port, in BIG-ENDIAN byte order
         struct in_addr sin_addr;  the IP address, big-endian too
         ...                        (padding to the generic sockaddr)
     };

   - htons(9000) on a little-endian machine stores the bytes reversed,
     so that the kernel (and the wire) sees 9000 correctly.
   - htonl(INADDR_ANY) means "accept connections on ALL interfaces."
   - inet_pton(AF_INET, "127.0.0.1", &addr.sin_addr) converts a
     dotted-quad string into the 32-bit address (the client's job).
\end{verbatim}

\begin{ruleofthumb}[The wire is big-endian; your struct is not]
Whatever your CPU does, the network format is fixed: most significant
byte first. Every integer that leaves your machine---ports, addresses,
and any data you serialize---must be converted to network byte order,
or the machine on the other end will read your numbers backwards. This
is Chapter~16's wire-format discipline with a byte order the protocol,
not you, chooses.
\end{ruleofthumb}

\section{A TCP Server, Step by Step}

TCP is the \emph{connection-oriented} protocol: before any data, the
two sides perform a handshake, and after the data, they close. The
server's setup has four calls, and each has one job:

\begin{verbatim}
   The server lifecycle:

     socket()   create the endpoint              (a descriptor)
     bind()     give it a name: IP + port        ("I live here")
     listen()   start accepting; the kernel       ("queue my calls")
                queues incoming connections
     accept()   take one connection off the       (returns a NEW
                queue; returns a NEW socket        socket for THIS
                for that one client                client)

   accept() is the confusing step, and it is the heart of the model:
   the LISTENING socket is not for talking -- it is for meeting.
   Each accepted connection is a separate socket, with its own buffer
   and its own conversation.
\end{verbatim}

\begin{lstlisting}[caption={tcp\_server.c: a complete one-client echo server.}]
#include <stdio.h>
#include <stdlib.h>
#include <string.h>
#include <unistd.h>
#include <arpa/inet.h>
#include <sys/socket.h>

#define PORT 9000

int main(void) {
    int listen_fd;
    int conn_fd;
    struct sockaddr_in addr;
    char buf[1024];
    ssize_t got;

    listen_fd = socket(AF_INET, SOCK_STREAM, 0);
    if (listen_fd < 0) {
        perror("socket");
        return 1;
    }

    memset(&addr, 0, sizeof(addr));
    addr.sin_family = AF_INET;
    addr.sin_port = htons(PORT);
    addr.sin_addr.s_addr = htonl(INADDR_ANY);

    if (bind(listen_fd, (struct sockaddr *) &addr, sizeof(addr)) < 0) {
        perror("bind");
        return 1;
    }

    if (listen(listen_fd, 16) < 0) {        /* backlog: 16 queued */
        perror("listen");
        return 1;
    }

    printf("listening on port %d\n", PORT);

    conn_fd = accept(listen_fd, NULL, NULL);  /* blocks until a client */
    if (conn_fd < 0) {
        perror("accept");
        return 1;
    }

    printf("client connected\n");

    /* Echo whatever the client sends until it closes. */
    while (1) {
        got = read(conn_fd, buf, sizeof(buf));
        if (got <= 0) {
            break;                  /* EOF: the client closed */
        }
        write(conn_fd, buf, (size_t) got);
    }

    close(conn_fd);
    close(listen_fd);
    return 0;
}
\end{lstlisting}

The \emph{listen backlog} is the kernel's welcome mat: it tells the
kernel how many pending connections it may queue while you are busy
with others. If the queue is full, new connection attempts are simply
delayed or refused---the kernel, not your program, holds the door.

\begin{verbatim}
   Why the backlog exists, drawn:

        kernel's listen queue (capacity 16)
        +------+------+------+...
        | conn | conn | ...         <- the kernel holds these WHILE
        +------+------+------+...      your program handles others
        ^
     clients connecting

     your program:   accept()  takes ONE from the queue
                     read/write over it
                     accept()  takes the next...
\end{verbatim}

\begin{osconcept}
The listening socket, the backlog queue, and the handshake are all
kernel state. Your \code{listen} call only configures the kernel's
welcoming; the kernel itself accepts the TCP handshake, queues the
connection, and hands it to your \code{accept}. This is why a server
can be slow to \code{accept} without dropping connections---the OS is
holding the door. Only the kernel can do this, because only the kernel
owns the protocol stack.
\end{osconcept}

\section{A TCP Client: The Three-Way Handshake}

The client's setup is two calls: \code{socket} to create the endpoint,
and \code{connect} to reach the server. The magic happens inside
\code{connect}: it drives the TCP \emph{three-way handshake} with the
kernel's help, so that when \code{connect} returns, both sides agree a
connection exists and data can flow.

\begin{lstlisting}[caption={tcp\_client.c: send a message, print the reply.}]
#include <stdio.h>
#include <stdlib.h>
#include <string.h>
#include <unistd.h>
#include <arpa/inet.h>
#include <sys/socket.h>

#define PORT 9000

int main(int argc, char *argv[]) {
    int fd;
    struct sockaddr_in addr;
    char buf[1024];
    ssize_t got;

    if (argc != 2) {
        fprintf(stderr, "usage: %s message\n", argv[0]);
        return 1;
    }

    fd = socket(AF_INET, SOCK_STREAM, 0);
    if (fd < 0) {
        perror("socket");
        return 1;
    }

    memset(&addr, 0, sizeof(addr));
    addr.sin_family = AF_INET;
    addr.sin_port = htons(PORT);
    if (inet_pton(AF_INET, "127.0.0.1", &addr.sin_addr) <= 0) {
        perror("inet_pton");
        return 1;
    }

    /* connect performs the three-way handshake. */
    if (connect(fd, (struct sockaddr *) &addr, sizeof(addr)) < 0) {
        perror("connect");
        return 1;
    }

    write(fd, argv[1], strlen(argv[1]));

    got = read(fd, buf, sizeof(buf));      /* blocks until the reply */
    if (got > 0) {
        printf("server echoed: %.*s\n", (int) got, buf);
    }

    close(fd);
    return 0;
}
\end{lstlisting}

\begin{verbatim}
   The three-way handshake, drawn (the SYN packets):

     server  (LISTEN)                 client
        |                              socket()
        |<---------- SYN ------------   connect(): "can we talk?"
        |---------- SYN + ACK ------->  "yes, and my handshake"
        |<---------- ACK -------------  "agreed"
        |                              connect() returns: ESTABLISHED
        |<--------- data -------------  now the read/write phase
        |--------- data ------------->

   Why THREE messages? Because each side must learn that the other is
   there AND willing: the SYN asks, the SYN+ACK answers and asks back,
   and the ACK closes the loop. A two-message version could not tell a
   dead peer from a busy one.

   Closing is a two-way FIN exchange: each side sends FIN when it is
   done writing, and the other sees EOF (read returns 0). Half-closing
   is legal: one side can finish sending while still reading.
\end{verbatim}

\subsection{How a Connection Works in Bytes}

The handshake above is real packets, and packets are real bytes. Every
TCP message---handshake or data---is an IP packet whose payload is a
\emph{TCP segment}: a header followed by the user's data. The header
is a wire format, exactly like the record format of Chapter~16, and it
is what carries the connection's state. Here is the header, byte by
byte:

\begin{verbatim}
   The TCP header (minimal form), drawn like a wire format:

     bytes    field          meaning
     ------   -------------  -----------------------------------------
     0-1      src port       where THIS segment comes from
     2-3      dst port       where it is going
     4-7      sequence       the byte's position in the outgoing stream
     8-11     ack number     the next byte expected from the peer
     12       data offset    header length in 32-bit words (4 bits)
     13       flags          (URG ACK PSH RST SYN FIN, among others)
     14-15    window         flow control: how many bytes the receiver
                             is willing to accept
     16-17    checksum       integrity check over header and payload
     18-19    urgent pointer (meaningful only when URG is set)
     20+      payload        the user's data bytes (for data segments)

   One segment is therefore: 20 bytes of protocol + the data.
   The connection itself has NO bytes on the wire -- it IS the
   sequence numbers and flags shared across segments.
\end{verbatim}

The two numbers that make the byte stream work are the
\emph{sequence} and \emph{acknowledgment} fields. Every byte of data
is numbered: the sender stamps each segment with the sequence number
of its first byte, and the receiver acknowledges by sending back
``I expect byte N next.'' This one mechanism does three jobs at once:
ordering (the receiver sorts out-of-order segments by their numbers),
retransmission (an unacknowledged byte is resent), and the byte
stream your \code{read} sees.

\begin{verbatim}
   The three-way handshake in bytes:

     client                            server
        |  SYN, seq=1000                |
        |------------------------------>|   "hello, start at byte 1000"
        |  SYN+ACK, seq=7000, ack=1001  |
        |<------------------------------|   "hello, I start at 7000;
        |                               |    I expect your byte 1001"
        |  ACK, seq=1001, ack=7001      |
        |------------------------------>|   "agreed; I'm sending 1001"
        |                               |
        |  data "ada", seq=1001         |   <-- the payload bytes
        |------------------------------>|       (sequence 1001 = first
        |                               |        byte of "ada")
        |  ACK, ack=1004                |
        |<------------------------------|   "got bytes 1001..1003;
        |                               |    next I expect 1004"

   Read it as a wire format conversation: each arrow is a real
   packet, each packet carries sequence numbers that let the kernel
   on each side turn an arbitrary stream of segments into the ordered,
   gap-free byte stream your read() returns.
\end{verbatim}

The lesson, tying the whole book together: a TCP connection is a
\emph{protocol}, and a protocol is bytes with rules---the same idea as
the MiniDB wire format of Chapter~16, scaled to two machines. The
kernel implements the rules; your program sees only the resulting byte
stream.

\begin{osconcept}
The byte-level machinery---sequence numbers, acknowledgments,
retransmission, flow control---is the kernel's own wire-format
implementation. Your program never sees segments or sequence numbers;
the kernel consumes the wire protocol on your behalf and delivers a
clean byte stream. This is the extreme of the book's theme: the OS
turns raw, unreliable bytes on a wire into the abstractions your code
was built to use.
\end{osconcept}

\begin{osconcept}
The TCP handshake and connection state machine live entirely in the
kernel. Your \code{connect} is one syscall; behind it the kernel
retransmits lost SYN packets, times out dead peers, and tracks the
connection's state (SYN-SENT, ESTABLISHED, FIN-WAIT, and so on)
throughout. Reliability on an unreliable network is exactly the kind
of tedious, time-sensitive machinery no application could duplicate--
and exactly what an operating system is for.
\end{osconcept}

\section{Reading and Writing on a Socket}

After the setup, a connected socket is a file descriptor, and every
lesson of Chapters~4 and~6 applies unchanged:

\begin{verbatim}
   The socket is a byte stream, exactly like a pipe:

     - read() may return FEWER bytes than requested: loop or accept
       short reads.
     - EOF (read returns 0) means the peer closed: its FIN arrived.
     - write() may write fewer bytes than asked: loop.
     - the data has NO message boundaries. "ab" then "cd" may arrive
       as "abcd" or "abc"+"d". A protocol that needs messages must
       frame them itself (length prefixes -- Chapter 16's wire format).
     - blocking reads park the process until data arrives, exactly
       like a pipe.

   The ONE socket-specific rule: a socket is FULL-DUPLEX. Both sides
   can read and write at the same time on the same socket -- unlike a
   pipe, which flows one way.
\end{verbatim}

For a server that must hold many connections, blocking reads are the
killer---the exact problem Chapter~10 solved with \code{poll}. The same
medicine works on sockets, because a socket is a descriptor: set it
\emph{non-blocking}, and reads that would block return \code{EAGAIN}
instead. The event loop of Chapter~10 was already the right shape; it
just needed a new source of descriptors.

\section{TCP vs UDP}

Not every conversation wants TCP's guarantees. The alternative,
UDP, trades reliability for speed and simplicity: a UDP socket sends
\emph{datagrams}---individual messages with boundaries, no connection,
no handshake, and no promise that a message arrives, arrives once, or
arrives in order.

\begin{verbatim}
   The two protocols, side by side:

   concern        TCP                              UDP
   ----------     -------------------------------  ---------------------------
   connection     connected, handshake first       connectionless: just send
   delivery       reliable, in order, no dupes     best effort: may lose,
                                                  duplicate, or reorder
   boundaries     byte STREAM: no message edges    DATAGRAM: one message,
                                                  one recvfrom
   overhead       higher (handshake, acks)         minimal
   use for        HTTP, files, databases,          DNS, games, live media,
                  anything needing correctness     broadcasts, fast probes
   API            read() / write()                 recvfrom() / sendto()

   The rule of thumb: if it must be correct, use TCP; if it must be
   fast and small and can tolerate loss, use UDP. Many systems use
   both -- TCP for control, UDP for media.
\end{verbatim}

A UDP server never accepts; it binds and then loops on
\code{recvfrom}, which returns the datagram \emph{and} the address it
came from, so the server can reply with \code{sendto}.

\begin{lstlisting}[caption={udp\_echo.c: a datagram echo server.}]
#include <stdio.h>
#include <stdlib.h>
#include <string.h>
#include <unistd.h>
#include <arpa/inet.h>
#include <sys/socket.h>

#define PORT 9001

int main(void) {
    int fd;
    struct sockaddr_in addr;
    struct sockaddr_in peer;
    socklen_t peer_len;
    char buf[1024];
    ssize_t got;

    fd = socket(AF_INET, SOCK_DGRAM, 0);     /* UDP */
    if (fd < 0) {
        perror("socket");
        return 1;
    }

    memset(&addr, 0, sizeof(addr));
    addr.sin_family = AF_INET;
    addr.sin_port = htons(PORT);
    addr.sin_addr.s_addr = htonl(INADDR_ANY);

    if (bind(fd, (struct sockaddr *) &addr, sizeof(addr)) < 0) {
        perror("bind");
        return 1;
    }

    printf("UDP echo on port %d\n", PORT);

    while (1) {
        peer_len = sizeof(peer);
        got = recvfrom(fd, buf, sizeof(buf), 0,
                       (struct sockaddr *) &peer, &peer_len);
        if (got < 0) {
            perror("recvfrom");
            break;
        }

        /* Each datagram carried its own source: reply to IT. */
        sendto(fd, buf, (size_t) got, 0,
               (struct sockaddr *) &peer, peer_len);
    }

    close(fd);
    return 0;
}
\end{lstlisting}

\begin{verbatim}
   recvfrom / sendto, and why they are different:

     read/write  on a CONNECTED socket: the peer is known, no address
     recvfrom    on an UNCONNECTED socket: each datagram brings its
                 own return address (peer), used by sendto to reply

   One UDP socket can talk to thousands of peers, because each
   datagram is self-contained. That is UDP's superpower and its
   curse: no connection means no one is keeping track of anyone.
\end{verbatim}

\begin{osconcept}
The choice of TCP or UDP is a contract about how much the kernel will
do for you. TCP asks the kernel to run a small protocol machine per
connection---acknowledging, retransmitting, ordering, and closing---
so your application can trust a byte stream. UDP asks the kernel for
almost nothing, and your application inherits the network's
unreliability. The OS provides both and the differences between them
are entirely kernel policy.
\end{osconcept}

\section{The Scalable Server: One Loop, Many Sockets}

The one-client echo server blocks on a single \code{accept}, then on a
single \code{read}---it can serve exactly one client at a time. The
fix is the chapter you already know: make every socket non-blocking,
put them all in a \code{poll} array, and let one loop serve thousands.
This is the event loop of Chapter~10, now watching sockets instead of
pipes, with one new move: when the \emph{listener} reports readiness,
we \code{accept} and add the new client to the set.

\begin{lstlisting}[caption={echo\_server.c: a poll-driven, multi-client echo server.}]
#include <stdio.h>
#include <stdlib.h>
#include <string.h>
#include <unistd.h>
#include <fcntl.h>
#include <poll.h>
#include <errno.h>
#include <arpa/inet.h>
#include <sys/socket.h>

#define PORT 9000
#define MAX_CLIENTS 64

static int set_nonblocking(int fd) {
    int flags;

    flags = fcntl(fd, F_GETFL, 0);
    if (flags < 0) {
        return -1;
    }
    return fcntl(fd, F_SETFL, flags | O_NONBLOCK);
}

int main(void) {
    int listen_fd;
    struct sockaddr_in addr;
    struct pollfd fds[MAX_CLIENTS];
    int nfds;
    char buf[1024];

    listen_fd = socket(AF_INET, SOCK_STREAM, 0);
    if (listen_fd < 0) {
        perror("socket");
        return 1;
    }

    memset(&addr, 0, sizeof(addr));
    addr.sin_family = AF_INET;
    addr.sin_port = htons(PORT);
    addr.sin_addr.s_addr = htonl(INADDR_ANY);

    if (bind(listen_fd, (struct sockaddr *) &addr, sizeof(addr)) < 0) {
        perror("bind");
        return 1;
    }
    if (listen(listen_fd, 16) < 0) {
        perror("listen");
        return 1;
    }
    set_nonblocking(listen_fd);

    printf("echo server on port %d\n", PORT);

    /* fds[0] is the listener; clients are appended after it. */
    fds[0].fd = listen_fd;
    fds[0].events = POLLIN;
    nfds = 1;

    while (1) {
        int r;
        int i;

        r = poll(fds, (nfds_t) nfds, -1);
        if (r < 0) {
            if (errno == EINTR) {
                continue;
            }
            perror("poll");
            break;
        }

        for (i = 0; i < nfds; i++) {
            if ((fds[i].revents & POLLIN) == 0) {
                continue;
            }

            if (fds[i].fd == listen_fd) {
                /* A NEW connection is ready to be accepted. */
                int cfd;
                int slot;

                cfd = accept(listen_fd, NULL, NULL);
                if (cfd < 0) {
                    continue;
                }
                set_nonblocking(cfd);

                slot = nfds;
                if (slot >= MAX_CLIENTS) {
                    close(cfd);        /* out of room: refuse politely */
                    continue;
                }
                fds[slot].fd = cfd;
                fds[slot].events = POLLIN;
                nfds = nfds + 1;
                printf("client %d connected\n", cfd);
            } else {
                /* A client has data, or closed. */
                ssize_t got;

                got = read(fds[i].fd, buf, sizeof(buf));
                if (got <= 0) {
                    /* EOF or error: close it, drop it from the set. */
                    printf("client %d gone\n", fds[i].fd);
                    close(fds[i].fd);
                    fds[i] = fds[nfds - 1];   /* swap-last-into-slot */
                    nfds = nfds - 1;
                    i = i - 1;
                } else {
                    write(fds[i].fd, buf, (size_t) got);   /* echo */
                }
            }
        }
    }

    return 0;
}
\end{lstlisting}

\begin{verbatim}
   The loop's two moves, drawn:

        poll(fds, nfds)  -- one wait for EVERY socket
              |
              +-- listener ready?  accept(); append to fds[]
              +-- client ready?    read(); echo
              +-- client EOF?      close(); remove (swap-last)

   - accept() is called only when the LISTENER is readable, so it
     never blocks.
   - reads happen only when a client is readable, and the sockets are
     non-blocking anyway: no client can stall the world.
   - removing a closed client by swapping the LAST slot into its place
     keeps the array dense in O(1) -- the same trick as open-address
     deletion with tombstones in Chapter 4.
\end{verbatim}

\section{Practical: Growing the Echo Server into a Database}

The echo server is a skeleton; its muscle comes from everything
earlier in the book. The same loop that echoes can \emph{parse} and
\emph{store}, turning the server into a networked database with the
command language of Chapter~14 and the durable storage of Chapter~16.

\begin{verbatim}
   The capstone architecture, all the pieces together:

     client                       server (the poll loop of ch.10/17)
     "SET name ada\n"                 |
        |                             v
        |  TCP ----------------->  read()  -- collect bytes into a
        |                             line buffer (may need SEVERAL
        |                             reads: TCP is a stream!)
        |                             v
        |                     lex() + parse()   (ch.14: the MiniDB
        |                             |            front end)
        |                             v
        |                     db_set()           (ch.16: the dblog,
        |                             |            durable on disk)
        |                             v
        |  <----------------------  write("ok\n")

   Two new lessons the network adds on top of the earlier chapters:
     1. A "line" may arrive split across reads, and several lines may
        arrive in one read. The server must buffer until a full
        record, using a length prefix (ch.16's wire format) to know
        when it has one.
     2. The event loop dispatches every connection's bytes to the
        same parser -- one loop, one parser, many clients, and the
        durable store behind it.
\end{verbatim}

Every capability in that picture was built in this book: descriptors
(Chapter~5), the event loop (Chapter~10), the parser (Chapter~14), the
durable log (Chapter~16). The network did not require a new idea; it
required reusing the old ones across a socket.

\begin{researchnote}[The philosophy of the internet protocols]
The finale deserves the field's founding argument. In 1988, David
Clark, one of the architects of the internet protocols, published
``The Design Philosophy of the DARPA Internet Protocols.'' It is the
clearest statement ever written of \emph{why} the network works the
way it does---and it explains every design choice in this chapter.

\begin{verbatim}
   Clark's first principle, condensed:

       "The internet is designed around a fundamental principle: the
        network itself is deliberately simple and makes no guarantees;
        the intelligence and reliability live at the ENDPOINTS, in
        the hosts, not in the middle of the network."

   In this chapter's terms: the network delivers packets best-effort,
   and TCP--which lives in your kernel, at the endpoint--supplies the
   reliability. The routers in the middle stay dumb and fast; the
   cleverness is at the edges. This is the "end-to-end argument," and
   it is why your socket's kernel does so much work.
\end{verbatim}

Clark's paper is short, opinionated, and remains the best one-page
answer to ``why is the network built this way?'' Reading it closes the
book's arc: from bytes and addresses on one machine, to a reliable
byte stream between two machines, designed by the same principles of
layering and explicit responsibility you have practiced all along.

\paperref{D.~D.~Clark, ``The Design Philosophy of the DARPA Internet
Protocols,'' ACM SIGCOMM 1988}{https://www.cs.princeton.edu/courses/archive/fall06/cos561/papers/clark88.pdf}
\end{researchnote}

\section{Structures in Bytes}

The TCP header you just read is one of the three byte structures of a
network conversation. The other two are the socket address you build
with \code{htons}/\code{htonl}, and the IP header that carries every
segment.

\begin{verbatim}
   struct sockaddr_in (IPv4): 16 bytes
   offset  size  field
   0       2     sin_family     (AF_INET = 2)
   2       2     sin_port       (network byte order!)
   4       4     sin_addr       (network byte order!)
   8       8     padding to the generic sockaddr

   the IPv4 header (no options): 20 bytes
   offset  size  field
   0       1     version + IHL      (4 + 4 bits; 5 words = 20 bytes)
   1       1     DSCP / ECN
   2       2     total length
   4       2     identification
   6       2     flags + fragment offset
   8       1     TTL
   9       1     protocol           (6 = TCP, 17 = UDP)
   10      2     header checksum
   12      4     source address
   16      4     destination address

   the UDP header: 8 bytes
   offset  size  field
   0       2     source port
   2       2     destination port
   4       2     length
   6       2     checksum
\end{verbatim}

A datagram on the wire is IP header + transport header + payload
(20 + 20 + data for TCP, 20 + 8 + data for UDP), every integer in
network byte order. The kernel builds and parses these; your program
only ever sets the two fields of the 16-byte \code{sockaddr\_in},
converted with \code{htons} and \code{htonl}.

\section{Summary}

\begin{verbatim}
   The networking toolbox, complete:

   socket(AF_INET, SOCK_STREAM|SOCK_DGRAM, 0)   create an endpoint
   sockaddr_in + htons/htonl/inet_pton          name it (byte order!)
   bind()      give the socket its address and port
   listen()    tell the kernel to queue connections
   accept()    take one connection -> a NEW socket for that client
   connect()   drive the three-way handshake (client side)
   read/write  the connected socket is a byte stream (ch.5 rules)
   recvfrom/sendto  UDP datagrams with per-message addresses
   poll + non-blocking  one loop, many clients (ch.10, reused)

   Rules that keep network code correct:
     1. put every integer on the wire in network byte order;
     2. a connected socket is a byte stream: buffer until you have a
        full record, using a length prefix;
     3. EOF (read == 0) means the peer closed; POLLHUP, too;
     4. never block on one socket when serving many: non-blocking +
        poll;
     5. accept() only when the listener is ready; it never blocks
        then.
\end{verbatim}

\begin{osconcept}
Networking is the operating system's grand finale abstraction: the
kernel owns the hardware, runs the protocols, holds the handshakes,
and presents every conversation as a descriptor. The end-to-end
principle says the network stays simple and the intelligence lives at
the edges---in the kernels of the two hosts, and in the discipline of
your program. You have now built, byte by byte, process by process,
every layer between a variable in RAM and a message on the wire.
\end{osconcept}

This is where the book's tour ends and your systems programming begins.
You hold all the tools: memory and pointers, structures and algorithms,
processes, signals, pipes, threads, filesystems, durable storage, and
the network. The machine is no longer a mystery; it is a set of
contracts you understand, and the rest is experience.

}{}

\backmatter

\end{document}
